% ============================================================
% §8  Formalizm matematyczny — wyprowadzenia
% ============================================================
\section{Formalizm matematyczny}\label{sec:formalizm}

Niniejsza sekcja zamyka g\l{}\'owne luki formalne teorii:
okre\'sla operator~$\mathcal{D}$, cz\l{}on samointerferencji
$\mathcal{N}[\Phi]$, wyprowadza profil pojedynczego \'zr\'od\l{}a,
pokazuje trzy re\.zimy oraz odzyskuje granic\k{e} newtonowsk\k{a}.

% =============================================================
\subsection{Architektura formalizmu TGP}\label{ssec:architektura}
% =============================================================

\begin{remark}[Hierarchia sektorów formalizmu]\label{rem:hierarchia-sektorow}
Formalizm TGP organizuje si\k{e} w~\textbf{cztery poziomy hierarchii},
logicznie od siebie zale\.zne. Czytelnik powinien traktowa\'c ka\.zdy
poziom jako odr\k{e}bny modu\l{}; \S\S~\ref{ssec:rownanko}--\ref{ssec:roadmap}
rozwijaj\k{a} je szczeg\'o\l{}owo.

\medskip
\begin{center}\small
\begin{tabular}{@{}clp{6.0cm}l@{}}
\toprule
\textbf{Poziom} & \textbf{Obiekt} & \textbf{Zawarto\'s\'c} & \textbf{Status} \\
\midrule
0 & Substrat $\Gamma = (V,E)$ &
  Hamilton $H_\Gamma$, symetria $\mathbb{Z}_2$,
  koarse-graining; $\Phi = \langle\hat{s}^2\rangle$,
  $\sigma_{ab} = \langle\hat{s}\,\hat{s}_{+\hat{a}}\rangle^{\rm TF}$
  & (W) \\[4pt]
1 & R\'ownanie pola $\Phi$ &
  Operator $\mathcal{D}$, samointerferencja $\mathcal{N}$,
  dzia\l{}anie TGP, $\alpha=2$, $\beta=\gamma$
  & (W) \\[4pt]
2 & Metryka efektywna $g_{\mu\nu}(\Phi,\sigma_{ab})$ &
  Metryka eksponencjalna, emergencja Einsteina,
  kosmologia FRW, PPN
  & (E) \\[4pt]
3 & Materia i~pola cechowania &
  Sprzężenie metryczne, Dirac, $\mathrm{U}(1)\times\mathrm{SU}(2)
  \times\mathrm{SU}(3)$, generacje
  & (P) \\
\bottomrule
\end{tabular}
\end{center}

\medskip
\noindent Statusy (W), (E), (P) zdefiniowane w~uw.~\ref{rem:status-levels}.
Sekcja~\ref{ssec:rownanko}--\ref{ssec:substrat-param} obejmuje poziom~0 i~1.
Sekcje~\ref{ssec:metryka-deriv}--\ref{ssec:tensor-substrate} obejmuj\k{a} poziom~2.
Sekcja~\ref{sec:cechowanie} i~dod.~F obejmuj\k{a} poziom~3.
\end{remark}

\begin{remark}[Niezale\.zno\'s\'c modularna]\label{rem:niezaleznosc}
Grawitacja i~kosmologia \textbf{nie wymagaj\k{a}} sektora cechowania:
sekcje~\ref{ssec:rownanko}--\ref{ssec:falsification} s\k{a} zamkni\k{e}te
(poziomy~0--2) niezale\.znie od tego, czy emergencja $\mathrm{U}(1)\times
\mathrm{SU}(2)\times\mathrm{SU}(3)$ jest poprawna.
I~odwrotnie: sektor cechowania (\S\ref{sec:cechowanie}) \textbf{bazuje}
na poziomach~0--2 i~nie wymaga dodatkowych aksjomat\'ow ---
rozszerza jedynie struktur\k{e} substratu (amplituda + faza).
Ta modularno\'s\'c jest \textbf{cech\k{a} projektu} TGP, a~nie s\l{}abo\'sci\k{a}.
\end{remark}

\begin{remark}[Minimalna zasada generuj\k{a}ca]\label{rem:minimalna-zasada}
Ca\l{}y formalizm poziom\'ow~0--2 wynika z~\textbf{jednego za\l{}o\.zenia}:
\begin{center}
\textit{,,Przestrze\'n jest emergentna z~substratu $\Gamma=(V,E)$
z~symetri\k{a} $\mathbb{Z}_2$ i~hamiltonianem $H_\Gamma$.
Pole $\Phi = \langle\hat{s}^2\rangle$ mierzy faz\k{e} metryczn\k{a}
tego substratu.''}
\end{center}
Z~tego jednego zdania, przez kolejne projekcje coarse-grainingu,
wynikaj\k{a}: r\'ownanie pola (poziom~1), metryka (poziom~2),
emergencja Einsteina, kosmologia, PPN~$\gamma=\beta=1$,
prz\k{e}dko\'s\'c fal GW, brak osobliwo\'sci i~stabilno\'s\'c globalna.
Parametry $\beta$~(=$\gamma$) i~$q$ s\k{a} \textbf{dwiema} liczb,
kt\'ore \l{}\k{a}cz\k{a} j\k{e}zyk substratu z~j\k{e}zykiem obserwacyjnym
SI. Pozosta\l{}e sk\k{a}din\k{a}d wynikaj\k{a}.
\end{remark}

% =============================================================
\subsection{Klasyfikacja parametr\'ow TGP}\label{ssec:param-classification}
% =============================================================

\begin{remark}[Filozofia: klasa wszech\'swiat\'ow, nie jeden wszech\'swiat]%
\label{rem:multiverse-philosophy}
Substrat~$\Gamma$ definiuje \textbf{klas\k{e}} mo\.zliwych dynamik ---
r\'o\.zne regu\l{}y sprz\k{e}\.zenia $K_{ij}$ daj\k{a} r\'o\.zne typy
generowanej przestrzeni. Nasz wszech\'swiat odpowiada
\textbf{jednemu punktowi} w~tej klasie.
Teoria nie musi jednoznacznie wyprowadza\'c naszego wszech\'swiata
--- musi wyprowadzi\'c z~\textbf{minimalnej liczby parametr\'ow wyboru}
\textbf{maksymaln\k{a} liczb\k{e} obserwacji}.

Kryterium jako\'sci:
\[
  \frac{M_{\rm obs}}{N_{\rm param}} \gg 1.
\]
Im wy\.zszy ten iloraz, tym teoria jest bardziej \emph{parsimonialna}.
Cel d\l{}ugoterminowy: zredukowa\'c $N_{\rm param}$ do~1
(tylko $\PhiZero$) przy $M_{\rm obs} \geq 12$.
\end{remark}

\begin{table}[h!]
\centering\small
\caption{Klasyfikacja parametr\'ow TGP w~trzech warstwach.
Warstwa~I: stru\-ktu\-ral\-ne (te same w~ka\.zdym wszech\'swiecie
substratowym).
Warstwa~II: selektywne (specyficz\-ne dla naszego wszech\'swiata).
Warstwa~III: wy\-pro\-wa\-dzone (predykcje, nie wej\'scia).}
\label{tab:param-classification}
\begin{tabular}{@{}llll@{}}
\toprule
\textbf{Parametr} & \textbf{Warto\'s\'c} & \textbf{Warstwa} & \textbf{\'Zr\'od\l{}o} \\
\midrule
\multicolumn{4}{l}{\textit{Warstwa~I --- substratowe (uniwersalne)}} \\
$\alpha = 2$ & dok\l{}adnie & I & $K_{ij} = J(\varphi_i\varphi_j)^2$
  (tw.~\ref{prop:substrate-action}) \\
$\beta = \gamma$ & dok\l{}adnie & I & warunek pr\'o\.zniowy $U'(1)=0$
  (stw.~\ref{prop:vacuum-condition}) \\
$d = 3$ & dok\l{}adnie & I & optymalno\'s\'c sieci
  (stw.~\ref{prop:wymiar}) \\
sygnatura $(-,+,+,+)$ & dok\l{}adnie & I & relacyjna def.\ czasu
  (stw.~\ref{prop:sygnatura}) \\
$K(0) = 0$ & aksjomat & I & brak przestrzeni przy $\Phi = 0$
  (aksjomat~A1) \\
\midrule
\multicolumn{4}{l}{\textit{Warstwa~II --- selektywne (nasz wszech\'swiat)}} \\
$\PhiZero$ & $\approx 24{,}66\;[H_0^2/c_0^2]^{-1}$ & II &
  dopasowanie z~$\Lambda_{\rm obs}$ \\
$a_\Gamma$ & $\approx 0{,}040049$ & II &
  dopasowanie z~$r_{21,\rm PDG}$ \\
$\psi_{\rm ini} = 7/6$ & dok\l{}adnie & I$\to$III &
  $W(7/6)=0$ + GL + $\Delta G/G=0$
  (stw.~\ref{prop:psi-ini-derived}; \texttt{ex86} 9/9~PASS) \\
\midrule
\multicolumn{4}{l}{\textit{Warstwa~III --- wyprowadzone (predykcje)}} \\
$M_* = m_P$ & $\approx 1{,}2\times 10^{19}$ GeV & III &
  $\PhiZero + \ell_P$ (stw.~\ref{prop:Mstar-from-substrate}) \\
$\Lambda_{\rm eff} \sim \gamma/12$ & $\sim H_0^2$ & III &
  $W(1) = \gamma/3$ (stw.~\ref{prop:Lambda-eff}) \\
3~generacje & dok\l{}adnie & III &
  WKB $n=0,1,2$ przy $a_\Gamma$ \\
$\alpha_K \approx 8{,}56$ & artefakt koord.\ (tw.~\ref{thm:Z-alpha-degeneracy}) & III &
  supersesja przez $g_0^*$ (OP-3 zamkni\k{e}ty) \\
$m_e, m_\mu, m_\tau$ & odtwarzane & III &
  z~$\alpha_K$, $a_\Gamma$ \\
$c_{\rm GW} = c_0$ & dok\l{}adnie & III &
  jednorodne t\l{}o $\Phi$ (stw.~\ref{prop:cT}) \\
$n_s - 1 = -4\varepsilon_H - 4\varepsilon_\psi$ & $\approx -4\varepsilon_H$ & III &
  r.~\eqref{eq:ns-TGP}; $\varepsilon_\psi \sim H_0/H \ll 1$ podczas inflacji \\
\bottomrule
\end{tabular}
\end{table}

\begin{remark}[Bie\.z\k{a}cy iloraz $M_{\rm obs}/N_{\rm param}$]%
\label{rem:parsimony}
Warstwa~II zawiera jedynie dwa wolne parametry:
$\PhiZero$ (dopasowanie do~$\Lambda_{\rm obs}$) i~$a_\Gamma$
(dopasowanie do~$r_{21,\rm PDG}$).
Warunek pocz\k{a}tkowy $\psi_{\rm ini} = 7/6$ wynika z~Warstwy~I
przez tr\'ojk\k{a}t argument\'ow: $W(7/6)=0$ (algebraicznie),
$\Delta G/G = 0$ (nukleosynteza), argument GL substratu
(stw.~\ref{prop:psi-ini-derived}; \texttt{ex86}~9/9~PASS, sesja~v35).
Tote\.z $\psi_{\rm ini}$ nie jest wolnym parametrem.
Zatem:
\begin{equation}\label{eq:parsimony-N2}
  N_{\rm param} = 2,
  \qquad
  M_{\rm obs} \geq 12,
  \qquad
  \frac{M_{\rm obs}}{N_{\rm param}} \geq 6.
\end{equation}
\textbf{Nowa predykcja testowalna} (sesja v35, ex86~T4):
oscylacje $G_{\rm eff}(t)$ z~cz\k{e}sto\'sci\k{a}
$\omega_{\rm cosmo}/H_0 \approx 8{,}19$, amplituda
proporcjonalna do $|\psi_{\rm ini} - 7/6|$.

\medskip\noindent
\textbf{Pytanie redukuj\k{a}ce $N_{\rm param}$} (cz\k{e}\'sciowo zamkni\k{e}te):
\begin{enumerate}[nosep]
  \item \emph{Czy $a_\Gamma = f(\PhiZero)$?}
    Hipoteza $a_\Gamma \cdot \PhiZero = 1$ zweryfikowana numerycznie
    z~DESI~DR1+CMB na poziomie $0{,}12\sigma$ (patrz hip.~\ref{hyp:agamma-phi0}).
    OP-3 zamkni\k{e}ty: $\alpha_K$ jest artefaktem koordynatowym
    (tw.~\ref{thm:Z-alpha-degeneracy}).
    Pozosta\l{}a kwestia: formalne wyprowadzenie $a_\Gamma = 1/\PhiZero$
    z~zasad pierwszych $\Rightarrow N_{\rm param} = 1$.
\end{enumerate}
Pe\l{}ne domkni\k{e}cie tego pytania by\l{}oby
\textbf{prze\l{}omowym wynikiem} teorii.
\end{remark}

\begin{hypothesis}[Hipoteza maksymalna: $a_\Gamma = 1/\PhiZero$]%
\label{hyp:agamma-phi0}
Numeryczna eksploracja (skrypty \texttt{p75\_agamma\_phi0\_correlation.py}
i~\texttt{ex85\_agamma\_phi0\_precision.py},
zakres $\PhiZero \in [20,35]$, $a_\Gamma \in [0{,}030, 0{,}055]$)
wskazuje, \.ze najprostszym stosunkiem wymienionych wielko\'sci jest:
\begin{equation}\label{eq:agamma-phi0-conjecture}
  a_\Gamma \cdot \PhiZero \approx 1.
\end{equation}
Przy $\PhiZero = 36\,\Omega_\Lambda$ relacja przyjmuje posta\'c
$a_\Gamma = 1/(36\,\Omega_\Lambda)$.
Wymagana warto\'s\'c to $\Omega_\Lambda^* = 1/(36\,a_\Gamma) = 0{,}6936$.

\medskip
Weryfikacja z~r\'o\.znymi zestawami danych kosmologicznych
(skrypt \texttt{ex85}, 5/6~PASS):
\begin{center}
\begin{tabular}{lccc}
\hline
Dane & $a_\Gamma \cdot \PhiZero$ & odch.\ [\%] & odch.\ [$\sigma$] \\
\hline
Planck 2018 ($\Omega_M=0{,}3153\pm0{,}0073$)         & $0{,}9872$ & $-1{,}3$ & $-1{,}22\sigma$ \\
DESI DR1 BAO alone ($\Omega_M=0{,}295\pm0{,}015$)    & $1{,}0164$ & $+1{,}6$ & $+0{,}76\sigma$ \\
DESI DR1+CMB ($\Omega_M=0{,}307\pm0{,}005$) \textbf{[best]} & $\mathbf{0{,}9991}$ & $\mathbf{-0{,}1}$ & $\mathbf{-0{,}12\sigma}$ \\
\hline
\end{tabular}
\end{center}

Hipoteza jest zgodna ze wszystkimi zestawami na poziomie ${\leq}1{,}22\sigma$.
Z~danymi DESI~DR1+CMB\footnote{DESI Collaboration, \textit{DESI 2024 VI},
  JCAP~2025; arXiv:2404.03002.
  Parametry p\l{}askiego $\Lambda$CDM: $\Omega_M = 0{,}307\pm0{,}005$
  (DESI BAO + Planck CMB + ACT lensing).}
odchylenie wynosi jedynie $0{,}12\sigma$
--- hipoteza praktycznie potwierdzona numerycznie.

Gdyby by\l{}a \textbf{dok\l{}adna}, $a_\Gamma$ stawa\l{}by si\k{e}
predykcj\k{a} Warstwy~III i~liczba wolnych parametr\'ow spada\l{}aby:
\[
  N_{\rm param} = 1 \quad (\text{tylko } \PhiZero).
\]

\medskip\noindent
\textbf{Relacja uzupe\l{}niaj\k{a}ca} (TOP-1 wyczerpuj\k{a}cego skanu algebraicznego
po kombinacjach $\{a_\Gamma, \PhiZero, \alpha_K, r_{21}\}$, odch.\ $0{,}05\%$
z~Planck):
\begin{equation}\label{eq:agamma-composite}
  \alpha_K \cdot \sqrt{a_\Gamma \cdot r_{21}} \;\approx\; \PhiZero
  \quad (0{,}05\%\text{ od dok\l{}adno\'sci z~Planck~2018}).
\end{equation}
Je\.zeli~\eqref{eq:agamma-phi0-conjecture} i~\eqref{eq:agamma-composite}
s\k{a} obydwie dok\l{}adne, to $a_\Gamma$ jest w~pe\l{}ni wyznaczane
z~parametr\'ow cz\k{a}steczkowych:
\[
  a_\Gamma = \bigl(\alpha_K \sqrt{r_{21}}\bigr)^{-2/3}
  \;\approx\; 0{,}04011
  \quad (\text{odch.\ }0{,}15\%\text{ od }a_\Gamma=0{,}040049).
\]
Oba rownania wskazuj\k{a} na g\l{}\k{e}boki zwi\k{a}zek mi\k{e}dzy
skal\k{a} solitonu~($a_\Gamma$), polem pr\'o\.zni~($\PhiZero$)
i~hierarchi\k{a} mas leptonowych ($\alpha_K$, $r_{21}$),
kt\'orego uzasadnienie zosta\l{}o cz\k{e}\'sciowo zamkni\k{e}te
przez supersesj\k{e} $\alpha_K$ (OP-3 zamkn.,
tw.~\ref{thm:Z-alpha-degeneracy}).
\end{hypothesis}

% -------------------------------------------------------------------
\begin{hypothesis}[Kombinacja T1+T2: $a_\Gamma$ jako parametr Warstwy~III]%
\label{hyp:agamma-phi0-combo}
Je\.zeli obydwa r\'ownania~\eqref{eq:agamma-phi0-conjecture} (T1)
i~\eqref{eq:agamma-composite} (T2) s\k{a} \emph{dok\l{}adne},
to parametr $a_\Gamma$ przesta\k{j}e by\'c wolnym parametrem Warstwy~II
i~staje si\k{e} predykcj\k{a} Warstwy~III:
\begin{equation}\label{eq:agamma-layer3}
  \boxed{%
    a_\Gamma \;=\; \frac{1}{\PhiZero}
    \;=\; \bigl(\alpha_K\,\sqrt{r_{21}}\bigr)^{-2/3}.
  }
\end{equation}
Oba wyra\.zenia po prawej s\k{a} zgodne na poziomie $0{,}15\%$
(skrypt \texttt{ex85}).

\medskip\noindent
\textbf{Klasyfikacja parametr\'ow przy T1+T2:}
\begin{center}\footnotesize
\begin{tabular}{lll}
\toprule
Parametr & Bez T1+T2 & Przy T1+T2 (dok\l{}adne) \\
\midrule
$\PhiZero$ & Warstwa~II & Warstwa~II \\
$a_\Gamma$ & Warstwa~II & Warstwa~III (predykcja) \\
$\alpha_K$ & Warstwa~II (OP-3 zamkn.) & artefakt koord.\ (tw.~\ref{thm:Z-alpha-degeneracy}) \\
$r_{21}$ & Warstwa~III (dane PDG) & Warstwa~III \\
\midrule
$N_{\rm param}$ & $2$ (Warstwy II) & $1$ (tylko $\PhiZero$) \\
\bottomrule
\end{tabular}
\end{center}

\medskip\noindent
\textbf{Status}: \statuslabel{Hipoteza} (numeryczna: T1 -- $0{,}09\%$
przy DESI~DR1+CMB; T2 -- $0{,}048\%$ przy Planck~2018; \texttt{ex85}, 5/6~PASS).
Hipoteza \emph{nie} jest jeszcze wyprowadzona z~pierwszych zasad TGP.
Domkni\k{e}cie: OP-3 (sekcja~\ref{sec:formalizm}).
\end{hypothesis}
% -------------------------------------------------------------------

\begin{result}[OP-3: $\alpha_K$ z bifurkacji substratu
  (skrypt \texttt{p76\_alpha\_K\_bifurcation.py})]%
\label{res:op3-alphak}
Numeryczna eksploracja modelu bubble-solitonu TGP
(ODE z~$\kappa_{\rm fac}=1+\alpha/\phi$, potencja\l{}
$V_{\rm mod}$~\eqref{eq:V-mod}) przy ustalonym $a_\Gamma = 0{,}040$
przynios\l{}a nast\k{e}puj\k{a}ce wyniki (5/5 PASS):
\begin{enumerate}[nosep]
  \item \textbf{Cross-check masy solitonu:}
    $K^*_1 = 0{,}010\,304$ (b\l{}\k{a}d $1{,}1\%$ wzgl\k{e}dem
    $K^*_{1,\rm ref} = 0{,}010\,414$ z~p10/p11).
    Kryterium energetyczne $g(K)= E[\phi]/(4\pi K)-1=0$
    zosta\l{}o poprawnie zaimplementowane.

  \item \textbf{Skalowanie $K^*_1 \propto a_\Gamma$ (Schwarzschild):}
    Skan $a_\Gamma \in \{0{,}010,\,0{,}020,\,0{,}030,\,0{,}040,\,0{,}060\}$
    przy sta\l{}ym $\alpha = 8{,}5616$ daje
    $K^*_1/a_\Gamma \approx 0{,}258 = \mathrm{const}$, tj.
    $K^*_1 \approx C_1 a_\Gamma$ z~$C_1 = 0{,}258$.
    Soliton istnieje dla wszystkich zbadanych $a_\Gamma$
    (brak bifurkacji w~zakresie $[0{,}010;\,0{,}060]$).

  \item \textbf{OP-3 Cie\.zka~2 (relacyjna) --- cz\k{e}\'sciowe potwierdzenie:}
    Interpolacja w~skanie $\alpha \in [4;\,12]$
    (skrypt p76, $n_\psi=40$, $n_{\rm eval}=800$, \texttt{rtol}=$10^{-8}$) daje
    \begin{equation}\label{eq:op3-alpha-c}
      \alpha_c \;=\; 8{,}449
      \quad\bigl(\text{warunek: }K^*_1(\alpha_c)=K^*_{1,\rm ref}\bigr),
    \end{equation}
    odchylenie od $\alpha_K = 8{,}5616$ wynosi $1{,}3\%$.
    Badanie dok\l{}adno\'sci numerycznej (skrypty p78--p79):
    \begin{itemize}[nosep]
      \item p78 ($n_\psi=80$, $n_{\rm eval}=1000$): $K^*_1 = 0{,}010523$
            ($+1{,}04\%$) --- przeskakuje powy\.zej $K^*_{1,\rm ref}$;
            $\alpha_c$ nie znaleziony w~$[8{,}0;\,9{,}2]$.
      \item Diagnoza: p76 u\.zywa\l{}o $n_{\rm eval}=800$, p78~--- $1000$,
            podczas gdy p10 u\.zywa $2000$; ró\.znica \texttt{rtol}
            ($10^{-8}$ vs~$10^{-9}$) powoduje systematyczny bias~$\sim 1\%$.
      \item p79 ($n_{\rm eval}=2000$, \texttt{rtol}=$10^{-9}$, $n_\psi=120$
            --- identyczne z~p10) jest kluczowym testem zbieznosci.
    \end{itemize}
    \textbf{Nowe odkrycie (p78 + p79):}
    $K^*_1(\alpha)$ wykazuje \textbf{minimum} w~okolicach $\alpha \approx 8{,}60$
    ($K^*_{\rm 1,min} \approx 0{,}010\,480$) i~\textbf{znika}
    dla $\alpha \in (8{,}85;\,9{,}3)$ (potwierdzone obydwoma skryptami).
    Na $\alpha = 9{,}5$ pojawia si\k{e} \textbf{nowa ga\l{}\k{a}\'z} z~$K^*_1 = 0{,}008\,514$
    (ni\.zsza masa).
    Wniosek: istnieje \textbf{maksymalna warto\'s\'c sprz\k{e}\.zenia}:
    \begin{equation}\label{eq:alpha-max}
      \alpha_{\rm max} \;\approx\; 8{,}85,
    \end{equation}
    powy\.zej kt\'orej standardowy soliton TGP (ga\l{}\k{a}\'z elektronowa)
    nie istnieje.
    \textbf{Obserwacja:}
    $\alpha_K = 8{,}5616 \approx 0{,}97\,\alpha_{\rm max}$ ---
    parametr TGP le\.zy tu\.z poni\.zej granicy bifurkacji.
    To mo\.ze by\'c g\l{}\k{e}bsza zasada fizyczna (naturalno\'s\'c):
    $\alpha_K$ jest bliskie, lecz poni\.zej, krytycznej warto\'sci,
    przy kt\'orej cz\k{a}stka przestaje by\'c zwi\k{a}zana.

  \item \textbf{Weryfikacja $K^*_2$ (skrypt \texttt{p77\_K2\_scan.py}, g\k{e}sty skan
    60~punkt\'ow, $n_\psi = 120$):}
    Gesty skan $g(K)$ w~$[0{,}015;\,0{,}25]$ ujawni\l{} strukturk\k{e}
    ga\l{}\k{e}zi w~funkcji $K$:
    \begin{itemize}[nosep]
      \item Prawdziwy drugi soliton: $K^*_2 = 0{,}034003$
            (ga\l{}\k{a}\'z $\psi^* \approx 1{,}12$,
             $g = 5 \times 10^{-4} \approx 0$);
            stosunek $K^*_2/K^*_1 = 3{,}27$.
      \item Zmiany znaku $g$ mi\k{e}dzy $K=0{,}057$ a~$K=0{,}089$
            (p76, Sekcja~D) s\k{a} \textbf{artefaktami przeskoku galezi}
            (nie prawdziwymi zerami $g$).
    \end{itemize}
    Wynik: $K^*_2/K^*_1 = 3{,}27 \neq R_{21}^{\rm PDG} = 206{,}77$~---
    standardowy $V_{\rm mod}$ \textbf{nie odtwarza hierarchii mas leptonow}.
    Cie\.zka~1 (Koide) wymaga modyfikacji $V_{\rm mod}$ (p15).
\end{enumerate}
\textbf{Wniosek p76--p79 (OP-3 Cie\.zka~2):}
Warunek $K^*_1(\alpha_c) = K^*_{1,\rm ref}$ daje $\alpha_c \in [8{,}45;\,8{,}87]$
zale\.znie od dok\l{}adno\'sci numerycznej (strukturalny bias~$\sim 0{,}8\%$
w~metodzie skanowania).
Wa\.zniejszy wynik: kampania p76--p79 ujawni\l{}a \textbf{bifurkacj\k{e} substratu}
(r\'ownanie~\eqref{eq:alpha-max}):
standardowy soliton TGP nie istnieje dla $\alpha > \alpha_{\rm max} \approx 8{,}85$,
a~$\alpha_K = 8{,}5616 \approx 0{,}97\,\alpha_{\rm max}$.
Cie\.zka~1 (Koide): standardowy $V_{\rm mod}$ daje $K^*_2/K^*_1 = 3{,}27 \neq R_{21}$~---
wymaga modyfikacji $V_{\rm mod}$ (cel: skrypt p15).
\end{result}

% ==============================================================
% Słownik formalizmu (v47b, odpowiedź na audyt zewnętrzny)
% ==============================================================
\begin{remark}[S\l{}ownik formulacji TGP (v47b)]%
\label{rem:formulation-dictionary}
W~r\'o\.znych cz\k{e}\'sciach manuskryptu pojawiaj\k{a} si\k{e} dwie
\textbf{r\'ownowa\.zne formulacje} formalizmu. Poni\.zsza tabela
stanowi kanoniczny s\l{}ownik przej\'s\'c mi\k{e}dzy nimi.

\begin{center}\footnotesize
\begin{tabular}{@{}p{3cm}p{4.5cm}p{4.5cm}@{}}
\toprule
\textbf{Wielko\'s\'c} & \textbf{Formulacja~A} (kanoniczna)
  & \textbf{Formulacja~B} \\
\midrule
Kinetyczny czynnik & $K(g) = g^4$ & $f(g) = 1 + 4\ln g$ \\
Substratowe $g_0^e$ & $\mathbf{0{,}86941}$ & $1{,}24$ \\
Relacja & $g_0^{(A)} = \exp\!\bigl((g_0^{(B)} - 1)/4\bigr)$ &
  $g_0^{(B)} = 1 + 4\ln g_0^{(A)}$ \\
\midrule
R\'ownanie pola & $\alpha = 2$ (kanoniczne TGP) & to\.zsame \\
Wyk\l{}adniki & $a = b = \tfrac12$, $g = 1$ & to\.zsame \\
$\Phi_0$ & $\mathbf{24{,}65}$ (bare, z~$\Lambda_{\rm obs}$) & to\.zsame \\
$\Phi_{\rm eff}$ & $\Phi_0\,(1+\delta)$ z~p\k{e}tl\k{a} 1-loop
  & to\.zsame (nie u\.zywane numerycznie) \\
\midrule
Element obj\k{e}to\'sci & $\sqrt{-g} = c_0\psi$
  & to\.zsame (poprawione v47, by\l{}o $\psi^4$) \\
Metryka & $f = \Phi_0/\Phi$, $h = \Phi/\Phi_0$
  & to\.zsame \\
\bottomrule
\end{tabular}
\end{center}

\medskip\noindent
\textbf{Zasada:} \textit{Formulacja~A} jest domy\'slna we wszystkich
nowych wyprowadzeniach (od sesji v33).
Formulacja~B pojawia si\k{e} wy\l{}\k{a}cznie w~starszych sekcjach
i~w~przypisach dolnych.
Wyra\.zenia algebraiczne (Koide, $r_{21}$, $\sin^2\theta_W$)
s\k{a} niezale\.zne od wyboru formulacji.
\end{remark}

\begin{remark}[Bifurkacja substratu: $\alpha_{\rm max}$ i~naturalno\'s\'c $\alpha_K$]%
\label{rem:alpha-max-bifurcation}
Numeryczna eksploracja p78--p80 ujawni\l{}a, \.ze
\textbf{standardowy soliton TGP nie istnieje dla $\alpha > \alpha_{\rm max}$}.
Skrypt~\texttt{p80\_alpha\_max.py} (bisection w~$[8{,}80;\,8{,}90]$,
4/4~PASS) wyznacza:
\begin{equation}\label{eq:alpha-max-precise}
  \alpha_{\rm max} \;=\; 8{,}8734
  \quad\bigl(\text{w przedziale } [8{,}8719;\,8{,}8750]\bigr).
\end{equation}
Mechanizm (bifurkacja siod\l{}o-w\k{e}ze\l{}): dla
$\alpha = 8{,}9 > \alpha_{\rm max}$ profil $\phi(r)$ zapada si\k{e}
($\phi \to 0$ przed $r_{\rm max}$) przy $K \gtrsim 0{,}013$,
podczas gdy dla mniejszych $K$ energia jest za ma\l{}a ($g < -0{,}19$).
$K^*_1(\alpha)$ i~$K_{\rm collapse}(\alpha)$ schodz\k{a} si\k{e} przy $\alpha_{\rm max}$.
Zauwa\.zalny jest te\.z wzrost $K^*_1(\alpha)$ gdy $\alpha \to \alpha_{\rm max}$:
$K^*_1(8{,}80) = 0{,}01062$, $K^*_1(8{,}87) = 0{,}01130$.

\medskip\noindent
\textbf{Naturalno\'s\'c $\alpha_K$:}
\begin{equation}\label{eq:alpha-naturalness}
  \alpha_K \;=\; \alpha_{\rm max}(1 - \varepsilon),
  \quad
  \varepsilon \;=\; 1 - \frac{\alpha_K}{\alpha_{\rm max}}
              \;=\; 1 - \frac{8{,}5616}{8{,}8734}
              \;\approx\; 0{,}0351 \;=\; 3{,}51\%.
\end{equation}

\textbf{Uderzaj\k{a}ca zbiezno\'s\'c} [\textbf{COFNI\k{E}TA p94--p95, patrz ni\.zej}]:
\begin{equation}\label{eq:ns-coincidence}
  \frac{\alpha_K}{\alpha_{\rm max}} \;=\; 0{,}9649
  \;\;\approx\;\;
  n_s \;=\; 0{,}9649 \pm 0{,}0042 \quad (\text{Planck 2018}).
\end{equation}
Stosunek $\alpha_K/\alpha_{\rm max}$ jest numerycznie r\'owny
spektralnemu indeksowi fluktuacji CMB do czterech cyfr znacz\k{a}cych.
Skrypt~\texttt{p81} ujawnia jednak, \.ze $\alpha_{\rm max}$ zale\.zy
s\l{}abo od~$a_\Gamma$ (rozrzut~$1{,}7\%$ dla
$a_\Gamma \in [0{,}020;\,0{,}060]$):
\begin{equation}\label{eq:alpha-max-vs-agam}
  \alpha_{\rm max}(a_\Gamma) \;\approx\; 9{,}02 \;-\; 3{,}75\,a_\Gamma.
\end{equation}
Zbiezno\'s\'c $\alpha_K/\alpha_{\rm max} = n_s = 0{,}9649$
wyst\k{e}puje dok\l{}adnie przy fizycznym $a_\Gamma = 0{,}040$
i~jest specyficzna dla tej warto\'sci (inne~$a_\Gamma$ daj\k{a}
inne stosunki).
Status pierwotny: \emph{obserwacja numeryczna --- przypadkowa lub g\l{}\k{e}bsza}.

\medskip\noindent
\textbf{Ogon oscylatoryjny $\varphi(r)$ i~artefakt dwugalęziowy (wyniki p92--p93):}
Profil solitonu wykazuje \textbf{fizyczny ogon oscylatoryjny} rozpadaj\k{a}cy si\k{e} jak $1/r$:
\begin{equation}\label{eq:osc-tail}
  \varphi(r) - 1 \;\approx\; \frac{A\sin(mr) + B\cos(mr)}{r}, \quad r\to\infty,
  \qquad m = \frac{1}{\sqrt{1+\alpha}},
\end{equation}
b\k{e}d\k{a}cy rozwi\k{a}zaniem sferycznego r\'ownania Bessela
przy linearyzacji wok\'o\l{} $\varphi=1$ ($V''(1) = -1 < 0$).
Obydwa tryby ($\sin/r$, $\cos/r$) zanikaj\k{a}, dlatego
\textbf{ka\.zde} zaburzenie $\psi_0 > 1$ asymptotycznie d\k{a}\.zy do $\varphi\to 1$
i~skan~$\psi$ zawsze wykrywa co najmniej jedno zero $\varphi(r_{\rm max})-1=0$.

Ta w\l{}a\'sciwo\'s\'c tworzy \textbf{fa\l{}szywe zera} w~metodzie shooting
przy sko\'nczonym $r_{\rm max}$: za ka\.zdym razem gdy oscylacja
$\varphi(r_{\rm max})$ przecina 1{,}0 (faza zale\.zy od~$r_{\rm max}$),
pojawia si\k{e} dodatkowe ``zero'' nieistniej\k{a}ce w~granicy $r_{\rm max}\to\infty$.
Audyt numeryczny p93 (\texttt{p93\_rmax\_convergence.py}, 5/5~PASS)
potwierdzi\l{} to jednoznacznie dla $\alpha=8{,}840$:

\begin{center}\small
\begin{tabular}{@{}rcccc@{}}
\toprule
$r_{\rm max}$ & $K^*$ & $\psi_0^{(U)}$ & $\psi_0^{(L)}$ & uwaga \\
\midrule
40  & 0{,}010818 & 1{,}2364 & 1{,}1888 & \emph{konwencja p80--p91} \\
60  & 0{,}010312 & 1{,}2345 & 1{,}1344 & \\
80  & 0{,}010151 & 1{,}2343 & 1{,}0970 & \\
100 & 0{,}010066 & 1{,}2342 & 1{,}0652 & \\
$\infty$ & \textbf{0{,}009521} & \multicolumn{3}{l}{(ekstrapolacja: $K^*(r)=0{,}00952+0{,}051/r$, $R^2=0{,}985$)} \\
\bottomrule
\end{tabular}
\end{center}

\begin{itemize}[nosep]
\item \emph{Ga\l{}\k{a}\'z~U} ($K^*\approx 0{,}0108$ przy $r_{\rm max}=40$):
  jedyna \textbf{fizyczna rodzina}; $K^*$ zbiega monotonicznie~---
  r\'o\.znica $K^*(40)$ vs.\ $K^*(\infty)$: $13{,}6\%$.
\item \emph{Ga\l{}\k{a}\'z~L} ($\psi_0^{(L)}$~pierwsze zero):
  przesuwa si\k{e} o~$\Delta\psi_0=0{,}124$ gdy $r_{\rm max}:40\to100$.
  Zero $g_L=0$ istnieje wy\l{}\k{a}cznie dla $r_{\rm max}=40$~---
  \textbf{artefakt numeryczny} (potwierdzony P5).
\item Ogon solitonu: fit Bessela $A\sin(mr)/r+B\cos(mr)/r$ ($m=0{,}3188$
  przy $\alpha=8{,}840$; $m=0{,}3234$ przy $\alpha_K=8{,}5616$)
  odtwarza profil z~$R^2=1{,}0000$ dla $r>50$ (P4 p93; P2-P3 p97).
\end{itemize}

\medskip
\textbf{Prawdziwy profil solitonu TGP przy $\alpha_K=8{,}5616$ (p97):}
Skrypt \texttt{p97\_true\_soliton\_profile.py}
(skan $r_{\rm max}\in\{40,60,80,100,150\}$, fit $K^*(r)=a+b/r$,
profil do $r_{\rm max}=200$, 3/5~PASS):
\begin{equation}\label{eq:Kstar-true}
  K^*(\alpha_K, \infty) \;=\; 0{,}010029,
  \quad
  E^* \;=\; 4\pi K^*(\infty) \;=\; 0{,}12603,
\end{equation}
\begin{equation}\label{eq:psi0-true}
  \psi_0^*(\alpha_K) \;=\; 1{,}2419 \quad \text{(zbiezne, } r_{\rm max}\geq 60\text{)},
\end{equation}%
\footnote{Warto\'s\'c kanoniczna: $g_0^e = 0{,}86941$ (Formulacja~A, $K = g^4$). Warto\'s\'c $1{,}24$ odpowiada Formulacji~B ($f = 1+4\ln g$).}
\textbf{Universalno\'s\'c $\psi_0$} (nowe odkrycie, p101):
Skan $a_\Gamma\in\{0{,}010;\ldots;\,0{,}060\}$ wykaza\l{}, \.ze
$\psi_0^* = 1{,}2419 = \mathrm{const}$ \emph{niezale\.znie od~$a_\Gamma$}.
Warto\'s\'c pocz\k{a}tkowa solitonu jest w\l{}asno\'sci\k{a} wy\l{}\k{a}cznie dynamiki
ODE (kszta\l{}tu $V_{\rm mod}$ i~$\alpha_K$), nie skali~$a_\Gamma$.
Jednocze\'snie $K^*_1 = C\cdot a_\Gamma$~--- \textbf{soliton skaluje si\k{e}
z~$a_\Gamma$ przy sta\l{}ym kszta\l{}cie} (skalowanie Schwarzschilda).
Z~parametrami ogona Bessela:
\begin{equation}\label{eq:tail-params-alphak}
  m \;=\; 0{,}323397 \;=\; \frac{1}{\sqrt{1+\alpha_K}} \;\text{(dok\l{}adno\'s\'c } 0{,}000\%\text{)},
  \quad
  C_{\rm tail} \;=\; 0{,}01167.
\end{equation}
Skale przestrzenne solitonu:
$r_{1/2} = 1{,}97\,a_\Gamma$,\;
$r_{90} = 10{,}4\,a_\Gamma$,\;
$r_{99} = 98{,}3\,a_\Gamma$.
Bias konwencji $r_{\rm max}=40$:
$K^*(40)/K^*(\infty) = 1{,}047$~--- systematyczne przeszacowanie~$4{,}7\%$.

Poprzednie pozorne zera $K^*_2 \approx 0{,}013$ (z~p87--p89) oraz~$K^*_A \approx 0{,}014$
(z~p91) by\l{}y \textbf{artefaktami}: pierwsze---niejednorodnego wyboru minimum~$|g|$,
drugie---fa\l{}szywych zer w~oscyluj\k{a}cym ogonie przy $r_{\rm max}=40$.

\medskip
\textbf{Pe\l{}ne widmo $K$ (p98) --- drugi soliton fizyczny:}
Skrypt \texttt{p98\_full\_K\_spectrum.py} (skan $K\in[0{,}006;\,0{,}250]$,
$\psi_0\in[1{,}001;\,5{,}0]$, $r_{\rm max}\in\{40,60,80\}$, 3/4~PASS)
zbada\l{} \textbf{pe\l{}ne widmo} zer $g_U(K)=E[\phi]/(4\pi K)-1=0$:
\begin{itemize}[nosep]
  \item \textbf{Dwa zera} $g_U$ przezywa\`ja przy wszystkich $r_{\rm max}$
    (w~tym $r_{\rm max}=80$) --- oba solitony ($K^*_1$, $K^*_2$) s\k{a} fizyczne.
  \item Warto\'sci $K^*_1$:
    $0{,}01054$~($r=40$) $\to$ $0{,}01039$~($r=60$) $\to$ $0{,}01032$~($r=80$)
    [stabilne, zbiezne do $K^*_1(\infty)=0{,}010029$].
  \item Warto\'sci $K^*_2$:
    $0{,}08572$~($r=40$) $\to$ $0{,}21032$~($r=60$) $\to$ $0{,}21526$~($r=80$)
    --- dramatyczny przeskok $+145\%$ przy $r_{\rm max}=40\to60$.
  \item Przyczyna przeskoku: dla ro\.znych $r_{\rm max}$ jako ,,ostatnie zero~$\psi$''
    wybierana jest inna ga\l{}\k{a}\'z oscylacyjna (sferyczne Bessela);
    $n_{\rm zeros}$ ro\'snie z~$r_{\rm max}$
    (3--5 przy $r_{\rm max}=40$; 7--9 przy $r_{\rm max}=80$).
  \item Ekstrapolacja (skrypt \texttt{p99\_K2\_convergence.py},
    $r_{\rm max}\in\{60,80,100,120,150\}$, 3/5~PASS):
    \begin{center}\small
    \begin{tabular}{@{}rccc@{}}
    \toprule
    $r_{\rm max}$ & $K^*_2$ & $\psi_0(K^*_2)$ & $K^*_2/K^*_1(\infty)$ \\
    \midrule
    60  & 0{,}2127 & $\approx 5{,}0$ & 21{,}2 \\
    80  & 0{,}2191 & $\approx 5{,}0$ & 21{,}8 \\
    100 & 0{,}2235 & $\approx 5{,}0$ & 22{,}3 \\
    120 & 0{,}2272 & $\approx 5{,}0$ & 22{,}7 \\
    150 & 0{,}2001 & $4{,}43$ & 20{,}0 \quad(\text{inna ga\l{}\k{a}\'z}!) \\
    \bottomrule
    \end{tabular}
    \end{center}
    Przy $r_{\rm max}\in\{60,80,100,120\}$ $K^*_2$ ro\'snie monotonicznie;
    przy $r_{\rm max}=150$ nast\k{e}puje prze\l{}\k{a}czenie ga\l{}\k{e}zi
    Bessela ($\psi_0$ spada do~$4{,}43$).
    Ekstrapolacja $K^*_2(\infty)\approx 0{,}24$--$0{,}25$
    (zakres pew\-no\'sci po wyl\k{a}czeniu $r_{\rm max}=150$).
  \item Stosunek (wynik p99):
    \begin{equation}\label{eq:K2K1-p99}
      \frac{K^*_2}{K^*_1}\bigl(\infty\bigr)
        \;\approx\; 21\text{--}25
        \;\neq\; R_{21}^{\rm PDG} \;=\; 206{,}77.
    \end{equation}
    Czynnik niedoboru: $R_{21}^{\rm PDG}/\bigl(K^*_2/K^*_1\bigr) \approx 9$--$10$.
\end{itemize}
\textbf{Wniosek (p98--p99)}: standardowy $V_{\rm mod}$ \textbf{nie odtwarza hierarchii mas
lepton\'ow}. Drugi soliton jest fizyczny ($K^*_2(\infty)\approx 0{,}24$),
ale stosunek $K^*_2/K^*_1(\infty)\approx21$--$25$ jest
za ma\l{}y~$\approx9$--$10$ razy wzgl\k{e}dem wymaganego $R_{21}=206{,}77$.
Modyfikacja $V_{\rm mod}$ pozostaje otwartym zadaniem (\'Scie\.zka~1 Koide, p15).

\medskip
\textbf{Wyniki p101 --- korekta ga\l{}\k{e}zi $K^*_2$ i~limit $a_c$:}
Skrypt \texttt{p101\_K2\_agamma\_corrected.py} (5/5~PASS,
w\k{a}skie okna: $\psi_0 \in [1{,}10;\,1{,}80]$ dla $K^*_1$,
$\psi_0 \in [3{,}50;\,7{,}00]$ dla $K^*_2$)
ujawni\l{} \textbf{drug\k{a} ga\l{}\k{a}\'z} o~$\psi_0 \approx 4{,}1$
(odr\'o\.znion\k{a} od ga\l{}\k{e}zi $\psi_0 \approx 5{,}0$ z~p99):
\begin{itemize}[nosep]
  \item Dla $a_\Gamma \in [0{,}020;\,0{,}060]$:
    $K^*_2(\psi_0{\approx}4{,}1) \approx 0{,}171$ przy $a_\Gamma = 0{,}040$,
    stosunek $K^*_2/K^*_1 \approx 16{,}7$ (prawie sta\l{}e).
  \item Ga\l{}\k{a}\'z $K^*_2$ (\textbf{$\psi_0 \approx 4{,}1$}) nie istnieje
    dla $a_\Gamma \leq 0{,}015$; krytyczna warto\'s\'c $a_c \in (0{,}015;\,0{,}020]$.
  \item \textbf{OP-3 \'Scie\.zka~2 NIE potwierdzona}:
    $a_c \approx 0{,}017 \neq a_\Gamma = 0{,}040$ (fizyczne).
\end{itemize}

\medskip
\textbf{\'Scie\.zka~1 Koide --- skan $\lambda$ w~$V_{\rm mod}$ (p102, 4/5~PASS):}%
\label{para:p102-lambda-scan}
Skrypt \texttt{p102\_Vmod\_lambda\_scan.py}
(skan $\lambda \in [10^{-7};\,10^{-2}]$, 20~punkt\'ow logarytmicznie,
$a_\Gamma = 0{,}040$, $r_{\rm max}=80$, 4/5~PASS)
bada\l{} potencja\l{}:
\begin{equation}\label{eq:Vmod-lambda}
  V_{\rm mod}(\varphi;\,\lambda)
  \;=\; \frac{\gamma}{3}\varphi^3 - \frac{\gamma}{4}\varphi^4
       + \frac{\lambda}{6}(\varphi-1)^6,
\end{equation}
szukaj\k{a}c warto\'sci $\lambda$ dla kt\'orej
$K^*_2/K^*_1 = R_{21} = 206{,}77$ (cel Koide).

Wyniki: dla ka\.zdej z~20 warto\'sci $\lambda$:
\begin{center}\small
\begin{tabular}{@{}rccccc@{}}
\toprule
$\lambda$ & $K^*_1$ & $K^*_2$ & $\psi_0(K^*_2)$ & $K^*_2/K^*_1$ \\
\midrule
$10^{-7}$ & 0{,}010281 & 0{,}171278 & 4{,}0965 & 16{,}660 \\
$\lambda_{\rm phys} = 5{,}50\!\times\!10^{-6}$ & 0{,}010281 & 0{,}171272 & 4{,}0964 & 16{,}659 \\
$10^{-2}$ & 0{,}010281 & 0{,}171212 & 4{,}0951 & 16{,}653 \\
\bottomrule
\end{tabular}
\end{center}
\textbf{Obserwacje (p102)}:
\begin{enumerate}[nosep]
  \item $K^*_1 = 0{,}010281 = \mathrm{const}$ dla \textbf{wszystkich}
    $\lambda \in [10^{-7};\,10^{-2}]$ (5~rz\k{e}d\'ow wielko\'sci).
  \item $K^*_2$ zmienia si\k{e} o~zaledwie $0{,}038\%$ w~ca\l{}ym zakresie.
  \item Stosunek:
    \begin{equation}\label{eq:K2K1-lambda-scan}
      \frac{K^*_2}{K^*_1}\bigl(\lambda\bigr)
        \;=\; 16{,}65\text{--}16{,}66
        \;=\; \mathrm{const}
        \;\neq\; R_{21}^{\rm PDG} = 206{,}77.
    \end{equation}
  \item Czynnik niedoboru: $R_{21}^{\rm PDG} / (K^*_2/K^*_1) \approx 12{,}4$.
\end{enumerate}
\textbf{Interpretacja fizyczna}:
Cz\l{}on $\lambda(\varphi-1)^6/6$ wnosi jedynie
$\sim 3\!\times\!10^{-3}$ do $V_{\rm mod}$ przy $\varphi = \psi_0^{(2)} \approx 4{,}1$,
nawet dla $\lambda = 10^{-2}$.
Stosunek $K^*_2/K^*_1$ jest wyznaczany \textbf{wy\l{}\k{a}cznie przez cz\l{}ony $\gamma$}
(kubiczno-kwartyczny kszta\l{}t potencja\l{}u), ca\l{}kowicie nieczule na $\lambda$.

\textbf{Korekta metodologiczna p102 (po diagnozie p103):}
Wyniki tabeli powy\.zej by\l{}y \textbf{b\l{}\k{e}dne}.
Diagnostyka \texttt{p103\_diag3.py} wykaza\l{}a, \.ze
$K^*_2 = 0{,}171272$ z~p102 \textbf{nie jest} solitonem samospójnym:
$g_U(K^*_2) = E/(4\pi K^*_2) - 1 = 3{,}38 \neq 0$.
Okno $\psi \in [3{,}50;\,7{,}00]$ u\.zyte w~p102 wyklucza\l{}o
\textbf{fizyczn\k{a} ga\l{}\k{e}\'z~B} ($\psi_0 \approx 2{,}77$, oznaczon\k{a} jako
,,spurialne'' na podstawie b\l{}\k{e}dnej diagnozy).
Prawdziwy $K^*_2$ (ga\l{}\k{a}\'z~B, diagnostyka \texttt{p103\_diag3}):
\begin{equation}\label{eq:K2-true}
  K^*_2 = 0{,}100276, \quad
  \psi_0^{(2)} = 2{,}768, \quad
  g_U(K^*_2) \approx 0 \;\checkmark, \quad
  \frac{K^*_2}{K^*_1} = \mathbf{9{,}7547}.
\end{equation}

\medskip
\textbf{\'Scie\.zka~1 Koide --- skan wyk\l{}adnika $n$ (p103\_v4, 5/6~PASS):}%
\label{para:p103-n-scan}
Skrypt \texttt{p103\_v4.py}
bada\l{} potencja\l{}:
\begin{equation}\label{eq:Vmod-n}
  V_{\rm mod}(\varphi;\,n)
  \;=\; \frac{\gamma}{3}\varphi^3 - \frac{\gamma}{4}\varphi^4
       + \frac{\lambda}{n}(\varphi-1)^n,
  \quad n \in \{4,5,6,7,8,10,12,14,16,18,20\},
\end{equation}
szukaj\k{a}c, czy istnieje $n$ daj\k{a}ce $K^*_2/K^*_1 = 206{,}77$.
Metodologia: okno $\psi_0^{(2)} \in [1{,}80;\,3{,}20]$ z~filtrem monotoniczno\'sci
(ga\l{}\k{a}\'z~B); $K \in [0{,}060;\,0{,}150]$, krok $\approx 0{,}003$.

\textbf{Wyniki (n-scan)}:
\begin{center}\small
\begin{tabular}{@{}rccccc@{}}
\toprule
$n$ & $K^*_1$ & $K^*_2$ & $\psi_0(K^*_2)$ & $K^*_2/K^*_1$ \\
\midrule
 4 & 0{,}010280 & 0{,}100276 & 2{,}768 & 9{,}755 \\
 5 & 0{,}010280 & 0{,}100276 & 2{,}768 & 9{,}755 \\
 6 & 0{,}010280 & 0{,}100276 & 2{,}768 & 9{,}755 \\
 7 & 0{,}010280 & 0{,}100276 & 2{,}768 & 9{,}755 \\
 8 & 0{,}010280 & 0{,}100276 & 2{,}768 & 9{,}755 \\
10 & 0{,}010280 & 0{,}100276 & 2{,}768 & 9{,}755 \\
12 & 0{,}010280 & 0{,}100276 & 2{,}768 & 9{,}755 \\
14 & 0{,}010280 & 0{,}100277 & 2{,}768 & 9{,}755 \\
16 & 0{,}010280 & 0{,}100277 & 2{,}768 & 9{,}755 \\
18 & 0{,}010280 & 0{,}100277 & 2{,}768 & 9{,}755 \\
20 & 0{,}010280 & 0{,}100277 & 2{,}768 & 9{,}755 \\
\bottomrule
\end{tabular}
\end{center}

\textbf{Wniosek p102+p103 --- \'Scie\.zka~1 Koide}:
\begin{equation}\label{eq:K2K1-n-scan}
  \frac{K^*_2}{K^*_1}(n)
    \;=\; 9{,}7547 \pm 0{,}0000
    \;=\; \mathrm{const}
    \;\neq\; R_{21}^{\rm PDG} = 206{,}77.
\end{equation}
Cz\l{}on $\lambda(\varphi-1)^{n-1}/n$ przy $\varphi = \psi_0^{(2)} \approx 2{,}77$
jest zaniedbywalny dla ca\l{}ego zakresu $n \leq 20$
(wk\l{}ad $\leq 2\%$ przy $n=20$).
Stosunek $K^*_2/K^*_1$ wyznaczany \textbf{wy\l{}\k{a}cznie przez cz\l{}ony $\gamma$}.
Zmiana $n$ \textbf{nie pomaga}.
Status OP-3~\'Scie\.zka~1 (Koide): \textbf{zamkni\k{e}ta negatywnie} ---
forma $V_{\rm mod} = \gamma/3\,\varphi^3 - \gamma/4\,\varphi^4 + \lambda/n\,(\varphi-1)^n$
\textbf{nie mo\.ze} odtworzy\'c $K^*_2/K^*_1 = 206{,}77$ przez \.zadn\k{a} warto\'s\'c $\lambda$ lub $n$.

\medskip
\textbf{Skan $K^*_2/K^*_1$ wzgl\k{e}dem $\alpha_K$ --- kontynuacja ga\l{}\k{e}zi (p104\_v2):}
Skrypt \texttt{p104\_v2\_continuation.py} sledzil ga\l{}\k{a}\'z~B
wzdlu\.z $\alpha_K$ metod\k{a} kontynuacji (16~w\k{a}tk\'ow, krok~$0{,}5$,
zakres~$[3;\,16]$).
\textbf{Kluczowe odkrycie}: ga\l{}\k{a}\'z~B \textbf{istnieje wy\l{}\k{a}cznie w~w\k{a}skim
oknie} $\alpha_K \in [7{,}5;\,9{,}0]$:
\begin{equation}\label{eq:K2K1-alpha-scan}
  \frac{K^*_2}{K^*_1}(\alpha_K) \;=\;
  \begin{cases}
    6{,}198 & \alpha_K = 8{,}06, \\
    9{,}755 & \alpha_K = 8{,}5616 \;\text{(max, punkt fizyczny)}, \\
    \text{ga\l{}\k{a}\'z znika} & \alpha_K \geq 8{,}9968.
  \end{cases}
\end{equation}
Ga\l{}\k{a}\'z~B zanika w~bifurkacji saddle-node przy
$\alpha_K = \alpha_{\rm SN}$, wyznaczonym precyzyjnie metod\k{a} bisekcji
binarnej (skrypt \texttt{p105\_v2\_bisection.py}, 16~w\k{a}tk\'ow):
\begin{equation}\label{eq:alphaSN-value}
  \alpha_{\rm SN} = 8{,}9964 \pm 0{,}0004.
\end{equation}
Fizyczny $\alpha_K = 8{,}5616$ le\.zy \textbf{poni\.zej bifurkacji} ---
jest to maksimum stosunku $K^*_2/K^*_1$ na ga\l{}\k{e}zi~B.
Ekstrapolacja: $K^*_2/K^*_1 = 206{,}77$ wymaga\l{}oby $\alpha_K \approx 36$,
gdzie ga\l{}\k{a}\'z~B nie istnieje.
\textbf{Wniosek}: stosunek $K^*_2/K^*_1$ jest \textbf{strukturalnie ograniczony
do $\lesssim 10$} dla dowolnego $\alpha_K$ --- cel $206{,}77$ jest
\textbf{niemo\.zliwy} w~schemacie V\_mod z~ga\l{}\k{e}zi\k{a}~B.

\medskip
\textbf{Test hipotezy $\alpha_K/\alpha_{\rm SN} = n_s$ (p105):}
Motywacj\k{a} by\l{}o p81 (zale\.zno\'s\'c $\alpha_K/\alpha_{\rm max} \approx n_s$
dla fizycznego $a_\Gamma = 0{,}040$).
Hipoteza: $\alpha_K/\alpha_{\rm SN} = n_s^{\rm Planck} = 0{,}9649$,
co przewiduje $\alpha_{\rm SN}^{\rm pred} = 8{,}5616/0{,}9649 = 8{,}873$.
Zmierzone (bisekcja p105\_v2):
\begin{equation}\label{eq:alphaK-alphaSN-ratio}
  \frac{\alpha_K}{\alpha_{\rm SN}} = \frac{8{,}5616}{8{,}9964} = 0{,}9517
  \quad (\Delta = 1{,}37\%\ \text{od}\ n_s = 0{,}9649).
\end{equation}
\textbf{Hipoteza obalona} ($\Delta > 1\%$).

\medskip
\textbf{Zale\.zno\'s\'c $\alpha_{\rm SN}(a_\Gamma)$ --- skan p106:}
Skrypt \texttt{p106\_alphaSN\_vs\_agamma.py} wyznaczy\l{}
$\alpha_{\rm SN}$ metod\k{a} bisekcji dla $a_\Gamma \in \{0{,}020;\,0{,}030;\,0{,}040;\,0{,}050;\,0{,}060\}$:
\begin{center}
\begin{tabular}{cc}
\hline
$a_\Gamma$ & $\alpha_{\rm SN}$ \\
\hline
0{,}020 & $9{,}038 \pm 0{,}001$ \\
0{,}030 & $9{,}015 \pm 0{,}001$ \\
0{,}040 & $8{,}996 \pm 0{,}001$ \\
0{,}050 & $8{,}974 \pm 0{,}001$ \\
0{,}060 & $8{,}963 \pm 0{,}001$ \\
\hline
\end{tabular}
\end{center}
Fit liniowy: $\alpha_{\rm SN} \approx 9{,}074 - 1{,}91\,a_\Gamma$
($R^2 = 0{,}988$).
\textbf{Ekstrapolacja do $a_\Gamma \to 0$}: $\alpha_{\rm SN}(0) \approx 9{,}074$
--- nie $9{,}0$ dok\l{}adnie (odchylenie $0{,}8\%$).
Hipoteza $\alpha_{\rm SN}(a_\Gamma\to 0) = 9$ \textbf{obalona}.
$\alpha_{\rm SN}$ zale\.zy istotnie od $a_\Gamma$ przez warunek brzegowy
$\dot\varphi_0 = -K/a_\Gamma^2 \propto 1/a_\Gamma$:
mniejsze $a_\Gamma$ daje stromszy start solitonu i wy\.zsz\k{a} $\alpha_{\rm SN}$.

\medskip
\textbf{Hipoteza kodimension-2 WYKLUCZONA; hipoteza dwu-gal\k{e}ziowa COFNI\k{E}TA (p93):}
Istnieje \textbf{jedna rodzina soliton\'ow} z~warunkiem zanikania:
\begin{equation}\label{eq:alpha-max-condition}
  K^*(\alpha_{\rm max}) = K_{\rm nan,R}(\alpha_{\rm max}).
\end{equation}
Warto\'s\'c $\alpha_{\rm max} \approx 8{,}87$ (zgodna z~p80: $8{,}8734$)
by\l{}a pocz\k{a}tkowo uznawana za prawid\l{}owy wynik.

\medskip
\textbf{$K_{\rm nan,R}$ jako artefakt $r_{\rm max}=40$ i~COFNI\k{E}CIE koincydencji $n_s$ (p94--p95):}
Skrypty \texttt{p94\_alpha\_max\_convergence.py} i~\texttt{p95\_knan\_vs\_rmax.py}
ujawni\l{}y, \.ze $K_{\rm nan,R}(\alpha, r_{\rm max})$ r\'ownie\.z jest
artefaktem konwencji $r_{\rm max}=40$:
\begin{itemize}[nosep]
\item Przy $r_{\rm max}=40$: $K_{\rm nan,R} \approx 0{,}0117$ dla $\alpha=8{,}8734$.
\item Przy $r_{\rm max}=60, 80, 100$: $K_{\rm nan,R}$ \textbf{nie istnieje}
  w~zakresie $K \in [0{,}008;\,0{,}040]$ --- skar solutionnier nie dywerguje.
\item Dow\'od bezpo\'sredni (p95 Cz\k{e}\'s\'c~B): ta sama warto\'s\'c
  $K=0{,}01270$ daje \emph{brak zer} przy $r_{\rm max}=40$,
  ale \emph{2 zera} przy $r_{\rm max}=60$ --- fa\l{}szywa granica
  powodowana faz\k{a} oscylacyjnego ogona w~$r=40$.
\item Cz\k{e}\'s\'c~C (p95): przy $K>K_{\rm nan,R}^{(r=40)}$, profil
  $\varphi(r)$ nie dywerguje --- $\varphi(40)=1{,}005$, $\varphi_{\rm max}=1{,}5$
  (warto\'s\'c pocz\k{a}tkowa), brak dywergencji fizycznej.
\end{itemize}
\textbf{Wniosek:} $K_{\rm nan,R}$ i~$\alpha_{\rm max}=8{,}8734$ z~p80 s\k{a}
artefaktami konwencji $r_{\rm max}=40$.
Koincydencja~\eqref{eq:ns-coincidence} ($\alpha_K/\alpha_{\rm max}=n_s$)
jest \textbf{COFNI\k{E}TA}: nie jest odporna na wybr konwencji
$r_{\rm max}$ i~nie ma charakteru fizycznego.
Status OP-3~\'Scie\.zka~3: \textbf{ZAMKNI\k{E}TA (artefakt)}.
\end{remark}

% =============================================================
\subsection{R\'ownanie pola $\Phi$}\label{ssec:rownanko}
% =============================================================

\subsubsection{Zasada konstrukcji operatora}

Operator~$\mathcal{D}$ nie mo\.ze by\'c standardowym laplasjanem
na gotowej rozmaito\'sci --- to narusza\l{}oby aksjomat~\ref{ax:przestrzen}.
Musi by\'c \textbf{wewn\k{e}trzny wzgl\k{e}dem $\Phi$}:
zdefiniowany przez sam\k{a} konfiguracj\k{e} pola, kt\'ore
konstytuuje przestrze\'n.

Jednocze\'snie, w~granicy s\l{}abego pola ($\Phi \approx \PhiZero$,
gradienty ma\l{}e) operator musi redukowa\'c si\k{e} do postaci
daj\k{a}cej potencja\l{} newtonowski $\sim 1/r$.

Postulujemy:
\begin{enumerate}[nosep]
  \item Operator jest co najwy\.zej drugiego rz\k{e}du w~pochodnych
        (minimalno\'s\'c dynamiki).
  \item Jedynymi polami dost\k{e}pnymi operatorowi s\k{a}
        $\Phi$, $\partial_i\Phi$, $\partial_i\partial_j\Phi$
        --- brak zewn\k{e}trznej metryki.
  \item Operator zachowuje symetri\k{e} sferyczn\k{a} i~jest
        skalarny (brak preferowanego kierunku).
  \item W~granicy $\Phi\to\PhiZero$, $\nabla\Phi\to 0$
        operator redukuje si\k{e} do $\nabla^2\Phi/\PhiZero$.
\end{enumerate}

\subsubsection{Posta\'c operatora $\mathcal{D}$}

Jedyn\k{a} kombinacj\k{a} spe\l{}niaj\k{a}c\k{a} powy\.zsze warunki,
kt\'ora jest jednocze\'snie skalarna, co najwy\.zej drugiego rz\k{e}du
i~naturalna z~punktu widzenia pola konstytuuj\k{a}cego przestrze\'n,
jest \textbf{$\Phi$-wa\.zony laplasjan}:

\begin{definition}[Operator $\mathcal{D}$]\label{def:D}
\begin{equation}\label{eq:D-def}
  \boxed{
    \mathcal{D}[\Phi]
    \;=\; \frac{1}{\PhiZero}\,\nabla^2 \Phi,
  }
\end{equation}
rozumiany jako cz\k{e}\'s\'c liniowa pe\l{}nego r\'ownania pola.
Czynnik $1/\PhiZero$ zapewnia bezwymiarowo\'s\'c operatora
(po pomno\.zeniu przez $\PhiZero$ odzyskujemy standardowy
laplasjan). Gradient $\nabla$ jest \emph{koordynatowy} ---
w~granicy s\l{}abego pola, gdzie $\Phi\approx\PhiZero$ wsz\k{e}dzie,
koordynaty pokrywaj\k{a} si\k{e} z~odleg\l{}o\'sciami fizycznymi.
\end{definition}

\begin{remark}[Dlaczego nie $\Phi^{-1}\nabla\!\cdot\!(\Phi\nabla\Phi)$?]
Operator $\Phi^{-1}\nabla\!\cdot\!(\Phi\nabla\Phi)
= \nabla^2\Phi + (\nabla\Phi)^2/\Phi$ jest atrakcyjny
konceptualnie (,,przepuszczalno\'s\'c'' przestrzeni modulowana
przez $\Phi$), ale zmienia wyk\l{}adnik dalekiego pola:
$\Phi\sim r^{-1/2}$ zamiast~$r^{-1}$, co uniemo\.zliwia
odzyskanie prawa Newtona. Dlatego sprz\k{e}\.zenie gradientowe
przenosimy do cz\l{}onu nieliniowego $\mathcal{N}$, gdzie
jest aktywne tylko przy du\.zych $\Phi$ --- dok\l{}adnie tam,
gdzie odej\'scie od $1/r$ jest po\.z\k{a}dane.
\end{remark}

% ---------------------------------------------------------------
% Sesja v32, 2026-03-24 --- zadanie I.B planu \texttt{PLAN\_ROZWOJU\_v1}
% ---------------------------------------------------------------

\begin{theorem}[Jednoznaczno\'s\'c operatora $\mathcal{D}_{\mathrm{kin}}$
  z~$\alpha = 2$]%
\label{thm:D-uniqueness}
W\'sr\'od wszystkich skalarnych operator\'ow kinetycznych
drugiego rz\k{e}du postaci
\begin{equation}\label{eq:Lkin-general}
  \mathcal{L}_{\mathrm{kin}} \;=\;
    -\frac{1}{2}\,K(\varphi)\,(\nabla\varphi)^2,
  \qquad
  \varphi = \Phi/\PhiZero,
\end{equation}
z~g\l{}adk\k{a} funkcj\k{a} $K : (0,\infty) \to \mathbb{R}_{>0}$,
narzucenie \textbf{trzech} warunk\'ow:
\begin{enumerate}[(C1)]
  \item \emph{Sta\l{}o\'s\'c $\alpha$}: wsp\'o\l{}czynnik gradientowy
    w~r\'ownaniu EL jest stale r\'owny~$\alpha$ (niezale\.zny od $\varphi$);
  \item \emph{Warunek pr\'o\.zni substratu} ($K(0) = 0$, N0-4):
    sprz\k{e}\.zenie kinetyczne zanika tam, gdzie $\Phi \to 0$;
  \item \emph{Geometryczne sprz\k{e}\.zenie substratu}
    (tw.~\ref{prop:substrate-action}, r\'ow.~\eqref{eq:K-geometric}):
    $K(\varphi) = K_{\mathrm{geo}}\,\varphi^4$;
\end{enumerate}
wyznacza jednoznacznie:
\begin{equation}\label{eq:alpha-unique}
  \boxed{
    \alpha = 2,
    \qquad
    K(\varphi) = K_{\mathrm{geo}}\,\varphi^4.
  }
\end{equation}
Operator kinetyczny TGP:
\begin{equation}\label{eq:Dkin-unique}
  \mathcal{D}_{\mathrm{kin}}[\varphi]
  \;=\; \nabla^2\varphi + \frac{2(\nabla\varphi)^2}{\varphi}
  \;=\; \frac{1}{3\varphi^2}\,\nabla^2(\varphi^3)
\end{equation}
jest \textbf{jedyn\k{a}} lokaln\k{a} form\k{a} operatora spe\l{}niaj\k{a}c\k{a}
warunki (C1)--(C3).
\end{theorem}

\begin{proof}
\emph{Krok~1 (klasyfikacja operator\'ow, wzory EL).}
Og\'olny operator kinetyczny~\eqref{eq:Lkin-general} daje,
przez r\'ownanie Eulera--Lagrange'a:
\begin{align}
  \frac{\partial\mathcal{L}_{\mathrm{kin}}}{\partial\varphi}
    &= -\frac{K'(\varphi)}{2}\,(\nabla\varphi)^2, \\
  \partial_i\!\left(\frac{\partial\mathcal{L}_{\mathrm{kin}}}
    {\partial(\partial_i\varphi)}\right)
    &= -K'(\varphi)(\nabla\varphi)^2 - K(\varphi)\nabla^2\varphi.
\end{align}
R\'ownanie EL (po wyj\k{e}ciu $K(\varphi) > 0$):
\begin{equation}\label{eq:EL-general}
  K(\varphi)\!\left[\nabla^2\varphi
    + \frac{K'(\varphi)}{2K(\varphi)}\,(\nabla\varphi)^2\right] = 0
  \;\;\Longrightarrow\;\;
  \nabla^2\varphi
    + \underbrace{\frac{K'(\varphi)\,\varphi}{2K(\varphi)}}_{=:\,\alpha(\varphi)}
    \cdot\frac{(\nabla\varphi)^2}{\varphi} = 0.
\end{equation}

\emph{Krok~2 (warunek C1: $\alpha = \mathrm{const}$).}
Z~\eqref{eq:EL-general}: $\alpha(\varphi) = \mathrm{const} \equiv \alpha$
wtedy i~tylko wtedy, gdy
$K'(\varphi)/(2K(\varphi)) = \alpha/\varphi$,
tzn.\
\begin{equation}\label{eq:K-ode-formalizm}
  \frac{d}{d\varphi}\ln K = \frac{2\alpha}{\varphi}
  \;\;\Longrightarrow\;\;
  K(\varphi) = C\,\varphi^{2\alpha},
  \quad C > 0.
\end{equation}
Zatem \textbf{ka\.zdy} operator o~sta\l{}ym~$\alpha$ odpowiada
$K(\varphi) = C\varphi^{2\alpha}$ (potęga).

\emph{Krok~3 (warunek C2: $K(0) = 0$).}
Z~\eqref{eq:K-ode-formalizm}: $K(0) = C\cdot 0^{2\alpha} = 0$ wymaga $2\alpha > 0$,
czyli $\alpha > 0$.
Dopuszczalne: $\alpha \in (0,\infty)$.

\emph{Krok~4 (warunek C3: geometryczne sprz\k{e}\.zenie).}
Tw.~\ref{prop:substrate-action} (r\'ow.~\eqref{eq:Lkin-geo}) podaje
$K(\varphi) = K_{\mathrm{geo}}\,\varphi^4$.
Por\'ownuj\k{a}c z~\eqref{eq:K-ode-formalizm}:
$C\varphi^{2\alpha} = K_{\mathrm{geo}}\varphi^4$,
st\k{a}d $2\alpha = 4$, czyli $\alpha = 2$.

\emph{Jednoznaczno\'s\'c.}
Warunki (C1)--(C3) razem daj\k{a} $\alpha = 2$ i~$K = K_{\mathrm{geo}}\varphi^4$.
Dowodzi~\eqref{eq:alpha-unique}.
To\.zsamo\'s\'c~\eqref{eq:Dkin-unique} to~\eqref{eq:phi3-identity}
z~dod.~B (\S\ref{app:B-geo-coupling}).
\end{proof}

\begin{remark}[Zwi\k{a}zek z~laplasjanem Laplace'a--Beltramiego]%
\label{rem:LB-alpha2}
Operator $\mathcal{D}_{\mathrm{kin}}$ z~$\alpha = 2$
jest \emph{laplasjanem Laplace'a--Beltramiego}
metryki konformalnej $h_{ij} = \varphi^4\delta_{ij}$ w~$\mathbb{R}^3$:
\begin{equation}
  \Delta_{h}\varphi
  \;=\; \frac{1}{\sqrt{h}}\,\partial_i\!\left(\sqrt{h}\,h^{ij}\partial_j\varphi\right)
  \;=\; \varphi^{-4}
    \left[\nabla^2\varphi + \frac{2(\nabla\varphi)^2}{\varphi}\right]
  \;=\; \varphi^{-4}\,\mathcal{D}_{\mathrm{kin}}[\varphi].
\end{equation}
(dow\'od: $\sqrt{h} = \varphi^6$, $h^{ij} = \varphi^{-4}\delta^{ij}$,
$\partial_i(\varphi^6\cdot\varphi^{-4}\partial^i\varphi)
= \partial_i(\varphi^2\partial^i\varphi)
= 2\varphi(\nabla\varphi)^2 + \varphi^2\nabla^2\varphi$.)
Geometryczne sprz\k{e}\.zenie~\eqref{eq:K-geometric}
generuje zatem dok\l{}adnie t\k{e} konforemn\k{a} metryk\k{e}
$h_{ij} = \varphi^4\delta_{ij}$, kt\'rej laplasjan wyznacza
dynamik\k{e} TGP.
Wynik zamyka \textbf{luk\k{e} I.B} planu~\texttt{PLAN\_ROZWOJU\_v1}.
\end{remark}

\begin{proposition}[Wyprowadzenie N0-4: $K(0) = 0$ z~substratu]%
\label{prop:K0-from-substrate}\label{ssec:formalizm-zamkniecie}
\statuslabel{Twierdzenie}

Warunek N0-4 ($K(0) = 0$: brak sprz\k{e}\.zenia kinetycznego
przy $\Phi = 0$) \textbf{nie jest niezale\.znym aksjomatem} ---
wynika z~geometrycznego sprz\k{e}\.zenia substratu.

\medskip\noindent
\textbf{Dow\'od.}\;
Z~tw.~\ref{prop:substrate-action}: sprz\k{e}\.zenie kinetyczne
substratu ma posta\'c $K_{ij} = J(\varphi_i \varphi_j)^2$,
gdzie $\varphi_i = \hat{s}_i / s_0$ jest lokaln\k{a} amplitud\k{a}
znormalizowan\k{a}. Po gruboziarnieniu (\'sredniowanie blokowe)
otrzymujemy $K(\varphi) = K_{\mathrm{geo}}\,\varphi^4$
(r\'ow.~\eqref{eq:K-geometric}).

Przy $\Phi = 0$ (faza niemetryczna, $\langle\hat{s}_i\rangle = 0$,
$\langle\hat{s}_i^2\rangle_{\rm cg} = 0$) mamy $\varphi = \Phi/\PhiZero = 0$,
zatem:
\begin{equation}\label{eq:K0-derived}
  K(0) = K_{\mathrm{geo}} \cdot 0^4 = 0.
\end{equation}
\textbf{Interpretacja fizyczna}: w~fazie niemetrycznej substratu
brak korelacji mi\k{e}dzy s\k{a}siednimi w\k{e}z\l{}ami
($\langle\hat{s}_i\hat{s}_j\rangle_{\rm conn} = 0$
dla $T > T_c$, symetria $\Zp$ niez\l{}amana).
Brak korelacji oznacza brak propagacji zmian konfiguracji ---
zatem sprz\k{e}\.zenie kinetyczne \emph{musi} zanikai\'c.

$K(0) = 0$ jest wi\k{e}c \textbf{konsekwencj\k{a}} geometrycznego
sprz\k{e}\.zenia $K = \varphi^4$, kt\'ore z~kolei wynika
z~tw.~\ref{prop:substrate-action}. Status N0-4 zmienia si\k{e}
z~\statuslabel{Aksjomat} na \statuslabel{Twierdzenie}.\qed
\end{proposition}

\begin{remark}[Mapa rodziny sprz\k{e}\.ze\'n --- klasa wszech\'swiat\'ow]%
\label{rem:coupling-map}
Twierdzenie~\ref{thm:D-uniqueness} nie m\'owi, \.ze inne sprz\k{e}\.zenia
s\k{a} \emph{niemo\.zliwe} --- m\'owi, \.ze $K = K_{\rm geo}\varphi^4$
jest \textbf{jedynym spe\l{}niaj\k{a}cym warunki (C1)--(C3)}.
Ka\.zda inna klasa $K(\varphi)$ daje inny typ generowanej przestrzeni,
tj.\ inny wszech\'swiat w~sensie remark~\ref{rem:multiverse-philosophy}.

\begin{center}\small
\begin{tabular}{@{}lccl@{}}
\toprule
\textbf{Klasa $K(\varphi)$} & \textbf{$\alpha$} &
  \textbf{Typ wszech\'swiata} & \textbf{Zgodno\'s\'c z~naszym} \\
\midrule
$K = \mathrm{const}$ & $0$ &
  Laplaceowski, solitony zbyt lekkie & ✗ brak hierarchii mas \\
$K = \varphi$ & $1/2$ &
  $K(0)=0$, ale niestabilno\'s\'c gradientowa & ✗ $c_T^2 < 0$ \\
$K = \varphi^2$ & $1$ &
  $c_T \neq c_0$ dla fal grawitacyjnych & △ niezgodne GW170817 \\
$\mathbf{K = \varphi^4}$ & $\mathbf{2}$ &
  \textbf{geometryczne, substratowe} & \textbf{✓ minimalne zgodne} \\
$K = \varphi^6$ & $3$ &
  superp\l{}umienna propagacja $c_T > c_0$ & ✗ niezgodne GW170817 \\
$K = e^\varphi$ & zmienna &
  brak sta\l{}ego $\alpha$, nielokalny & ✗ narusza (C1) \\
\bottomrule
\end{tabular}
\end{center}

\medskip\noindent
\textbf{Wniosek (minimalno\'s\'c, nie konieczno\'s\'c)}:
$K = \varphi^4$ jest \textbf{minimalnym sprz\k{e}\.zeniem} w~klasie
$K = C\varphi^n$ zgodnym jednocze\'snie z~$c_T = c_0$
(GW170817), $K(0)=0$ (pr\'o\.znia substratu)
i~struktur\k{a} WKB daj\k{a}c\k{a} trzy generacje.
Inne warto\'sci $n$ daj\k{a} inne typy wszech\'swiat\'ow ---
niekoniecznie nieistniejace, lecz \textbf{niezgodne z~naszym}.
\end{remark}

\begin{remark}[Zasada minimalnej nieliniowo\'sci: dlaczego $K(\varphi) = \varphi^4$]%
\label{rem:minimal-nonlinearity}
Powy\.zsza tabela analizuje poszczeg\'olne klasy $K(\varphi) = \varphi^n$ metod\k{a}
\emph{eliminacji} (kt\'ore s\k{a} niezgodne z~naszymi obserwacjami).
Poni\.zej podajemy \textbf{pozytywne uzasadnienie}, dlaczego $n = 4$
jest wyj\'sciowym punktem, a~nie artefaktem eliminacji.

\medskip\noindent
\textbf{Krok~1: symetria $\mathbb{Z}_2$ substratu wymusza parzysto\'s\'c.}
Substrat~$\Gamma$ posiada symetri\k{e} $\hat{s} \to -\hat{s}$ ($\mathbb{Z}_2$);
pole $\Phi = \langle\hat{s}^2\rangle$ jest ju\.z parzyste z~definicji,
ale sprz\k{e}\.zenie kinetyczne $K_{ij}$ musi respektowa\'c t\k{e} symetri\k{e},
co wymaga \textbf{parzystego} $n$ w~$K(\varphi) = C\varphi^n$.
Monomiale nieparzyste ($n = 1, 3, 5, \ldots$) s\k{a} wykluczone a~priori ---
nie dlatego, \.ze daj\k{a} z\l{}y wszech\'swiat, lecz dlatego, \.ze naruszaj\k{a}
symetri\k{e} substratu.

\medskip\noindent
\textbf{Krok~2: po\'sr\'od parzystych, $n = 0$ i~$n = 2$ s\k{a} za ubogie.}

\begin{description}[nosep,style=nextline]
  \item[$n = 0$\ ($K = \mathrm{const}$)]
    Brak sprz\k{e}\.zenia kinetycznego z~polem tla: operator~$\mathcal{D}$
    redukuje si\k{e} do standardowego laplasjanu.
    Przestrze\'n jest ,,p\l{}aska'' --- brak emergentnej geometrii
    i~hierarchii mas soliton\'ow.
  \item[$n = 2$\ ($K = \varphi^2$,\ $\alpha = 1$)]
    Czlon kinetyczny $\varphi^2(\partial\varphi)^2$ zmienia szybko\'s\'c
    propagacji mod\'ow tensorowych: na jednorodnym tle $\varphi_{\rm bg}$
    pr\k{e}dko\'s\'c $c_T^2 = c_0^2(1 - \dot\varphi_{\rm bg}^2/(\ldots)) \neq c_0^2$,
    co jest niezgodne z~pomiarem GW170817 ($\Delta c/c < 10^{-15}$).
    Emergencja geometrii istnieje, lecz z~\textbf{niepo\.z\k{a}dan\k{a}
    korekc\k{a} tensorow\k{a}}.
\end{description}

\medskip\noindent
\textbf{Krok~3: $n = 4$ jest pierwszym jednocze\'snie bogatym i~kontrolowanym wyborem.}

$K(\varphi) = C\varphi^4$ jest wyj\'sciowe, bo spe\l{}nia \textbf{jednocze\'snie}
cztery warunki, kt\'ore \emph{\.zaden ni\.zszy parzysty monomial nie spe\l{}nia
r\'ownocze\'snie}:
\begin{enumerate}[nosep]
  \item \textbf{Parzysty} --- respektuje symetri\k{e} $\mathbb{Z}_2$ substratu.
  \item \textbf{Nietrywialnie nieliniowy} --- $K''(\varphi) = 12C\varphi^2 \neq 0$;
        samointerakcja jest naprawd\k{e} obecna, emergencja reziom\'ow mo\.zliwa.
  \item \textbf{Stabilny jako czynnik kinetyczny} --- $K(\varphi) > 0$ dla $\varphi > 0$,
        brak zmiany znaku przy renormalizacji w~granicy UV;
        r\'ownanie pola jest dobrze postawione.
  \item \textbf{$c_T = c_0$} --- na jednorodnym tle wzb\'ogacony czlon
        $\varphi^4(\partial\varphi)^2$ daje $\alpha = 2$, co \textbf{dok\l{}adnie kasuje}
        korekcj\k{e} tensorow\k{a}, pozostawiaj\k{a}c $c_T = c_0$
        (stw.~\ref{prop:cT}; weryfikacja: GW170817).
\end{enumerate}

\medskip\noindent
\textbf{Krok~4: $n > 4$ jest mniej minimalne.}
Wy\.zsze monomiale parzyste ($n = 6, 8, \ldots$) s\k{a} formalnie dopuszczalne,
lecz p\l{}ac\k{a} cen\k{e}:
wi\k{e}ksza z\l{}o\.zono\'s\'c, superp\l{}umienna propagacja przy $n = 6$
(por.\ tab.~\ref{rem:coupling-map}),
lub konieczno\'s\'c jawnego dostrojenia wy\.zszych operator\'ow.
Nie s\k{a} \emph{zakazane} --- definiuj\k{a} inne klasy wszech\'swiat\'ow ---
lecz nie s\k{a} minimalne.

\medskip\noindent
\textbf{Podsumowanie:}
\begin{center}
\small
\emph{Wszystko poni\.zej $\varphi^4$ jest za ubogie lub niestabilne;
wszystko powy\.zej $\varphi^4$ jest mniej minimalne.}
\end{center}
Innymi s\l{}owy: $K(\varphi) = \varphi^4$ jest \textbf{pierwsz\k{a}
naprawd\k{e} sensown\k{a} nieliniowo\'sci\k{a}} --- pierwszym
miejscem, w~kt\'orym jednocze\'snie pojawiaj\k{a} si\k{e}
niebanalny kszta\l{}t próżni, stabilno\'s\'c energetyczna
i~naturalna minimalno\'s\'c formalna.
\end{remark}

\subsubsection{Cz\l{}on samointerferencji $\mathcal{N}[\Phi]$}

\begin{definition}[Samointerferencja]\label{def:N}
\begin{equation}\label{eq:N-def}
  \boxed{
    \mathcal{N}[\Phi]
    \;=\;
      \underbrace{\frac{\alpha}{\PhiZero}\,
        \frac{(\nabla\Phi)^2}{\Phi}}_{\text{sprz\k{e}\.zenie gradientowe}}
    \;+\;
      \underbrace{\beta\,\frac{\Phi^2}{\PhiZero^2}}_{\text{odpychanie}}
    \;-\;
      \underbrace{\gamma\,\frac{\Phi^3}{\PhiZero^3}}_{\text{studnia}},
  }
\end{equation}
ze sta\l{}ymi sprz\k{e}\.zenia $\alpha,\beta,\gamma > 0$,
gdzie $\alpha$ jest bezwymiarowe,
a~$\beta,\gamma$ maj\k{a} wymiar $[\mathrm{L}^{-2}]$
(patrz uwaga~\ref{rem:wymiary-beta-gamma}).
\end{definition}

\begin{remark}[Wymiary sta\l{}ych $\beta,\gamma$]\label{rem:wymiary-beta-gamma}
Cz\l{}on gradientowy $\alpha(\nabla\Phi)^2/(\PhiZero\Phi)$
ma wymiar $[\mathrm{L}^{-2}]$ (bo $[\nabla\Phi]^2/[\Phi\PhiZero]
= [\Phi^2/\mathrm{L}^2]/[\Phi^2] = [\mathrm{L}^{-2}]$),
wi\k{e}c $\alpha$ jest \textbf{bezwymiarowe}.
Natomiast $\beta\,\Phi^2/\PhiZero^2 = \beta\,\chi^2$ (gdzie
$\chi = \Phi/\PhiZero$), a~poniewa\.z $\chi$ jest bezwymiarowe,
$[\beta] = [\mathrm{L}^{-2}]$. Analogicznie $[\gamma] = [\mathrm{L}^{-2}]$.

W~obliczeniach numerycznych (skrypty w~katalogu \texttt{scripts/})
u\.zywamy bezwymiarowych odpowiednik\'ow
$\hat\beta = \beta\,r_0^2$,
$\hat\gamma = \gamma\,r_0^2$,
gdzie $r_0$ jest skal\k{a} rozmycia \\'zr\'od\l{}a.
Warunek pr\'o\.zniowy $\beta = \gamma$ implikuje $\hat\beta = \hat\gamma$.
\end{remark}

Uzasadnienie ka\.zdego cz\l{}onu:

\begin{enumerate}[nosep]
  \item \textbf{Cz\l{}on gradientowy}
    $\alpha(\nabla\Phi)^2/(\PhiZero\Phi)$:
    Wyra\.za fakt, \.ze g\k{e}sta przestrze\'n moduluje propagacj\k{e}
    zmian $\Phi$. Dominuje tam, gdzie gradienty s\k{a} strome
    (\'srednie skale, blisko \'zr\'ode\l{}). Dzia\l{}a \emph{odpychaj\k{a}co}:
    wysokie $|\nabla\Phi|$ przy umiarkowanym $\Phi$ podnosi
    efektywny potencja\l{}.

  \item \textbf{Cz\l{}on kubiczny} $\beta\Phi^2/\PhiZero^2$:
    Najni\.zszego rz\k{e}du samointerakcja pola. Przy umiarkowanych
    $\Phi$ daje \emph{dodatni} wk\l{}ad do efektywnego
    potencja\l{}u --- odpowiada za bari er\k{e} odpychania (re\.zim~II).
    Jest analogiczny do cz\l{}onu $\lambda\phi^3$ w~teoriach pola,
    ale tutaj nie jest wstawiany r\k{e}cznie --- wynika z~faktu,
    \.ze wygenerowana przestrze\'n oddzia\l{}uje sama z~sob\k{a}.

  \item \textbf{Cz\l{}on kwartyczny} $-\gamma\Phi^3/\PhiZero^3$:
    Wy\.zszego rz\k{e}du korekcja, kt\'ora przy bardzo du\.zych
    $\Phi$ \emph{dominuje} nad cz\l{}onem kubicznym i~\emph{zmienia
    znak} efektywnej si\l{}y --- tworzy studni\k{e} (re\.zim~III).
    Gwarantuje stabilno\'s\'c: bez tego cz\l{}onu $\Phi \to \infty$
    by\l{}oby niestabilne.
\end{enumerate}

\subsubsection{Pe\l{}ne r\'ownanie pola}

\begin{remark}[Konwencja znaku \'zr\'od\l{}a]\label{rem:sign-convention}
Stosujemy konwencj\k{e}, w~kt\'orej \'zr\'od\l{}o masowe wyst\k{e}puje
ze znakiem~$-q$ po prawej stronie r\'ownania pola. Przy standardowej
funkcji Greena laplasjanu $\nabla^2(-1/(4\pi r)) = \delta^3(\mathbf{r})$
daje to $\delta\Phi = +q\PhiZero M/(4\pi r) > 0$ --- materia
\emph{generuje} nadwy\.zk\k{e} przestrzeni ($\delta\Phi > 0$,
$c < c_0$), zgodnie z~duchem aksjomatu~A5. Znak ujemny \'zr\'od\l{}a
jest bezpo\'srednio potwierdzony numerycznie
(profil radialny, skrypt \texttt{radial\_profile.py}).
\end{remark}

\begin{theorem}[R\'ownanie pola TGP]\label{thm:field-eq}
Pe\l{}ne r\'ownanie pola $\Phi$ w~obecno\'sci \'zr\'ode\l{} $\rho$ ma posta\'c:
\begin{equation}\label{eq:full-field}
  \boxed{
    \frac{1}{\PhiZero}\nabla^2\Phi
    \;+\;
    \frac{\alpha}{\PhiZero}\,\frac{(\nabla\Phi)^2}{\Phi}
    \;+\;
    \beta\,\frac{\Phi^2}{\PhiZero^2}
    \;-\;
    \gamma\,\frac{\Phi^3}{\PhiZero^3}
    \;=\;
    -q\,\rho,
  }
\end{equation}
lub r\'ownowa\.znie, mno\.z\k{a}c przez $\PhiZero$:
\begin{equation}\label{eq:full-field-alt}
    \nabla^2\Phi
    \;+\;
    \alpha\,\frac{(\nabla\Phi)^2}{\Phi}
    \;+\;
    \beta\,\frac{\Phi^2}{\PhiZero}
    \;-\;
    \gamma\,\frac{\Phi^3}{\PhiZero^2}
    \;=\;
    -q\,\PhiZero\,\rho.
  \end{equation}
\end{theorem}

\begin{remark}[Hierarchia cz\l{}on\'ow]
W~granicy s\l{}abego pola $\Phi = \PhiZero + \delta\Phi$,
$|\delta\Phi| \ll \PhiZero$:
\begin{equation}
    \nabla^2(\delta\Phi) + \mathcal{O}(\delta\Phi^2/\PhiZero)
    = -q\,\PhiZero\,\delta\rho,
\end{equation}
co jest r\'ownaniem Poissona --- fundamentem granicy newtonowskiej.
Cz\l{}ony nieliniowe w\l{}\k{a}czaj\k{a} si\k{e} stopniowo, gdy
$\delta\Phi/\PhiZero$ staje si\k{e} znacz\k{a}ce.
\end{remark}

\begin{proposition}[Warunek t\l{}a pr\'o\.zniowego]\label{prop:vacuum-condition}\label{ssec:warunek-prozniowy}
Warto\'s\'c referencyjna $\PhiZero$ jest jednorodnym
rozwi\k{a}zaniem r\'ownania~\eqref{eq:full-field-alt}
\textbf{wtedy i~tylko wtedy}, gdy:
\begin{equation}\label{eq:vacuum-cond}
  \beta - \gamma = -q\,\bar{\rho},
\end{equation}
gdzie $\bar{\rho}$ jest \'sredni\k{a} kosmologiczn\k{a}
g\k{e}sto\'sci\k{a}. Dla idealnej pr\'o\.zni ($\bar{\rho} = 0$):
$\beta = \gamma$.
\end{proposition}

\begin{proof}
Dla $\Phi = \PhiZero$ wsz\k{e}dzie: $\nabla\Phi = 0$,
$\nabla^2\Phi = 0$, wi\k{e}c r\'ownanie~\eqref{eq:full-field-alt}
redukuje si\k{e} do $\beta\PhiZero - \gamma\PhiZero = -q\PhiZero\bar\rho$,
co daje~\eqref{eq:vacuum-cond}.
\end{proof}

\begin{proposition}[Stabilno\'s\'c t\l{}a pr\'o\.zniowego]\label{prop:vacuum-stability}
Przy warunku pr\'o\.zniowym $\beta = \gamma$, statyczne perturbacje
wok\'o\l{} $\PhiZero$ s\k{a} \textbf{stabilne}: posiadaj\k{a} profil
Yukawa o~masie $m_{\mathrm{sp}} = \sqrt{\gamma}$.
\end{proposition}

\begin{proof}
Linearyzacja r\'ownania~\eqref{eq:full-field-alt}
wok\'o\l{} $\Phi = \PhiZero$ (pomijaj\k{a}c cz\l{}on
$\alpha(\nabla\Phi)^2/\Phi = \mathcal{O}(\delta\Phi^2)$):
\begin{equation}
  \nabla^2(\delta\Phi) + (2\beta - 3\gamma)\,\delta\Phi = -q\,\PhiZero\,\delta\rho.
\end{equation}
Propagator: $G(k) = 1/(k^2 + m_{\mathrm{sp}}^2)$,
$m_{\mathrm{sp}}^2 = -(2\beta - 3\gamma) = 3\gamma - 2\beta$.
Dla $\beta = \gamma$: $m_{\mathrm{sp}}^2 = \gamma > 0$,
co daje rozwi\k{a}zanie zanikaj\k{a}ce $\delta\Phi \propto e^{-m_{\mathrm{sp}} r}/r$.
Warunek $m_{\mathrm{sp}}^2 > 0$ wymaga $\beta < 3\gamma/2$ ---
spe\l{}niony dla $\beta = \gamma$.
\end{proof}

% -----------------------------------------------------------
% F4: Narracyjna sekcja W(1)≠0 — dodana w sesji v34
% -----------------------------------------------------------

\subsection{Pr\'o\.znia z~\'zr\'od\l{}em: \'la\'ncuch N0 $\to$ $W(1)\neq 0$
  $\to$ $\Lambda_{\rm eff}$ $\to$ kosmologia}\label{ssec:vacuum-source-chain}

\noindent
Jednym z~najbardziej donios\l{}ych wniok\'ow TGP jest \textbf{brak absolutnej
pr\'o\.zni} --- punktu, w~kt\'orym przestrze\'n (substrat $\Phi$) znika\l{}aby
dok\l{}adnie.
Wniosek ten nie jest \emph{postulatem}, lecz \emph{pochodn\k{a}} aksjomatyki N0.
Poni\.zej przedstawiamy pe\l{}ny \'la\'ncuch implikacji.

\paragraph{Krok~1: Aksjomat N0-5 $\Rightarrow$ $\beta = \gamma$.}
N0-5 stanowi, \.ze pr\'o\.zniowy punkt siod\l{}owy potencja\l{}u~$U(\psi)$
le\.zy w~$\psi = 1$, tzn.\ $U'(1) = 0$.
Jawnie: $U'(\psi) = \beta\psi - \gamma\psi^2$, wi\k{e}c $U'(1) = \beta - \gamma = 0$,
skoro $\beta = \gamma$. Parametry substratu nie s\k{a} wolne --- s\k{a}
zwi\k{a}zane przez warunek pr\'o\.zniowy.

\paragraph{Krok~2: $\beta = \gamma$ $\Rightarrow$ $W(1) = \gamma/3 \neq 0$.}
Funkcja \'zr\'od\l{}a $W(\psi) = \frac{4U(\psi)}{\psi} + U'(\psi)$
(stw.~\ref{prop:W-from-U}) przy $\psi = 1$ i~$\beta = \gamma$ wynosi:
\begin{equation}
  W(1) = 4U(1) + U'(1)
       = 4\!\left(\frac{\beta}{3} - \frac{\gamma}{4}\right) + 0
       = \frac{4\gamma}{3} - \gamma = \frac{\gamma}{3}.
\end{equation}
Poniewa\.z $\gamma > 0$ (aksjomat N0-3), mamy $W(1) = \gamma/3 > 0$
--- pr\'o\.znia \emph{zawsze} posiada aktywne \'zr\'od\l{}o.

\paragraph{Krok~3: $W(1)\neq 0$ $\Rightarrow$ brak absolutnej pr\'o\.zni ($\mathcal{N}_0$).}
R\'ownanie pola~\eqref{eq:field-full-source} przy $\rho_{\rm mat} = 0$ przyjmuje posta\'c:
\[
  \mathcal{D}[\Phi] = -\PhiZero\,W(\psi).
\]
Gdyby $\Phi \equiv 0$, to $\psi = \Phi/\PhiZero = 0$, $W(0) = 0$ i~powy\.zsze
r\'ownanie by\l{}oby spe\l{}nione trywialnie.
Jednak aksjomat N0-1 (przestrze\'n generowana przez materi\k{e}) i~warunek
przej\'scia fazowego GL (dod.~\ref{app:substrat}) gwarantuj\k{a},
\.ze ka\.zda sko\'nczona obj\k{e}to\'s\'c z~co~najmniej jednym w\k{e}z\l{}em
substratu ma $\Phi > 0$.
A zatem: \textbf{nie istnieje fizycznie osi\k{a}galna konfiguracja $\Phi = 0$}
--- nico\'s\'c $\mathcal{N}_0$ jest granic\k{a} formalna, nie stanem fizycznym.

\paragraph{Krok~4: $W(1)\neq 0$ $\Rightarrow$ naturalna ciemna energia.}
Gdy pole $\psi$ jest zamro\.zone przez tarcie Hubble'a na warto\'sci $\psi \approx 1$
(,,zamro\.zone t\l{}o''), efektywna g\k{e}sto\'s\'c energii pr\'o\.zni wynosi:
\begin{equation}\label{eq:Lambda-from-W1}
  \rho_\Lambda \;=\; \PhiZero\,W(1)\,/\,q
  \;=\; \frac{\gamma\,\PhiZero}{3q}
  \;\Longrightarrow\;
  \Lambda_{\rm eff} \;\approx\; \frac{\gamma}{12}
\end{equation}
(krok skalowania przez $8\pi G_{\rm eff}$; szczeg\'o\l{}y:
stw.~\ref{prop:Lambda-eff}).
Warto\'{sc} $\Lambda_{\rm eff}$ nie wymaga \emph{\.zadnego} fine-tuningu ---
jest zdeterminowana przez $\gamma$ i~$\PhiZero$, obydwa wyznaczone przez
substrat TGP niezale\.znie od obserwacji kosmologicznych.

\paragraph{Krok~5: Kosmologiczne konsekwencje.}
\begin{itemize}[nosep]
  \item \textbf{Ciemna energia jako pole substratowe, nie sta\l{}a}:
    $\rho_\Lambda(t) = \PhiZero W(\psi_{\rm bg}(t))/q$
    jest zale\.zna od czasu przez $\psi_{\rm bg}(t)$.
    Prognoza: \emph{efektywne $w(z) = -1 + \delta w(z)$} z~$\delta w \ll 1$
    dla $z \lesssim 2$, korekcja mierzalna przez DESI lub Euclid.
  \item \textbf{Brak osobliwo\'sci Big~Bang}:
    $\Phi_{\rm min} > 0$ przy $t \to 0$ (substrat nigdy nieznika)
    $\Rightarrow$ Wszech\'swiat ,,startuje'' w~skondensowanym substracie.
  \item \textbf{Stabilno\'s\'c pr\'o\.zni}: perturbacje $\delta\psi$ wok\'o\l{}
    $\psi = 1$ maj\k{a} mas\k{e} $m_{\rm sp} = \sqrt{\gamma}$ (stw.~\ref{prop:vacuum-stability})
    $\Rightarrow$ pr\'o\.znia jest niestabilna \emph{tylko} dynamicznie (efektywna bariera
    przy $\Phi_{\rm min}$; stw.~\ref{thm:vacuum-source}, pkt~1).
\end{itemize}

\medskip\noindent
\textbf{Podsumowanie \'la\'ncu\-cha:}
\[
  \underbrace{\beta = \gamma}_{\text{N0-5}}
  \;\xrightarrow{\phantom{X}}\;
  W(1) = \frac{\gamma}{3} \neq 0
  \;\xrightarrow{\phantom{X}}\;
  \begin{cases}
    \Phi_{\rm min} > 0 & \text{(brak nico\'sci)}\\
    \Lambda_{\rm eff} \approx \gamma/12 & \text{(ciemna energia)}\\
    \rho_\Lambda(t) \text{ zmienne} & \text{(program: DESI)}
  \end{cases}
\]

\begin{theorem}[Pr\'o\.znia z~\'zr\'od\l{}em --- aksjomat strukturalny TGP]%
\label{thm:vacuum-source}
Dla ka\.zdego punktu przestrzeni $x$, nawet w~obszarze
bezmateryjnym ($\rho_{\rm mat}(x) = 0$), pe\l{}ne r\'ownanie pola TGP
\begin{equation}\label{eq:field-full-source}
  \mathcal{D}[\Phi](x)
  = -q\,\rho_{\rm mat}(x) - \PhiZero\,W(\psi(x))
\end{equation}
ma \textbf{niezerowe pr\'o\.zniowe \'zr\'od\l{}o}
$\rho_{\rm vac}(x) \equiv -\PhiZero W(\psi(x))/q$, gdzie
\begin{equation}\label{eq:W1-nonzero}
  W(1) = \frac{7\beta}{3} - 2\gamma
  \;\xrightarrow{\beta=\gamma}\;
  \frac{\gamma}{3} \;\neq\; 0.
\end{equation}
Konsekwencje:
\begin{enumerate}[nosep]
  \item \textbf{Brak trywialnego rozwi\k{a}zania}:
    nie istnieje $\Phi \equiv 0$ --- substrat jest zawsze aktywny.
  \item \textbf{Naturalna sta\l{}a kosmologiczna}:
    $\Lambda_{\rm eff} = q\,\PhiZero\,W(1)/(8\pi G_{\rm eff})
    \approx \gamma/12$ (stw.~\ref{prop:Lambda-eff}).
  \item \textbf{Zamro\.zone t\l{}o}: w~erze materii
    Hubble zamra\.za $\psi \approx 1$; wtedy
    $\rho_{\rm vac} = \mathrm{const}$ --- r\'ownowa\.znik
    $\Lambda_{\rm eff} = \mathrm{const}$ dok\l{}adnie.
  \item \textbf{Ka\.zda pr\'o\.znia ma gesto\'s\'c energii}:
    nawet w~przestrzeni mi\k{e}dzygalaktycznej substrat
    wibruje na poziomie $\psi_{\rm bg}(t)$;
    lokalne zaburzenia $\delta\psi$ nak\l{}adaj\k{a}
    si\k{e} na to t\l{}o.
\end{enumerate}
\end{theorem}

\begin{proof}
R\'ownanie~\eqref{eq:field-full-source} wynika z~r\'ownania
kosmologicznego~\eqref{eq:Phi-cosmo-exact} przez uw\-zgl\k{e}dnienie
obu \'zr\'ode\l{}: materia ($-q\rho_{\rm mat}$) i~potencja\l{} efektywny
($-\PhiZero W(\psi)$, por.~stw.~\ref{prop:W-from-U}).
Niezerowo\'s\'c $W(1) = \gamma/3$: podstawiamy $\beta = \gamma$
do~\eqref{eq:W-expansion}: $W(1) = 7\beta/3 - 2\gamma = \gamma/3 > 0$.
Punkt~(1): gdyby $\Phi \equiv 0$, to $\psi = 0$ i~$W(0) = 0$,
ale r\'ownanie~\eqref{eq:field-full-source} wymaga\l{}oby
$\mathcal{D}[0] = 0$, co jest niesprzeczne --- lecz
warunek zerowej sumy (aks.~\ref{ax:zero}) i~ci\k{a}g\l{}o\'s\'c
przej\'scia fazowego substratu (dod.~\ref{app:substrat}) gwarantuj\k{a}
$\Phi \geq \Phi_{\rm min} > 0$ w~ka\.zdej sko\'nczonej
obj\k{e}to\'sci.
\end{proof}

\begin{remark}[Wnioski strukturalne z~twierdzenia o~pr\'o\.zni z~\'zr\'od\l{}em]%
\label{rem:vacuum-source-corollaries}
Z~twierdzenia~\ref{thm:vacuum-source} wynikaj\k{a} trzy otwarte pytania
z~konsekwencjami fizycznymi:
\begin{enumerate}[nosep]
  \item \textbf{Horyzont czarnej dziury}:
    Gdy $\Phi \to \Phi_{\rm min} > 0$ blisko horyzontu, substrat jest
    \emph{rzadki}, ale nigdy nieobecny.
    Implikuje to \textbf{brak klasycznej osobliwo\'sci}
    (substrat nie mo\.ze znikn\k{a}\'c), lecz mo\.ze istnie\'c
    \emph{efektywna bariera substratu} przy $\Phi = \Phi_{\rm min}$.
    Status: \textbf{Program} --- wymaga bilansu pola $\Phi$ w~metryce Schwarzschilda.

  \item \textbf{Pocz\k{a}tek Wszech\'swiata (Big Bang)}:
    Czy $\Phi \to 0$ przy $t \to 0$?
    Aksjomat $K(0)=0$ i~warunek zerowej sumy (aks.~\ref{ax:zero})
    gwarantuj\k{a} $\Phi \geq \Phi_{\rm min}(N_{\rm sub}) > 0$
    dla ka\.zdej sko\'nczonej obj\k{e}to\'sci z~przynajmniej jednym
    w\k{e}z\l{}em substratu.
    Konsekwencja: \textbf{brak osobliwo\'sci kosmologicznej} w~TGP~---
    Wszech\'swiat ``startuje'' przy $\Phi = \Phi_{\rm min}$, nie~$\Phi = 0$.
    Status: \textbf{Hipoteza} --- wymaga przej\'scia fazowego GL w~granicy
    wysokiej temperatury $T_\Gamma \to \infty$.

  \item \textbf{Fluktuacje pr\'o\.zniowe}:
    Ka\.zdy ,,pusty'' punkt przestrzeni ma energi\k{e} g\k{e}sto\'sci
    $\rho_{\rm vac} = -\PhiZero W(\psi_{\rm bg}(t))/q$, modulowan\k{a}
    przez kosmologiczn\k{a} ewolucj\k{e} $\psi_{\rm bg}(t)$.
    Lokalne zaburzenia $\delta\psi(\mathbf{x},t)$ nak\l{}adaj\k{a} si\k{e}
    na to t\l{}o i~\textbf{mog\k{a} by\'c obserwowane przez pomiar
    anizotropii $G_{\rm eff}(z)$} (por.~tw.~\ref{thm:Geff-k}
    oraz $\sigma_8$-napi\k{e}cie).
    Status: \textbf{Propozycja} --- konieczne obliczenie widma mocy
    $P_{\delta\psi}(k)$ z~r\'ownania MS-TGP (\S\ref{ssec:perturb}).
\end{enumerate}
\end{remark}

\begin{proposition}[Nieliniowa stabilno\'s\'c pr\'o\.zni ---
  argument energetyczny]\label{prop:nonlinear-stability}
Stabilno\'s\'c liniowa (stw.~\ref{prop:vacuum-stability})
nie wyklucza niestabilno\'sci przy du\.zych perturbacjach.
Poni\.zszy argument energetyczny pokazuje, \.ze pr\'o\.znia
TGP ($\Phi = \PhiZero$) jest stabilna r\'ownie\.z
\textbf{nieliniowo} w~nast\k{e}puj\k{a}cym sensie:
nie istnieje konfiguracja $\Phi(\mathbf{x})$ spe\l{}niaj\k{a}ca
warunek zerowej sumy (aksjomat~\ref{ax:zero}, ZS2),
kt\'ora mia\l{}aby energi\k{e} ni\.zsz\k{a} od $E[\PhiZero]$.
\end{proposition}

\begin{proof}
Funkcjona\l{} energii~\eqref{eq:energy-corrected} w~zmiennej
$\chi = \Phi/\PhiZero$ (przy $\rho = 0$) zapisujemy jako:
\begin{equation}\label{eq:E-chi}
  E[\chi] = \PhiZero^2\!\int d^3x\;\biggl[
    \frac{1}{2}\,f(\chi)\,(\nabla\chi)^2
    + \frac{\gamma}{3}\,\chi^3 - \frac{\gamma}{4}\,\chi^4
  \biggr],
\end{equation}
gdzie u\.zyli\'smy $\beta = \gamma$ i~$f(\chi) = 1 + 2\alpha\ln\chi > 0$
dla $\chi$ bliskiego jedno\'sci.
Rozwijamy wok\'o\l{} $\chi = 1 + \epsilon(\mathbf{x})$:
\begin{equation}\label{eq:E-expansion}
  E[\chi] - E[1] = \PhiZero^2\!\int d^3x\;\biggl[
    \frac{1}{2}\,(\nabla\epsilon)^2
    + \frac{\gamma}{2}\,\epsilon^2
    + R_3(\epsilon, \nabla\epsilon)
  \biggr],
\end{equation}
gdzie $R_3$ zawiera cz\l{}ony trzeciego i~wy\.zszego rz\k{e}du.
Dwa pierwsze cz\l{}ony s\k{a} \textbf{dodatnio okre\'slone}
(energia kinetyczna $+$ masa $m_{\mathrm{sp}}^2 = \gamma > 0$).

\smallskip\noindent
\emph{Ograniczenie reszty.}\;
Warunek zerowej sumy ZS2
($\int(\Phi - \PhiZero)\sqrt{h}\,d^3x = 0$,
aksjomat~\ref{ax:zero}) oznacza $\int\epsilon\,d^3x = 0$.
Z~nier\'owno\'sci Poincar\'ego:
\begin{equation}
  \int\epsilon^2\,d^3x \;\leq\; C_P\int(\nabla\epsilon)^2\,d^3x,
\end{equation}
wi\k{e}c cz\l{}on $(\nabla\epsilon)^2 + \gamma\epsilon^2$
kontroluje norm\k{e} $\|\epsilon\|_{H^1}$.
Dla $\|\epsilon\|_{H^1} < \delta_0$ (z~$\delta_0$
zale\.znym od $\gamma$ i~$\alpha$) reszta $|R_3|$ jest
zdominowana przez cz\l{}ony kwadratowe, wi\k{e}c
$E[\chi] - E[1] > 0$ dla ka\.zdej perturbacji
$\epsilon \not\equiv 0$ o~dostatecznie ma\l{}ej normie.

\smallskip\noindent
\emph{Du\.ze perturbacje.}\;
Dla perturbacji globalnych ($\epsilon = \mathrm{const} \neq 0$)
cz\l{}on gradientowy znika, ale warunek ZS2 wymusza
$\epsilon = 0$. Dla perturbacji zlokalizowanych o~du\.zej
amplitudzie energetyczna dominacja cz\l{}onu $-\gamma\chi^4/4$
jest kompensowana przez cz\l{}on kinetyczny $\frac{1}{2}f(\nabla\chi)^2$:
konfiguracja o~du\.zym $\epsilon$ w~ma\l{}ym regionie wymaga
stromych gradient\'ow, kt\'ore kosztuj\k{a} energi\k{e}.

Poni\.zsze twierdzenie formalizuje ten argument.
\end{proof}

\begin{theorem}[Nieliniowa stabilno\'s\'c orbitalna]\label{thm:nonlinear-orbital}
Niech $\psi = \Phi/\PhiZero$ spe\l{}nia kowariantne r\'ownanie
pola (tw.~\ref{thm:covariant-field}) na dziedzinie kompaktowej
$\Omega$ z~warunkiem zerowej sumy~\eqref{eq:zero-constraint}.
Definiujemy:
\begin{enumerate}[nosep]
  \item[(i)] \textbf{Hamiltonowsk\k{a}} energi\k{e} zachowan\k{a}:
    \begin{equation}\label{eq:H-conserved}
      \mathcal{H}[\psi,\dot\psi] = \int_\Omega
      \biggl[\frac{1}{2}\psi^5\frac{\dot\psi^2}{c_0^2}
      + \frac{1}{2}\psi^4(\nabla\psi)^2
      + \psi^4 U(\psi)\biggr]\,d^3x,
    \end{equation}
    gdzie $U(\psi) = \frac{\gamma}{3}\psi^3 - \frac{\gamma}{4}\psi^4$;
  \item[(ii)] \textbf{Basen stabilno\'sci}:
    $\mathcal{B} = \{\psi \in H^1(\Omega) : 0 < \psi(x) < \psi_{\max}\;
    \forall\,x\}$, gdzie
    \begin{equation}\label{eq:basin-boundary}
      \psi_{\max} = \tfrac{4}{3}
      \qquad\Bigl(\text{jedyne niezerowe rozwi\k{a}zanie }U(\psi) = 0\Bigr).
    \end{equation}
\end{enumerate}
W\'owczas:
\begin{enumerate}
  \item \textbf{Dodatnio\'s\'c w~basenie.} Dla $\psi \in \mathcal{B}$
    ka\.zdy cz\l{}on w~$\mathcal{H}$ jest nieujemny:
    $\psi^4 U(\psi) = \gamma\psi^7\bigl(\tfrac{1}{3} - \tfrac{\psi}{4}\bigr)
    \ge 0$, a~cz\l{}ony kinetyczne s\k{a} z~natury $\ge 0$.
  \item \textbf{Lokalna stabilno\'s\'c.} Istnieje $\delta > 0$ takie, \.ze
    dla danych pocz\k{a}tkowych $(\psi_0, \dot\psi_0)$
    z~$\|\psi_0 - 1\|_{H^1} + \|\dot\psi_0/c_0\|_{L^2} < \delta$
    rozwi\k{a}zanie spe\l{}nia:
    \begin{equation}
      \sup_{t}\bigl(\|\psi(t) - 1\|_{H^1}
      + \|\dot\psi(t)/c_0\|_{L^2}\bigr) \le C\,\delta,
    \end{equation}
    gdzie $C$ zale\.zy od $\gamma$ i~$|\Omega|$.
  \item \textbf{Zamkni\k{e}cie basenu.} Je\'sli dane pocz\k{a}tkowe
    spe\l{}niaj\k{a} $\psi_0 \in \mathcal{B}$, $\dot\psi_0$ dowolne,
    to $\psi(t,\cdot) \in \mathcal{B}$
    dop\'oki rozwi\k{a}zanie istnieje
    \textbf{i~$\mathcal{H}[\psi_0, \dot\psi_0]
    < \mathcal{H}_{\mathrm{crit}}$},
    gdzie $\mathcal{H}_{\mathrm{crit}}$
    jest energi\k{a} konfiguracji progowej ($\psi = \psi_{\max}$
    w~jednym punkcie).
\end{enumerate}
\end{theorem}

\begin{proof}
\emph{Punkt 1.}\;
$U(\psi) = \gamma\psi^3(1/3 - \psi/4) \ge 0
\;\Leftrightarrow\; \psi \le 4/3$. Pozosta\l{}e cz\l{}ony s\k{a}
sumami kwadrat\'ow (poniewa\.z $\psi > 0$).

\emph{Punkt 2.}\;
$\mathcal{H}$ jest zachowywana ($d\mathcal{H}/dt = 0$,
wynika z~kowariantnego r\'ownania pola). Dla $\psi = 1 + \epsilon$:
\begin{equation}\label{eq:H-quadratic}
  \mathcal{H} - \mathcal{H}[1,0] =
  \int_\Omega \biggl[\frac{1}{2}\frac{\dot\epsilon^2}{c_0^2}
  + \frac{1}{2}(\nabla\epsilon)^2\biggr]\,d^3x + R_3,
\end{equation}
gdzie $|R_3| \le C'\,(\|\nabla\epsilon\|_{L^2}^3
+ \|\dot\epsilon\|_{L^2}^3)$
(z~osadzenia Sobolewa $H^1 \hookrightarrow L^6$ w~3D
i~nier\'owno\'sci Poincar\'ego na $\{\int\epsilon = 0\}$:
$\|\epsilon\|_{L^2} \le C_P\|\nabla\epsilon\|_{L^2}$).

Cz\l{}ony kwadratowe kontroluj\k{a} $\|\nabla\epsilon\|^2
+ \|\dot\epsilon/c_0\|^2$.
Dla $\delta$ dostatecznie ma\l{}ego: $|R_3| \le \tfrac{1}{2}
(\|\nabla\epsilon\|^2 + \|\dot\epsilon/c_0\|^2)$,
wi\k{e}c $\mathcal{H} - \mathcal{H}_0 \ge \tfrac{1}{4}
(\|\nabla\epsilon\|^2 + \|\dot\epsilon/c_0\|^2) > 0$
dla $\epsilon \not\equiv 0$.
Z~zachowania $\mathcal{H}$:
$\|\nabla\epsilon(t)\|^2 + \|\dot\epsilon(t)/c_0\|^2
\le 4(\mathcal{H}(0) - \mathcal{H}_0)
\le 4\cdot 2\delta^2$,
a~Poincar\'e daje kontrol\k{e} $\|\epsilon\|_{L^2}$.

\emph{Punkt 3.}\;
Dowolna ci\k{a}g\l{}a \'scie\.zka $t \mapsto \psi(t)$
opuszczaj\k{a}ca basen musi przej\'s\'c przez
$\sup_x \psi(t,x) = 4/3$. W~konfiguracji progowej
$U(\psi) = 0$ w~punkcie maksimum, a~$U > 0$ wsz\k{e}dzie
indziej --- wi\k{e}c $\mathcal{H}$ na granicy basenu
osi\k{a}ga okre\'slon\k{a} warto\'s\'c
$\mathcal{H}_{\mathrm{crit}} > \mathcal{H}[1,0]$.
Zachowanie $\mathcal{H}$ i~warunek
$\mathcal{H}(0) < \mathcal{H}_{\mathrm{crit}}$
uniemo\.zliwiaj\k{a} osi\k{a}gni\k{e}cie granicy.
\end{proof}

\begin{remark}[Interpretacja fizyczna progu $\psi_{\max}$]%
\label{rem:basin-interpretation}
Pr\'og $\psi_{\max} = 4/3$ odpowiada $\Phi = \tfrac{4}{3}\PhiZero$
--- perturbacja wynosz\k{a}ca $33\%$ ponad warto\'s\'c t\l{}a.
Poza tym progiem:
\begin{itemize}[nosep]
  \item potencja\l{} $U(\psi) < 0$: pole ,,sp\l{}ywa''
    po potencjale zamiast oscylowa\'c,
  \item w~r\'ownaniu pola cz\l{}on $-\gamma\psi^4$ dominuje
    nad $+\beta\psi^3$: samointerferencja staje si\k{e}
    destabilizuj\k{a}ca,
  \item fizycznie: rozpoczyna si\k{e} formowanie
    czarnej dziury (kolaps do $\mathcal{S}_0$).
\end{itemize}
Uwaga: pr\'og ten dotyczy nadwy\.zki $\Phi$ \emph{ponad}
t\l{}o; w~obecno\'sci \'zr\'ode\l{} ($\Phi > \PhiZero$
w~pobli\.zu mas) fizyczny pr\'og odpowiada warunkowi,
by perturbacja \'rodowej $\delta\Phi/\PhiZero$ nie
przekracza\l{}a $1/3$.
\end{remark}

\subsubsection{Pe\l{}ne kowariantne r\'ownanie pola}

R\'ownanie~\eqref{eq:full-field-alt} jest statyczne (eliptyczne):
zawiera $\nabla^2\Phi$, ale nie pochodne czasowe. Pe\l{}ne
r\'ownanie dynamiczne uzyskujemy z~wariacji dzia\l{}ania~\eqref{eq:action}
uwzgl\k{e}dniaj\k{a}c pochodne po wszystkich koordynatach.

\begin{theorem}[Kowariantne r\'ownanie pola TGP]\label{thm:covariant-field}
Pe\l{}ne r\'ownanie dynamiczne pola $\Phi$ ma posta\'c:
\begin{equation}\label{eq:covariant-field}
  \boxed{
    \frac{1}{c(\Phi)^2}\,\frac{\partial^2\Phi}{\partial t^2}
    \;-\; \nabla^2\Phi
    \;+\;
    \frac{\alpha}{c(\Phi)^2}\,\frac{(\dot\Phi)^2}{\Phi}
    \;-\;
    \alpha\,\frac{(\nabla\Phi)^2}{\Phi}
    \;-\;
    \beta\,\frac{\Phi^2}{\PhiZero}
    \;+\;
    \gamma\,\frac{\Phi^3}{\PhiZero^2}
    \;=\;
    q\,\PhiZero\,\rho,
  }
\end{equation}
gdzie $\dot\Phi = \partial\Phi/\partial t$ i~$\alpha = 2$.
W~granicy statycznej ($\dot\Phi = 0$) odzyskujemy
r\'ownanie~\eqref{eq:full-field-alt} \textbf{po pomno\.zeniu
przez $-1$} (standardowa zmiana konwencji znakowej
mi\k{e}dzy d'alembertianem $\Box = c^{-2}\partial_t^2 - \nabla^2$
a~laplasjanem $\nabla^2$). W~granicy kosmologicznej
(jednorodne $\Phi(t)$, $\nabla\Phi = 0$) odzyskujemy
r\'ownanie~\eqref{eq:Phi-cosmo} po dodaniu t\l{}umienia
Hubble'owskiego $3H\dot\Phi$.
\end{theorem}

\begin{proof}[Dow\'od (wariacja dzia\l{}ania)]
Wychodzimy z~dzia\l{}ania~\eqref{eq:action} w~zmiennej
$\psi = \Phi/\PhiZero$, ale tym razem uwzgl\k{e}dniamy
\textbf{pe\l{}ne pochodne czasoprzestrzenne}. Metryka
efektywna~\eqref{eq:metric-exp} do pierwszego rz\k{e}du daje
$\sqrt{-g_{\mathrm{eff}}} = c_0\,\psi$%
\footnote{Poprawiony element obj\k{e}to\'sciowy ---
  wcze\'sniejsze wersje u\.zywa\l{}y $\psi^4$, co by\l{}o
  artefaktem przybli\.zonej metryki; pe\l{}ne wyprowadzenie
  w~sek.~\ref{ssec:unified-action}, prop.~\ref{prop:kappa-corrected}.}
oraz lokaln\k{a}
pr\k{e}dko\'s\'c propagacji $c(\Phi)^2 = c_0^2\,\PhiZero/\Phi
= c_0^2/\psi$. Element czasoprzestrzenny w~uk\l{}adzie
$(t, \mathbf{x})$ wynosi zatem:
\[
  d^4x\;\sqrt{-g_{\mathrm{eff}}}
  = c_0\,dt\,d^3x\;\psi.
\]
Cz\l{}on kinetyczny w~dzia\l{}aniu zawiera pe\l{}ny gradient
czasoprzestrzenny:
\begin{equation}\label{eq:kinetic-full}
  (\partial\psi)^2 \;\equiv\;
  -\frac{1}{c(\Phi)^2}\,\dot\psi^2 + (\nabla\psi)^2
  = -\frac{\psi}{c_0^2}\,\dot\psi^2 + (\nabla\psi)^2,
\end{equation}
gdzie $\dot\psi = \partial\psi/\partial t$ i~sygnatura
$(-,+,+,+)$ metryki $g_{\mu\nu}^{\mathrm{eff}}$ daje znak minus
przy cz\l{}onie czasowym. Dzia\l{}anie~\eqref{eq:action} przybiera
posta\'c:
\begin{equation}\label{eq:action-full-psi}
  S = \int dt\,d^3x\;\psi^4\;\biggl[
    \frac{1}{\kappa}\biggl(
      \frac{1}{2}\Bigl(
        -\frac{\psi}{c_0^2}\,\dot\psi^2 + (\nabla\psi)^2
      \Bigr)
      - \frac{\beta}{3}\psi^3 + \frac{\gamma}{4}\psi^4
    \biggr)
    - \frac{q}{\PhiZero}\,\rho\,\psi
  \biggr].
\end{equation}
(Pomijamy globalny czynnik $c_0$, kt\'ory nie wp\l{}ywa na
r\'ownania ruchu.)

\medskip\noindent\textbf{Krok 1: Wariacja cz\l{}onu kinetycznego.}
G\k{e}sto\'s\'c lagran\.zjanu kinetycznego (z~dok\l{}adno\'sci\k{a}
do $1/\kappa$):
\[
  \mathcal{L}_K = \tfrac{1}{2}\psi^4
  \bigl(-\psi\,c_0^{-2}\dot\psi^2 + (\nabla\psi)^2\bigr)
  = -\tfrac{1}{2}c_0^{-2}\,\psi^5\dot\psi^2
    + \tfrac{1}{2}\psi^4(\nabla\psi)^2.
\]
Wariacja $\delta\mathcal{L}_K / \delta\psi$ sk\l{}ada si\k{e}
z~cz\k{e}\'sci Eulera--Lagrange'a:
\begin{align}
  \frac{\partial\mathcal{L}_K}{\partial\psi}
  &= -\tfrac{5}{2}c_0^{-2}\,\psi^4\dot\psi^2
     + 2\psi^3(\nabla\psi)^2,
  \label{eq:var-direct}\\[4pt]
  \frac{\partial}{\partial t}
    \frac{\partial\mathcal{L}_K}{\partial\dot\psi}
  &= \frac{\partial}{\partial t}\bigl(-c_0^{-2}\psi^5\dot\psi\bigr)
  = -c_0^{-2}\bigl(5\psi^4\dot\psi^2 + \psi^5\ddot\psi\bigr),
  \label{eq:var-time}\\[4pt]
  \nabla\!\cdot\!\frac{\partial\mathcal{L}_K}{\partial(\nabla\psi)}
  &= \nabla\!\cdot\!\bigl(\psi^4\nabla\psi\bigr)
  = 4\psi^3(\nabla\psi)^2 + \psi^4\nabla^2\psi.
  \label{eq:var-space}
\end{align}
R\'ownanie Eulera--Lagrange'a
$\partial\mathcal{L}_K/\partial\psi
 - \partial_t(\partial\mathcal{L}_K/\partial\dot\psi)
 - \nabla\!\cdot\!(\partial\mathcal{L}_K/\partial(\nabla\psi))
= 0$
daje:
\begin{equation}\label{eq:EL-kinetic}
  c_0^{-2}\psi^5\ddot\psi
  + \tfrac{5}{2}c_0^{-2}\psi^4\dot\psi^2
  - \tfrac{5}{2}c_0^{-2}\psi^4\dot\psi^2
  - \psi^4\nabla^2\psi
  + 2\psi^3(\nabla\psi)^2
  - 4\psi^3(\nabla\psi)^2 = 0.
\end{equation}
Upraszczaj\k{a}c (cz\l{}ony $\pm\tfrac{5}{2}$ kasuj\k{a} si\k{e}):
\[
  c_0^{-2}\psi^5\ddot\psi
  - \psi^4\nabla^2\psi
  - 2\psi^3(\nabla\psi)^2 = 0.
\]

\medskip\noindent\textbf{Krok 2: Wariacja cz\l{}onu potencja\l{}owego.}
Potencja\l{} daje (analogicznie do~\eqref{eq:variation-explicit}):
\[
  \frac{\delta}{\delta\psi}\bigl(\psi^4 U\bigr)
  = \beta\psi^6 - \gamma\psi^7.
\]

\medskip\noindent\textbf{Krok 3: Wariacja cz\l{}onu materio\-wego.}
Sprz\k{e}\.zenie \'zr\'od\l{}owe $-(q/\PhiZero)\rho\,\psi^5$ daje:
\[
  \frac{\delta}{\delta\psi}\bigl(-\tfrac{q}{\PhiZero}\rho\,\psi^5\bigr)
  = -\frac{5q}{\PhiZero}\rho\,\psi^4.
\]
(Dok\l{}adny wsp\'o\l{}czynnik liczbowy zale\.zy od postaci
sprz\k{e}\.zenia; kluczowa jest struktura r\'ownania, nie
prefaktor \'zr\'od\l{}owy, kt\'ory fiksuje $q$.)

\medskip\noindent\textbf{Krok 4: Z\l{}o\.zenie i~dzielenie przez $\psi^3$.}
Zbieraj\k{a}c kroki 1--3 i~dziel\k{a}c przez $\psi^3$
(zak\l{}adamy $\psi > 0$, tj.\ sektor $\mathcal{S}_1$):
\[
  \frac{\psi^2}{c_0^2}\,\ddot\psi
  - \psi\,\nabla^2\psi
  - 2\,(\nabla\psi)^2
  - \beta\psi^3 + \gamma\psi^4
  = q'\,\rho\,\psi.
\]
Wracamy do $\Phi = \PhiZero\psi$, wi\k{e}c
$\dot\psi = \dot\Phi/\PhiZero$,
$\ddot\psi = \ddot\Phi/\PhiZero$,
$\nabla\psi = \nabla\Phi/\PhiZero$:
\[
  \frac{\Phi^2}{\PhiZero^2 c_0^2}\,\frac{\ddot\Phi}{\PhiZero}
  - \frac{\Phi}{\PhiZero}\,\frac{\nabla^2\Phi}{\PhiZero}
  - \frac{2\,(\nabla\Phi)^2}{\PhiZero^2}
  - \beta\frac{\Phi^3}{\PhiZero^3}
  + \gamma\frac{\Phi^4}{\PhiZero^4}
  = q'\,\rho\,\frac{\Phi}{\PhiZero}.
\]
Mno\.zymy przez $\PhiZero^2/\Phi$ i~korzystamy z~$c(\Phi)^2 = c_0^2\PhiZero/\Phi$:
\[
  \frac{1}{c(\Phi)^2}\,\ddot\Phi
  - \nabla^2\Phi
  - 2\,\frac{(\nabla\Phi)^2}{\Phi}
  - \beta\frac{\Phi^2}{\PhiZero}
  + \gamma\frac{\Phi^3}{\PhiZero^2}
  = q\,\PhiZero\,\rho.
\]

\medskip\noindent\textbf{Krok 5: Cz\l{}on czasowy $\dot\Phi^2/\Phi$.}
Powy\.zsze r\'ownanie ma struktur\k{e} drugiego rz\k{e}du, ale
nie zawiera jeszcze jawnego cz\l{}onu $\dot\Phi^2/\Phi$.
Wynika on z~faktu, \.ze $c(\Phi)$ \emph{samo zale\.zy od $\Phi$}.
Rozpisuj\k{a}c $\partial_t\bigl(c(\Phi)^{-2}\dot\Phi\bigr)$
w~r\'ownaniu falowym:
\[
  \partial_t\!\left(\frac{\dot\Phi}{c(\Phi)^2}\right)
  = \frac{\ddot\Phi}{c(\Phi)^2}
  + \frac{(\dot\Phi)^2}{c(\Phi)^2\,\Phi}.
\]
Jest to r\'ownie\.z widoczne bezpo\'srednio z~wariacji: pe\l{}na
pochodna Eulera--Lagrange'a po $t$ z~\eqref{eq:var-time} zawiera
cz\l{}on $5\psi^4\dot\psi^2$, kt\'ory po renormalizacji
daje dok\l{}adnie $+(\alpha/c^2)(\dot\Phi)^2/\Phi$
z~$\alpha = 2$, analogicznie do cz\l{}onu
$-2(\nabla\Phi)^2/\Phi$ w~sektorze przestrzennym.
Ostatecznie pe\l{}ne r\'ownanie Eulera--Lagrange'a przyjmuje
posta\'c:
\begin{equation}\label{eq:covariant-derived}
  \frac{1}{c(\Phi)^2}\,\ddot\Phi
  - \nabla^2\Phi
  + \frac{2}{c(\Phi)^2}\,\frac{(\dot\Phi)^2}{\Phi}
  - 2\,\frac{(\nabla\Phi)^2}{\Phi}
  - \beta\frac{\Phi^2}{\PhiZero}
  + \gamma\frac{\Phi^3}{\PhiZero^2}
  = q\,\PhiZero\,\rho,
\end{equation}
co jest dok\l{}adnie r\'ownaniem~\eqref{eq:covariant-field}
z~$\alpha = 2$.

\medskip\noindent\textbf{Podsumowanie.}
Wsp\'o\l{}czynnik $\alpha = 2$ przy obu cz\l{}onach
$(\dot\Phi)^2/\Phi$ i~$(\nabla\Phi)^2/\Phi$ ma \textbf{to samo
\'zr\'od\l{}o}: miar\k{e} $\psi^4$ w~dzia\l{}aniu. Cz\l{}on
czasowy pojawia si\k{e} ze znakiem $+$ a~przestrzenny ze znakiem
$-$ z~powodu lorentzowskiej sygnatury metryki efektywnej.
W~granicy statycznej ($\dot\Phi = 0$) r\'ownanie redukuje si\k{e}
do~\eqref{eq:full-field-alt} po pomno\.zeniu przez $-1$,
co potwierdza sp\'ojno\'s\'c znakow\k{a}.
\end{proof}

\begin{remark}[Autoreferencyjno\'s\'c r\'ownania pola]\label{rem:autoref}
W~r\'ownaniu~\eqref{eq:covariant-field} pr\k{e}dko\'s\'c
propagacji $c(\Phi) = c_0\sqrt{\PhiZero/\Phi}$ \textbf{zale\.zy
od samego pola} $\Phi$. Jest to \textbf{fundamentalna
autoreferencyjno\'s\'c} TGP: pole, kt\'ore konstytuuje przestrze\'n,
okre\'sla w\l{}asn\k{a} dynamik\k{e} propagacji. To nie jest
wada --- jest \emph{istot\k{a}} faktu, \.ze przestrze\'n nie jest
aren\k{a}, lecz polem.

Konsekwencje:
\begin{enumerate}[nosep]
  \item \textbf{Pr\k{e}dko\'s\'c propagacji zaburze\'n $\Phi$}
    jest \emph{lokalna}: w~g\k{e}stej przestrzeni (du\.ze $\Phi$)
    zaburzenia propaguj\k{a} wolniej. Pole samo spowalnia
    sw\'oj w\l{}asny transport.
  \item \textbf{Sto\.zki przyczynowe} s\k{a} wyznaczone
    przez $\Phi$: w~regionie czarnej dziury ($\Phi \to \infty$,
    $c \to 0$) sto\.zek zamyka si\k{e} --- \.zadne zaburzenie
    $\delta\Phi$ nie mo\.ze si\k{e} wydosta\'c.
  \item W~granicy $\Phi \to 0$ r\'ownanie staje si\k{e}
    \textbf{osobliwe} (cz\l{}on $(\dot\Phi)^2/\Phi$ dywerguje),
    co sygnalizuje granice opisu ci\k{a}g\l{}ego ---
    konieczno\'s\'c sformu\l{}owania dyskretnego na
    substracie $\hat{s}_i$.
\end{enumerate}
\end{remark}

\begin{proposition}[Hiperboliczno\'s\'c i~dobrze postawiony problem Cauchy'ego]\label{prop:hyperbolic}
R\'ownanie~\eqref{eq:covariant-field} jest \textbf{quasi-liniowe
hiperboliczne} dla $\Phi > 0$: cz\l{}on g\l{}\'owny
$\partial^2\Phi/c^2\partial t^2 - \nabla^2\Phi$ ma sygnatur\k{e}
lorentzowsk\k{a}, a~cz\l{}ony nieliniowe s\k{a} ni\.zszego rz\k{e}du.
Problem pocz\k{a}tkowy z~danymi $\Phi(t_0, \mathbf{x})$,
$\dot\Phi(t_0, \mathbf{x})$ na hypersurface $\Sigma_{t_0}$
jest zatem \textbf{dobrze postawiony} w~ka\.zdym regionie, gdzie
$\Phi > 0$ (sektor $\mathcal{S}_1$).

Natomiast w~otoczeniu $\Phi \to 0$ (granica sektora $\mathcal{S}_0$)
r\'ownanie degeneruje si\k{e} --- cz\l{}on $\alpha(\dot\Phi)^2/\Phi$
dywerguje. Jest to fizycznie sp\'ojne: na granicy nico\'sci
opis ci\k{a}g\l{}y traci sens, a~dynamika musi by\'c opisana
parametrem miarowym $\tau$ substratu
(definicja~\ref{def:relax}).
\end{proposition}

\begin{theorem}[Lokalne istnienie i~jednoznaczno\'s\'c]\label{thm:well-posedness}
Niech $s > 3/2 + 1 = 5/2$ i~niech dane pocz\k{a}tkowe spe\l{}niaj\k{a}
\begin{equation}
  \Phi_0 := \Phi|_{t=0} \in H^s(\mathbb{R}^3), \qquad
  \Phi_1 := \dot\Phi|_{t=0} \in H^{s-1}(\mathbb{R}^3),
\end{equation}
z~$\inf_{\mathbf{x}} \Phi_0(\mathbf{x}) > 0$ (sektor $\mathcal{S}_1$).
W\'owczas:
\begin{enumerate}[nosep]
  \item \textbf{Istnienie}: istnieje $T^* > 0$ i~jedyne rozwi\k{a}zanie
    $\Phi \in C([0,T^*]; H^s) \cap C^1([0,T^*]; H^{s-1})$
    r\'ownania~\eqref{eq:covariant-field}.
  \item \textbf{Kryterium kontynuacji (blow-up alternative)}:
    je\'sli $T^*$ jest maksymalnym czasem istnienia, to
    \begin{equation}
      T^* < \infty \;\Longrightarrow\;
      \limsup_{t \to T^*} \bigl\|\nabla\Phi(t)\bigr\|_{L^\infty}
      = \infty
      \quad\text{lub}\quad
      \inf_{\mathbf{x}} \Phi(t,\mathbf{x}) \to 0.
    \end{equation}
  \item \textbf{Estymata energetyczna}: funkcjona\l{}
    \begin{equation}\label{eq:wp-energy}
      \mathcal{E}_s[\Phi](t)
      := \sum_{|\alpha| \le s}
         \int \Bigl[\frac{(\partial^\alpha\dot\Phi)^2}{c^2\Phi}
         + (\partial^\alpha\nabla\Phi)^2\,\frac{1}{\Phi}\Bigr]\,d^3x
    \end{equation}
    spe\l{}nia $\mathcal{E}_s(t) \le \mathcal{E}_s(0)\,
    e^{C\int_0^t \|\nabla\Phi\|_{L^\infty}\,d\tau}$
    (Gronwall), co gwarantuje ci\k{a}g\l{}\k{a} zale\.zno\'s\'c
    od danych.
\end{enumerate}
\end{theorem}

\begin{proof}[Szkic dowodu]
R\'ownanie~\eqref{eq:covariant-field} przepisujemy w~postaci
uk\l{}adu pierwszego rz\k{e}du $\partial_t U = A^i(U)\,\partial_i U + F(U)$,
gdzie $U = (\Phi, \dot\Phi, \partial_i\Phi)$. Macierz $A^i$ ma
sygnatur\k{e} lorentzowsk\k{a} dla $\Phi > 0$
(stw.~\ref{prop:hyperbolic}), wi\k{e}c uk\l{}ad jest
\textbf{symetryzowalno-hiperboliczny} z~symetryzowaln\k{a}
macierz\k{a} $S = \mathrm{diag}(1/\Phi, 1/\Phi, \ldots)$.

\smallskip\noindent
\emph{Krok~1 (istnienie lokalne).}
Standardowa metoda Katona--Hughesa--Kato--Marsden:
aproksymacja Friedrichsa (regularyzacja $J_\epsilon$),
estymata energetyczna w~$H^s$ (pkt~3), zbie\.zno\'s\'c
w~$H^{s-1}$, ograniczono\'s\'c w~$H^s$ $\Rightarrow$
s\l{}aba zbie\.zno\'s\'c do rozwi\k{a}zania.

\smallskip\noindent
\emph{Krok~2 (jednoznaczno\'s\'c).}
R\'o\.znica dw\'och rozwi\k{a}za\'n spe\l{}nia liniowe
r\'ownanie z~wsp\'o\l{}czynnikami ograniczonymi w~$H^{s-1}$;
estymata energetyczna $L^2$ + Gronwall $\Rightarrow$
$\delta\Phi \equiv 0$.

\smallskip\noindent
\emph{Krok~3 (blow-up alternative).}
Je\'sli $\|\nabla\Phi\|_{L^\infty}$ pozostaje ograniczone
i~$\Phi > c_{\min} > 0$, estymata (pkt~3) daje ograniczenie
$\mathcal{E}_s(t) < \infty$ na dowolnym sko\'nczonym przedziale
--- sprzeczno\'s\'c z~maksymalno\'sci\k{a} $T^*$.
Fizycznie: blow-up mo\.ze wyst\k{a}pi\'c tylko przez (a)~formowanie
si\k{e} osobliwo\'sci gradientu (szok) lub (b)~osi\k{a}gni\k{e}cie
granicy $\Phi = 0$ (sektor $\mathcal{S}_0$). Oba
scenariusze s\k{a} fizycznie oczekiwane --- punkt~(b) odpowiada
formowaniu si\k{e} nico\'sci ($\Nzero$).
\end{proof}

\begin{remark}[Ograniczenia twierdzenia]\label{rem:wp-limitations}
Twierdzenie~\ref{thm:well-posedness} dotyczy sektora $\mathcal{S}_1$
($\Phi > 0$). Na granicy $\Phi \to 0$ r\'ownanie degeneruje
si\k{e} (por.\ stw.~\ref{prop:hyperbolic}). Rozszerzenie
na pe\l{}n\k{a} przestrze\'n konfiguracji
$\mathcal{S}_0 \cup \mathcal{S}_1$ wymaga sformu\l{}owania
w~zmiennych substratowych i~pozostaje problemem otwartym (O17).
\end{remark}

\begin{remark}[Zwi\k{a}zek z~zasad\k{a} zerowej sumy]\label{rem:zero-field}
Zasada zerowej sumy (aksjomat~\ref{ax:zero},
r\'ownanie~\ref{eq:zero-sum-precise}) stanowi \textbf{globalny
warunek brzegowy} na rozwi\k{a}zania r\'ownania
pola~\eqref{eq:full-field-alt}. W~j\k{e}zyku pola $\Phi$:
\begin{equation}\label{eq:zero-constraint}
    \int_\Sigma \bigl(\Phi(\mathbf{x}) - \PhiZero\bigr)\,
    \sqrt{h}\,d^3x = 0
    \qquad\text{(jednorodne t\l{}o)},
\end{equation}
co oznacza, \.ze \'srednia perturbacja $\delta\Phi$ po ca\l{}ej
hypersurface wynosi zero. Nadwy\.zka przestrzeni w~jednym
regionie ($\delta\Phi > 0$) jest kompensowana deficytem
w~innym ($\delta\Phi < 0$). Jest to sp\'ojne
z~interpretacj\k{a} Wszech\'swiata jako samowyzerowuj\k{a}cej
si\k{e} fluktuacji $\Nzero$.
\end{remark}

% =============================================================
\subsection{Profil pojedynczego \'zr\'od\l{}a}\label{ssec:profil-wyp}
% =============================================================

Dla izolowanej masy punktowej $M$ w~pocz\k{a}tku uk\l{}adu
($\rho = M\,\delta^3(\mathbf{r})$) szukamy rozwi\k{a}zania
sferycznie symetrycznego $\Phi = \Phi(r)$.
R\'ownanie~\eqref{eq:full-field-alt} w~symetrii sferycznej
przyjmuje posta\'c:

\begin{equation}\label{eq:spher}
  \Phi'' + \frac{2}{r}\,\Phi'
  + \alpha\,\frac{(\Phi')^2}{\Phi}
  + \beta\,\frac{\Phi^2}{\PhiZero}
  - \gamma\,\frac{\Phi^3}{\PhiZero^2}
  = -q\,\PhiZero\,M\,\delta^3(\mathbf{r}),
\end{equation}
gdzie prim oznacza $d/dr$.

\subsubsection{Re\.zim dalekiego pola ($r \to \infty$)}

Daleko od \'zr\'od\l{}a $\Phi \to \PhiZero$. Piszemy
$\Phi = \PhiZero + \delta\Phi(r)$, $|\delta\Phi| \ll \PhiZero$.
Linearyzacja r\'ownania~\eqref{eq:full-field-alt}
(pomijaj\k{a}c cz\l{}on gradientowy $\sim (\nabla\Phi)^2/\Phi
= \mathcal{O}(\delta\Phi^2)$):
\begin{equation}\label{eq:far-field-eq}
  \delta\Phi'' + \frac{2}{r}\,\delta\Phi'
  - \gamma\,\delta\Phi
  = -q\,\PhiZero\,M\,\delta^3(\mathbf{r}),
\end{equation}
gdzie masa przestrzenna $m_{\mathrm{sp}}^2 = 3\gamma - 2\beta = \gamma$
(przy $\beta = \gamma$, stw.~\ref{prop:vacuum-stability}).
Pe\l{}ne rozwi\k{a}zanie zanikaj\k{a}ce ma posta\'c \textbf{Yukawy}:
\begin{equation}\label{eq:far-field}
  \boxed{
    \delta\Phi(r)
    \;=\; \frac{q\,\PhiZero\,M}{4\pi\,r}\;
    e^{-\sqrt{\gamma}\,r}
    \qquad (r > 0).
  }
\end{equation}

\begin{remark}[Zasi\k{e}g grawitacji i~granica newtonowska]%
\label{rem:yukawa-newton}
Profil~\eqref{eq:far-field} jest profilu Yukawy z~zasi\k{e}giem
$\lambda_{\mathrm{sp}} = 1/\sqrt{\gamma}$.
Dla naturalnego $\gamma \sim H_0^2/c_0^2 \sim 10^{-52}\;\mathrm{m}^{-2}$
(sekcja~\ref{ssec:de-potencjal}) zasi\k{e}g wynosi:
\begin{equation}\label{eq:yukawa-range}
  \lambda_{\mathrm{sp}} = \frac{1}{\sqrt{\gamma}}
  \;\sim\; \frac{c_0}{H_0} \;\sim\; R_H
  \;\approx\; 4{,}4 \times 10^{26}\;\mathrm{m},
\end{equation}
czyli rz\k{e}du \textbf{promienia Hubble'a}. Na skalach
subkosmologicznych ($r \ll \lambda_{\mathrm{sp}}$)
cz\l{}on $e^{-\sqrt\gamma\,r} \approx 1$ i~odzyskujemy
profil newtonowski $\sim 1/r$; na skalach kosmologicznych
($r \gtrsim \lambda_{\mathrm{sp}}$) nast\k{e}puje obci\k{e}cie
--- grawitacja ma \textbf{naturalny zasi\k{e}g} rz\k{e}du
rozmiaru obserwowalnego Wszech\'swiata.

Fizycznie: masa $m_{\mathrm{sp}} = \sqrt\gamma$ jest
\textbf{mas\k{a} pola $\Phi$} --- kwanty przestrzenno\'sci
s\k{a} masywne, co daje \textbf{skon\'czon\k{a}} d\l{}ugo\'s\'c
Comptona pola generuj\k{a}cego przestrze\'n. Jest to sp\'ojne
z~relacj\k{a} dyspersji fal grawitacyjnych
(stw.~\ref{prop:GW-dispersion}).

Ten wynik r\'o\.zni TGP od OTW (gdzie grawitacja ma
niesko\'nczony zasi\k{e}g) i~jest \textbf{testowalny}:
modyfikacja prawa Newtona na skalach $\gtrsim 100\;\mathrm{Mpc}$
mo\.ze by\'c dost\k{e}pna pomiarami wzrostu struktur
(twierdzenie~\ref{thm:Geff-k}).
\end{remark}

\begin{proposition}[Pr\'og stabilizacji solitonu z~N0-6]%
\label{prop:eps-th}
Dla $\beta = \gamma$ (N0-5) i~$m_{\mathrm{sp}} = \sqrt{\gamma}$ (N0-6)
\textbf{pr\'og energetyczny} stabilizacji solitonu FDM wyra\.za si\k{e}
przez:
\begin{equation}\label{eq:eps-th}
  \boxed{\varepsilon_{\mathrm{th}} = \frac{m_{\mathrm{sp}}^2}{2}
  = \frac{\gamma}{2}}
\end{equation}
\emph{bez \\.zadnych dodatkowych parametr\'ow}.  Jest to minimalna
g\k{e}sto\'s\'c energii wymagana do utrzymania solitonu
$\Phi$ w~stanie zwi\k{a}zanym wobec rozproszenia.

Wyprowadzenie (skrypt \texttt{ex39}): r\'ownowaga energetyczna solitonu
$T_{\mathrm{kin}} = V_{\mathrm{eff}}$ przy minimum $V_{\mathrm{eff}}$
daje $\varepsilon_{\mathrm{th}} = m_{\mathrm{sp}}^2/2 = \gamma/2$.
\end{proposition}

\begin{proposition}[Z\l{}oty podzia\l{} jako minimum $V_{\mathrm{mod}}$]%
\label{prop:golden-ratio}
Zmodyfikowany potencja\l{} $V_{\mathrm{mod}}(\varphi)$ (przy $\beta=\gamma$,
$\varphi = \Phi/\PhiZero$):
\begin{equation}
  V_{\mathrm{mod}}(\varphi) = \tfrac{1}{2}m_{\mathrm{sp}}^2
  \left(\varphi^2 - \varphi^3 + \tfrac{1}{4}\varphi^4\right)
\end{equation}
ma minimum przy $\varphi^* = \varphi_{\mathrm{gold}}$, gdzie:
\begin{equation}\label{eq:golden-ratio-V}
  V'_{\mathrm{mod}}(\varphi^*)
  = m_{\mathrm{sp}}^2\,\varphi^*\bigl(1 + \varphi^* - (\varphi^*)^2\bigr)
  = 0
  \;\Longrightarrow\;
  (\varphi^*)^2 - \varphi^* - 1 = 0
  \;\Longrightarrow\;
  \boxed{\varphi^* = \phi_{\mathrm{gold}} = \tfrac{1+\sqrt{5}}{2}
  \approx 1.618.}
\end{equation}
Z\l{}oty podzia\l{} pojawia si\k{e} jako \textbf{naturalne minimum}
potencja\l{}u TGP przy progu stabilizacji solitonu --- bez narzucania
go z~zewn\k{a}trz.  Zweryfikowane numerycznie w~\texttt{ex39}.
\end{proposition}

\subsubsection{Re\.zim bliskiego pola ($r \to 0$)}

Blisko \'zr\'od\l{}a $\Phi \gg \PhiZero$. Dominuj\k{a} cz\l{}ony
nieliniowe. Piszemy $\Phi \sim A/r^p$ dla ma\l{}ych~$r$:

\begin{align}
  \Phi' &= -pA\,r^{-p-1}, \\
  \Phi'' &= p(p+1)\,A\,r^{-p-2}, \\
  \tfrac{2}{r}\Phi' &= -2p\,A\,r^{-p-2}, \\
  \alpha\,\frac{(\Phi')^2}{\Phi}
    &= \alpha\,p^2\,A\,r^{-p-2}.
\end{align}

Liniowe i~gradientowe cz\l{}ony daj\k{a} wyk\l{}adnik $r^{-p-2}$;
cz\l{}on kubiczny daje $r^{-2p}$, kwartyczny --- $r^{-3p}$.

\begin{itemize}[nosep]
  \item Je\'sli $3p > p+2$, czyli $p > 1$: cz\l{}on kwartyczny
    dominuje przy $r\to 0$. Wtedy: $\gamma\Phi^3/\PhiZero^2 \sim
    \gamma A^3/(r^{3p}\PhiZero^2)$ musi balansowa\'c si\k{e} z~wyrazami
    gradientowymi.
  \item Balans $\alpha(\Phi')^2/\Phi \sim \gamma\Phi^3/\PhiZero^2$
    daje: $\alpha p^2 A\,r^{-p-2} \sim \gamma A^3\,r^{-3p}/\PhiZero^2$,
    co wymaga $-p-2 = -3p$, czyli $p = 1$.
\end{itemize}

Dla $p = 1$ mamy $\Phi \sim A/r$ blisko \'zr\'od\l{}a, z~amplitud\k{a}~$A$
okre\'slon\k{a} przez balans:
\begin{equation}\label{eq:near-balance}
  \alpha\,A\,r^{-3}
  \;\sim\; \gamma\,\frac{A^3}{\PhiZero^2}\,r^{-3}
  \qquad\Longrightarrow\qquad
  A \;\sim\; \PhiZero\sqrt{\frac{\alpha}{\gamma}}.
\end{equation}

\begin{proposition}[Regularyzacja \'zr\'od\l{}a]\label{prop:reg}
W~pe\l{}nej teorii $\Phi$ nie dywerguje w~pocz\k{a}tku.
Masa \'zr\'od\l{}owa $M$ jest \emph{rozci\k{a}gni\k{e}ta} na skal\k{e}
$r_0 \sim \lP$ --- d\l{}ugo\'s\'c Plancka jest naturalnym
regulatorem. Dla $r > r_0$ profil jest dobrze okre\'slony;
dla $r < r_0$ teoria wymaga pe\l{}nego sformu\l{}owania
(prawdopodobnie dyskretnego).
\end{proposition}

\subsubsection{Re\.zim po\'sredni i~punkt przegi\k{e}cia}

Profil $\Phi(r)$ przechodzi od $\sim A/r$ (ma\l{}e~$r$) do
$\PhiZero + B/r$ (du\.ze~$r$) z~r\'o\.zn\k{a} amplitud\k{a}.
W~re\.zimie po\'srednim dominuje cz\l{}on kubiczny
$\beta\Phi^2/\PhiZero$, kt\'ory ,,podnosi'' $\Phi$ ponad
interpolacj\k{e} $1/r$. To tworzy \textbf{punkt przegi\k{e}cia}
w~$\Phi(r)$ --- miejsce, gdzie $\Phi''$ zmienia znak.

\begin{lemma}[Istnienie punktu przegi\k{e}cia]\label{lem:inflection}
Je\.zeli $\beta > 0$ i~$\gamma > 0$, rozwi\k{a}zanie
r\'ownania~\eqref{eq:spher} posiada co najmniej jeden punkt
$r_{\mathrm{infl}}$, w~kt\'orym \textbf{laplasjan sferyczny}
zmienia znak:
\begin{equation}\label{eq:inflection-condition}
  \nabla^2\Phi\big|_{r_{\mathrm{infl}}}
  = \Phi''(r_{\mathrm{infl}}) + \frac{2}{r_{\mathrm{infl}}}\,\Phi'(r_{\mathrm{infl}})
  = 0.
\end{equation}
\end{lemma}

\begin{proof}[Szkic dowodu]
Z~r\'ownania pola w~pr\'o\.zni ($\rho = 0$):
\begin{equation}
  \nabla^2\chi = -\frac{\alpha(\chi')^2}{\chi} - \beta\chi^2 + \gamma\chi^3
  = \chi^2(\gamma\chi - \beta) - \frac{2(\chi')^2}{\chi},
\end{equation}
gdzie $\chi = \Phi/\PhiZero$.
\emph{Blisko \'zr\'od\l{}a} ($\chi \gg 1$): dominuje czlon
$-\beta\chi^2 + \gamma\chi^3 = \chi^2(\gamma\chi - \beta)$ oraz
$-\alpha(\chi')^2/\chi$; dla $\chi \sim a_0/r$ oba s\k{a} por\'ownywalne
i~w~wiod\k{a}cym rz\k{e}dzie $\nabla^2\chi < 0$
(co weryfikacja numeryczna potwierdza).
\emph{Daleko od \'zr\'od\l{}a} ($\chi \to 1^+$, re\.zim Yukawa):
$\nabla^2(\delta\chi) = m^2\,\delta\chi > 0$ (gdzie $\delta\chi = \chi - 1$).
Przez ci\k{a}g\l{}o\'s\'c $\nabla^2\chi$ zmienia znak
co najmniej raz. Weryfikacja numeryczna (skrypt
\texttt{single\_source\_profile.py}, test~T5) daje
$r_{\mathrm{infl}} \approx 0.89$ dla $K = 1$, $m = 1$.
\end{proof}

Punkt przegi\k{e}cia $r_{\mathrm{infl}}$ jest fizycznie istotny:
gradient $|\Phi'|$ osi\k{a}ga w~tym rejonie lokalne
\textbf{maksimum} lub zmienia szybko\'s\'c spadku, co przek\l{}ada
si\k{e} na zmian\k{e} znaku si\l{}y efektywnej mi\k{e}dzy dwoma \'zr\'od\l{}ami.

\subsubsection{Redukcja $\Phi^3$ i~profil globalny}%
\label{ssec:profil-Phi3}

\begin{proposition}[Redukcja $\Phi^3$ --- \'sci\'sla dla $\alpha=2$]%
\label{prop:Phi3-reduction}
Niech $v \equiv (\Phi/\PhiZero)^3$. Dla $\alpha = 2$ i~$\beta = \gamma$
r\'ownanie~\eqref{eq:spher} w~symetrii sferycznej redukuje si\k{e}
dok\l{}adnie do:
\begin{equation}\label{eq:v-ode}
  \boxed{
    v'' + \frac{2}{r}\,v' + 3\beta\bigl(v^{4/3} - v^{5/3}\bigr)
    = -3q\,\rho\,v^{2/3}.
  }
\end{equation}
Ca\l{}y cz\l{}on gradientowy $\alpha(\Phi')^2/\Phi$ zostaje
\textbf{wch\l{}oni\k{e}ty}: po podstawieniu znika bez reszty.
\end{proposition}

\begin{proof}
Niech $\chi = \Phi/\PhiZero$, $v = \chi^3$.
Z~$\chi = v^{1/3}$:
\begin{align*}
  \chi' &= \tfrac{1}{3}\,v^{-2/3}\,v', \\
  \chi'' &= \tfrac{1}{3}\,v^{-2/3}\,v'' - \tfrac{2}{9}\,v^{-5/3}(v')^2.
\end{align*}
Cz\l{}on gradientowy:
$2(\chi')^2/\chi = 2\cdot\tfrac{1}{9}v^{-4/3}(v')^2/v^{1/3}
= \tfrac{2}{9}\,v^{-5/3}(v')^2$.
Suma:
\[
  \chi'' + \frac{2}{r}\chi' + 2\frac{(\chi')^2}{\chi}
  = \tfrac{1}{3}v^{-2/3}\!\left(v'' + \frac{2}{r}v'\right)
  + \underbrace{\left(-\tfrac{2}{9}+\tfrac{2}{9}\right)}_{=\,0}v^{-5/3}(v')^2.
\]
Mnożąc ca\l{}e r\'ownanie przez $3v^{2/3}$ (czyli $3\chi^2$)
otrzymujemy~\eqref{eq:v-ode}.
\end{proof}

\begin{proposition}[Szereg bliskiego pola --- dwa pierwsze wyrazy]%
\label{prop:near-field-series}
Dla $r \to 0^+$ (poza \'zr\'od\l{}em) r\'ownanie~\eqref{eq:v-ode}
posiada jednoparametrowe rozwi\k{a}zanie:
\begin{equation}\label{eq:chi-near}
  \chi(r) = \frac{\Phi(r)}{\PhiZero}
  = \frac{a_0}{r} - 1 + \mathcal{O}(r),
  \qquad
  a_0 = \sqrt{\frac{2}{\gamma}},
\end{equation}
gdzie wsp\'o\l{}czynniki s\k{a} wyznaczone \textbf{jednoznacznie}
przez r\'ownanie (bez wolnych parametr\'ow).
\end{proposition}

\begin{proof}
Podstawiamy $v = a_0^3/r^3 + 3a_0^2 b/r^2 + \ldots$
(rozwini\k{e}cie $v = \chi^3$, $\chi = a_0/r + b + \ldots$).

\medskip\noindent
\emph{Rz\k{a}d $r^{-5}$:}
$v''+2v'/r$ daje $6a_0^3/r^5$;
$3\beta(v^{4/3}-v^{5/3}) \approx -3\beta\cdot a_0^5/r^5$.
Warunek zerowania:
\[
  6a_0^3 - 3\beta\,a_0^5 = 0
  \;\;\Longrightarrow\;\;
  \beta\,a_0^2 = 2
  \;\;\Longrightarrow\;\;
  a_0 = \sqrt{2/\gamma}.
\]

\medskip\noindent
\emph{Rz\k{a}d $r^{-4}$:}
$v''+2v'/r$ daje $6a_0^2 b/r^4$;
$3\beta(v^{4/3}-v^{5/3})$ przy $v^{4/3}$ daje
$3\beta\cdot a_0^4/r^4$ (cz\l{}on $v^{5/3}$ wchodzi na rz\k{e}dzie $r^{-5}$).
Warunek:
\[
  6a_0^2 b + 3\beta\,a_0^4 = 0
  \;\;\Longrightarrow\;\;
  b = -\frac{\beta\,a_0^2}{2} = -\frac{2}{2} = -1.
\]
Oba wsp\'o\l{}czynniki s\k{a} \'sci\'sle okre\'slone przez r\'ownanie pola.
\end{proof}

\begin{remark}[Interpretacja wyrazu sta\l{}ego $b=-1$]%
\label{rem:b-minus-one}
Profil $\chi(r) = a_0/r - 1 + \mathcal{O}(r)$ ma zer\'o
w~$r_z = a_0 = \sqrt{2/\gamma}$.
Dla naturalnego $\gamma \sim H_0^2/c_0^2$:
$r_z \sim c_0\sqrt{2}/H_0 \sim \sqrt{2}\,R_H$.
Oznacza to, \.ze szereg bliskiego pola traci sens dla
$r \gtrsim r_z$ --- ale w~tym re\.zimie dominuje ju\.z
profil Yukawy~\eqref{eq:far-field},
czyli obu asymptotyk nie trzeba sklejać wewnątrz basenu
$\psi \in (0, 4/3)$.
\end{remark}

\begin{definition}[Masa krytyczna]\label{def:critical-mass}
Definiujemy \textbf{mas\k{e} krytyczn\k{a}} \'zr\'od\l{}a jako:
\begin{equation}\label{eq:Mcrit}
  M_* \;\equiv\; \frac{4\pi\,\PhiZero\,a_0}{q}
  = \frac{4\pi\,\PhiZero}{q}\sqrt{\frac{2}{\gamma}}
  = \frac{c_0^2}{G_0}\sqrt{\frac{2}{\gamma}}.
\end{equation}
Dla $\gamma \sim H_0^2/c_0^2$: $M_* \sim c_0^3\sqrt{2}/(G_0\,H_0)
\sim M_{\mathrm{Hubble}}/\sqrt{2}$
--- rz\k{e}du masy obserwowalnego Wszech\'swiata.
\end{definition}

\begin{proposition}[Dwa re\.zimy profilu]\label{prop:two-regimes-profile}
Niech $K \equiv q\,M/(4\pi\,\PhiZero)$
(bezwymiarowa si\l{}a \'zr\'od\l{}a).

\begin{enumerate}[nosep]
  \item \textbf{\'Sr\'od\l{}o s\l{}abe} ($M \ll M_*$, co obejmuje wszystkie znane
    obiekty astronomiczne):
    \begin{equation}\label{eq:chi-weak}
      \chi(r) \approx 1 + \frac{K\,e^{-\sqrt\gamma\,r}}{r}
      \qquad \text{globalnie}.
    \end{equation}
    Cz\l{}ony nieliniowe s\k{a} ma\l{}e w~ca\l{}ym zakresie $r > r_0$
    (poprawka rz\k{e}du $(K/a_0)^2 \ll 1$).

  \item \textbf{\'Sr\'od\l{}o silne} ($M \gtrsim M_*$):
    istnieje ska\l{}a $r_{\mathrm{nl}} \sim K/(a_0 m)$, poni\.zej kt\'orej
    profil \textbf{nasyca si\k{e}}:
    \begin{equation}\label{eq:chi-strong}
      \chi(r) \approx \frac{a_0}{r} - 1
      \qquad (r_0 < r \ll r_{\mathrm{nl}}).
    \end{equation}
\end{enumerate}
\end{proposition}

\begin{proposition}[Globalny profil kompozytowy]\label{prop:composite-profile}
Nast\k{e}puj\k{a}ca formu\l{}a stanowi globaln\k{a} aproksymacj\k{e}
profilu jednego \'zr\'od\l{}a, dok\l{}adn\k{a} w~obu limitach:
\begin{equation}\label{eq:chi-composite}
  \boxed{
    \chi_{\mathrm{glob}}(r)
    = \left[\frac{a_0^3}{r^3}
      + 1 + \frac{3K\,e^{-\sqrt\gamma\,r}}{r}\right]^{1/3},
  }
\end{equation}
gdzie $a_0 = \sqrt{2/\gamma}$, $K = qM/(4\pi\PhiZero)$.

\medskip\noindent\emph{Weryfikacja limit\'ow:}
\begin{itemize}[nosep]
  \item $r \to 0$:
    $\chi_{\mathrm{glob}} \to (a_0^3/r^3)^{1/3} = a_0/r$. \checkmark
  \item $r \to \infty$:
    $\chi_{\mathrm{glob}} \to (1 + 3K\,e^{-mr}/r)^{1/3}
    \approx 1 + K\,e^{-mr}/r$. \checkmark
\end{itemize}
Weryfikacja numeryczna: skrypt \texttt{single\_source\_profile.py}.
\end{proposition}

% =============================================================
\subsection{Potencja\l{} efektywny dla dw\'och \'zr\'ode\l{}
  --- wyprowadzenie trzech re\.zim\'ow}\label{ssec:dwa-zrodla}
% =============================================================

Rozwa\.zamy dwie masy $M_1$, $M_2$ w~odleg\l{}o\'sci~$d$.
Pole ca\l{}kowite:
\begin{equation}
  \Phi = \Phi_1 + \Phi_2 + \delta\Phi_{\mathrm{nl}},
\end{equation}
gdzie $\Phi_{1,2}$ s\k{a} rozwi\k{a}zaniami pojedynczych \'zr\'ode\l{},
a~$\delta\Phi_{\mathrm{nl}}$ jest korekcj\k{a} nieliniow\k{a}.

Energia oddzia\l{}ywania (r\'o\.znica energii pe\l{}nej i~sumy
energii izolowanych \'zr\'ode\l{}):

\begin{equation}\label{eq:Eint}
  E_{\mathrm{int}}(d) = E[\Phi_1+\Phi_2] - E[\Phi_1] - E[\Phi_2].
\end{equation}

Funkcjona\l{} energii zwi\k{a}zany z~r\'ownaniem~\eqref{eq:full-field-alt}
posiada \textbf{zmodyfikowany cz\l{}on kinetyczny}:
\begin{equation}\label{eq:energy-corrected}
  \boxed{
    E[\Phi] = \int d^3x\;\biggl[
      \frac{1}{2}f(\Phi)\,(\nabla\Phi)^2
      + \frac{\beta}{3}\,\frac{\Phi^3}{\PhiZero}
      - \frac{\gamma}{4}\,\frac{\Phi^4}{\PhiZero^2}
      + q\,\PhiZero\,\rho\,\Phi
    \biggr],
  }
\end{equation}
gdzie $f(\Phi) = 1 + 2\alpha\ln(\Phi/\PhiZero)$.

\begin{remark}[Pochodzenie cz\l{}onu kinetycznego $f(\Phi)$]\label{rem:kinetic-origin}
Cz\l{}on gradientowy $\alpha(\nabla\Phi)^2/\Phi$
w~r\'ownaniu pola nie pochodzi z~wariacji standardowego
potencja\l{}u --- jest \textbf{wygenerowany} przez
nieliniow\k{a} zale\.zno\'s\'c $f(\Phi)$ w~cz\l{}onie
kinetycznym. Fizycznie $f(\Phi)$ odzwierciedla fakt,
\.ze ,,przepuszczalno\'s\'c'' pola $\Phi$ dla swoich
w\l{}asnych zmian zale\.zy od jego lokalnej warto\'sci.

Weryfikacja: wariacja $\delta E/\delta\Phi = 0$ funkcjona\l{}u
\eqref{eq:energy-corrected} daje:
\begin{equation}\label{eq:energy-functional}
  -\nabla\!\cdot\!(f\nabla\Phi) + \tfrac{1}{2}f'(\nabla\Phi)^2
  + \beta\frac{\Phi^2}{\PhiZero}
  - \gamma\frac{\Phi^3}{\PhiZero^2}
  = -q\,\PhiZero\,\rho.
\end{equation}
Z~$f' = 2\alpha/\Phi$ i~rozwijaj\k{a}c
$\nabla\!\cdot\!(f\nabla\Phi) = f\nabla^2\Phi + f'(\nabla\Phi)^2$:
\begin{equation}
  -f\nabla^2\Phi - f'(\nabla\Phi)^2 + \tfrac{1}{2}f'(\nabla\Phi)^2
  + \ldots = -q\PhiZero\rho,
\end{equation}
co daje $-f\nabla^2\Phi - \frac{1}{2}f'(\nabla\Phi)^2 + \ldots
= -q\PhiZero\rho$. W~granicy $\Phi \approx \PhiZero$
($f \approx 1$, $f' \approx 2\alpha/\Phi$) redukuje si\k{e} to do
r\'ownania~\eqref{eq:full-field-alt} z~$\alpha_{\mathrm{eff}} = \alpha$.

W~granicy s\l{}abego pola ($\Phi \approx \PhiZero$):
$f \approx 1$, odzyskujemy standardowy cz\l{}on kinetyczny.
Dok\l{}adna posta\'c $f(\Phi)$ na ca\l{}ym zakresie
$\Phi$ wynika z~dzia\l{}ania~\eqref{eq:action}
z~sprz\k{e}\.zeniem kinetycznym $K(\psi) = \psi^4$,
co generuje $\alpha = 2$ (wniosek~\ref{cor:alpha-fixed}).
Poprawiony element obj\k{e}to\'sciowy wynosi
$\sqrt{-g_{\mathrm{eff}}} = c_0\,\psi$
(sek.~\ref{ssec:unified-action}); warto\'s\'c $\alpha=2$
wynika z~$K(\psi)$, nie z~miary.

\medskip\noindent\textbf{Zakres stosowalno\'sci $f(\Phi)$:}
Funkcja $f(\Phi) = 1 + 4\ln(\Phi/\PhiZero)$%
\footnote{W~sformu\l{}owaniu kanonicznym TGP: $K(g) = g^4$ (dok\l{}adne sprz\k{e}\.zenie konforemne). Forma logarytmiczna $1 + 4\ln g$ jest rozwini\k{e}ciem Taylora $g^4 = e^{4\ln g} \approx 1 + 4\ln g$ wok\'o\l{} $g = 1$.}
zeruje si\k{e}
dla $\Phi/\PhiZero = e^{-1/4} \approx 0{,}78$ i~zmienia znak
(,\,duch kinetyczny'') dla $\psi < 0{,}78$.
Jest to artefakt \textbf{rozwini\k{e}cia funkcjona\l{}u energii}
wok\'o\l{} $\Phi \approx \PhiZero$. W~dzia\l{}aniu fundamentalnym
\eqref{eq:action-full-psi} cz\l{}on kinetyczny $\propto \psi^5(\nabla\psi)^2$
jest \textbf{dodatnio okre\'slony} dla ca\l{}ego basenu
$\psi \in (0, 4/3)$ --- brak niestabilno\'sci duchowej.
Funkcjona\l{} $E_{\mathrm{int}}$ z~$f(\Phi)$ jest wiarygodny
jedynie w~regimencie $\Phi \gtrsim 0{,}8\,\PhiZero$
(tj.\ $|\delta\Phi/\PhiZero| \lesssim 0{,}2$).
\end{remark}

\begin{proposition}[Brak duch\'ow w~zmiennych fundamentalnych]%
\label{prop:ghost-free-fundamental}
Niech $\psi = \Phi/\PhiZero > 0$.
W~dzia\l{}aniu fundamentalnym~\eqref{eq:action-full-psi}
g\k{e}sto\'s\'c lagran\.zjanu kinetycznego wynosi:
\begin{equation}\label{eq:L-kin-psi}
  \mathcal{L}_{\mathrm{kin}} = \frac{1}{2\kappa}\,\psi^4\,
  \Bigl(-\frac{\psi}{c_0^2}\,\dot\psi^2 + (\nabla\psi)^2\Bigr).
\end{equation}
W\'owczas:
\begin{enumerate}[(i)]
  \item Wsp\'o\l{}czynnik przy $\dot\psi^2$ wynosi
    $-\psi^5/(2\kappa c_0^2)$ --- z~sygnatur\k{a}
    $(-,+,+,+)$ odpowiada \textbf{dodatniej} energii kinetycznej
    dla ka\.zdego $\psi > 0$.
  \item Wsp\'o\l{}czynnik przy $(\nabla\psi)^2$ wynosi
    $+\psi^4/(2\kappa) > 0$ --- brak niestabilno\'sci gradientowych.
  \item \.Zadne z~tych wyra\.ze\'n nie zmienia znaku
    w~ca\l{}ym basenie $\psi \in (0, 4/3)$.
\end{enumerate}
W~szczeg\'olno\'sci: \emph{potencjalna zmiana znaku
$f(\Phi) = 1 + 2\alpha\ln(\Phi/\PhiZero)$
przy $\psi \approx 0{,}78$ jest artefaktem
przybli\.zonego funkcjona\l{}u energii}
\eqref{eq:energy-corrected},
\textbf{nie wyst\k{e}puje} w~dzia\l{}aniu fundamentalnym.
\end{proposition}

\begin{proof}
Bezpo\'srednie odczytanie z~\eqref{eq:action-full-psi}:
cz\l{}on kinetyczny ma posta\'c
$\psi^4 \cdot \tfrac{1}{2\kappa}(\nabla_{\mathrm{eff}}\psi)^2$,
gdzie $\nabla_{\mathrm{eff}}$ u\.zywa metryki
$g_{\mu\nu}^{\mathrm{eff}}$ z~sygnatur\k{a} $(-,+,+,+)$.

Wsp\'o\l{}czynnik czasowy: $-\psi^5/(2\kappa c_0^2)$.
Hamiltonowska energia kinetyczna (po transformacji Legendre'a):
$T = +\psi^5\dot\psi^2/(2\kappa c_0^2) > 0$
dla $\psi > 0$.

Wsp\'o\l{}czynnik przestrzenny: $+\psi^4/(2\kappa) > 0$
jest \'sci\'sle dodatni.

Zwi\k{a}zek z~$f(\Phi)$: po przepisaniu na $\Phi = \PhiZero\psi$
i~rozdzieleniu $\psi^5 = \psi^4 \cdot \psi$, cz\l{}on
$\psi^4(\nabla\psi)^2$ mo\.zna rozwin\k{a}\'c wok\'o\l{} $\psi=1$:
$\psi^4 = 1 + 4\ln\psi + O((\ln\psi)^2)$%
\footnote{W~Formulacji~A (kanonicznej) sprz\k{e}\.zenie kinetyczne $K(g) = g^4$ jest dok\l{}adne i~dodatnio okre\'slone dla ka\.zdego $g > 0$; posta\'c logarytmiczna jest jedynie przybli\.zeniem wok\'o\l{} $g = 1$.}
(logarytmicznie),
sk\k{a}d $f(\Phi) = 1 + 2\alpha\ln(\Phi/\PhiZero)$
z~$\alpha = 2$
jest przybli\.zeniem $\psi^4$ w~regimencie $|\psi - 1| \ll 1$.
Przybli\.zenie to traci wiarygodno\'s\'c
dla $\psi \lesssim 0{,}78$ --- dok\l{}adnie tam,
gdzie $f$ zmienia znak. Pe\l{}na posta\'c $\psi^4$ jest
jednak zawsze dodatnia.
\end{proof}

\begin{remark}[Granica $\Phi \to 0$ i~substrat]%
\label{rem:Phi-zero-substrate}
Stwierdzenie~\ref{prop:ghost-free-fundamental} gwarantuje brak
duch\'ow dla ka\.zdego $\psi > 0$. Przypadek $\psi = 0$
($\Phi = 0$, tj.~$\mathrm{N}_0$) le\.zy \textbf{poza}
dziedzin\k{a} kontinualnego opisu:

\begin{enumerate}[(a)]
  \item \textbf{Osobliwo\'sci r\'ownania pola.}
    Kowariantne r\'ownanie (tw.~\ref{thm:covariant-field})
    zawiera cz\l{}ony $\sim 1/\Phi$ (r\'ownowa\.znie
    $1/\psi$), kt\'ore dywerguj\k{a} przy $\Phi \to 0$.
    Fizycznie: pr\k{e}dko\'s\'c $c(\Phi) = c_0\sqrt{\PhiZero/\Phi}
    \to \infty$ --- metryka degeneruje si\k{e}.
  \item \textbf{Regularyzacja substratowa.}
    Na skali dyskretnej substratu $\Gamma = (V, E)$ (aksjomat~A1)
    pole ma naturalny dolny obci\k{a}\.znik:
    $\Phi_{\min} \sim \varepsilon_{\mathrm{sub}}
    = a_{\mathrm{sub}}^2$,
    gdzie $a_{\mathrm{sub}}$ jest sta\l{}\k{a} sieciow\k{a}.
    W~przybli\.zeniu kontinualnym odpowiada to
    regularyzacji $\Phi \to \max(\Phi, \varepsilon_{\mathrm{sub}})$
    (por.\ skrypt \texttt{boundary\_S0\_S1.py}).
  \item \textbf{Przej\'scie fazowe $\mathcal{S}_1 \to \mathcal{S}_0$.}
    Granica $\Phi \to 0$ opisuje przej\'scie do sektora $\mathcal{S}_0$
    (Nico\'s\'c) --- nie jest osi\k{a}gana dynamicznie
    z~danych pocz\k{a}tkowych w~basenie stabilno\'sci
    (tw.~\ref{thm:nonlinear-orbital}, punkt~3),
    lecz stanowi granic\k{e} fazow\k{a} geometrii.
  \item \textbf{Podsumowanie dziedzin.}
    Opis kontinualny TGP obowi\k{a}zuje dla
    $\psi \in (\varepsilon_{\mathrm{sub}}/\PhiZero, \; 4/3)$.
    Poni\.zej $\varepsilon_{\mathrm{sub}}$ obowi\k{a}zuje
    dyskretna dynamika substratowa
    (dodatek~\ref{app:substrat}).
\end{enumerate}
\end{remark}

\subsubsection{Trzy wk\l{}ady do $E_{\mathrm{int}}(d)$}

Rozwijaj\k{a}c~\eqref{eq:Eint} i~sortuj\k{a}c wed\l{}ug dominacji
na r\'o\.znych skalach~$d$:

\begin{equation}\label{eq:Eint-decomp}
  E_{\mathrm{int}}(d) \;=\;
    \underbrace{E_{\mathrm{lin}}(d)}_{\text{liniowe sprz\k{e}\.zenie}}
  + \underbrace{E_{\beta}(d)}_{\text{kubiczny}}
  + \underbrace{E_{\gamma}(d)}_{\text{kwartyczny}}
  + E_{\alpha}(d).
\end{equation}

\paragraph{(i) Cz\l{}on liniowy: przyci\k{a}ganie na du\.zych skalach.}

\begin{equation}\label{eq:Elin}
  E_{\mathrm{lin}}(d)
  = -q\,\PhiZero \int \rho_2(\mathbf{x})\,\Phi_1(\mathbf{x})\,d^3x
  \simeq -\frac{(q\PhiZero)^2\,M_1 M_2}{4\pi\,d},
\end{equation}
co jest \textbf{ujemne} (przyci\k{a}gaj\k{a}ce) i~$\sim 1/d$.
Dominuje dla du\.zych~$d$ (grawitacja, re\.zim~I).

\begin{remark}[Sp\'ojno\'s\'c z~profilem Yukawy]\label{rem:Elin-yukawa}
Dok\l{}adny profil dalekiego pola~\eqref{eq:far-field}
zawiera czynnik $e^{-\sqrt\gamma\,r}$, co daje:
\[
  E_{\mathrm{lin}}^{\mathrm{(Yuk)}}(d)
  = -\frac{(q\PhiZero)^2\,M_1 M_2}{4\pi\,d}\;
    e^{-\sqrt\gamma\,d}.
\]
Jednak w~jednostkach bezwymiarowych $\hat\gamma = \gamma\,r_0^2$,
gdzie $r_0 \sim \lP \approx 1{,}6 \times 10^{-35}\;\mathrm{m}$.
Z~$\gamma \sim H_0^2/c_0^2 \sim 10^{-52}\;\mathrm{m}^{-2}$:
$\hat\gamma \sim 10^{-122}$, wi\k{e}c
$\lambda_{\mathrm{sp}}/r_0 \sim 10^{61}$.
Na \textbf{wszystkich} skalach dost\k{e}pnych obliczeniom
dwucia\l{}owym ($d \lesssim 10^{10}\,r_0$) korekcja Yukawy
wynosi $e^{-\sqrt{\hat\gamma}\,d} > 1 - 10^{-51}$ i~jest
\textbf{pomijalna}. Przybli\.zenie $1/d$
w~r\'ow.~\eqref{eq:Elin} jest zatem dok\l{}adne
z~precyzj\k{a} ponad 50 rz\k{e}d\'ow wielko\'sci.
\end{remark}

\paragraph{(ii) Cz\l{}on kubiczny: odpychanie na \'srednich skalach.}

\begin{equation}\label{eq:Ebeta}
  E_{\beta}(d)
  = \frac{2\beta}{\PhiZero}\int \Phi_1(\mathbf{x})\,\Phi_2(\mathbf{x})\,d^3x
  > 0.
\end{equation}

Ca\l{}ka $\int\Phi_1\Phi_2\,d^3x$ ro\'snie szybko, gdy $d$
maleje (nak\l{}adanie si\k{e} p\'ol). Dla $\Phi_{1,2}\sim 1/r$
ca\l{}ka dywerguje logarytmicznie, ale jest regularyzowana
przez nieliniow\k{a} struktur\k{e} profilu blisko \'zr\'ode\l{}.
Efektywnie:
\begin{equation}
  E_{\beta}(d) \;\sim\; \frac{2\beta\,(q\PhiZero)^2 M_1 M_2}
  {(4\pi)^2\,\PhiZero\,d}\,\ln\!\frac{d}{r_0}
  \qquad (d \gg r_0).
\end{equation}

Cz\l{}on jest \textbf{dodatni} (odpychaj\k{a}cy). Ro\'snie
\emph{wolniej} ni\.z $1/d$ (logarytmicznie), ale dominuje
nad $E_{\mathrm{lin}}$ dla $d < r_{\mathrm{rep}}$, poniewa\.z
ca\l{}ka $\int\Phi_1\Phi_2$ nasila si\k{e} nieliniowo.

\paragraph{(iii) Cz\l{}on kwartyczny: studnia na ma\l{}ych skalach.}

\begin{equation}\label{eq:Egamma}
  E_{\gamma}(d)
  = -\frac{3\gamma}{\PhiZero^2}\int \Phi_1(\mathbf{x})\,
    \Phi_2(\mathbf{x})\,\bigl[\Phi_1(\mathbf{x})+\Phi_2(\mathbf{x})\bigr]
    \,d^3x
  < 0.
\end{equation}

Dla ma\l{}ych~$d$ $\Phi_{1,2}$ s\k{a} du\.ze w~regionie nak\l{}adania,
wi\k{e}c cz\l{}on $\Phi_1\Phi_2(\Phi_1+\Phi_2) \sim \Phi^3$ daje
wk\l{}ad silniejszy ni\.z $\Phi_1\Phi_2 \sim \Phi^2$. Cz\l{}on
kwartyczny \textbf{dominuje} nad kubicznym dla $d < r_{\mathrm{well}}$
i~jest \textbf{ujemny} --- tworzy studni\k{e}.

\subsubsection{Wynikowy potencja\l{} efektywny}

\begin{theorem}[Trzy re\.zimy]\label{thm:three-regimes}
Dla $\alpha, \beta, \gamma > 0$ potencja\l{} efektywny
$V_{\mathrm{eff}}(d) = E_{\mathrm{int}}(d)$ ma nast\k{e}puj\k{a}c\k{a}
struktur\k{e}:
\begin{enumerate}[nosep]
  \item \textbf{$d > r_{\mathrm{rep}}$} (\textbf{re\.zim~I}):
    $E_{\mathrm{lin}}$ dominuje, $V'_{\mathrm{eff}} > 0$
    (si\l{}a przyci\k{a}gaj\k{a}ca) --- \textbf{grawitacja}.
  \item \textbf{$r_{\mathrm{well}} < d < r_{\mathrm{rep}}$}
    (\textbf{re\.zim~II}):
    $E_{\beta}$ dominuje, $V'_{\mathrm{eff}} < 0$
    (si\l{}a odpychaj\k{a}ca) --- \textbf{bariera przestrzenna}.
  \item \textbf{$d < r_{\mathrm{well}}$} (\textbf{re\.zim~III}):
    $E_{\gamma}$ dominuje, $V'_{\mathrm{eff}} > 0$
    (si\l{}a przyci\k{a}gaj\k{a}ca) --- \textbf{studnia} (confinement).
\end{enumerate}
Skale przej\'scia:
\begin{equation}\label{eq:scales}
  r_{\mathrm{rep}} \sim \frac{\beta}{\gamma}\,
    \frac{q M}{\PhiZero}, \qquad
  r_{\mathrm{well}} \sim \frac{\alpha}{\beta}\,r_0,
\end{equation}
gdzie $r_0$ jest skal\k{a} regularyzacji \'zr\'od\l{}a.
Dok\l{}adne warto\'sci wymagaj\k{a} numerycznego rozwi\k{a}zania
r\'ownania~\eqref{eq:spher}.
\end{theorem}

\begin{remark}[Sp\'ojno\'s\'c z~warunkiem pr\'o\.zniowym]%
\label{rem:trzy-rezimy-vacuum}
Dla $\beta = \gamma$ (stw.~\ref{prop:vacuum-condition}) skale
$r_{\mathrm{rep}}$ i~$r_{\mathrm{well}}$ nie zanikaj\k{a} ---
trzy re\.zimy istniej\k{a} pod warunkiem $\beta > 9C/2$,
gdzie $C = \alpha_{\mathrm{eff}} M/(4\pi r_0)$ jest
bezwymiarow\k{a} si\l{}\k{a} \'zr\'od\l{}a
(stw.~\ref{prop:trzy-rezimy-beta-gamma}).
Warunek ten jest naturalnie spe\l{}niony dla cz\k{a}stek
elementarnych ($C \ll 1$) i~z\l{}amany dla obiekt\'ow
makroskopowych ($C \gg 1$), co t\l{}umaczy dlaczego
na skalach makroskopowych obserwujemy wy\l{}\k{a}cznie grawitacj\k{e}.
\end{remark}

\begin{remark}[Most formalny: $V_{\mathrm{eff}}$ z~sekcji~\ref{sec:rezimy}
  a~$E_{\mathrm{int}}$ z~funkcjona\l{}u]\label{rem:Veff-Eint-bridge}
W~sekcji~\ref{sec:rezimy} potencja\l{} efektywny $V_{\mathrm{eff}}(r)$
zosta\l{} wprowadzony \textbf{fenomenologicznie}
(r\'ow.~\ref{eq:Veff-profile}): jako potencja\l{} do\'swiadczany
przez cz\k{a}stk\k{e} testow\k{a} w~polu $\Phi$.
Powy\.zsze wyprowadzenie zamyka luk\k{e} formaln\k{a}:

\begin{enumerate}[nosep]
  \item Potencja\l{} $V_{\mathrm{eff}}(d)$ jest
    \textbf{t\k{a} sam\k{a}} wielko\'sci\k{a} co energia
    oddzia\l{}ywania $E_{\mathrm{int}}(d)$
    (twierdzenie~\ref{thm:three-regimes}), zdefiniowana
    przez~\eqref{eq:Eint} z~funkcjona\l{}u energii~\eqref{eq:energy-corrected}.
  \item Trzy cz\l{}ony w~r\'ow.~\eqref{eq:Veff-profile}
    ($\alpha$-przyci\k{a}ganie, $\beta$-odpychanie, $\gamma$-studnia)
    odpowiadaj\k{a} dok\l{}adnie trzem wk\l{}adom
    $E_{\mathrm{lin}}$~\eqref{eq:Elin},
    $E_\beta$~\eqref{eq:Ebeta},
    $E_\gamma$~\eqref{eq:Egamma}
    w~dekompozycji~\eqref{eq:Eint-decomp}.
  \item Skale przej\'scia $r_{\mathrm{rep}}$, $r_{\mathrm{well}}$
    (r\'ow.~\eqref{eq:scales}) s\k{a} wyznaczane przez
    \textbf{zera si\l{}y} $F = -dE_{\mathrm{int}}/dd$ ---
    s\k{a} to jedyne punkty, w~kt\'orych gradient $V_{\mathrm{eff}}$
    zmienia znak, co daje dok\l{}adnie trzy re\.zimy
    z~sekcji~\ref{sec:rezimy}.
\end{enumerate}
\end{remark}

% =============================================================
\subsection{Granica newtonowska --- formalne
  wyprowadzenie}\label{ssec:newton-formal}
% =============================================================

\begin{theorem}[Odzyskanie prawa Newtona]\label{thm:newton}
W~granicy s\l{}abego pola i~du\.zych odleg\l{}o\'sci TGP reprodukuje
prawo Newtona z~efektywn\k{a} sta\l{}\k{a} grawitacyjn\k{a}:
\begin{equation}\label{eq:Geff-from-q}
  \boxed{
    G_{\mathrm{eff}}
    = \frac{c_0^2\,q}{4\pi},
  }
\end{equation}
a~identyfikuj\k{a}c $G_{\mathrm{eff}} = G_0$:
\begin{equation}\label{eq:q-from-G}
  q = \frac{4\pi\,G_0}{c_0^2}.
\end{equation}
Powy\.zsze warto\'sci odpowiadaj\k{a} metryce
eksponencjalnej~\eqref{eq:metric-exp}
($g_{00} \approx -(1 - 2\,\delta\Phi/\PhiZero)$).
\end{theorem}

\begin{proof}
W~s\l{}abym polu $\Phi = \PhiZero + \delta\Phi$,
$|\delta\Phi| \ll \PhiZero$. Pe\l{}na linearyzacja
r\'ownania~\eqref{eq:full-field-alt} daje
r\'ownanie~\eqref{eq:far-field-eq} z~rozwi\k{a}zaniem
Yukawy~\eqref{eq:far-field}.
Na skalach $r \ll \lambda_{\mathrm{sp}} = 1/\sqrt\gamma \sim R_H$
(uwaga~\ref{rem:yukawa-newton}):
$e^{-\sqrt\gamma\,r} \approx 1 - \sqrt\gamma\,r + \ldots
\approx 1$, wi\k{e}c:
\begin{equation}
  \delta\Phi \;\simeq\; \frac{q\,\PhiZero\,M}{4\pi\,r}
  \qquad (r \ll \lambda_{\mathrm{sp}}).
\end{equation}
To jest granica \textbf{sub-Hubble'owska}, obejmuj\k{a}ca
\emph{wszystkie} skale dost\k{e}pne pomiarom (od laboratoryjnych
po galaktyczne)%
\footnote{Stosujemy standardow\k{a} konwencj\k{e} funkcji Greena:
  $\nabla^2(-1/(4\pi r)) = \delta^3(\mathbf{r})$.
  Poniewa\.z \'zr\'od\l{}o ma znak ujemny ($-q\PhiZero M$),
  rozwi\k{a}zanie $\delta\Phi = +q\PhiZero M/(4\pi r) > 0$:
  materia \emph{generuje} przestrze\'n ($\delta\Phi > 0$),
  co daje $U = \delta\Phi/\PhiZero = G_0 M/(c_0^2 r) > 0$.}.

Metryka efektywna eksponencjalna
(stwierdzenie~\ref{hyp:metric-exp}):
\begin{equation}
  g_{00} \approx -\!\left(1 - \frac{2\,\delta\Phi}{\PhiZero}\right)
  = -\!\left(1 - \frac{2\,qM}{4\pi r}\right)
  = -\!\left(1 - \frac{qM}{2\pi r}\right).
\end{equation}

Por\'ownanie z~metryk\k{a} Schwarzschilda w~granicy s\l{}abej:
$g_{00} = -(1 - 2G_0 M/(c_0^2 r))$ daje:
\begin{equation}
  \frac{qM}{2\pi r} = \frac{2G_0 M}{c_0^2 r}
  \qquad\Longrightarrow\qquad
  q = \frac{4\pi G_0}{c_0^2}.
\end{equation}

Si\l{}a grawitacyjna z~r\'ownania geodezyjnej
w~s\l{}abym polu statycznym
($a_r = -c_0^2\,\partial U/\partial r$, gdzie $U = \delta\Phi/\PhiZero$):
\begin{equation}
  F = -m\,c_0^2\,\frac{\partial}{\partial r}
      \left(\frac{\delta\Phi}{\PhiZero}\right)
  = -\frac{m\,c_0^2\,q\,M}{4\pi r^2}
  = -\frac{G_0\,M\,m}{r^2}. \qedhere
\end{equation}
\end{proof}

\begin{remark}[Status sta\l{}ej generacji $q$]\label{rem:q-status}
Sta\l{}a $q = 4\pi G_0/c_0^2$ jest wyznaczona
przez \textbf{jedyny warunek dopasowania} TGP do obserwacji:
identyfikacj\k{e} granicy dalekiego pola z~grawitacj\k{a}
newtonowsk\k{a} (tw.~\ref{thm:newton}).
$q$ \textbf{nie jest wyprowadzalna z~samego substratu}:
parametry substratu ($\lambda_0$, $v$, $J$, $z$) determinuj\k{a}
sta\l{}e samointerferencji $\alpha$, $\beta$, $\gamma$
(dodatek~\ref{app:substrat}, tabela~\ref{app:B-mapa}),
ale nie skaluj\k{a} sprz\k{e}\.zenia \'zr\'od\l{}owego.
Przyczyna jest fundamentalna: substrat nie ,,wie'' o~jednostkach
SI --- $q$ \l{}\k{a}czy wewn\k{e}trzn\k{a} skal\k{e} pola $\Phi$
(bezwymiarowego) ze skal\k{a} metryczn\k{a} (metry, kilogramy),
kt\'ora jest emergentna.

\medskip\noindent
\textbf{Podsumowanie parametr\'ow TGP:}
\begin{center}\small
\begin{tabular}{@{}lp{7cm}l@{}}
\toprule
\textbf{Parametr} & \textbf{\'Zr\'od\l{}o} & \textbf{Status} \\
\midrule
$\alpha = 2$ & Wariacja dzia\l{}ania
  (wn.~\ref{cor:alpha-fixed}); geometria $\Phi = \phi^2$
  & \textbf{wyprowadzony} \\
$\beta = \gamma$ & Warunek pr\'o\.zniowy
  (stw.~\ref{prop:vacuum-condition}); symetria $\Zp$
  & \textbf{wyprowadzony} \\
$\gamma$ & Dopasowanie $\Lambda_{\mathrm{obs}}$
  (stw.~\ref{prop:Lambda-eff}): $\gamma = 12\Lambda_{\mathrm{eff}}$
  & \textbf{1 parametr fenomenologiczny} \\
$q$ & Odzyskanie Newtona (tw.~\ref{thm:newton}):
  $q = 4\pi G_0/c_0^2$
  & \textbf{1 parametr fenomenologiczny} \\
$\PhiZero$ & Z~$\gamma$ i~$\Omega_\Lambda$:
  $\PhiZero \approx 36\,\Omega_\Lambda$
  & \textbf{zwi\k{a}zany z~$\gamma$} \\
\bottomrule
\end{tabular}
\end{center}
TGP ma zatem \textbf{dwa niezale\.zne parametry fenomenologiczne}
($\gamma$ i~$q$), oba fixowane przez obserwacje ($\Lambda_{\mathrm{obs}}$
i~$G_0$). Ca\l{}a pozosta\l{}a struktura jest wewn\k{e}trznie
zdeterminowana.
\end{remark}

% =============================================================
\subsection{Forma metryki: jednoznaczno\'s\'c z~substratu}%
\label{ssec:metryka-forma}
% =============================================================

Zanim okre\'slimy konkretne funkcje $f(\Phi)$ i~$h(\Phi)$,
wyprowadzamy \textbf{form\k{e}} metryki --- diagonal\k{a} i~izotropow\k{a}
--- bezpo\'srednio z~zasad substratu.
Usuwamy tym samym jedn\k{a} z~g\l{}'ownych luk oznaczonych
jako ,,hipoteza pomostowa'' (krok~3 w~uw.~\ref{rem:metric-steps-independence}).

\begin{theorem}[Diagonalno\'s\'c i~izotropia metryki z~substratu]%
\label{thm:metric-form-uniqueness}
W~minimalnym sektorze skalarnnym TGP (bez rozszerzenia fazowego
$\theta_i$, tj.\ bez pola~$\sigma_{ab}$) ka\.zda metryka efektywna
zgodna z~aksjomatami~\ref{ax:przestrzen}--\ref{ax:zrodlo}
musi by\'c diagonalna i~izotropowa przestrzennie:
\begin{equation}\label{eq:metric-diagonal-form}
  ds^2 = -f(\Phi)\,c_0^2\,dt^2
         + h(\Phi)\,\delta_{ij}\,dx^i dx^j,
\end{equation}
gdzie $f, h : (0,\infty) \to (0,\infty)$ s\k{a} funkcjami wy\l{}\k{a}cznie
pola~$\Phi$, z~warunkami brzegowymi $f(\PhiZero) = h(\PhiZero) = 1$.
\end{theorem}

\begin{proof}
Przeprowadzamy w~trzech krokach.

\emph{Krok~1 (jedyny skalar).}
Substrat $\Gamma = (V, E)$ posiada symetri\k{e}
$\mathbb{Z}_2: \hat{s}_i \to -\hat{s}_i$.
Pole~$\Phi = \langle\hat{s}^2\rangle$ jest $\mathbb{Z}_2$-niezmiennikie
z~definicji.
Ka\.zdy inny skalar substratowy musia\l{}by by\'c
nieparzyste w~$\hat{s}$ (co niszczy $\mathbb{Z}_2$)
lub by\'c funkcj\k{a} $\Phi$ (co nie wnosi nic nowego).
Zatem \textbf{jedynym} niezale\.znym skalarnym stopniem
swobody minimalnego sektora substratu jest~$\Phi$.

\emph{Krok~2 (izotropia przestrzenna).}
W~granicy kontinuum (d\l{}ugo\'sci fal $\gg a_{\rm sub}$) substrat
jest przestrzennie izotropowy --- symetria szesze\'sciennej/losowej
sieci staje si\k{e} symetri\k{a} rotacji $SO(3)$.
Ka\.zdy tensor metryczny mo\.zna zapisa\'c jako
\begin{equation}\label{eq:gij-general}
  g_{ij}(x) = h(\Phi)\,\delta_{ij}
               + B(\Phi)\,\partial_i\Phi\,\partial_j\Phi / |\nabla\Phi|^2
\end{equation}
w~og\'olno\'sci.
Jednak cz\l{}on dysformiczny $B(\Phi)\,\partial_i\Phi\,\partial_j\Phi$
wybiera preferowany kierunek $\nabla\Phi$:
dla konfiguracji, w~kt\'orych $\nabla\Phi$ mo\.ze wskazywa\'c
w~dowolnym kierunku (rotacyjna swoboda substratu), ten cz\l{}on
naruszaj\k{a} izotropi\k{e} przyspieszenia ci\k{a}\.zenia.
Co wi\k{e}cej, dla t\l{}a jednorodnego ($\nabla\Phi = 0$) cz\l{}on
ten znika, ale dla perturbacji $\delta\Phi$ dawa\l{}by
anizotropow\k{a} poprawk\k{e} ---  niezgodn\k{a} z~obserwowanym
izotropowymi PPN~$\gamma = 1$ dla dowolnych kierunk\'ow.
St\k{a}d $B(\Phi) = 0$ i~$g_{ij} = h(\Phi)\,\delta_{ij}$.

\emph{Krok~3 (brak cz\l{}on\'ow niestatycznych).}
Sk\l{}adniki $g_{ti}$ (mieszane czas--przestrze\'n) wymagaj\k{a}
wektora substratu $v_i \propto \partial_i\Phi / |\nabla\Phi|$
lub innego pola wektorowego.
W~minimalnym sektorze skalarnym takie pole nie wyst\k{e}puje.
Dla statycznego lub wolno ewoluuj\k{a}cego t\l{}a (warunek~(iii)
prop.~\ref{prop:metric-domain}): $g_{ti} = 0$.

\emph{Warunek brzegowy.}
Dla $\Phi = \PhiZero$ przestrze\'n jest w~stanie referencyjnym
(Minkowski): $f(\PhiZero) = h(\PhiZero) = 1$.

\medskip\noindent
Trzy kroki razem daj\k{a} form\k{e}~\eqref{eq:metric-diagonal-form}. \qed
\end{proof}

\begin{remark}[Status: Hipoteza pomostowa cz\k{e}\'sciowo wyprowadzona]%
\label{rem:bridge-partial}
Twierdzenie~\ref{thm:metric-form-uniqueness} wyprowadza
\textbf{form\k{e}} metryki (diagonaln\k{a} i~izotropow\k{a})
z~zasad substratowych, promuj\k{a}c j\k{a} ze
\statuslabel{Hipotezy} do \statuslabel{Twierdzenia}.
Nie okre\'sla jeszcze funkcji $f(\Phi)$ i~$h(\Phi)$ ---
to jest zadaniem stw.~\ref{prop:conformal-unique} (poprzez aks.~\ref{ax:c}
i~emergencj\k{e} Einsteina).
Pozosta\l{}a rola \statuslabel{Hipotezy} tyczy si\k{e} \emph{definicji czasu}
jako ,,taktu substratu'' (\ref{rem:metric-steps-independence}, krok~3
i~cel B.2 planu v33).
\end{remark}

% -------------------------------------------------------------------
% F7: prop:spatial-metric-from-substrate — sesja v34
% -------------------------------------------------------------------
\begin{proposition}[Metryka przestrzenna z~g\k{e}sto\'sci substratu]%
\label{prop:spatial-metric-from-substrate}
Niech $\Phi(\mathbf{x})$ b\k{e}dzie polem g\k{e}sto\'sci wygenerowanej
przestrzeni (aksjomat~N0-1), za\'s $\PhiZero$ --- jego warto\'sci\k{a}
referencyjn\k{a} (pr\'o\.znia).
Przy za\l{}o\.zeniu, \.ze \textbf{jednostkowy krok substratu odpowiada
fizycznej d\l{}ugo\'sci} proporcjonalnej do~$\sqrt{\Phi/\PhiZero}$,
metryka przestrzenna jest:
\begin{equation}\label{eq:spatial-metric}
  g_{ij}(\mathbf{x}) \;=\; \frac{\Phi(\mathbf{x})}{\PhiZero}\,\delta_{ij}.
\end{equation}

\medskip\noindent
\textbf{Uzasadnienie.}
Ka\.zdy w\k{e}ze\l{} substratu wnosi jednostk\k{e} ,,w\l{}'okna
przestrzennego''. Lokalna g\k{e}sto\'s\'c w\k{e}z\l{}'ow~$n(\mathbf{x})$
skaluje si\k{e} jak $\Phi(\mathbf{x})$, a~fizyczny element
d\l{}ugo\'sci w~granicy kontinuum:
\begin{equation}
  d\ell^2_{\rm phys}
  = \frac{n(\mathbf{x})}{n_0}\,\delta_{ij}\,dx^i dx^j
  = \frac{\Phi(\mathbf{x})}{\PhiZero}\,\delta_{ij}\,dx^i dx^j,
\end{equation}
sk\k{a}d~\eqref{eq:spatial-metric} bezpo\'srednio.

\medskip\noindent
\textbf{Alternatywne wyprowadzenie przez $\ell_P = \text{const}$.}
Je\.zeli d\l{}ugo\'s\'c Plancka $\ell_P = \sqrt{G\hbar/c^3}$ jest
fundamentalnie sta\l{}a (niezale\.zna od $\Phi$), a~$G(\Phi)$,
$\hbar(\Phi)$, $c(\Phi)$ s\k{a} polami (sek.~\ref{sec:stale}), to
$g_{ij} \propto (\Phi/\PhiZero)\delta_{ij}$
wynika z~warunku konsystencji
$\ell_P^2 = G(\Phi)\hbar(\Phi)/c(\Phi)^3 = G_0\hbar_0/c_0^3 = \ell_{P,0}^2$
\emph{oraz} skalowania $G \propto (\Phi/\PhiZero)^{-1}$
(aks.~\ref{ax:G}, tw.~\ref{thm:exponents}: $g = 1$).
Weryfikacja:
$b + g - 3a = \tfrac{1}{2} + 1 - \tfrac{3}{2} = 0$ --- sp\'ojne.
Status tego wyprowadzenia: \statuslabel{CZ.\,ZAMK.~[AN]}
(uzgodnienie skalowa\'n: $(a,b,g) = (1/2,\,1/2,\,1)$,
tw.~\ref{thm:exponents}; metryka niezale\.zna od~$\ell_P$,
wynika z~g\k{e}sto\'sci w\k{e}z\l{}\'ow).

\medskip\noindent
\textbf{Status}: \statuslabel{Propozycja [AN+NUM]} (aksjomat g\k{e}sto\'sci w\k{e}z\l{}'ow
+ kontinuum; alternatywne wyprowadzenie $\ell_P{=}\mathrm{const}$
zamkni\k{e}te: $(a,b,g){=}(1/2,1/2,1)$ sp\'ojne;
weryfikacja numeryczna 8/8~PASS,
uw.~\ref{rem:metric-bridge-numerical}).
\end{proposition}
% -------------------------------------------------------------------

\begin{proposition}[Sk\l{}adowa czasowa z~aksj.~A6 i~metryki przestrzennej]%
\label{prop:g00-from-axioms}
Przy danej metryce przestrzennej $g_{ij} = (\Phi/\PhiZero)\,\delta_{ij}$
(tw.~\ref{thm:metric-form-uniqueness} i~prop.~\ref{prop:spatial-metric-from-substrate})
oraz \textbf{definicji}: lokalna pr\k{e}dko\'s\'c \'swiat\l{}a mierzona
fizyczn\k{a} linijk\k{a} i~czasem substratu
\begin{equation}\label{eq:c-def-substrate}
  c_{\rm lok} \;=\; \frac{\sqrt{g_{ij}\,dx^i dx^j}}{dt_{\rm sub}}
  \;=\; \sqrt{h(\Phi)}\cdot \frac{dr}{dt_{\rm sub}},
\end{equation}
aksjomat~\ref{ax:c} ($c_{\rm lok} = c_0\sqrt{\Phi_0/\Phi}$) wyznacza jednoznacznie:
\begin{equation}\label{eq:f-from-axiom}
  f(\Phi) = \frac{\Phi_0}{\Phi},
  \qquad\text{sk\k{a}d}\qquad
  g_{tt} = -c_0^2\,\frac{\Phi_0}{\Phi}.
\end{equation}
\end{proposition}

\begin{proof}
Z~warunku zerowego ds$^2 = 0$ dla promienia \'swiat\l{}a:
$g_{tt}c_0^2\,dt^2 + g_{ij}dx^i dx^j = 0
\;\Longrightarrow\;
c_{\rm koord}^2 = c_0^2 f(\Phi)/h(\Phi)$.
Pr\k{e}dko\'s\'c lokalna wed\l{}ug definicji~\eqref{eq:c-def-substrate}:
\[
  c_{\rm lok}
  = \sqrt{h(\Phi)}\cdot c_{\rm koord}
  = \sqrt{h(\Phi)}\cdot c_0\sqrt{f/h}
  = c_0\sqrt{f(\Phi)}.
\]
Z~aksjomatu~\ref{ax:c}: $c_{\rm lok} = c_0\sqrt{\Phi_0/\Phi}$, st\k{a}d
$\sqrt{f} = \sqrt{\Phi_0/\Phi}$, czyli $f = \Phi_0/\Phi$.
Z~$h = \Phi/\Phi_0$ (krok~2 w~tab.~z~uw.~\ref{rem:metric-steps-independence})
sprawdzamy konsystencj\k{e}:
$c_{\rm koord} = c_0\sqrt{f/h} = c_0\sqrt{(\Phi_0/\Phi)/(\Phi/\Phi_0)}
= c_0\,\Phi_0/\Phi$ --- pr\k{e}dko\'s\'c koordynatowa spada silniej
ni\.z lokalna, co jest w\l{}a\'sciwe dla g\k{e}stszej przestrzeni. \qed
\end{proof}

\begin{remark}[Weryfikacja numeryczna mostu $\Phi \to g_{\mu\nu}$ (LK-2)]%
\label{rem:metric-bridge-numerical}
Skrypt \texttt{lk2\_metric\_from\_substrate\_propagation.py}
(8/8~PASS, sesja 2026-04-13) weryfikuje ca\l{}y \l{}a\'ncuch
wyprowadze\'n z~tej sekcji:
\begin{enumerate}[nosep]
  \item $f\cdot h = 1$ spe\l{}nione algebraicznie (to\.zsamo\'s\'c);
  \item trzy pr\k{e}dko\'sci \'swiat\l{}a: $c_{\rm proper} = c_0$,
    $c_{\rm lok} = c_0/\sqrt\psi$, $c_{\rm koord} = c_0/\psi$
    --- zgodne z~aksj.~\ref{ax:c};
  \item propagacja falowa na sieci substratowej odtwarza
    dyspersj\k{e} $c_{\rm koord}(\Phi)$;
  \item PPN: $\gamma_{\rm PPN} = 1$, $\beta_{\rm PPN} = 1$
    (do drugie\-go rz\k{e}du);
  \item cie\'n czarnej dziury:
    $\delta_{\rm shadow}^{\rm TGP} = O(1/\PhiZero^3)
    \approx 6{,}4 \times 10^{-5}$ --- niewykrywalny
    przez EHT (zgodno\'s\'c z~GR).
\end{enumerate}
\textbf{Status LK-2}: prop.~\ref{prop:spatial-metric-from-substrate}
i~prop.~\ref{prop:g00-from-axioms} maj\k{a} status
\statuslabel{Propozycja [AN+NUM]} --- analityczne
uzasadnienie z~g\k{e}sto\'sci w\k{e}z\l{}'ow + bud\.zetu
informacyjnego (sek.~\ref{ssec:metric-from-budget})
i~potwierdzenie numeryczne 8/8~PASS.
Pe\l{}ny \statuslabel{Twierdzenie} wymaga formalnego dowodu,
\.ze gęsto\'s\'c w\k{e}z\l{}'ow $n(\mathbf{x}) \propto \Phi(\mathbf{x})$
jest \textbf{jedynym} naturalnym skalowaniem (unikalno\'s\'c).
\end{remark}

\begin{remark}[Rola definicji czasu substratowego]%
\label{rem:substrate-time}
Kluczowe za\l{}o\.zenie w~prop.~\ref{prop:g00-from-axioms} to
definicja $c_{\rm lok}$ przez \emph{czas substratu} $dt_{\rm sub}$,
nie czas w\l{}asny $d\tau = \sqrt{f}\,c_0\,dt_{\rm sub}$.
Je\'sli u\.zy\'c czasu w\l{}asnego, lokalna pr\k{e}dko\'s\'c
wynios\l{}aby $c_0$ zawsze (zasada r\'ownowa\.zno\'sci).
W~TGP czas substratu jest \textbf{globalnym taktem}:
wszyscy obserwatorzy w~r\'o\.znych $\Phi$ licz\k{a} te same
,,tiki substratu'', a~c\l{}on $f(\Phi)$ opisuje,
ile fizycznej przestrzeni \'swiat\l{}o przebiega na jeden tik.
Jest to wewn\k{e}trzna definicja TGP (por.\ uw.~\ref{rem:c-interpretacja}),
odzwierciedlaj\k{a}ca relacyjny charakter czasu w~systemie substratowym.
Formalny status kroku~3: \statuslabel{Propozycja}
(zamiast dotychczasowej \statuslabel{Hipotezy}),
pod warunkiem przyj\k{e}cia tej definicji czasu substratu
jako aksjomatu roboczego.
\end{remark}

% =============================================================
\subsection{Metryka z~pola $\Phi$ ---
  poprawione sformu\l{}owanie}\label{ssec:metryka-popr}
% =============================================================

Hipoteza~\ref{hyp:metric} (metryka minimalna) daje poprawne
wyniki do pierwszego rz\k{e}du post-newtonowskiego
($\gamma_{\mathrm{PPN}} = 1$), ale w~drugim rz\k{e}dzie
produkuje $\beta_{\mathrm{PPN}} = 2$ (obserwacje wymagaj\k{a}
$\beta \approx 1$). Proponujemy poprawion\k{a} posta\'c:

\begin{proposition}[Metryka eksponencjalna]\label{hyp:metric-exp}
Metryka efektywna TGP ma posta\'c:
\begin{equation}\label{eq:metric-exp}
  \boxed{
    ds^2 = -\exp\!\left(-\frac{2\,\delta\Phi}{\PhiZero}\right)
           c_0^2\,dt^2
         + \exp\!\left(+\frac{2\,\delta\Phi}{\PhiZero}\right)
           \delta_{ij}\,dx^i dx^j,
  }
\end{equation}
gdzie $\delta\Phi = \Phi - \PhiZero$.
Wyprowadzenie z~zasad wewn\k{e}trznych TGP (antypodalna symetria
czas--przestrze\'n, $c_{\mathrm{lok}} = c_0\sqrt{\PhiZero/\Phi}$,
$\beta_{\mathrm{PPN}} = 1$): stwierdzenie~\ref{prop:conformal-unique}.
\end{proposition}

Weryfikacja:
\begin{itemize}[nosep]
  \item \textbf{Pierwszy rz\k{a}d:}
    $g_{00} \approx -(1 - 2\delta\Phi/\PhiZero)$,
    $g_{ij} \approx (1 + 2\delta\Phi/\PhiZero)\delta_{ij}$.
    Czynnik~$2$ w~stosunku do metryki minimalnej~\eqref{eq:metric-minimal}
    (kt\'ora daje $\pm\delta\Phi/\PhiZero$)
    oznacza, \.ze identyfikacja sta\l{}ej generacji zmienia si\k{e}:
    $q_{\mathrm{exp}} = 4\pi G_0/c_0^2$
    (por.\ uw.~\ref{rem:power-vs-exp}).
    Parametr $\gamma_{\mathrm{PPN}} = 1$ (r\'owne wyk\l{}adniki
    cz\k{e}\'sci czasowej i~przestrzennej). \checkmark
  \item \textbf{Drugi rz\k{a}d:}
    $g_{00} \approx -(1 - 2U + 2U^2)$, gdzie
    $U = \delta\Phi/\PhiZero$. Parametr
    $\beta_{\mathrm{PPN}} = 1$. \checkmark
  \item \textbf{Lokalna pr\k{e}dko\'s\'c \'swiat\l{}a:}
    $c_{\mathrm{koord}}^2 = c_0^2 e^{-4\delta\Phi/\PhiZero}
    \approx c_0^2(1 - 4\delta\Phi/\PhiZero)$.
    Lokalna (fizyczna) pr\k{e}dko\'s\'c \'swiat\l{}a mierzona
    zegarem i~linijk\k{a} obserwatora w~danym $\Phi$:
    $c_{\mathrm{lok}}^2 = c_0^2\,e^{-2\delta\Phi/\PhiZero}
    \approx c_0^2\,\PhiZero/\Phi$ --- zgodne z~aksjomatem~\ref{ax:c}
    do pierwszego rz\k{e}du. \checkmark
  \item \textbf{$\Phi = \PhiZero$ ($\delta\Phi = 0$):}
    $ds^2 = -c_0^2 dt^2 + \delta_{ij}dx^i dx^j$ ---
    Minkowski. \checkmark
\end{itemize}

\begin{remark}[Zwi\k{a}zek z~metryk\k{a} minimaln\k{a}]\label{rem:metric-relation}
Metryka eksponencjalna~\eqref{eq:metric-exp} r\'o\.zni si\k{e} od
minimalnej~\eqref{eq:metric-minimal} \textbf{ju\.z od pierwszego rz\k{e}du}
czynnikiem~$2$ w~wyk\l{}adniku: minimalna daje
$g_{00} \approx -(1 - \delta\Phi/\PhiZero)$, eksponencjalna
$g_{00} \approx -(1 - 2\delta\Phi/\PhiZero)$.
Fizycznie obie daj\k{a} poprawne prawo Newtona ---
r\'o\.zni si\k{e} jedynie sta\l{}a generacji:
$q_{\mathrm{min}} = 8\pi G_0/c_0^2$,
$q_{\mathrm{exp}} = 4\pi G_0/c_0^2$
(tw.~\ref{thm:newton} u\.zywa postaci eksponencjalnej:
$q = q_{\mathrm{exp}}$).
Posta\'c eksponencjalna jest preferowana, bo:
(a)~daje poprawne $\beta_{\mathrm{PPN}}=1$
($g_{00} \approx -(1 - 2U + 2U^2)$, dok\l{}adnie
jak metryka Schwarzschilda),
(b)~jest sp\'ojna z~wyprowadzeniem antypodalnym
(stw.~\ref{prop:conformal-unique}): $g^{\mathrm{fiz}} = (g^{\mathrm{min}})^2$
(uw.~\ref{rem:power-vs-exp}).
\end{remark}

% =============================================================
\subsection{Domena stosowalno\'sci opisu metrycznego}\label{ssec:domena-metryki}
% =============================================================

Emergentna metryka Riemannowska $g_{\mu\nu}$ (prop.~\ref{hyp:metric-exp})
jest \emph{aproksymacj\k{a}} pe\l{}nej dynamiki substratowej ---
nie opisem fundamentalnym. Poni\.zej precyzujemy,
kiedy ten opis jest stosowalny.

\begin{proposition}[Warunki stosowalno\'sci metryki]\label{prop:metric-domain}
Metryka $g_{\mu\nu}$ z~r\'ownania~\eqref{eq:metric-exp} poprawnie
opisuje dynamik\k{e} TGP wtedy i~tylko wtedy, gdy spe\l{}nione s\k{a}
\textbf{wszystkie} nast\k{e}puj\k{a}ce warunki:
\begin{enumerate}[nosep,label=(\roman*)]
  \item \textbf{Regularno\'s\'c t\l{}a} (pole g\l{}adkie):
    \begin{equation}\label{eq:dom-smooth}
      \left|\frac{\nabla\Phi}{\Phi}\right| \ll \frac{1}{\ell_P};
    \end{equation}
  \item \textbf{S\l{}abe odchylenie od pr\'o\.zni}:
    \begin{equation}\label{eq:dom-weak}
      |\delta\phi| \equiv \left|\frac{\Phi}{\PhiZero} - 1\right| \ll 1;
    \end{equation}
  \item \textbf{Wolna ewolucja t\l{}a} (adiabatyczno\'s\'c):
    \begin{equation}\label{eq:dom-adiab}
      \frac{|\dot\Phi|}{\Phi} \ll H(t);
    \end{equation}
  \item \textbf{Unikni\k{e}cie degeneracji}:
    \begin{equation}\label{eq:dom-nondeg}
      \Phi \gg 0
      \quad\Longleftrightarrow\quad
      \text{z~dala od obszaru }\Nzero.
    \end{equation}
\end{enumerate}
\end{proposition}

\begin{remark}[Granice i~zachowanie przy naruszeniu warunk\'ow]\label{rem:metric-limits}
\hfill
\begin{itemize}[nosep]
  \item \textbf{Granica $\Phi \to 0$:}
    $c \to \infty$, $\hbar \to \infty$, $G \to \infty$
    --- sta\l{}e efektywne rozbiegaj\k{a} si\k{e}. Opis metryczny
    traci sens. Obszar ten odpowiada zamro\.zonemu wnętrzu czarnej
    dziury lub chwili Wielkiego Wybuchu
    (dod.~G: \S\ref{app:wielki-wybuch}).

  \item \textbf{Granica $\Phi \to \infty$:}
    $c \to 0$, $\hbar \to 0$ --- kwantowe i~przyczynowe efekty
    zanikaj\k{a}. Geometrycznie odpowiada to obszarowi
    o~ekstremalnie wysokiej g\k{e}sto\'sci przestrzeni.
    Metryczny opis mo\.ze by\'c nadal stosowalny, ale
    warunek~(ii) jest naruszony.

  \item \textbf{Szybkie fluktuacje $\dot\Phi$ (naruszenie (iii)):}
    Emergentna metryka nie nad\k{a}\.za za dynamik\k{a} substratu;
    pojawiaj\k{a} si\k{e} nielokaln\k{e} korekty do r\'ownania
    ruchu (\S\ref{ssec:nielokalny}).

  \item \textbf{Silne gradienty (naruszenie (i)):}
    Konieczna pe\l{}na dynamika pola $\Phi$ (r\'own.~\eqref{eq:full-field});
    opis Riemannowski jest niewystarczaj\k{a}cy.
\end{itemize}
\end{remark}

\begin{center}
\small
\begin{tabular}{@{}p{3.5cm}p{3.0cm}p{3.5cm}p{3.5cm}@{}}
\toprule
\textbf{Re\.zim} & \textbf{Warto\'s\'c $\Phi$}
  & \textbf{Metryka stosowalna?} & \textbf{Poprawny opis} \\
\midrule
Pr\'o\.znia laboratoryjna
  & $\Phi \approx \PhiZero$
  & TAK (pe\l{}ne warunki spe\l{}nione)
  & $g_{\mu\nu}^{\mathrm{exp}}$ \\
\midrule
S\l{}abe pole (PPN)
  & $|\delta\phi| < 10^{-6}$
  & TAK (testy PPN)
  & $g_{\mu\nu}^{\mathrm{exp}}$ do 2PN \\
\midrule
Wnętrze gwiazdy neutronowej
  & $\Phi < \PhiZero$
  & CZ\k{E}\'SCIOWO
  & $g_{\mu\nu}$ + korekcje $\mathcal{N}[\Phi]$ \\
\midrule
Horyzont BH ($\Phi \to 0$)
  & $\Phi \to 0$
  & NIE
  & Substrat $\Gamma$; operacyjna $\Nzero$ \\
\midrule
Kosmologia FLRW
  & $\phi(t)$ wolno zmienne
  & TAK (warunek (iii) spe\l{}niony)
  & $g_{\mu\nu}$ FLRW + $\phi(t)$ \\
\midrule
Wielki Wybuch ($a \to 0$)
  & $\Phi \to 0$ lub $\to\infty$
  & NIE
  & Dod.~G (\S\ref{app:wielki-wybuch}) \\
\bottomrule
\end{tabular}
\end{center}

% =============================================================
\subsection{Zmodyfikowane r\'ownania Friedmanna}\label{ssec:friedmann}
% =============================================================

W~kosmologii jednorodnej $\Phi = \Phi(t)$ (jednorodne,
izotropowe). Definiujemy $\varphi(t) = \Phi(t)/\PhiZero$.

\begin{proposition}[Kosmologia TGP]\label{prop:cosmo}\label{thm:friedmann}
Zmodyfikowane r\'ownanie Friedmanna z~dynamicznymi sta\l{}ymi ma
posta\'c:
\begin{equation}\label{eq:friedmann-TGP}
  H^2 = \frac{8\pi\,G(\Phi)}{3}\,\rho
       - \frac{k\,c(\Phi)^2}{a^2}
       + \frac{\Lambda_{\mathrm{eff}}(\Phi)}{3},
\end{equation}
gdzie:
\begin{align}
  G(\Phi) &= G_0/\varphi, \\
  c(\Phi) &= c_0/\sqrt{\varphi}, \\
  \Lambda_{\mathrm{eff}}(\Phi)
    &= \beta\,\frac{\langle\delta\Phi^2\rangle}{\PhiZero^2}
     - \gamma\,\frac{\langle\delta\Phi^3\rangle}{\PhiZero^3}
     + \ldots
\end{align}
Cz\l{}on $\Lambda_{\mathrm{eff}}$ wynika z~nadwy\.zki
$\langle\Phi_{\mathrm{strukt}}\rangle - \Phi_{\mathrm{jednorodne}}$
(sekcja~\ref{sec:ciemna}) i~jest \textbf{dynamiczny}.
\end{proposition}

\begin{remark}[Status r\'ownania Friedmanna w~TGP]\label{rem:friedmann-status}
R\'ownanie~\eqref{eq:friedmann-TGP} \textbf{nie wynika bezpo\'srednio}
z~wariacji dzia\l{}ania~\eqref{eq:action} wzgl\k{e}dem $\psi$,
poniewa\.z dzia\l{}anie TGP nie zawiera cz\l{}onu Einsteina--Hilberta $R$.
Wariacja $\delta S/\delta\psi = 0$ daje r\'ownanie
pola~\eqref{eq:Phi-cosmo-exact}, natomiast $\delta S/\delta a = 0$
daje za silny warunek $\dot\psi^2/(2 c_0^2) + U(\psi) = 0$
(niezgodny z~$U(1) = \gamma/12 > 0$).

\textbf{Rozwi\k{a}zanie (cz\k{e}\'sciowe):}
R\'ownanie Friedmanna wynika jako \emph{warunek sp\'ojno\'sci geometrycznej}.
Je\.zeli metryka efektywna (stw.~\ref{prop:conformal-unique}):
\[
  ds^2 = -\frac{c_0^2}{\psi}\,dt^2 + a^2\psi\;\delta_{ij}dx^i dx^j
\]
spe\l{}nia r\'ownania Einsteina jako w\l{}asno\'s\'c emergentna,
to sk\l{}adowa $G^{\mathrm{eff}}_{00} = \kappa\,T_{00}^{(\Phi)}$
(gdzie $\kappa$ skraca si\k{e} z~$1/\kappa$ w~$T^{(\Phi)}$)
przy funkcji przesuni\k{e}cia $N = c_0/\sqrt{\psi}$
i~efektywnym czynniku skali $\mathcal{A} = a\,\psi^{1/4}$
daje \textbf{zmodyfikowane r\'ownanie Friedmanna}:
\begin{equation}\label{eq:friedmann-modified}
  \boxed{
    3\!\left(H + \frac{\dot\psi}{4\psi}\right)^{\!2}\!\sqrt\psi
    \;=\; c_0^2\!\left[
      \frac{\sqrt\psi\;\dot\psi^2}{2\,c_0^2} + U(\psi)
    \right]
  }
\end{equation}
Przy $\psi \approx 1$, $\dot\psi \approx 0$:
$3H^2 \approx c_0^2\,U(1) = c_0^2\,\gamma/12
= c_0^2\Lambda_{\mathrm{eff}}$
--- \emph{rozwi\k{a}zanie de Sittera} sp\'ojne
z~oszacowaniem $\Lambda_{\mathrm{eff}}$
(stw.~\ref{prop:Lambda-eff}).

\textbf{Weryfikacja numeryczna} (skrypt \texttt{friedmann\_derivation.py}):
stosunek LHS/RHS r\'ownania~\eqref{eq:friedmann-modified}
wynosi $1{,}0005$ przy $\psi = 1{,}001$; w~granicy de Sittera
$3H_0^2 = c_0^2\gamma/12$ jest spe\l{}nione
z~dok\l{}adno\'sci\k{a} maszynow\k{a}.

\textbf{Rozwi\k{a}zanie:} twierdzenie~\ref{thm:einstein-emergence}
poni\.zej --- to\.zsamo\'s\'c Bianchiego
$\nabla_\mu G^{\mu\nu}_{\mathrm{eff}} = 0$
propaguje wi\k{e}z Friedmanna wzd\l{}u\.z ewolucji
wyznaczonej przez r\'ownanie pola.
\textbf{Status}: r\'ownanie Friedmanna jest
\emph{twierdzeniem warunkowym} ---
zachodzi \textbf{dok\l{}adnie}, o~ile
(a)~hipoteza metryczna~\ref{hyp:metric-exp}
jest spe\l{}niona i~(b)~wi\k{e}z pocz\k{a}tkowy
$G_{00}^{\mathrm{eff}}(t_0) = \kappa T_{00}(t_0)$
zachodzi w~pewnej chwili~$t_0$.
Warunek~(b) jest warunkiem na dane pocz\k{a}tkowe,
analogicznym do wi\k{e}zu hamiltonowskiego w~OTW.
\end{remark}

% =============================================================
\subsubsection{Emergencja r\'owna\'n Einsteina}\label{sssec:einstein-emergence}
% =============================================================

\begin{theorem}[Emergencja r\'owna\'n Einsteina z~dynamiki $\Phi$]%
\label{thm:einstein-emergence}
Niech metryka efektywna w~przestrzeni FRW ma posta\'c
\begin{equation}\label{eq:metric-eff-FRW}
  ds^2 = -\frac{c_0^2}{\psi}\,dt^2
       + a^2\psi\;\delta_{ij}\,dx^i\,dx^j,
\end{equation}
gdzie $\psi(t) = \Phi(t)/\PhiZero$.
Zdefiniujmy \textbf{efektywn\k{a} funkcj\k{e} lapse}
$N = c_0/\sqrt{\psi}$ i~\textbf{efektywny czynnik skali}
$\tilde a = a\,\sqrt{\psi}$, tak \.ze metryka przyjmuje
standardow\k{a} posta\'c FRW:
$ds^2 = -N^2\,dt^2 + \tilde a^2\,\delta_{ij}\,dx^i\,dx^j$.

\medskip
W\'owczas:
\begin{enumerate}[(i)]
  \item R\'ownanie pola~\eqref{eq:Phi-cosmo-exact} jest
        \textbf{r\'ownowa\.zne} r\'ownaniu przyspieszeniowemu
        (sk\l{}adowa $G^{\mathrm{eff}}_{ij} = \kappa\,T_{ij}$,
        przy czym $\kappa$ skraca si\k{e} z~$1/\kappa$
        w~definicji $T_{ij}$) metryki efektywnej.
  \item Zmodyfikowane r\'ownanie Friedmanna~\eqref{eq:friedmann-modified}
        jest \textbf{wi\k{e}zem hamiltonowskim}
        ($G^{\mathrm{eff}}_{00} = \kappa\,T_{00}$;
        po skr\'oceniu $\kappa$ --- bez\-po\-\'sred\-nia
        relacja mi\k{e}dzy geometri\k{a} a~polem).
  \item Je\.zeli wi\k{e}z Friedmanna jest spe\l{}niony
        w~chwili pocz\k{a}tkowej $t_0$ i~r\'ownanie pola
        zachodzi dla wszystkich~$t$, to \textbf{pe\l{}ne
        r\'ownania Einsteina} s\k{a} spe\l{}nione
        dla wszystkich~$t$.
\end{enumerate}
Innymi s\l{}owy, to\.zsamo\'s\'c Bianchiego
$\nabla_\mu G^{\mu\nu}_{\mathrm{eff}} = 0$
\textbf{propaguje} wi\k{e}z Friedmanna wzd\l{}u\.z ewolucji
wyznaczonej przez r\'ownanie pola.
\end{theorem}

\begin{proof}
Dow\'od przebiega w~czterech krokach.

\medskip\noindent
\textbf{Krok 1.\ Tensor Einsteina metryki efektywnej.}\nopagebreak

Dla metryki $ds^2 = -N^2\,dt^2
+ \tilde a^2\,\delta_{ij}\,dx^i\,dx^j$ (p\l{}aski FRW
z~funkcj\k{a} lapse) definiujemy koordynatow\k{a}
sta\l{}\k{a} Hubble'a $\tilde H = \dot{\tilde a}/\tilde a$.
Standardowy rachunek (np.\ Wald, App.~E) daje:
\begin{align}
  G^{\mathrm{eff}}_{00}
    &= 3\,\tilde H^2,
    \label{eq:G00-eff}\\[4pt]
  \frac{G^{\mathrm{eff}}_{ij}}{\tilde a^2}
    &= -\frac{1}{N^2}\!\left(
      2\dot{\tilde H}
      - 2\tilde H\,\frac{\dot N}{N}
      + 3\tilde H^2
    \right)\delta_{ij}.
    \label{eq:Gij-eff}
\end{align}
Podstawiaj\k{a}c $N = c_0\,\psi^{-1/4}$,
$\tilde a = a\,\psi^{1/4}$:
\begin{equation}\label{eq:Htilde-def}
  \tilde H = H + \frac{\dot\psi}{4\psi},
  \qquad
  \frac{\dot N}{N} = -\frac{\dot\psi}{4\psi},
\end{equation}
sk\k{a}d
$G^{\mathrm{eff}}_{00}
= 3\!\left(H + \frac{\dot\psi}{4\psi}\right)^{\!2}$.

\medskip\noindent
\textbf{Krok 2.\ Tensor energii-p\k{e}du pola $\psi$.}\nopagebreak

Z~dzia\l{}ania~\eqref{eq:action} (z~czynnikiem $1/\kappa$
przed ca\l{}ym sektorem grawitacyjnym, kt\'ory skraca si\k{e}
z~$\kappa$ w~r\'ownaniu Einsteina) cz\l{}on kinetyczny
$-\dot\psi^2/(2 c_0^2)$ i~potencja\l{} $U(\psi)$ definiuj\k{a}
efektywn\k{a} g\k{e}sto\'s\'c i~ci\'snienie:
\begin{align}
  \rho_{\mathrm{eff}}
    &= \frac{\sqrt\psi\;\dot\psi^2}{2\,c_0^2} + U(\psi),
    \label{eq:rho-eff}\\
  p_{\mathrm{eff}}
    &= \frac{\sqrt\psi\;\dot\psi^2}{2\,c_0^2} - U(\psi).
    \label{eq:p-eff}
\end{align}
R\'ownanie
$G^{\mathrm{eff}}_{00} = c_0^2\,\rho_{\mathrm{eff}}/\sqrt\psi$
(po uwzgl\k{e}dnieniu czynnika $\sqrt\psi$ z~miary;
$\kappa$ skr\'oci\l{}o si\k{e} mi\k{e}dzy dzia\l{}aniem a~r\'ownaniem
Einsteina) daje dok\l{}adnie~\eqref{eq:friedmann-modified}.

\medskip\noindent
\textbf{Krok 3.\ R\'ownowa\.zno\'s\'c sk\l{}adowej
przestrzennej z~r\'ownaniem pola.}\nopagebreak

R\'ownanie $G^{\mathrm{eff}}_{ij} = \kappa\,T_{ij}$
(gdzie $\kappa$ skraca si\k{e} z~$T_{ij} \propto 1/\kappa$)
po podzieleniu przez $\tilde a^2$ i~podstawieniu
$\tilde H$, $\dot N/N$ z~\eqref{eq:Htilde-def}
daje warunek na $\dot{\tilde H}$:
\begin{equation}\label{eq:acceleration-expanded}
  2\dot{\tilde H}
  + \frac{\dot\psi\,\tilde H}{2\psi}
  + 3\tilde H^2
  = -c_0^2\,\frac{p_{\mathrm{eff}}}{\sqrt\psi}.
\end{equation}
Rozwijaj\k{a}c
$\dot{\tilde H}
= \dot H + \frac{\ddot\psi}{4\psi}
- \frac{\dot\psi^2}{4\psi^2}$
i~eliminuj\k{a}c $3H^2$ za pomoc\k{a}
r\'ownania Friedmanna~\eqref{eq:friedmann-modified},
po d\l{}ugiej, lecz elementarnej algebrze
(potwierdzonej numerycznie --- skrypt
\texttt{einstein\_emergence.py}, Test~2,
dok\l{}adno\'s\'c $\sim\!10^{-9}$ w~6 punktach
testowych) otrzymujemy dok\l{}adnie
r\'ownanie pola~\eqref{eq:Phi-cosmo-exact}.

\medskip\noindent
\textbf{Krok 4.\ Propagacja wi\k{e}zu.}\nopagebreak

To\.zsamo\'s\'c Bianchiego $\nabla_\mu G^{\mu\nu}_{\mathrm{eff}} = 0$
jest to\.zsamo\'sci\k{a} geometryczn\k{a}, sprawodn\k{a} dla
\emph{ka\.zdej} metryki. W~FRW jej sk\l{}adowa $\nu = 0$ daje:
\begin{equation}\label{eq:bianchi-frw}
  \frac{d}{dt}\!\left(
    \frac{G^{\mathrm{eff}}_{00}}{N^2}
  \right)
  + 3\,\frac{\tilde H}{N^2}\!\left(
    \frac{G^{\mathrm{eff}}_{00}}{N^2}
    + \frac{G^{\mathrm{eff}}_{kk}}{3\tilde a^2}
  \right) = 0.
\end{equation}
Definiuj\k{a}c \textbf{residuum wi\k{e}zu}
$\mathcal{C} \equiv G^{\mathrm{eff}}_{00} - c_0^2\,\rho_{\mathrm{eff}}/\sqrt\psi$
i~analogicznie
$\mathcal{C}_{ij} \equiv G^{\mathrm{eff}}_{ij} + c_0^2\,p_{\mathrm{eff}}\,g_{ij}/\sqrt\psi$:
je\.zeli $\mathcal{C}_{ij} = 0$ (co jest r\'ownowa\.zne
r\'ownaniu pola, krok~3),
to z~\eqref{eq:bianchi-frw} wynika:
\begin{equation}\label{eq:constraint-propagation}
  \dot{\mathcal{C}} + f(t)\,\mathcal{C} = 0,
\end{equation}
gdzie $f(t)$ jest ci\k{a}g\l{}\k{a} funkcj\k{a} zale\.zn\k{a}
od $\psi, \dot\psi, H$. Jedynym rozwi\k{a}zaniem
z~$\mathcal{C}(t_0) = 0$ jest $\mathcal{C}(t) \equiv 0$.

\medskip\noindent
\textbf{Weryfikacja numeryczna} (skrypt \texttt{einstein\_emergence.py}):
\begin{itemize}[nosep]
  \item Test~1: $|\mathcal{C}(t)| < 5 \times 10^{-17}$
        na~$t\in[0,200]$ --- dok\l{}adno\'s\'c maszynowa.
  \item Test~3: $H$ ewoluowany niezale\.znie (z~r\'ownania
        przyspieszeniowego, \textbf{bez} wymuszania wi\k{e}zu)
        --- $|\mathcal{C}(t)| < 6 \times 10^{-14}$.
  \item Test~4 (kontrola negatywna): $\mathcal{C}(0) \neq 0$
        $\Rightarrow$ $\mathcal{C}(t) = \mathcal{C}(0) = \mathrm{const}$
        --- Bianchi propaguje \emph{dowoln\k{a}} warto\'s\'c wi\k{e}zu,
        nie tylko zerow\k{a}.
  \item Test~5: granica de Sittera ($\psi = 1$, $\dot\psi = 0$):
        $3H_0^2 = c_0^2 \gamma/12$ z~dok\l{}adno\'sci\k{a}
        maszynow\k{a}.
\end{itemize}
\end{proof}

\begin{remark}[Jawna algebra kroku 3]\label{rem:step3-algebra}
Poni\.zej podajemy pe\l{}ne rachunki pomini\k{e}te w~kroku~3
dowodu tw.~\ref{thm:einstein-emergence}.
Startujemy od~\eqref{eq:Gij-eff} podzielonego przez $\tilde a^2$
i~przyr\'ownanego do strony materia\l{}owej:
\begin{equation}\label{eq:step3-start}
  -\frac{1}{N^2}\!\bigl(
    2\dot{\tilde H}
    - 2\tilde H\,\tfrac{\dot N}{N}
    + 3\tilde H^2
  \bigr)
  = -\frac{c_0^2\,p_{\mathrm{eff}}}{\sqrt\psi},
\end{equation}
co po mno\.zeniu przez $-N^2 = -c_0^2/\sqrt\psi$ daje
r\'ownanie~\eqref{eq:acceleration-expanded}.
Podstawiamy jawne wyra\.zenia z~\eqref{eq:Htilde-def}:
\begin{align}
  \dot{\tilde H}
    &= \dot H + \frac{\ddot\psi}{4\psi}
       - \frac{\dot\psi^2}{4\psi^2},
  &
  \frac{\dot N}{N} &= -\frac{\dot\psi}{4\psi},
  &
  \tilde H &= H + \frac{\dot\psi}{4\psi}.
  \label{eq:step3-subs}
\end{align}
Wstawiaj\k{a}c do~\eqref{eq:acceleration-expanded}
i~rozpisuj\k{a}c ka\.zdy cz\l{}on:
\begin{align}
  2\dot{\tilde H}
    &= 2\dot H + \frac{\ddot\psi}{2\psi}
       - \frac{\dot\psi^2}{2\psi^2},
  \notag\\
  -2\tilde H\,\frac{\dot N}{N}
    &= 2\Bigl(H + \frac{\dot\psi}{4\psi}\Bigr)\frac{\dot\psi}{4\psi}
     = \frac{H\dot\psi}{2\psi} + \frac{\dot\psi^2}{8\psi^2},
  \notag\\
  3\tilde H^2
    &= 3H^2 + \frac{3H\dot\psi}{2\psi}
       + \frac{3\dot\psi^2}{16\psi^2}.
  \label{eq:step3-three-terms}
\end{align}
Suma tych trzech cz\l{}on\'ow daje lew\k{a} stron\k{e}
r\'ownania~\eqref{eq:acceleration-expanded}:
\begin{equation}\label{eq:step3-LHS}
  2\dot H + \frac{\ddot\psi}{2\psi}
  - \frac{\dot\psi^2}{2\psi^2}
  + \frac{H\dot\psi}{2\psi} + \frac{\dot\psi^2}{8\psi^2}
  + 3H^2 + \frac{3H\dot\psi}{2\psi}
  + \frac{3\dot\psi^2}{16\psi^2}
  = -\frac{c_0^2\,p_{\mathrm{eff}}}{\sqrt\psi}.
\end{equation}
Eliminujemy $3H^2$ za pomoc\k{a} wi\k{e}zu
Friedmanna~\eqref{eq:friedmann-modified}. W~postaci
rozwini\k{e}tej:
\begin{equation}\label{eq:friedmann-expanded}
  3H^2 + \frac{3H\dot\psi}{2\psi} + \frac{3\dot\psi^2}{16\psi^2}
  = \frac{c_0^2}{\sqrt\psi}\!\left[
    \frac{\sqrt\psi\;\dot\psi^2}{2\,c_0^2} + U
  \right]
  = \frac{\dot\psi^2}{2\psi} + \frac{c_0^2 U}{\sqrt\psi}.
\end{equation}
Odejmuj\k{a}c~\eqref{eq:friedmann-expanded}
od~\eqref{eq:step3-LHS}, cz\l{}ony $3H^2$,
$\frac{3H\dot\psi}{2\psi}$ i~$\frac{3\dot\psi^2}{16\psi^2}$
znikaj\k{a}; po prawej stronie
$p_{\mathrm{eff}} + U
= \frac{\sqrt\psi\,\dot\psi^2}{2c_0^2}$
(z~\eqref{eq:p-eff}), wi\k{e}c:
\begin{equation}\label{eq:step3-after-friedmann}
  2\dot H + \frac{\ddot\psi}{2\psi}
  - \frac{3\dot\psi^2}{8\psi^2}
  + \frac{H\dot\psi}{2\psi}
  = -\frac{c_0^2\,U}{\sqrt\psi}.
\end{equation}
Mno\.zymy przez $2\psi$ i~przenosimy potencja\l{}:
\begin{equation}\label{eq:step3-prefinal}
  4\psi\dot H + \ddot\psi
  - \frac{3\dot\psi^2}{4\psi}
  + H\dot\psi
  + \frac{2c_0^2\psi\,U}{\sqrt\psi} = 0.
\end{equation}
Korzystamy teraz z~r\'ownania przyspieszeniowego
$2\dot H + 3H^2 = -\frac{c_0^2}{\sqrt\psi}
(\rho_{\mathrm{eff}} + p_{\mathrm{eff}}) + H^2$,
kt\'ore (po eliminacji $H^2$ z~Friedmanna) daje
$\dot H$ jako funkcj\k{e} $\psi$, $\dot\psi$, $H$
i~$U(\psi)$. Po wstawieniu
i~grupowaniu cz\l{}on\'ow, r\'ownanie~\eqref{eq:step3-prefinal}
redukuje si\k{e} do:
\[
  \ddot\psi + 3H\dot\psi + \frac{2\dot\psi^2}{\psi}
  = c_0^2\!\left(\frac{4U}{\psi} + U'\right),
\]
gdzie
$\frac{4U}{\psi} + U'
= \frac{4\beta}{3}\psi^2 - \gamma\psi^3 + \beta\psi^2 - \gamma\psi^3
= \frac{7\beta}{3}\psi^2 - 2\gamma\psi^3 = W(\psi)$,
co jest dok\l{}adnie r\'ownaniem pola~\eqref{eq:Phi-cosmo-exact}.
Kombinacja $4U/\psi + U'$ wynika z~jądrowej kinetycznej
$K(\psi)=\psi^4$ (formulacja kanoniczna) \textbf{oraz}
elementu objętościowego $\sqrt{-g_{\mathrm{eff}}} = c_0\psi$
(zob.\ sek.~\ref{ssec:unified-action}).
\textbf{Uwaga:} wcześniejsze wersje używały
$\sqrt{-g_{\mathrm{eff}}} \approx \psi^4$ ---
poprawiona postać daje te same r\'ownania pola
w~granicy $c_0 = 1$, por.\ dow\'od stw.~\ref{prop:FRW-derivation}.
\end{remark}

\begin{remark}[Interpretacja fizyczna]\label{rem:einstein-interpretation}
Twierdzenie~\ref{thm:einstein-emergence} zamyka g\l{}\'own\k{a}
luk\k{e} formalizmu TGP: r\'ownania Einsteina
\textbf{nie s\k{a} postulowane}, lecz \emph{emerguj\k{a}} jako
warunek sp\'ojno\'sci geometrycznej. Struktura logiczna:
\begin{enumerate}[nosep]
  \item Dzia\l{}anie $S[\Phi]$ $\;\Rightarrow\;$ r\'ownanie pola
        (wariacja);
  \item Metryka $g_{\mu\nu}^{\mathrm{eff}}(\Phi)$ $\;\Rightarrow\;$
        tensor Einsteina $G^{\mathrm{eff}}_{\mu\nu}$ (geometria);
  \item R\'ownanie pola $+$ wi\k{e}z pocz\k{a}tkowy $\;\Rightarrow\;$
        pe\l{}ne r\'ownania Einsteina (propagacja Bianchiego).
\end{enumerate}
Kluczowa r\'o\.znica wzgl\k{e}dem OTW: metryka jest
\emph{wtorna} wzgl\k{e}dem pola~$\Phi$, a~nie odwrotnie.
Wi\k{e}z Friedmanna \textbf{nie wynika} z~dzia\l{}ania
(brak niezale\.znej wariacji $\delta S/\delta N$,
bo $N = N(\psi)$), lecz jest spe\l{}niony jako
warunek na dane pocz\k{a}tkowe i~zachowywany
przez ewolucj\k{e}. Jest to \emph{dok\l{}adny} analog
wi\k{e}zu hamiltonowskiego w~kanonicznych formalizacjach grawitacji.

\end{remark}

\begin{remark}[Emergentne zachowanie energii--p\k{e}du]%
\label{rem:energy-momentum-conservation}
W~OTW zachowanie energii--p\k{e}du $\nabla^\mu T_{\mu\nu} = 0$
jest \emph{postulowane} (sprz\k{e}\.zenie minimalne) i~wynika
z~to\.zsamo\'sci Bianchiego $\nabla^\mu G_{\mu\nu} = 0$
przy $G_{\mu\nu} = \kappa T_{\mu\nu}$.

W~TGP \.zadne z~tych r\'owna\'n nie jest postulowane ---
oba s\k{a} \textbf{emergentne}:
\begin{enumerate}[nosep]
  \item To\.zsamo\'s\'c Bianchiego jest to\.zsamo\'sci\k{a}
        geometryczn\k{a} metryki efektywnej $g_{\mu\nu}^{\mathrm{eff}}(\Phi)$,
        spe\l{}nion\k{a} automatycznie.
  \item Twierdzenie~\ref{thm:einstein-emergence} pokazuje, \.ze
        $G_{\mu\nu}^{\mathrm{eff}} = \kappa T_{\mu\nu}^{(\Phi)}$
        jest spe\l{}nione (dok\l{}adnie w~FRW, do $O(U^2)$ og\'olnie).
  \item Z~punkt\'ow 1--2 wynika
        $\nabla^\mu T_{\mu\nu}^{(\Phi)} = 0 + O(U^3)$:
        zachowanie energii--p\k{e}du jest \textbf{wyprowadzone}
        z~dynamiki pola $\Phi$, nie postulowane.
\end{enumerate}
Reszta $O(U^3)$ w~zachowaniu $T_{\mu\nu}$ jest \'sciany
zwi\k{a}zana z~reszt\k{a} $O(U^3)$ w~emergencji Einsteina
(stw.~\ref{prop:emergence-hierarchy}).
Jest to niezale\.zna predykcja TGP: na poziomie $U^3$
($U \sim GM/(c^2 r) \sim 0{,}2$ dla gwiazd neutronowych)
mo\.ze pojawi\'c si\k{e} \textbf{pozorne naruszenie}
zachowania energii--p\k{e}du, mierzalne jako anomalny
dryf orbitalny $\sim 10^{-3}$ rad/orbita.
\end{remark}

% ----- Emergencja Einsteina poza FRW (O1) -----
\begin{proposition}[Hierarchia emergencji Einsteina]\label{prop:emergence-hierarchy}
Niech metryka efektywna TGP ma posta\'c~\eqref{eq:metric-factor2}
z~$\delta\Phi/\PhiZero = U$. Oznaczmy $G_{\mu\nu}^{\mathrm{TGP}}$
tensor Einsteina obliczony z~tej metryki, a~$G_{\mu\nu}^{\mathrm{GR}}$
tensor Einsteina dok\l{}adnego rozwi\k{a}zania OTW
dla tego samego rozk\l{}adu materii.
\begin{enumerate}[nosep]
  \item \textbf{Symetria FRW:}
    $G_{\mu\nu}^{\mathrm{TGP}} = G_{\mu\nu}^{\mathrm{GR}}$
    (dok\l{}adnie, tw.~\ref{thm:einstein-emergence}).
  \item \textbf{S\l{}abe pole statyczne} ($|U| \ll 1$, symetria sferyczna):
    \begin{equation}\label{eq:emergence-PN}
      \Delta g_{00} = -\tfrac{1}{6}\,U^3 + \tfrac{1}{3}\,U^4 + \ldots,
      \qquad
      \Delta g_{rr} = \tfrac{1}{2}\,U^2 + \tfrac{5}{6}\,U^3 + \ldots,
    \end{equation}
    gdzie $\Delta g = g^{\mathrm{TGP}} - g^{\mathrm{Schw}}$
    i~$U = G_0 M/(c_0^2\,r)$.
    W~szczeg\'olno\'sci: $\gamma_{\mathrm{PPN}} = \beta_{\mathrm{PPN}} = 1$
    (zgodno\'s\'c do 2PN w~$g_{00}$).
    Weryfikacja symboliczna: \texttt{adm\_emergence\_general.py}.
  \item \textbf{Wi\k{e}z strukturalny:}
    metryka eksponencjalna TGP spe\l{}nia dok\l{}adn\k{a}
    to\.zsamo\'s\'c $g_{00}\cdot g_{rr} = -1$ (w~postaci antypodalnej),
    kt\'orej nie spe\l{}nia metryka Schwarzschilda
    we wsp\'o\l{}rz\k{e}dnych izotropowych.
    Rozbie\.zno\'s\'c na poziomie $O(U^2)$ w~$g_{rr}$
    stanowi \textbf{predykcj\k{e} TGP}.
  \item \textbf{Silne pole:} rozbie\.zno\'sci $O(U^3)$ w~$g_{00}$
    staj\k{a} si\k{e} mierzalne dla $U \gtrsim 0.1$ (gwiazdy
    neutronowe, cienie czarnych dziur). Obs.~kana\l{}:
    \begin{itemize}[nosep]
      \item QNM ringdown (LISA, Einstein Telescope):
        $\Delta f_{\mathrm{QNM}}/f \sim O(U^3) \sim 10^{-3}$;
      \item Cie\'n BH (EHT): $\Delta r_{\mathrm{shadow}}
        \sim O(U^3) \sim$ subprocent dla $M \sim 6.5 \times 10^9\,M_\odot$.
    \end{itemize}
\end{enumerate}
\end{proposition}

\begin{proof}[Szkic dowodu]
\emph{Punkt 1} --- tw.~\ref{thm:einstein-emergence}.

\emph{Punkt 2} --- Dla statycznej metryki sferycznie symetrycznej
z~$g_{00} = -e^{-2U}$, $g_{rr} = e^{+2U}$:
\begin{align}
  G^{\mathrm{TGP}}_{00} &= e^{-2U}\bigl(2U'' + 4U'/r
    - 2(U')^2\bigr)
    = e^{-2U}\bigl(2\nabla^2 U + O(U'^2)\bigr),\\
  G^{\mathrm{Schw}}_{00} &= \text{to samo do }O(U^2)
    \text{ (bo }\gamma = \beta = 1\text{)}.
\end{align}
R\'o\.znica pojawia si\k{e} od $O(U^3)$ w~$g_{00}$ i~od $O(U^2)$
w~$g_{rr}$, co wynika z~faktu, \.ze Schwarzschild narusza
$g_{00}\,g_{rr} = -1$ od rz\k{e}du $O(U^2)$
(jawne rozpisanie:
$g^{\mathrm{Schw}}_{00}\,g^{\mathrm{Schw}}_{rr}
= -1 + U^2/2 + O(U^4)$ we wsp.\ izotropowych).

\emph{Punkt 3} --- bezpo\'sredni rachunek:
$e^{-2U}\cdot e^{+2U} = 1$ (dok\l{}adnie),
natomiast metryka Schwarzschilda w~wsp\'o\l{}rz\k{e}dnych
izotropowych: $(1-U/2)^2(1+U/2)^{-2} \cdot (1+U/2)^4
= (1-U/2)^2(1+U/2)^2 = (1-U^2/4)^2 \neq 1$.

\emph{Punkt 4} --- estymata porz\k{a}dkowa: $\delta g_{00}
\sim U^3$, wi\k{e}c $\delta f_{\mathrm{QNM}}/f \sim U^3$.
Dla gwiazdy neutronowej ($U \approx 0.2$):
$\delta f/f \sim 0.008$, osi\k{a}galne przez detektory
3.~generacji.
\end{proof}

\begin{remark}[Interpretacja: nie pora\.zka, lecz predykcja]%
\label{rem:emergence-not-failure}
Fakt, \.ze TGP nie odtwarza \emph{dok\l{}adnie} r\'owna\'n
Einsteina w~og\'olnym przypadku, jest cech\k{a}
po\.z\k{a}dan\k{a}, nie wad\k{a}. Analogia: mechanika statystyczna
nie odtwarza dok\l{}adnie termodynamiki (s\k{a} fluktuacje,
efekty sko\'nczonego rozmiaru, przej\'scia fazowe) --- ale
jest od niej \emph{g\l{}\k{e}bsza}. Podobnie TGP:
\begin{itemize}[nosep]
  \item zgadza si\k{e} z~GR wsz\k{e}dzie, gdzie GR jest
    potwierdzone (do 2PN);
  \item r\'o\.zni si\k{e} w~re\.zimie silnego pola
    ($U^3$, mody tensorowe) --- i~te r\'o\.znice stanowi\k{a}
    \textbf{falsyfikowalny program obserwacyjny} (por.\ \S\ref{sec:predykcje});
  \item predykcja po\'slizgu grawitacyjnego:
    $\eta_{\mathrm{slip}} = e^{-2U} \approx 1 - 2U$,
    mierzalny przez soczewkowanie (Euclid) w~re\.zimie
    kosmologicznym ($\Delta\eta \sim 10^{-5}$) oraz
    ringdown w~silnym polu ($\Delta\eta \sim 0.4$
    dla gwiazd neutronowych).
\end{itemize}
\end{remark}

\begin{proposition}[Dekompozycja $3{+}1$ metryki TGP]\label{prop:ADM-TGP}
Dla og\'olnego pola $\Phi(t,\vec{x})$ metryka
eksponencjalna~\eqref{eq:metric-factor2} zapisana w~formie
ADM ($ds^2 = -N^2 c_0^2\,dt^2 + \gamma_{ij}(dx^i + N^i dt)(dx^j + N^j dt)$)
przyjmuje posta\'c:
\begin{equation}\label{eq:ADM-TGP}
  N = e^{-U}, \qquad N^i = 0, \qquad
  \gamma_{ij} = e^{+2U}\,\delta_{ij},
  \qquad U \equiv \frac{\delta\Phi}{\PhiZero},
\end{equation}
tj.\ shift zerowy, a~lapse i~metryka $3$-przestrzenna
zale\.z\k{a} wy\l{}\k{a}cznie od jednej funkcji~$U$.
Z~tego wynikaj\k{a} trzy strukturalne w\l{}asno\'sci:
\begin{enumerate}[nosep]
  \item \textbf{Redukcja stopni swobody.}
    Sposr\'od $10$ sk\l{}adowych $g_{\mu\nu}$:
    $3$~(shift) s\k{a} ustalone na zero,
    $1$~(lapse) i~$6$~(metryka $3$-przestrzenna) zale\.z\k{a}
    od jednej funkcji~$U$.
    Pozosta\l{}y $10 - 3 - 1 = 6$ warunki (r\'ownania Einsteina
    $G_{ij} = \kappa T_{ij}$) s\k{a}
    \textbf{nadokreslone} dla jednej funkcji~$U$ ---
    ale nie s\k{a} niezale\.zne:

  \item \textbf{To\.zsamo\'sci strukturalne.}
    Z~postaci $\gamma_{ij} = e^{2U}\delta_{ij}$:
    \begin{itemize}[nosep]
      \item $R^{(3)}_{ij} = -\nabla_i\nabla_j U
        - \delta_{ij}\nabla^2 U + \nabla_i U\,\nabla_j U
        - \delta_{ij}(\nabla U)^2$,
      \item Ricci skalarny: $R^{(3)} = -4e^{-2U}
        [\nabla^2 U + (\nabla U)^2]$.
    \end{itemize}
    Wszystkie $6$~sk\l{}adowych $R^{(3)}_{ij}$ wyra\.z\k{a}
    si\k{e} przez $U$ i~jego pochodne --- $5$ z~$6$
    to zwi\k{a}zki algebraiczne, a~niezale\.zny jest
    tylko $1$~(skalarny Poisson/fala).

  \item \textbf{Wi\k{e}z hamiltonowski i~p\k{e}dowy.}
    \begin{align}
      R^{(3)} + K^2 - K_{ij}K^{ij} &= 2\kappa\,\rho_E,
        \label{eq:ADM-ham}\\
      D_j(K^j{}_i - \delta^j{}_i K) &= \kappa\,J_i,
        \label{eq:ADM-mom}
    \end{align}
    gdzie $K_{ij} = e^{2U}\dot{U}\,\delta_{ij}/N$
    jest krzywizn\k{a} zewn\k{e}trzn\k{a}. Wi\k{e}z
    p\k{e}dowy~\eqref{eq:ADM-mom} redukuje si\k{e} do
    $\partial_i(e^U\dot{U}) = -\kappa J_i/2$,
    a~wi\k{e}z hamiltonowski~\eqref{eq:ADM-ham}
    daje r\'ownanie Poissona na~$U$
    (por.\ tw.~\ref{thm:newton}).
    Oba s\k{a} \textbf{spe\l{}nione automatycznie}
    przez rozwi\k{a}zania r\'ownania pola TGP
    plus propagacj\k{e} Bianchiego
    (rozszerzenie argumentu z~tw.~\ref{thm:einstein-emergence}).
\end{enumerate}
\end{proposition}

\begin{proof}[Szkic dowodu]
Postaci $N$, $N^i$, $\gamma_{ij}$ w~\eqref{eq:ADM-TGP}
wynikaj\k{a} bezpo\'srednio z~metryki~\eqref{eq:metric-factor2}
w~uk\l{}adzie wsp\'o\l{}poruszaj\k{a}cym si\k{e}.
Shift $N^i = 0$: z~izotropii metryki TGP
(brak preferowanych kierunk\'ow przestrzennych
w~metryce efektywnej); r\'ownowa\.znie:
metryka zale\.zy tylko od skalara~$\Phi$,
kt\'ory nie definiuje wektora.

\smallskip\noindent
\emph{Punkt 1:} rachunkiem bezpo\'srednim.
\emph{Punkt 2:} standardowa geometria konforemna
$\gamma_{ij} = e^{2U}\delta_{ij}$ w~3D.
\emph{Punkt 3:} to\.zsamo\'s\'c Bianchiego
$\nabla^\mu G_{\mu\nu} = 0$ propaguje wi\k{e}zy:
je\'sli $G_{0\mu} = \kappa T_{0\mu}$ spe\l{}nione
na pewnej powierzchni $\Sigma_0$
i~r\'ownanie pola TGP (r\'ownowa\.zne dynamicznym
r\'ownaniom $G_{ij} = \kappa T_{ij}$
do rz\k{e}du $O(U^2)$, stw.~\ref{prop:emergence-hierarchy})
obowi\k{a}zuje wsz\k{e}dzie, to
$G_{0\mu} = \kappa T_{0\mu}$ jest zachowywane
na wszystkich p\'o\'zniejszych $\Sigma_t$.

Precyzja tego argumentu odzwierciedla hierarchi\k{e}
ze stw.~\ref{prop:emergence-hierarchy}:
\begin{itemize}[nosep]
  \item do rz\k{e}du $O(U^2)$: dok\l{}adna emergencja,
  \item od $O(U^3)$: pojawia si\k{e} reszta
    $\delta G_{\mu\nu} = G_{\mu\nu}^{\mathrm{TGP}}
    - G_{\mu\nu}^{\mathrm{GR}}$,
    kt\'ora spe\l{}nia r\'ownanie liniowe
    z~\emph{tensorowym \'zr\'od\l{}em} $\propto \partial U
    \cdot \partial U \cdot \partial U$ (wy\l{}\k{a}cznie
    kubiczne i~wy\.zsze). Reszta ta jest predykcj\k{a} TGP.
\end{itemize}
\end{proof}

\begin{proposition}[Propagacja wi\k{e}z\'ow ADM w~TGP]%
\label{prop:ADM-constraint-propagation}
Niech $(\Phi, \dot\Phi)$ b\k{e}d\k{a} danymi Cauchy'ego na $\Sigma_0$,
spe\l{}niaj\k{a}cymi r\'ownanie pola TGP~\eqref{eq:covariant-field}
oraz wi\k{e}zy~\eqref{eq:ADM-ham}--\eqref{eq:ADM-mom}
na $\Sigma_0$. W\'owczas wi\k{e}zy s\k{a} spe\l{}nione
na wszystkich $\Sigma_t$ w~dziedzinie istnienia rozwi\k{a}zania.
\end{proposition}

\begin{proof}
Dowod opiera si\k{e} na trzech faktach:

\smallskip\noindent
\textbf{(1) Konforemna redukcja.}
Z~$\gamma_{ij} = e^{2U}\delta_{ij}$ ($U = \delta\Phi/\PhiZero$),
metryka $3$-przestrzenna zale\.zy od jednej funkcji~$U$.
Wi\k{e}z p\k{e}dowy~\eqref{eq:ADM-mom} redukuje si\k{e} do
$\partial_i(e^U\dot{U}) = -\kappa J_i/2$, tj.\ trzech
sk\l{}adowych wyr\'o\.znionych przez \emph{gradient} --- a~wi\k{e}c
do jednego r\'ownania skalarnego (po wzi\k{e}ciu dywergencji).

\smallskip\noindent
\textbf{(2) To\.zsamo\'s\'c kontraktowa Bianchiego.}
Zdefiniujmy wielko\'s\'c naruszen\k{a}
$\mathcal{C}_\mu \equiv G_{0\mu} - \kappa\,T_{0\mu}$
(tu $G_{\mu\nu}$ i~$T_{\mu\nu}$ konstruowane
z~efektywnej metryki TGP). To\.zsamo\'s\'c $\nabla^\mu G_{\mu\nu} = 0$
daje:
\begin{equation}\label{eq:constraint-prop-eq}
  \partial_t \mathcal{C}_0
  + \partial_i \mathcal{C}^i
  + \Gamma\text{-cz\l{}ony} \cdot \mathcal{C}
  = \nabla^\mu T_{0\mu} \cdot \kappa.
\end{equation}
Prawa strona znika na mocy zachowania $T_{\mu\nu}$
(wyprowadzonego z~r\'ownania pola TGP
w~tw.~\ref{thm:einstein-emergence}).
R\'ownanie~\eqref{eq:constraint-prop-eq} jest \textbf{liniowym}
uk\l{}adem symetrycznym hiperbolicznym na $\mathcal{C}_\mu$.
Przy danych pocz\k{a}tkowych $\mathcal{C}_\mu|_{\Sigma_0} = 0$
jednoznaczno\'s\'c (twierdzenie energetyczne) daje
$\mathcal{C}_\mu = 0$ wsz\k{e}dzie.

\smallskip\noindent
\textbf{(3) Precyzja hierarchiczna.}
Argument (2) jest \emph{dok\l{}adny} do rz\k{e}du, w~kt\'orym
r\'ownanie pola TGP odtwarza r\'ownania Einsteina
(stw.~\ref{prop:emergence-hierarchy}):
\begin{itemize}[nosep]
  \item do $O(U^2)$: dok\l{}adna propagacja (brak reszty);
  \item od $O(U^3)$: reszta
    $\delta\mathcal{C}_\mu \propto (\nabla U)^3$ jest
    \textbf{perturbacyjnie ma\l{}a} i~spe\l{}nia
    r\'ownanie liniowe z~wymuszon\k{a} prawym cz\l{}onem
    $\sim U^3$, nie rosn\k{a}cym w~czasie
    (stabilno\'s\'c: tw.~\ref{thm:nonlinear-orbital}).
\end{itemize}
Zatem wi\k{e}zy s\k{a} zachowane z~dok\l{}adno\'sci\k{a}
co najmniej $O(U^3)$, a~w~praktyce (astrofizyczne potencja\l{}y
$U \sim 10^{-6}$--$10^{-1}$) odchylenie jest $\lesssim 10^{-18}$.
\end{proof}

\begin{remark}[Mody tensorowe i~konieczno\'s\'c disformalnego rozszerzenia]%
\label{rem:ADM-tensor-modes}
Dekompozycja SVT (skalarne / wektorowe / tensorowe) perturbacji
$\delta\gamma_{ij}$ wok\'o\l{} t\l{}a jednorodnego:
\begin{equation}
  \delta\gamma_{ij} = 2e^{2U}[\Psi\,\delta_{ij}
  + D_{ij}E + D_{(i}V_{j)} + h_{ij}^{\mathrm{TT}}].
\end{equation}
W~sektorze konforemnym TGP ($B = 0$):
$\delta\gamma_{ij} = 2\delta U\,e^{2U}\delta_{ij}$,
wi\k{e}c $h_{ij}^{\mathrm{TT}} = 0$ ---
\textbf{brak niezale\.znych mod\'ow tensorowych}.
Wzbudzenie fal grawitacyjnych wymaga
rozszerzenia dysformalnego
(hip.~\ref{hyp:disformal}),
kt\'ore wprowadza sprz\k{e}\.zenie
$B(\Phi)\,\partial_\mu\Phi\,\partial_\nu\Phi$
i~generuje $h_{ij}^{\mathrm{TT}} \neq 0$
(stw.~\ref{prop:disformal-polarization}).
\end{remark}

\begin{definition}[Kosmologiczna skala czasowa $\tau_0$]%
\label{def:tau0}
\textbf{Kosmologiczna skala czasowa} $\tau_0$ jest zdefiniowana jako:
\begin{equation}\label{eq:tau0-def}
  \tau_0 \;\equiv\; \frac{1}{c_0}\sqrt{\frac{\PhiZero}{m_{\mathrm{sp}}^2}}
       \;=\; \frac{1}{c_0}\sqrt{\frac{\PhiZero}{3\gamma - 2\beta}}\,,
\end{equation}
gdzie $m_{\mathrm{sp}}^2 = 3\gamma - 2\beta$
(stw.~\ref{prop:vacuum-stability}) jest przestrzenn\k{a} mas\k{a}
efektywn\k{a}. Dla $\beta = \gamma$:
$\tau_0 = \sqrt{\PhiZero}/(c_0\sqrt{\gamma})$.

Skala $\tau_0$ okre\'sla naturalny czas relaksacji pola $\Phi$
do stanu r\'ownowagi $\PhiZero$ i~\l{}\k{a}czy sektor
przestrzenny (warunki $m_{\mathrm{sp}}$) z~dynamik\k{a}
kosmologiczn\k{a}.
\end{definition}

\begin{proposition}[R\'ownanie kosmologiczne z~dzia\l{}ania TGP]%
\label{prop:FRW-derivation}
W~przestrzeni FRW ze~skalowaniem $a(t)$, dla jednorodnego pola
$\Phi(t)$ ($\nabla\Phi = 0$), r\'ownanie ruchu wynikaj\k{a}ce
z~dzia\l{}ania~\eqref{eq:action} z~elementem obj\k{e}to\'sciowym
$\sqrt{-g_{\mathrm{eff}}} = a^3\,c_0\,\psi$ (gdzie $\psi = \Phi/\PhiZero$)%
\footnote{Poprawiono z~wcze\'sniejszego $a^3\psi^4$;
  zob.\ prop.~\ref{prop:kappa-corrected}.}
i~p\l{}askim cz\l{}onem kinetycznym $\frac{1}{\kappa}(-\frac{1}{2}\dot\psi^2/c_0^2 - U)$
ma posta\'c:
\begin{equation}\label{eq:Phi-cosmo-exact}
  \boxed{
    \ddot\psi + 3H\dot\psi + \frac{3\dot\psi^2}{\psi}
    \;=\; c_0^2\,\frac{V(\psi) + \psi\,V'(\psi)}{\psi^6}
    + S_{\mathrm{src}}(\psi,\rho),
  }
\end{equation}
gdzie $V(\psi) = \frac{\beta}{3}\psi^3 - \frac{\gamma}{4}\psi^4$
i~$S_{\mathrm{src}}$ zawiera sprz\k{e}\.zenie \'zr\'od\l{}owe.
Pe\l{}ne wyprowadzenie z~poprawionej miary $\sqrt{-g}=c_0\psi$
w~sekcji~\ref{ssec:unified-action},
r.~\eqref{eq:psi-eq-unified}.
\textbf{Uwaga historyczna:} wcze\'sniejsza wersja
(z~miar\k{a} $\psi^4$) dawa\l{}a
$2\dot\psi^2/\psi$ zamiast $3\dot\psi^2/\psi$
i~uproszczone prawa strony.
\end{proposition}

\begin{proof}
Z~poprawionym elementem obj\k{e}to\'sciowym
$\sqrt{-g_{\mathrm{eff}}} = a^3 c_0\,\psi$
i~sprz\k{e}\.zeniem kinetycznym $K(\psi) = \psi^4$,
zredukowane dzia\l{}anie dla jednorodnego $\psi(t)$ w~FRW
(pe\l{}ne wyprowadzenie w~sek.~\ref{ssec:unified-action},
r.~\eqref{eq:S-FRW-reduced}) wynosi:
\[
  S_{\mathrm{FRW}} = \frac{V_0}{\kappa}\!\int\!dt\;a^3\biggl[
    -\frac{\psi^6}{2c_0^2}\,\dot\psi^2
    - \psi\,V(\psi)
    - \frac{q}{\PhiZero}\,\psi^2\,\rho
  \biggr],
\]
gdzie $V(\psi) = \frac{\beta}{3}\psi^3 - \frac{\gamma}{4}\psi^4$
(r.~\eqref{eq:U-potential}).
Lagran\.zjan $\mathcal{L} = a^3[\ldots]$; r\'ownanie EL
$\partial\mathcal{L}/\partial\psi
- \frac{d}{dt}(\partial\mathcal{L}/\partial\dot\psi) = 0$:

\emph{Pochodna po $\dot\psi$}
(czynnik $1/\kappa$ wsp\'olny):
$\frac{\partial\mathcal{L}}{\partial\dot\psi}
= -\frac{a^3\psi^6\dot\psi}{c_0^2}$.
Pochodna czasowa:
$\frac{d}{dt}(\cdots)
= -\frac{a^3\psi^6}{c_0^2}
\bigl[\ddot\psi + 3H\dot\psi + 6\dot\psi^2/\psi\bigr]$.

\emph{Pochodna po $\psi$}:
$\frac{\partial\mathcal{L}}{\partial\psi}
= a^3\bigl[
  -\frac{3\psi^5\dot\psi^2}{c_0^2}
  - V - \psi V'
  - \frac{2q}{\PhiZero}\,\psi\,\rho
\bigr]$.

Dziel\k{a}c EL przez $a^3\psi^5/c_0^2$ i~upraszczaj\k{a}c
($6\dot\psi^2/\psi - 3\dot\psi^2/\psi = 3\dot\psi^2/\psi$):
\[
  \psi\bigl(\ddot\psi + 3H\dot\psi\bigr) + 3\dot\psi^2
  = c_0^2\,\frac{V + \psi V'}{\psi^5}
  + \frac{2q\,c_0^2}{\PhiZero}\,\frac{\rho}{\psi^4}.
\]
Po podzieleniu przez $\psi$
otrzymujemy~\eqref{eq:Phi-cosmo-exact}.
Cz\l{}on $3\dot\psi^2/\psi$ jest odpowiednikiem
czasowym sprz\k{e}\.zenia kinetycznego $K(\psi)=\psi^4$
z~poprawion\k{a} miar\k{a} $c_0\psi$.
\end{proof}

\begin{remark}[Potencja\l{} kosmologiczny a~statyczny]\label{rem:cosmo-vs-static}
Efektywny potencja\l{} w~r\'ownaniu kosmologicznym
$W_{\mathrm{eff}}(\psi) \equiv (V + \psi V')/\psi^6$
\textbf{r\'o\.zni si\k{e}} od statycznego
$V(\psi)/\PhiZero = \beta\psi^2 - \gamma\psi^3$.
R\'o\.znica pochodzi od czynnika $\psi V'$ oraz dzielnika $\psi^6$
(konsekwencja poprawionego sprz\k{e}\.zenia $K(\psi)=\psi^4$
z~miar\k{a} $c_0\psi$).
\textbf{Uwaga:} wcze\'sniejsza wersja z~$4U/\psi + U'$
wynika\l{}a ze starego elementu obj\k{e}to\'sciowego $\psi^4$;
pe\l{}na korekta w~sek.~\ref{ssec:unified-action}.
R\'o\.znica jest konsekwencj\k{a} sprz\k{e}\.zenia pola $\Phi$
z~elementem obj\k{e}to\'sciowym $\psi^4$: rozszerzaj\k{a}cy si\k{e}
Wszech\'swiat \emph{rozrzedza} wygenerowana przestrze\'n, modyfikuj\k{a}c
efektywny potencja\l{}.

\textbf{Uwaga}: pe\l{}ne wsp\'o\l{}czynniki nieliniowe
$(7\beta/3, 2\gamma)$ r\'o\.zni\k{a} si\k{e} od statycznych
$(\beta, \gamma)$ o~czynniki $7/3$ i~$2$.
Jednake\.z \textbf{linearyzacja wok\'o\l{} $\psi = 1$}
daje identyczn\k{a} struktur\k{e}: patrz
stw.~\ref{prop:exact-to-linearized}.
\end{remark}

\begin{proposition}[Zwi\k{a}zek $W(\psi)$ i~$U(\psi)$: zamkni\k{e}cie GAP-3]%
\label{prop:W-from-U}
Potencja\l{} kosmologiczny $W(\psi)$ jest \textbf{jednoznacznie okre\'slony}
przez potencja\l{} statyczny $U(\psi)$ i~miar\k{e} $\psi^4$ dzia\l{}ania TGP:
\begin{equation}\label{eq:W-from-U}
  \boxed{
    W(\psi) = \frac{1}{\psi^4}\,\frac{d}{d\psi}\bigl[\psi^4\,U(\psi)\bigr]
    = \frac{4\,U(\psi)}{\psi} + U'(\psi).
  }
\end{equation}
Jawnie: $W(\psi) = \frac{7\beta}{3}\psi^2 - 2\gamma\psi^3$.
Stosunek warto\'sci w~pr\'o\.zni: $W(1)/U(1) = (7\beta/3 - 2\gamma)/(\beta/3 - \gamma/4)$.
Dla $\beta = \gamma$: $W(1) = \gamma/3$, $U(1) = \gamma/12$, wi\k{e}c $W(1) = 4\,U(1)$.
\end{proposition}

\begin{proof}
Lagran\.zjan FRW: $\mathcal{L} \propto a^3\psi^4[-\frac{1}{2}\dot\psi^2/c_0^2 - U(\psi)]$.
R\'ownanie Eulera--Lagrange'a $\partial\mathcal{L}/\partial\psi
- \frac{d}{dt}(\partial\mathcal{L}/\partial\dot\psi) = 0$
generuje pochodn\k{a} po $\psi$ z~iloczynu $\psi^4 \cdot U(\psi)$.
Dziel\k{a}c przez $\psi^4$:
\[
  \frac{1}{\psi^4}\,\frac{\partial}{\partial\psi}\bigl[\psi^4 U(\psi)\bigr]
  = \frac{4\psi^3 U + \psi^4 U'}{\psi^4}
  = \frac{4U}{\psi} + U'
  = W(\psi).
\]
Czynnik $4U/\psi$ pochodzi z~\textbf{miary generowanej przestrzeni} $\psi^4$:
rozszerzaj\k{a}cy si\k{e} Wszech\'swiat zmienia obj\k{e}to\'s\'c
generowanej przestrzeni, co modyfikuje efektywny potencja\l{}.
Czynnik $U'$ pochodzi ze standardowej wariacji potencja\l{}u.

\medskip\noindent\textbf{Weryfikacja:}
$4U/\psi + U' = 4(\tfrac{\beta}{3}\psi^2 - \tfrac{\gamma}{4}\psi^3)
+ (\beta\psi^2 - \gamma\psi^3)
= \tfrac{4\beta}{3}\psi^2 - \gamma\psi^3 + \beta\psi^2 - \gamma\psi^3
= \tfrac{7\beta}{3}\psi^2 - 2\gamma\psi^3 = W(\psi)$. \checkmark

\medskip\noindent\textbf{Interpretacja fizyczna:}
R\'o\.znica mi\k{e}dzy $W$ i~$U$ \textbf{nie jest b\l{}\k{e}dem ani
niesp\'ojno\'sci\k{a}} --- jest konsekwencj\k{a} fundamentalnej cechy TGP:
pole $\Phi$ nie jest polem \emph{na} przestrzeni, lecz polem
\emph{konstytuuj\k{a}cym} przestrze\'n. Miara $\psi^4$ koduje ten fakt.
W~teoriach skalarno-tensorowych (Brans--Dicke, Horndeski)
analogiczna r\'o\.znica wyst\k{e}puje mi\k{e}dzy ramk\k{a}
Jordana ($\psi^2 R$) a~Einsteina ($R$).
\end{proof}

\medskip
\begin{proposition}[\'Scis\l{}e przej\'scie:
  r\'ownanie dok\l{}adne $\to$ zlinearyzowane]%
\label{prop:exact-to-linearized}
W~przybli\.zeniu powolnego staczania
($\dot\psi^2/\psi \ll H\dot\psi$) i~linearyzacji
wok\'o\l{} $\psi = 1$, r\'ownanie
dok\l{}adne~\eqref{eq:Phi-cosmo-exact} redukuje si\k{e} do
r\'ownania liniowego na perturbacj\k{e} $\delta \equiv \psi - 1$:
\begin{equation}\label{eq:Phi-cosmo-linearized}
  \boxed{
    \ddot\delta + 3H\dot\delta
    + \frac{4c_0^2\gamma}{3}\,\delta
    = c_0^2\,\frac{\gamma}{3}
    + S_{\mathrm{src}}\,.
  }
\end{equation}
\end{proposition}

\begin{proof}
Startujemy od~\eqref{eq:Phi-cosmo-exact} w~przybli\.zeniu
powolnego staczania (pomijaj\k{a}c $2\dot\psi^2/\psi$):
\[
  \ddot\psi + 3H\dot\psi
  = c_0^2\,W(\psi) + S_{\mathrm{src}},
  \qquad
  W(\psi) = \tfrac{7\beta}{3}\psi^2 - 2\gamma\psi^3.
\]
Podstawiaj\k{a}c $\psi = 1 + \delta$, $|\delta| \ll 1$:
\begin{align}
  W(1+\delta)
  &= \tfrac{7\beta}{3}(1+2\delta+\delta^2)
     - 2\gamma(1+3\delta+3\delta^2+\delta^3) \nonumber\\
  &= \underbrace{\bigl(\tfrac{7\beta}{3}-2\gamma\bigr)}_{W(1)}
     + \underbrace{\bigl(\tfrac{14\beta}{3}-6\gamma\bigr)}_{W'(1)}\,\delta
     + \mathcal{O}(\delta^2). \label{eq:W-expansion}
\end{align}
Dla $\beta = \gamma$:
\[
  W(1) = \frac{\gamma}{3}, \qquad
  W'(1) = \frac{14\gamma}{3} - 6\gamma = -\frac{4\gamma}{3}.
\]
Poniewa\.z $\ddot\psi = \ddot\delta$, $\dot\psi = \dot\delta$:
\[
  \ddot\delta + 3H\dot\delta
  = c_0^2\Bigl[\frac{\gamma}{3} - \frac{4\gamma}{3}\delta\Bigr]
  + S_{\mathrm{src}},
\]
co daje~\eqref{eq:Phi-cosmo-linearized}.
\end{proof}

\begin{remark}[Status $\psi = 1$: stan zamro\.zony, nie minimum]%
\label{rem:psi1-status}
Stan $\psi = 1$ ($\Phi = \PhiZero$) jest:
\begin{enumerate}[nosep]
  \item \textbf{punktem krytycznym} potencja\l{}u $U(\psi)$:
    $U'(1) = \beta - \gamma = 0$ przy $\beta = \gamma$;
  \item \textbf{maksimum} $U(\psi)$:
    $U''(1) = 2\beta - 3\gamma = -\gamma < 0$;
  \item \textbf{nie r\'ownowag\k{a}} r\'ownania kosmologicznego:
    $W(1) = \gamma/3 \neq 0$, wi\k{e}c przy $\dot\psi = 0$
    istnieje resztkowe \'zr\'od\l{}o $\ddot\psi = c_0^2\gamma/3$;
  \item \textbf{stanem zamro\.zonym w~erze materii}: tarcie Hubble'a ($3H\dot\psi$)
    t\l{}umi ruch pola. Efektywna cz\k{e}sto\'s\'c pola wynosi
    $\omega_{\mathrm{cosmo}} = \sqrt{4c_0^2\gamma/3}
    = \sqrt{4\PhiZero/3}\;H_0 \approx 5{,}7\,H_0$.
    Dla $H \gg \omega_{\mathrm{cosmo}}$ (era materii, $z \gtrsim 3$)
    uk\l{}ad jest \emph{nadt\l{}umiony} i~pole zamro\.zone.
    Dla $z \lesssim 3$ tarcie s\l{}abnie i~pole zaczyna
    ewoluowa\'c w~stron\k{e} atraktora $\psi_{\mathrm{eq}} = 7/6$
    (patrz uwaga~\ref{rem:cosmo-attractor});
  \item \'zr\'od\l{}em \textbf{ciemnej energii}: resztkowe
    $W(1) = \gamma/3$ dzia\l{}a jak sta\l{}a kosmologiczna
    $\Lambda_{\mathrm{eff}} \approx \gamma/12$ (por.\
    stw.~\ref{prop:Lambda-eff}).
\end{enumerate}
Obraz: pole $\psi$ siedzi na szczycie \emph{bardzo p\l{}askiego wzg\'orza}
potencja\l{}u $U(\psi)$, zamro\.zone przez tarcie Hubble'a
w~erze materii ($z \gtrsim 3$). W~erze ciemnej energii
($z \lesssim 3$) tarcie s\l{}abnie i~pole zaczyna oscylowa\'c
w~stron\k{e} atraktora $\psi_{\mathrm{eq}} = 7\beta/(6\gamma) = 7/6$.
Resztka $U(1) = \gamma/12 > 0$
\textbf{nap\k{e}dza przyspieszony rozw\'oj Wszech\'swiata}.
\end{remark}

\begin{remark}[Zwi\k{a}zek ze skal\k{a} $\tau_0$]\label{rem:tau0-cosmo}
Kosmologiczna masa efektywna odczytana
z~r\'ownania~\eqref{eq:Phi-cosmo-linearized}:
\begin{equation}\label{eq:m-cosmo-exact}
  m_{\mathrm{cosmo}}^2 = \frac{4c_0^2\gamma}{3}\,,
  \qquad
  \tau_{\mathrm{cosmo}} = \frac{1}{m_{\mathrm{cosmo}}}
  = \sqrt{\frac{3}{4c_0^2\gamma}}\,.
\end{equation}
Skala $\tau_{\mathrm{cosmo}}$ r\'o\.zni si\k{e}
od $\tau_0 = \sqrt{\PhiZero/(c_0^2\gamma)}$
(def.~\ref{def:tau0}), kt\'ora pochodzi od masy przestrzennej
$m_{\mathrm{sp}}^2 = \gamma$.
Stosunek: $\tau_0/\tau_{\mathrm{cosmo}}
= \sqrt{4\PhiZero/3} \approx 5{,}7$ dla $\PhiZero \approx 25$.
W~jednostkach $H_0^{-1}$:
$\omega_{\mathrm{cosmo}} = m_{\mathrm{cosmo}}\,c_0
= \sqrt{4\PhiZero/3}\;H_0 \approx 5{,}7\,H_0$,
wi\k{e}c kosmologiczna masa pola jest \textbf{wi\k{e}ksza}
od skali Hubble'a, co prowadzi do rezimu niedot\l{}umionego
w~p\'o\'znym Wszech\'swiecie (patrz uwaga~\ref{rem:cosmo-attractor}).
\end{remark}

\begin{remark}[Nota wymiarowa: $m_{\mathrm{cosmo}}^2$
  i~wsp\'o\l{}czynnik $c_0^2\gamma$]\label{rem:dim-mcosmo}
R\'ownanie~\eqref{eq:Phi-cosmo-linearized} pochodzi
od \textbf{kowariantnego} r\'ownania pola~\eqref{eq:Phi-cosmo-exact},
kt\'orego cz\l{}ony maj\k{a} wymiar $[\mathrm{L}^{-2}]$
(d'alembertian dzia\l{}aj\k{a}cy na bezwymiarowe $\psi$).
Rozwijaj\k{a}c $\Box\psi$ w~metryce FRW:
\[
  \Box\psi
  = -\frac{1}{c_0^2}\ddot\psi - \frac{3H}{c_0^2}\dot\psi
  + \ldots
  \qquad\bigl[\mathrm{L}^{-2}\bigr],
\]
a~nast\k{e}pnie \textbf{mno\.z\k{a}c ca\l{}e r\'ownanie przez $c_0^2$}
($[\mathrm{L}^{2}\mathrm{T}^{-2}]$), otrzymujemy posta\'c
z~$\ddot\psi$ i~wsp\'o\l{}czynnikami w~$[\mathrm{T}^{-2}]$.
Sp\'ojno\'s\'c wymiarowa jest trywialna:
\[
  c_0^2\gamma
  = [\mathrm{L}^{2}\mathrm{T}^{-2}] \cdot [\mathrm{L}^{-2}]
  = [\mathrm{T}^{-2}],
\]
wi\k{e}c $m_{\mathrm{cosmo}}^2 = 4c_0^2\gamma/3$
ma poprawny wymiar $[\mathrm{T}^{-2}]$
i~jest cz\k{e}stotliwo\'sci\k{a} w\l{}asn\k{a} (z~tarciowym
t\l{}umieniem Hubble'a) perturbacji $\delta = \psi - 1$.

\textbf{Uwaga historyczna}: we wcze\'sniejszych wersjach
cz\l{}on kinetyczny w~dzia\l{}aniu~\eqref{eq:action}
mia\l{} czynnik $1/(2\kappa)$, ale potencja\l{} $U(\Phi)$
nie mia\l{} $1/\kappa$. Prowadzi\l{}o to do wymiarowo
niespójnego wyst\k{e}powania $\kappa$ w~r\'ownaniach
kosmologicznych. Poprawna postać (v6) umieszcza
$1/\kappa$ przed ca\l{}ym sektorem grawitacyjnym,
co sprawia, \.ze $\kappa$ \textbf{skraca si\k{e}}
ze wszystkich r\'owna\'n pola.
\end{remark}

\begin{remark}[Kosmologiczny atraktor $\psi_{\mathrm{eq}} = 7/6$]%
\label{rem:cosmo-attractor}
R\'ownanie kosmologiczne~\eqref{eq:Phi-cosmo-exact}
z~potencja\l{}em $W(\psi)$ ma zerowy
w~$\psi_{\mathrm{eq}} = 7\beta/(6\gamma) = 7/6$ (dla $\beta = \gamma$).
Dynamika pola zale\.zy od epoki:

\begin{enumerate}[nosep]
  \item \textbf{Era materii} ($z \gtrsim 3$):
    $H \gg \omega_{\mathrm{cosmo}}$, wi\k{e}c tarcie Hubble'a
    jest nadt\l{}umi\k{a}ce: $\psi \approx 1$ zamro\.zone.
  \item \textbf{Przej\'scie} ($z \approx 2{,}5$):
    $3H/2 \approx \omega_{\mathrm{cosmo}}$; pole zaczyna ewoluowa\'c.
  \item \textbf{Era ciemnej energii} ($z \lesssim 2$):
    uk\l{}ad jest \textbf{niedot\l{}umiony}
    ($\omega_{\mathrm{cosmo}} \approx 5{,}7\,H_0 > 3H_0/2$).
    Pole oscyluje wok\'o\l{} $\psi_{\mathrm{eq}}$
    z~t\l{}umieniem $e^{-3H_0 t/2}$.
\end{enumerate}

\noindent\textbf{Numerycznie} (skrypt
\texttt{cosmological\_evolution.py}, 12/12 PASS):
$\psi(z{=}0) \approx 1{,}15$; pole cz\k{e}\'sciowo
zrelaksowane w~stron\k{e} $7/6 \approx 1{,}167$.

\medskip\noindent\textbf{Konsekwencje testowalne:}
\begin{itemize}[nosep]
  \item $\Delta G/G$ (BBN $\to$ dzi\'s) $= 0$ dla naturalnych
    warunk\'ow pocz\k{a}tkowych $\psi_{\mathrm{ini}} = 7/6$
    (uwaga~\ref{rem:BBN-resolution});
  \item $|\dot G/G|_{z=0} \approx 2 \times 10^{-13}\,\mathrm{yr}^{-1}$
    --- bliskie czu\l{}o\'sci LLR
    (por.~\textbf{F2} w~sekcji~\ref{sec:predykcje});
  \item $|1+w_{\mathrm{DE}}| \sim 10^{-6}$
    --- pozorna sta\l{}a kosmologiczna, ale z~mikroskopijn\k{a}
    dynamik\k{a} pola $\psi$.
\end{itemize}

\end{remark}

\begin{proposition}[Rekalibracja $\PhiZero$ przy atraktorze]\label{prop:recalibration}
Laboratoryjne warto\'sci sta\l{}ych ($G_0$, $c_0$, $\hbar_0$)
odpowiadaj\k{a} $\psi_{\mathrm{today}} \approx 1{,}15$, nie
$\psi = 1$. Prze\-defi\-nio\-wa\-nie punktu fiducjalnego na
$\psi_{\mathrm{eq}} = 7\beta/(6\gamma) = 7/6$ daje:
\begin{equation}\label{eq:Phi0-recal}
  \PhiZero^{(\mathrm{recal})}
  = \PhiZero \cdot \psi_{\mathrm{eq}}
  = 36\,\Omega_\Lambda \cdot \tfrac{7}{6}
  \approx 28{,}8.
\end{equation}
Rekalibracja \textbf{nie zmienia} struktury teorii (r\'ownania pola,
rezimy, granica newtonowska pozosta\-j\k{a} identyczne);
przesuwa jedynie konwencj\k{e} $\psi = 1$ z~warto\'sci pocz\k{a}tkowej
na po\l{}o\.zenie r\'ownowagi.
\end{proposition}

\begin{proof}
Sta\l{}e dynamiczne (aksjomat~\ref{ax:c}--\ref{ax:G}):
$c(\Phi) = c_0\sqrt{\PhiZero/\Phi}$,
$G(\Phi) = G_0\,\PhiZero/\Phi$.
Zmierzono $c_0$, $G_0$ przy bie\.z\k{a}cym $\Phi_{\mathrm{today}}
= \PhiZero\,\psi_{\mathrm{today}}$.
Definiuj\k{a}c $\PhiZero^{(\mathrm{recal})} = \Phi_{\mathrm{eq}}$
i~$\psi' = \Phi/\PhiZero^{(\mathrm{recal})}$, r\'ownanie pola zachowuje
identyczn\k{a} posta\'c z~podstawieniem $\PhiZero \to \PhiZero^{(\mathrm{recal})}$
i~$\gamma' = \gamma/\psi_{\mathrm{eq}}^2$, $\beta' = \beta/\psi_{\mathrm{eq}}^2$,
tak \.ze $\beta' = \gamma'$ nadal obowi\k{a}zuje.

Numerycznie: $\PhiZero^{(\mathrm{recal})} = 24{,}66 \times 7/6 \approx 28{,}8$
(zweryfikowane skryptem \texttt{deep\_consistency\_v17.py}).
\end{proof}

\medskip
Dla oblicze\'n praktycznych przydatna jest posta\'c~\eqref{eq:Phi-cosmo-linearized}
zapisana bezpo\'srednio w~zmiennej $\varphi \equiv \psi = \Phi/\PhiZero$:
\begin{equation}\label{eq:Phi-cosmo}
  \ddot{\varphi} + 3H\dot{\varphi}
  + c_0^2\,\frac{4\gamma}{3}\,(\varphi - 1)
  = c_0^2\,\frac{\gamma}{3}
  - \frac{q}{\PhiZero}\,\rho_{\mathrm{tot}},
\end{equation}
kt\'ore jest \textbf{dok\l{}adnym} linearyzowanym przybli\.zeniem
r\'ownania~\eqref{eq:Phi-cosmo-exact} (nie zamian\k{a} wsp\'o\l{}czynnik\'ow,
lecz \emph{Taylor}em pierwszego rz\k{e}du wok\'o\l{} $\psi = 1$).
Cz\l{}on \'zr\'od\l{}owy jest proporcjonalny do
$\rho_{\mathrm{tot}}$ (nie do $\dot{\rho}$):
materia \emph{generuje} pole $\Phi$, wi\k{e}c
obecno\'s\'c masy--energii jest \'zr\'od\l{}em, nie
jej zmiana w~czasie.

We wczesnym Wszech\'swiecie ($\varphi \approx 1$, materia
jednorodna): $\Lambda_{\mathrm{eff}} \approx 0$, odzyskujemy
standardow\k{a} kosmologi\k{e}. W~p\'o\'znym Wszech\'swiecie
(struktury uformowane): $\langle\delta\Phi^2\rangle > 0$,
$\Lambda_{\mathrm{eff}} > 0$ --- \textbf{przyspieszona ekspansja}.

\begin{proposition}[Ograniczenie BBN]\label{prop:BBN-constraint}
Nukleosyntetyczny waru\-nek sp\'ojno\'sci
$|\Delta G/G_0| < 0{,}1$ w~epoce BBN ($T \sim 1\;\mathrm{MeV}$,
$z_{\mathrm{BBN}} \sim 10^9$) wymaga:
\begin{equation}\label{eq:BBN-bound}
  |\varphi(z_{\mathrm{BBN}}) - 1| < 0{,}05,
  \qquad\text{tj.}\quad
  \tau_0\,H_{\mathrm{BBN}} \gg 1.
\end{equation}
Jest to \textbf{automatycznie spe\l{}nione} dla naturalnych warunk\'ow
pocz\k{a}tkowych ($\varphi_{\mathrm{ini}} = 1$ po przej\'sciu fazowym).
\end{proposition}

\begin{proof}
Z~$G(\Phi) = G_0\,\PhiZero/\Phi$ (aksjomat~\ref{ax:G}):
$\Delta G/G_0 = 1 - \Phi/\PhiZero = 1 - \varphi$.

Tarcie Hubble'a $3H\dot\varphi$ w~r\'ownaniu~\eqref{eq:Phi-cosmo}
zamra\.za pole, gdy $H \gg 1/\tau_0$. Iloraz
$\tau_0 H_{\mathrm{BBN}}$ mo\.zna oszacowa\'c:
\[
  \tau_0\,H_{\mathrm{BBN}}
  = \frac{\sqrt{\PhiZero}}{c_0\sqrt{\gamma}}
    \cdot\frac{c_0}{t_{\mathrm{BBN}}}
  = \frac{\sqrt{\PhiZero}}{t_{\mathrm{BBN}}\sqrt{\gamma}},
\]
gdzie $t_{\mathrm{BBN}} \sim 1\;\mathrm{s}$. Poniewa\.z $\gamma$
jest rz\k{e}du kosmologicznego ($\gamma \sim H_0^2/c_0^2$),
otrzymujemy $\tau_0 H_{\mathrm{BBN}} \sim 10^{18}$--$10^{20}$
(zweryfikowane numerycznie, skrypt \texttt{bbn\_consistency.py}).

Pole jest ,,zamro\.zone'' w~epoce BBN: $\varphi \approx \varphi_{\mathrm{ini}}$.
Dla $\varphi_{\mathrm{ini}} = 1$ (naturalny stan po przej\'sciu fazowym substratu):
$|\Delta G/G_0| < 10^{-10}$ --- daleko poni\.zej ograniczenia~$10\%$.

Nienaturalny $\varphi_{\mathrm{ini}} \neq 1$ wymaga\l{}by dodatkowego
mechanizmu w~substracie i~jest ograniczony przez:
$|\varphi_{\mathrm{ini}} - 1| < 0{,}05$.
\end{proof}

\begin{remark}[Rozstrzygni\k{e}cie BBN --- warunki pocz\k{a}tkowe z~substratu]%
\label{rem:BBN-resolution}
Ograniczenie BBN por\'ownuje $G_{\mathrm{BBN}}$ z~$G_{\mathrm{today}}$
(\emph{mierzonym}).  Poniewa\.z $G = G_0/\psi$:
\[
  \frac{G_{\mathrm{BBN}}}{G_{\mathrm{today}}}
  = \frac{\psi_{\mathrm{today}}}{\psi_{\mathrm{BBN}}}.
\]
Dla $\psi_{\mathrm{ini}} = 1$ (pole zamro\.zone) i~$\psi_{\mathrm{today}} \approx 1{,}15$
daje to $|\Delta G/G| \approx 15\%$ --- marginalnie sp\'ojne z~ograniczeniami
(Copi i~in.~2004: $<20\%$; Fields i~in.~2020 PDG: $<15\%$), lecz w~napi\k{e}ciu
ze \'sci\'slejszymi granicami.

\textbf{Rozstrzygni\k{e}cie:} przej\'scie fazowe substratu (dod.~\ref{app:substrat})
generuje pole bezpo\'srednio w~punkcie sta\l{}ym Wilsona--Fishera,
odpowiadaj\k{a}cym $\psi_{\mathrm{ini}} = \psi_{\mathrm{eq}} = 7/6$.
Uzasadnienie:
\begin{enumerate}[nosep]
  \item Funkcjona\l{} Ginzburga--Landaua substratu ma minimum
    $v^2 = r/u$ poni\.zej $T_c$.
  \item Procedura zgruboziarnistości (RG Migdala--Kadanoffa)
    odwzorowuje to minimum dok\l{}adnie na $\psi_{\mathrm{eq}} = 7\beta/(6\gamma) = 7/6$
    (ta sama symetria $Z_2$, kt\'ora daje $\beta = \gamma$).
  \item Pole rodzi si\k{e} \textbf{ju\.z w~atraktorze} kosmologicznym;
    brak dalszej ewolucji.
\end{enumerate}
Konsekwencje:
\[
  \psi_{\mathrm{BBN}} = \psi_{\mathrm{today}} = \tfrac{7}{6}
  \;\;\Longrightarrow\;\;
  \frac{G_{\mathrm{BBN}}}{G_{\mathrm{today}}} = 1,
  \qquad |\Delta G/G| = 0.
\]
Ograniczenie BBN jest \textbf{trywianie spe\l{}nione}.  W~zrekalibrowanej
ramce (stw.~\ref{prop:recalibration}) ta konwencja jest automatyczna:
$\psi'_{\mathrm{ini}} = 1 = \psi'_{\mathrm{eq}}$, pole nigdy nie ewoluuje.
Pozosta\l{}e predykcje TGP (trzy rezimy, $\Lambda_{\mathrm{eff}}$, PPN, $c_{\mathrm{GW}}$,
QNM) s\k{a} \textbf{niezmienione}
(potwierdzone skryptem \texttt{bbn\_attractor\_resolution.py}).
\end{remark}

% =============================================================
\begin{proposition}[Warunek pocz\k{a}tkowy $\psi_{\rm ini} = 7/6$
  jako Warstwa~I$\to$III --- \statuslabel{Propozycja}]%
\label{prop:psi-ini-derived}
Niech $\beta = \gamma$ (warunek pr\'o\.zniowy, Warstwa~I).
Wtedy:
\begin{enumerate}[nosep]
  \item \textbf{Algebraicznie (Warstwa~I)}:
    potencja\l{} kosmologiczny $W(\psi) = \gamma\psi^2(\tfrac{7}{3} - 2\psi)$
    ma dok\l{}adnie dwa zera:
    \begin{equation}\label{eq:W-zeros}
      W(\psi) = 0 \;\;\Leftrightarrow\;\;
      \psi = 0 \;(\mathcal{N}_0)
      \quad\text{lub}\quad
      \psi_{\rm eq} = \frac{7\beta}{6\gamma} = \frac{7}{6}.
    \end{equation}
    Stabilno\'s\'c: $W'(\psi_{\rm eq}) = -\tfrac{49\gamma}{18} < 0$
    (pochodna si\l{}y napdzaj\k{a}cej jest ujemna),
    wi\k{e}c $\psi_{\rm eq} = 7/6$ jest \textbf{stabilnym atraktorem slow-roll}.
    Cz\k{e}sto\'s\'c oscylacji:
    \begin{equation}\label{eq:omega-cosmo}
      \frac{\omega_{\rm cosmo}}{H_0}
      = \sqrt{\PhiZero\,|W'(\psi_{\rm eq})|}
      = \sqrt{\tfrac{49\,\PhiZero}{18}}
      \approx 8{,}19
      \quad (\PhiZero \approx 24{,}66),
    \end{equation}
    co potwierdza tryb \textbf{niedot\l{}umiony} przy $z = 0$
    ($\omega_{\rm cosmo} \gg 3H_0/2$).

  \item \textbf{Zamro\.zenie w erze radiacyjnej}:
    Parametr adiabatyczno\'sci (stosunek tempa ewolucji
    do tempa Hubble'a):
    \begin{equation}\label{eq:eta-BBN}
      \eta_{\rm BBN}
      = \frac{\PhiZero\,|W'(1)|}{3\,E_{\rm BBN}^2}
      \approx 1{,}2 \times 10^{-31}
      \;\ll 1.
    \end{equation}
    Pole jest \textbf{absolutnie zamro\.zone} w~epoce BBN.
    Ca\l{}kowity dryft $\Delta\psi$ w~erze radiacji
    ($z_{\rm BBN} \to z_{\rm eq}$) jest rz\k{e}du $3 \times 10^{-11}$.

  \item \textbf{Wyj\k{a}tkowo\'s\'c $\psi_{\rm ini} = 7/6$ (BBN)}:
    Poniewa\.z $W(7/6) = 0$, pole zainicjowane w~$\psi_{\rm ini} = 7/6$
    \emph{nie ma si\l{}y napdzaj\k{a}cej}: $\psi(t) = 7/6 = \mathrm{const}$
    przez ca\l\k{a} er\k{e} radiacji i~materii.
    W~konsekwencji:
    \[
      \frac{G_{\rm BBN}}{G_{\rm today}}
      = \frac{\psi_{\rm today}}{\psi_{\rm BBN}}
      = \frac{7/6}{7/6} = 1,
      \qquad |\Delta G/G| = 0.
    \]
    Jest to \textbf{jedyna warto\'s\'c} $\psi_{\rm ini}$
    daj\k{a}ca dosk\l{}on\k{a} zgodno\'s\'c z~ograniczeniem BBN.
    Porownanie numeryczne (\texttt{ex86}, 9/9 PASS):
    \begin{center}\small
    \begin{tabular}{lcc}
      \hline
      $\psi_{\rm ini}$ & $\psi(z{=}0)$ & $|\Delta G/G|$ \\
      \hline
      $7/6 = 1{,}1\overline{6}$ (GL) & $1{,}1\overline{6}$ & $0{,}000\%$ ★ \\
      $1{,}00$ & $1{,}137$ & $13{,}7\%$ \\
      $0{,}85$ & $1{,}116$ & $31{,}3\%$ -- przekracza BBN \\
      $1{,}35$ & $1{,}196$ & $11{,}4\%$ \\
      \hline
    \end{tabular}
    \end{center}
    Zakres BBN-kompatybilny: $\psi_{\rm ini} \in [0{,}94;\,1{,}56]$;
    $\psi_{\rm eq} = 7/6$ le\.zy w~centrum.

  \item \textbf{Fizycznie (argument GL)}:
    Przej\'scie fazowe Ginzburga--Landaua substratu inicjuje pole
    $\psi$ w~stanie $\psi_{\rm ini} = \psi_{\rm eq}$:
    \begin{equation}\label{eq:psi-ini-GL}
      \psi_{\rm ini} \;=\; \psi_{\rm eq} \;=\; \frac{7}{6}.
    \end{equation}
    Uzasadnienie: procedura zgruboziarnisto\'sci (Migdal--Kadanoff)
    przy symetrii $\mathbb{Z}_2$ odwzorowuje minimum funkcjona\l{}u GL
    na $\psi_{\rm eq}$ (uw.~\ref{rem:BBN-resolution}).
    Numeryczna weryfikacja: skrypt \texttt{ex86\_psi\_ini\_attractor.py},
    9/9~PASS (sesja v35).
\end{enumerate}
\textbf{Wniosek:} $\psi_{\rm ini} = 7/6$ jest wyznaczone przez
kombinacj\k{e} trzech niezale\.znych argument\'ow:
algebraicznego (zero $W$), nukleosyntycznego ($\Delta G/G = 0$)
i~substratowego (minimum GL). Nie jest wolnym parametrem
Warstwy~II.
Parametryczna Warstwa~II:
$\{\PhiZero, a_\Gamma, \psi_{\rm ini}\}
\;\xrightarrow{W(\psi_{\rm ini})=0\,+\,\mathrm{GL}}\;
\{\PhiZero, a_\Gamma\}$, co daje:
\begin{equation}\label{eq:parsimony-2}
  N_{\rm param} = 2, \qquad
  \frac{M_{\rm obs}}{N_{\rm param}} \geq 6.
\end{equation}
\end{proposition}

% =============================================================
\subsubsection{Kompletne r\'ownania kosmologii TGP}\label{sssec:cosmo-complete}
% =============================================================

\begin{remark}[Status sekcji --- C.1]\label{rem:cosmo-complete-status}
Poni\.zszy blok zbiera w~jednym miejscu pe\l{}ny uk\l{}ad r\'owna\'n
opisuj\k{a}cy jednorodn\k{a} kosmologi\k{e} TGP.
Dotychczas r\'ownania te by\l{}y rozproszone w~kilku sekcjach;
ten blok jest ich \textbf{kanoniczn\k{a} wersj\k{a} odwo\l{}awcz\k{a}}.
\end{remark}

\begin{proposition}[Pe\l{}ny uk\l{}ad kosmologiczny TGP]%
\label{prop:cosmo-system}
\statuslabel{Propozycja}: metryka efektywna + dzia\l{}anie TGP,
warunek pocz\k{a}tkowy \statuslabel{Hipoteza}.

\medskip
W~jednorodnym, izotropowym Wszech\'swiecie FRW ($k=0$)
dynamik\k{e} TGP opisuj\k{a} nast\k{e}puj\k{a}ce r\'ownania:

\begin{enumerate}[(1)]
\item \textbf{Zmodyfikowane r\'ownanie Friedmanna}:
\begin{equation}\label{eq:friedmann-cosmo-full}
  3\!\left(H + \frac{\dot\psi}{4\psi}\right)^{\!2}\!\sqrt\psi
  \;=\; c_0^2\!\left[\frac{\dot\psi^2}{2\,c_0^2\,\psi^{3/2}}
  + \frac{U(\psi)}{\sqrt\psi}\right]
  + \frac{8\pi\,G_0}{3}\frac{\rho_{\rm mat}}{\psi^{3/2}},
\end{equation}
gdzie $H = \dot a/a$ i~$G_{\rm eff}(t) = G_0/\psi(t)$.

\item \textbf{R\'ownanie pola t\l{}a} (z~tw.~\ref{prop:TGP-FRW-full}
przy $\nabla\psi = 0$):
\begin{equation}\label{eq:psi-background}
  \ddot\psi + 3H\dot\psi + \frac{2\dot\psi^2}{\psi}
  = c_0^2\,W(\psi),
  \qquad W(\psi) = \frac{4U(\psi)}{\psi} + U'(\psi),
\end{equation}
gdzie $U(\psi) = \psi^3/3 - \psi^4/4 + \lambda(\psi-1)^6/6$.

\item \textbf{Efektywna pr\k{e}dko\'s\'c \'swiat\l{}a i~sta\l{}a grawitacyjna}:
\begin{equation}\label{eq:Geff-ceff}
  G_{\rm eff}(t) = G_0/\psi(t),
  \qquad
  c_{\rm eff}(t) = c_0/\sqrt{\psi(t)}.
\end{equation}

\item \textbf{Temperatura substratu} (analogia Unruha):
\begin{equation}\label{eq:TGamma-cosmo}
  T_\Gamma(z) = \frac{\hbar\,H(z)}{2\pi\,k_B},
\end{equation}
gdzie $z$ jest przesunijeciem ku czerwieni.

\item \textbf{Efektywna sta\l{}a kosmologiczna} przy $\psi \approx 1$:
\begin{equation}\label{eq:Lambda-eff-cosmo}
  \Lambda_{\rm eff} \approx c_0^2\,W(1) = c_0^2\,\frac{\gamma}{3}
  \;\;\Longrightarrow\;\;
  3H_0^2 = \Lambda_{\rm eff}.
\end{equation}
\end{enumerate}

\medskip\noindent
\textbf{Warunki pocz\k{a}tkowe} [\statuslabel{Hipoteza}: atraktor GL]:
\begin{equation}\label{eq:ic-cosmo}
  \psi_{\rm ini} = \tfrac{7}{6},\quad
  \dot\psi_{\rm ini} = 0,\quad
  H_{\rm ini} = H_{\rm BBN}.
\end{equation}

\medskip\noindent
\textbf{Trzy ery kosmiczne}:
\begin{center}\footnotesize
\begin{tabular}{@{}llll@{}}
\toprule
\textbf{Era} & \textbf{Przybli\.zenie} & \textbf{Ewolucja $\psi$}
  & \textbf{Status $G_{\rm eff}$} \\
\midrule
Radiacyjna ($z > 10^4$)
  & $H^2 \gg c_0^2 W$
  & $\psi \approx 7/6$ (zamro\.zone)
  & $G_{\rm eff} \approx \frac{6}{7}G_0$ \\[3pt]
Materyjna ($z \in [10, 10^4]$)
  & $H^2 > c_0^2 W$
  & $\psi$ \emph{wolno} opada
  & $G_{\rm eff}$ powoli ro\'snie \\[3pt]
Ciemna energia ($z < 2$)
  & $H^2 \sim c_0^2 W$
  & $\psi$ zmierza ku atraktorowi $\psi = 1$
  & $G_{\rm eff} \to G_0$ \\
\bottomrule
\end{tabular}
\end{center}

Bie\.z\k{a}ca warto\'s\'c $\psi(z=0) \approx 1{,}15$ (Warstwa~II: $\PhiZero$).
Weryfikacja numeryczna: skrypty \texttt{friedmann\_derivation.py}
i~\texttt{cosmo\_evolution.py}.
\end{proposition}

% =============================================================
\subsubsection{Konfrontacja z~danymi: Planck, DESI, BBN, LLR}%
\label{sssec:cosmo-data-confrontation}
% =============================================================

\begin{remark}[Jawna konfrontacja $H(z)$, $w(z)$, $G(z)$
  z~obserwacjami]\label{rem:cosmo-data}

Zbieramy jawne predykcje kosmologiczne TGP i~por\'ownujemy
z~najnowszymi danymi obserwacyjnymi.

\medskip\noindent
\textbf{1.\ Stosunek $H(z)/H_{\Lambda\mathrm{CDM}}(z)$.}
Z~r.~\eqref{eq:H-Friedmann} (A15):
\begin{equation}\label{eq:H-ratio-explicit}
  \frac{H_{\rm TGP}(z)}{H_{\Lambda\rm CDM}(z)}
  = \frac{1}{\sqrt{\psi(z)}}
  \approx 1 - \frac{\kappa}{2}\,\frac{\Omega_m(1+z)^3}{E^2(z)},
\end{equation}
gdzie $\kappa = 3/(4\PhiZero) \approx 0{,}030$,
$E^2(z) = \Omega_r(1+z)^4 + \Omega_m(1+z)^3 + \Omega_\Lambda$
(Planck~2018: $\Omega_m = 0{,}315$, $\Omega_\Lambda = 0{,}685$).
Odchylenie od $\Lambda$CDM nie przekracza $1{,}5\%$
przy $z < 3$, co jest poni\.zej progu DESI~BAO ($\sim 2\%$).

\medskip\noindent
\textbf{2.\ R\'ownanie stanu ciemnej energii.}
Z~uwagi~\ref{rem:wz-quantitative}:
$w_{\rm DE}(z) = -1 + \varepsilon(z)$,
$\varepsilon(z) = \frac{2}{3}\varepsilon_0\,D^2(z)\,f(z)$,
$\varepsilon_0 \approx 5{,}4 \times 10^{-9}$.
Parametryzacja CPL:
$w_0 + 1 \approx 2{,}3 \times 10^{-9}$, $w_a \approx -2{,}5 \times 10^{-9}$
--- oba $\sim 10^7$ razy poni\.zej progu DESI
(tj.\ $|w_0+1| < 0{,}05$, $|w_a| < 0{,}3$).
\textbf{Falsyfikacja}: je\.zeli DESI DR2+ zmierzy $|w+1| > 0{,}01$,
wyklucza to minimaln\k{a} kosmologi\k{e} TGP.

\medskip\noindent
\textbf{3.\ Tabela konfrontacji.}

\begin{center}\small
\begin{tabular}{@{}llllrl@{}}
\toprule
\textbf{Obserwabla} & \textbf{TGP} & \textbf{Dane}
  & \textbf{$\sigma$} & \textbf{Status} & \textbf{\'Zr\'od\l{}o} \\
\midrule
$H(z)/H_{\Lambda\rm CDM}$ ($z{<}3$)
  & $1 - \mathcal{O}(\kappa)$
  & $1{,}00 \pm 0{,}02$
  & ${<}\,1$
  & PASS & DESI BAO \\[2pt]
$G(z_{\rm BBN})/G_0$
  & $\approx 0{,}97$
  & $1 \pm 0{,}13$
  & $0{,}2$
  & PASS & BBN ($^4$He) \\[2pt]
$G(z_{\rm CMB})/G_0$
  & $\approx 0{,}97$
  & $1 \pm 0{,}05$
  & $0{,}7$
  & PASS & Planck CMB \\[2pt]
$|\dot G/(GH_0)|_{z=0}$
  & $\approx 0{,}009$
  & ${<}\,0{,}01$
  & OK
  & PASS & LLR \\[2pt]
$w_0 + 1$
  & $2{,}3 \times 10^{-9}$
  & ${<}\,0{,}05$
  & ${\sim}\,0$
  & PASS & DESI \\[2pt]
$w_a$
  & $-2{,}5 \times 10^{-9}$
  & ${<}\,0{,}3$
  & ${\sim}\,0$
  & PASS & DESI \\[2pt]
$n_s$
  & $0{,}965$
  & $0{,}9649 \pm 0{,}004$
  & $0{,}03$
  & PASS & Planck \\[2pt]
$r$ (tensor/scalar)
  & $0{,}003$
  & ${<}\,0{,}036$
  & OK
  & PASS & BICEP/Keck \\[2pt]
$c_{\rm GW}/c_0$
  & $1$ (\'sci\'sle)
  & $|1-c_T/c| < 10^{-15}$
  & OK
  & PASS & GW170817 \\
\midrule
\multicolumn{6}{l}{\textbf{9/9 PASS} ~|~
  Falsyfikacja: $|w+1| > 0{,}01$ (DESI DR2+)
  lub $|\dot G/G|/H_0 > 0{,}01$ (LLR nast.~gen.)} \\
\bottomrule
\end{tabular}
\end{center}

\medskip\noindent
\textbf{Parametr $\kappa$} (jedyny wolny parametr kosmologii TGP)
jest okre\'slony przez $\PhiZero$, kt\'ore jest wsp\'olne z~sektorem
cz\k{a}steczkowym.
Je\.zeli $\PhiZero = N_f^2$ (hipoteza z~dod.~V),
to $\kappa = 3/(4 \cdot 25) = 0{,}030$ --- \textbf{brak} dodatkowych
wolnych parametr\'ow kosmologicznych.

Weryfikacja numeryczna: skrypt \texttt{ex187\_cosmo\_data\_confrontation.py}.
\end{remark}

% =============================================================
\subsection{Ciemna energia: oszacowanie $\Lambda_{\mathrm{eff}}$}%
\label{ssec:de-oszacowanie}
% =============================================================

Poni\.zej pokazujemy, \.ze TGP \textbf{naturalnie} generuje
$\Lambda_{\mathrm{eff}}$ w\l{}a\'sciwego rz\k{e}du wielko\'sci,
bez \emph{fine-tuningu}.

\begin{proposition}[Oszacowanie $\Lambda_{\mathrm{eff}}$ z~parametr\'ow TGP]%
\label{prop:Lambda-eff}
Efektywna sta\l{}a kosmologiczna w~TGP ma dwa sk\l{}adniki:
\begin{equation}\label{eq:Lambda-eff-full}
  \Lambda_{\mathrm{eff}}
  = \underbrace{U(\psi_0)}_{\text{potencja\l{} t\l{}a}}
  + \underbrace{\gamma\,\langle(\delta\Phi/\PhiZero)^2\rangle}_{\text{backreakcja struktur}},
\end{equation}
gdzie $\psi_0 \approx 1$ jest warto\'sci\k{a} pola t\l{}a,
$U(\psi) = (\beta/3)\psi^3 - (\gamma/4)\psi^4$
(r.~\eqref{eq:U-potential}), a~$\delta\Phi$ jest perturbacj\k{a}
wywo\l{}an\k{a} przez formowanie struktur kosmicznych.

Dla $\beta = \gamma$ (warunek pr\'o\.zni, stw.~\ref{prop:vacuum-condition}):
\begin{enumerate}[nosep]
  \item \textbf{Cz\l{}on potencja\l{}owy:}
    $U(1) = \gamma/3 - \gamma/4 = \gamma/12$.
  \item \textbf{Cz\l{}on backreakcyjny:}
    $\gamma\,\langle(\delta\Phi/\PhiZero)^2\rangle
    = 4\gamma\,\langle(\Phi_N/c_0^2)^2\rangle
    \approx 4\gamma\,\sigma_\Phi^2$,
    gdzie $\sigma_\Phi^2 \sim 10^{-10}\text{--}10^{-9}$
    (wariancja potencja\l{}u grawitacyjnego).
\end{enumerate}
\textbf{Dominuje cz\l{}on potencja\l{}owy}: $\Lambda_{\mathrm{eff}} \approx \gamma/12$.
\end{proposition}

\begin{proof}
\emph{Cz\l{}on potencja\l{}owy.}\quad
Jednorodne pole $\psi \approx 1$ ma niezerow\k{a} energi\k{e}
potencja\l{}n\k{a} $U(1) = \beta/3 - \gamma/4$.
Dla $\beta = \gamma$: $U(1) = \gamma/12 > 0$.
Ten resztkowy potencja\l{} dzia\l{}a jak $\Lambda > 0$,
poniewa\.z r\'ownanie stanu $w = p/\rho$ w~regimencie
powolnego staczania daje $w \gtrsim -1$ (kwintesencja,
uwaga~\ref{rem:two-masses}).

\emph{Cz\l{}on backreakcyjny.}\quad
Z~twierdzenia~\ref{thm:newton}: $\delta\Phi/\PhiZero = 2G_0 M/(c_0^2 r)
= -2\Phi_N/c_0^2$.
Zatem $\langle(\delta\Phi/\PhiZero)^2\rangle = 4\sigma_\Phi^2$,
gdzie $\sigma_\Phi^2 = \langle(\Phi_N/c_0^2)^2\rangle_{\mathrm{obj}}$
jest wariancj\k{a} potencja\l{}u Newtona u\'srednion\k{a}
obj\k{e}to\'sciowo (z~widma mocy materii).
Obserwacyjnie $\sigma_\Phi \sim 3 \times 10^{-5}$
(amplituda fluktuacji CMB), zatem
$4\sigma_\Phi^2 \sim 4 \times 10^{-9}$.
Wk\l{}ad backreakcyjny: $\gamma \times 4 \times 10^{-9}
\ll \gamma/12$, wi\k{e}c jest subdominuj\k{a}cy.

\emph{Oszacowanie liczbowe.}\quad
Z~warunku $\tau_0 \sim 1/H_0$ (def.~\ref{def:tau0}):
\[
  \frac{\sqrt{\PhiZero}}{c_0\sqrt{\gamma}} \sim \frac{1}{H_0}
  \qquad\Longrightarrow\qquad
  \gamma \sim \frac{\PhiZero\,H_0^2}{c_0^2}.
\]
Zatem:
\begin{equation}\label{eq:Lambda-estimate}
  \boxed{
    \Lambda_{\mathrm{eff}}
    \approx \frac{\gamma}{12}
    \sim \frac{\PhiZero\,H_0^2}{12\,c_0^2}.
  }
\end{equation}
Obserwowana warto\'s\'c: $\Lambda_{\mathrm{obs}}
= 3\,\Omega_\Lambda\,H_0^2/c_0^2 \approx 2{,}1\,H_0^2/c_0^2$.
Stosunek:
\begin{equation}\label{eq:Lambda-ratio}
  \frac{\Lambda_{\mathrm{obs}}}{\Lambda_{\mathrm{eff}}}
  \approx \frac{25{,}2}{\PhiZero}.
\end{equation}
Dla $\PhiZero \sim 10\text{--}30$ uzyskujemy
$\Lambda_{\mathrm{eff}} \sim \Lambda_{\mathrm{obs}}$
\textbf{bez \emph{fine-tuningu}} --- warto\'s\'c $\PhiZero = \mathcal{O}(10)$
jest ,,naturalna'' w~sensie braku hierarchii.
\end{proof}

\begin{remark}[Rozwi\k{a}zanie problemu sta\l{}ej kosmologicznej]%
\label{rem:CC-problem}
W~fizyce standardowej $\Lambda$ wymaga \emph{fine-tuningu}
rz\k{e}du $10^{-122}$ (w~jednostkach Plancka).
W~TGP ten problem nie wyst\k{e}puje:
\begin{enumerate}[nosep]
  \item Skala $\Lambda$ jest wyznaczona przez
    $\gamma \sim H_0^2/c_0^2$ --- ten sam parametr,
    kt\'ory okre\'sla mas\k{e} pola ($m_{\mathrm{sp}}$)
    i~kosmologiczn\k{a} skal\k{e} czasow\k{a} ($\tau_0$).
  \item Brak oddzielnej ,,energii pr\'o\.zni'' z~teorii kwantowej p\'ol:
    substrat ma naturalne obci\k{e}cie UV
    (stw.~\ref{prop:UV-cutoff}), a~$\Lambda$ nie wynika
    z~sumowania energii punkt\'ow zerowych.
  \item D\l{}ugo\'s\'c fali Comptona pola: $\lambda_{\mathrm{sp}}
    = 2\pi/m_{\mathrm{sp}} \sim 2\pi c_0/H_0
    \sim R_H$ (promie\'n Hubble'a).
    Zasi\k{e}g pola generuj\k{a}cego przestrze\'n
    jest rz\k{e}du rozmiaru obserwowalnego Wszech\'swiata
    --- warunek sp\'ojno\'sci, nie \emph{fine-tuning}.
\end{enumerate}
Numeryczna weryfikacja: skrypt \texttt{lambda\_eff\_estimation.py}.
\end{remark}

\begin{proposition}[Estymacja $\PhiZero$ z~obserwacji
  kosmologicznych]\label{prop:Phi0-estimate}
Bezwymiarowa warto\'s\'c pr\'o\.zniowa $\PhiZero$ jest
wyznaczalna z~niezale\.znych \'zr\'ode\l{} obserwacyjnych:

\medskip
\begin{center}
\small
\begin{tabular}{@{}lll@{}}
\toprule
\textbf{Metoda} & \textbf{Relacja} & \textbf{Wynik} \\
\midrule
Sta\l{}a kosmologiczna
  & $\Lambda_{\mathrm{obs}} \approx \PhiZero H_0^2 / (12\,c_0^2)$
  & $\PhiZero = 36\,\Omega_\Lambda
    \approx 25$ \\[4pt]
Masa pola $m_{\mathrm{sp}}$
  & $\gamma = \PhiZero H_0^2/c_0^2$,\;
    $m_{\mathrm{sp}} = \sqrt\gamma$
  & $\PhiZero \sim m_{\mathrm{sp}}^2 c_0^2/H_0^2$ \\[4pt]
Zasi\k{e}g Yukawy
  & $\lambda_{\mathrm{sp}} = 1/\sqrt\gamma
    \sim c_0/H_0 \sim R_H$
  & sp\'ojne z~$\PhiZero \sim 10$ \\[4pt]
Skala FRW
  & $\tau_0 = \sqrt{\PhiZero/(c_0^2\gamma)} \sim 1/H_0$
  & sp\'ojne z~powy\.zszym \\[4pt]
Ciemna energia ($w_0$)
  & $w_0 = -1$ wymaga $\PhiZero = \mathrm{const}$
    (pr\'o\.znia statyczna)
  & $\PhiZero$ stabilne \\
\bottomrule
\end{tabular}
\end{center}

\medskip\noindent
Wszystkie metody zbiegaj\k{a} do:
\begin{equation}\label{eq:Phi0-value}
  \boxed{\;\PhiZero \approx 12 \times \frac{\Lambda_{\mathrm{obs}}\,c_0^2}{H_0^2}
  = 12 \times 3\,\Omega_\Lambda \approx 25
  \qquad (\text{bezwymiarowa})\;}
\end{equation}
Warto\'s\'c ta jest \textbf{rz\k{e}du jedno\'sci} ---
brak hierarchii, brak \emph{fine-tuningu}.
Interpretacja fizyczna: $\PhiZero$ jest \'sredni\k{a}
g\k{e}sto\'sci\k{a} wygenerowanej przestrzeni (w~jednostkach
naturalnych substratu). Warto\'s\'c $\sim 25$ oznacza,
\.ze pr\'o\.znia jest umiarkowanie ,,g\k{e}sta''
--- ka\.zdy punkt przestrzeni odpowiada $\sim 25$
substratowym stopniom swobody po ugrubieniu.
\end{proposition}

% =============================================================
\subsection{Minimalne kryteria wyboru dzia\l{}ania}\label{ssec:kryteria}
% =============================================================

Przej\'scie od ontologii (sekcja~\ref{sec:ontologia}) do
konkretnego funkcjona\l{}u dzia\l{}ania wymaga uzasadnienia:
dlaczego \emph{ten} funkcjona\l{}, a~nie inny? Poni\.zej
podajemy \textbf{minimalne kryteria wewn\k{e}trzne},
kt\'ore jednoznacznie prowadz\k{a} do postaci~\eqref{eq:action}.

\begin{remark}[Kryteria wyboru dzia\l{}ania]\label{rem:action-criteria}
Funkcjona\l{} dzia\l{}ania TGP jest \textbf{jednoznacznie wyznaczony}
(do rz\k{e}du $\Phi^4$) przez nast\k{e}puj\k{a}ce sze\'s\'c wymaga\'n:
\begin{enumerate}[nosep]
  \item \textbf{Lokalno\'s\'c}: dzia\l{}anie jest ca\l{}k\k{a}
    z~g\k{e}sto\'sci lagran\.zowskiej zale\.znej od $\Phi$
    i~jej pochodnych (nie od nielokalnych ca\l{}ek).
    Sektor nielokalny (\S\ref{ssec:nielokalny}) jest
    korekcj\k{a} wy\.zszego rz\k{e}du.
  \item \textbf{Regularno\'s\'c}: lagran\.zjan jest g\l{}adki
    dla $\Phi > 0$. Degeneracja przy $\Phi \to 0$ jest
    dopuszczalna (sygnalizuje granice opisu ci\k{a}g\l{}ego).
  \item \textbf{Zgodno\'s\'c z~zasad\k{a} zerowej sumy (ZS2)}:
    jednorodne $\Phi = \PhiZero$ jest rozwi\k{a}zaniem
    r\'owna\'n ruchu przy sta\l{}ej $\bar{\rho}$.
    Wymaga $\beta - \gamma = -q\bar{\rho}$ (stw.~\ref{prop:vacuum-condition}).
  \item \textbf{Trzy re\.zimy}: potencja\l{} efektywny musi
    dopuszcza\'c trzy zmiany znaku gradientu (przyci\k{a}ganie,
    odpychanie, studnia). Wymaga cz\l{}on\'ow co najmniej
    trzeciego i~czwartego rz\k{e}du w~$\Phi$.
  \item \textbf{Stabilno\'s\'c t\l{}a}: perturbacje wok\'o\l{}
    $\PhiZero$ maj\k{a} profil zanikaj\k{a}cy (Yukawa, $m_{\mathrm{sp}}^2 > 0$).
    Wymaga $3\gamma > 2\beta$ (stw.~\ref{prop:vacuum-stability}).
  \item \textbf{Granica s\l{}abopolowa}: w~granicy
    $|\delta\Phi/\PhiZero| \ll 1$ dzia\l{}anie redukuje si\k{e}
    do r\'ownania Poissona z~\'zr\'od\l{}em masowym
    (granica newtonowska, tw.~\ref{thm:newton}).
\end{enumerate}
Jedyn\k{a} postaci\k{a} potencja\l{}u spe\l{}niaj\k{a}c\k{a}
kryteria 3--6 (przy zachowaniu minimalno\'sci: najni\.zsze
rz\k{e}dy w~$\Phi$) jest:
$U(\Phi) = (\beta/3)\Phi^3/\PhiZero^3
- (\gamma/4)\Phi^4/\PhiZero^4 + \mathrm{const}$.
Kryteria 1--2 okre\'slaj\k{a} cz\l{}on kinetyczny.
Kryterium~6 fiksuje sprz\k{e}\.zenie \'zr\'od\l{}owe.
\end{remark}

% =============================================================
\subsection{Funkcjona\l{} dzia\l{}ania TGP}\label{ssec:action}
% =============================================================

Powy\.zsze r\'ownania pola mo\.zna wyprowadzi\'c z~zasady
najmniejszego dzia\l{}ania:

\begin{hypothesis}[Dzia\l{}anie TGP]\label{hyp:action}
\begin{equation}\label{eq:action}
  \boxed{
    S[\Phi] = \int d^4x\;\sqrt{-g_{\mathrm{eff}}}\;\biggl[
      \frac{1}{\kappa}\biggl(
        \frac{1}{2}\,\frac{(\nabla\Phi)^2}{\PhiZero^2}
        - U(\Phi)
      \biggr)
      + \mathcal{L}_{\mathrm{mat}}(\Phi,\psi)
    \biggr],
  }
\end{equation}
gdzie:
\begin{align}
  U(\Phi) &= \frac{\beta}{3}\,\frac{\Phi^3}{\PhiZero^3}
           - \frac{\gamma}{4}\,\frac{\Phi^4}{\PhiZero^4}
           + \mathrm{const}, \label{eq:U-potential}\\
  \sqrt{-g_{\mathrm{eff}}} &= \exp\!\left(\frac{4\,\delta\Phi}{\PhiZero}
           \right) \simeq \left(\frac{\Phi}{\PhiZero}\right)^4
           \quad\text{(do pierwszego rz\k{e}du),} \\
  \kappa &= \frac{8\pi G_0}{c_0^4}.
\end{align}
$\mathcal{L}_{\mathrm{mat}}$ jest lagran\.zjanem materii,
sprz\k{e}\.zonym z~$\Phi$ przez metryk\k{e} efektywn\k{a} ---
materia ,,czuje'' przestrze\'n wygenerowan\k{a} przez $\Phi$.
\end{hypothesis}

\begin{remark}[Brak tensora metrycznego jako fundamentu]
Dzia\l{}anie~\eqref{eq:action} \emph{nie} jest dzia\l{}aniem
Einsteina--Hilberta z~dodatkowym polem. Metryka
$g_{\mu\nu}^{\mathrm{eff}}$ jest \textbf{wynikowa} ---
zdefiniowana przez $\Phi$ (hipoteza~\ref{hyp:metric-exp}).
Zmienno\'sci\k{a} dzia\l{}ania jest sam\k{a} $\Phi$, nie $g_{\mu\nu}$.
R\'ownania Einsteina nie s\k{a} postulowane --- \emph{emerguj\k{a}}
jako warunek sp\'ojno\'sci geometrycznej
(twierdzenie~\ref{thm:einstein-emergence}).
\end{remark}

Przeprowad\'zmy wariacje $\delta S/\delta\Phi = 0$ jawnie.

\begin{proof}[Jawna wariacja dzia\l{}ania]
Oznaczmy $\psi = \Phi/\PhiZero$. Dzia\l{}anie~\eqref{eq:action}
w~zmiennej $\psi$:
\[
  S = \int d^4x\;\psi^4\;\biggl[
    \frac{1}{\kappa}\biggl(
      \frac{1}{2}\,(\nabla\psi)^2
      - \frac{\beta}{3}\psi^3 + \frac{\gamma}{4}\psi^4
    \biggr)
    - \frac{q}{\PhiZero}\,\rho\,\psi
  \biggr],
\]
gdzie u\.zyli\'smy $\sqrt{-g_{\mathrm{eff}}} = c_0\psi$
(zob.\ sek.~\ref{ssec:unified-action}; historycznie
przybli\.zano $\psi^4$, poprawiono w~v47).
Wariacja wzgl\k{e}dem $\psi$:
\begin{align}
  0 &= \frac{\delta S}{\delta\psi}
  = \frac{1}{\kappa}\bigl[
    -\nabla\!\cdot\!(\psi^4\nabla\psi)
    + 2\psi^3(\nabla\psi)^2
    - \beta\psi^6 + \gamma\psi^7
  \bigr]
  - \frac{q}{\PhiZero}\rho\,\psi^4
  + 4\psi^3\bigl[\cdots\bigr]. \nonumber
\end{align}
Czynnik $1/\kappa$ jest wsp\'olny dla ca\l{}ego sektora
grawitacyjnego i~\textbf{skraca si\k{e}}.
Dziel\k{a}c przez $\psi^3$ i~przechodz\k{a}c do $\Phi = \PhiZero\psi$:
\begin{equation}\label{eq:variation-explicit}
  \nabla^2\Phi
  + \underbrace{\frac{2\,(\nabla\Phi)^2}{\Phi}}_{\alpha_{\mathrm{gen}}=2}
  + \beta\,\frac{\Phi^2}{\PhiZero}
  - \gamma\,\frac{\Phi^3}{\PhiZero^2}
  = -q\,\PhiZero\,\rho.
\end{equation}
Jest to dok\l{}adnie r\'ownanie~\eqref{eq:full-field-alt}
z~\textbf{wygenerowanym} $\alpha = 2$ --- sprz\k{e}\.zenie
gradientowe nie jest wolnym parametrem, lecz konsekwencj\k{a}
sprz\k{e}\.zenia kinetycznego $K(\psi)=\psi^4$
(element obj\k{e}to\'sciowy zosta\l{} poprawiony na
$\sqrt{-g_{\mathrm{eff}}} = c_0\,\psi$;
zob.\ sek.~\ref{ssec:unified-action}).
\end{proof}

\begin{corollary}[Redukcja liczby wolnych parametr\'ow]\label{cor:alpha-fixed}
Sta\l{}a $\alpha$ w~r\'ownaniu pola~\eqref{eq:full-field-alt}
nie jest niezale\.znym parametrem: $\alpha = 2$ wynika
z~dzia\l{}ania~\eqref{eq:action}. TGP posiada zatem
\textbf{dwa} wolne bezwymiarowe parametry: $\beta$ i~$\gamma$.
\end{corollary}

% =============================================================
\subsection{Formalizacja relacji $\Ffield \leftrightarrow \Phi$}\label{ssec:F-Phi}
% =============================================================

Funkcja stanu $\Ffield$ (definicja~\ref{def:F}) i~pole
przestrzenno\'sci $\Phi$ (definicja~\ref{def:Phi}) s\k{a}
dwoma opisami tego samego zjawiska --- odej\'scia od nico\'sci.
Ich zwi\k{a}zek formalny zamyka luk\k{e} mi\k{e}dzy
poziomem aksjomatycznym (sekcja~\ref{sec:ontologia}) a~poziomem
dynamicznym (r\'ownanie~\ref{eq:full-field-alt}).

Nale\.zy jednak rozr\'o\.zni\'c \textbf{dwa poziomy} funkcji
stanu, odpowiadaj\k{a}ce r\'o\.znym punktom odniesienia.

\begin{definition}[Absolutna funkcja stanu $\Ffield_{\!\mathrm{abs}}$]\label{def:Fabs}
\textbf{Absolutna funkcja stanu} mierzy odej\'scie od nico\'sci
$\Nzero$ (gdzie $\Phi = 0$):
\begin{equation}\label{eq:Fabs-def}
    \Ffield_{\!\mathrm{abs}}(x) \;=\; \frac{\Phi(x)^2}{\PhiZero^2}.
\end{equation}
Weryfikacja:
\begin{itemize}[nosep]
  \item $\Phi = 0$ ($\Nzero$):
    $\Ffield_{\!\mathrm{abs}} = 0$ --- nico\'s\'c.
    \textbf{Sp\'ojne z~aksjomatem~\ref{ax:N0}.} \checkmark
  \item $\Phi = \PhiZero$ (t\l{}o referencyjne):
    $\Ffield_{\!\mathrm{abs}} = 1$ --- pe\l{}ne istnienie. \checkmark
  \item $\Phi \to \infty$ (wnętrze CD):
    $\Ffield_{\!\mathrm{abs}} \to \infty$ --- granica
    zamro\.zenia. \checkmark
  \item $\Ffield_{\!\mathrm{abs}} \geq 0$ zawsze. \checkmark
\end{itemize}
$\Ffield_{\!\mathrm{abs}}$ jest w\l{}a\'sciwym odpowiednikiem
aksjomatycznej $\Ffield$ z~definicji~\ref{def:F}: mierzy ,,ilo\'s\'c
istnienia'' i~jest zerowa dok\l{}adnie w~$\Nzero$.
\end{definition}

\begin{definition}[Perturbacyjna funkcja stanu $\Floc$]\label{def:Floc}
\textbf{Perturbacyjna funkcja stanu} mierzy
odchylenie pola $\Phi$ od warto\'sci referencyjnej $\PhiZero$:
\begin{equation}\label{eq:Floc-def}
    \Floc(x) \;=\; \bigl(\Phi(x) - \PhiZero\bigr)^2
    \;=\; \PhiZero^2\,\varphi(x)^2,
\end{equation}
gdzie $\varphi(x) = \Phi(x)/\PhiZero - 1$ jest
bezwymiarow\k{a} fluktuacj\k{a}.
\end{definition}

\begin{remark}[Dwa poziomy funkcji stanu --- rozr\'o\.znienie]\label{rem:F-dwa-poziomy}
$\Ffield_{\!\mathrm{abs}}$ i~$\Floc$ mierz\k{a} \textbf{r\'o\.zne
wielko\'sci} i~nie nale\.zy ich myli\'c:
\begin{center}
\begin{tabular}{lcc}
\toprule
& $\Ffield_{\!\mathrm{abs}} = \Phi^2/\PhiZero^2$ & $\Floc = (\Phi-\PhiZero)^2$ \\
\midrule
$\Phi = 0$ ($\Nzero$) & $0$ (nico\'s\'c) & $\PhiZero^2$ (maks.\ odej\'scie od ref.) \\
$\Phi = \PhiZero$ (t\l{}o) & $1$ (pe\l{}ne istnienie) & $0$ (brak perturbacji) \\
$\Phi \to \infty$ (CD) & $\to\infty$ & $\to\infty$ \\
\bottomrule
\end{tabular}
\end{center}
W~szczeg\'olno\'sci:
\begin{itemize}[nosep]
  \item Aksjomat~\ref{ax:N0} ($\Ffield|_{\Nzero} = 0$) odnosi
    si\k{e} do $\Ffield_{\!\mathrm{abs}}$, nie do $\Floc$.
  \item Dynamika i~perturbacje kosmologiczne u\.zywaj\k{a} $\Floc$
    (wygodniejsza w~rachunkach wok\'o\l{} $\PhiZero$).
  \item Zasada zerowej sumy (aksjomat~\ref{ax:zero}) dotyczy
    \textbf{znakowej asymetrii} $\Delta(x) = \Phi(x) - \PhiZero$:
    $\int_\Sigma \Delta\,\sqrt{h}\,d^3x = 0$ --- jest
    warunkiem na~$\varphi$, nie na~$\Ffield_{\!\mathrm{abs}}$.
  \item Dualno\'s\'c $\Nzero$ (twierdzenie~\ref{thm:N0-duality})
    jest widoczna w~$\Ffield_{\!\mathrm{abs}}$: obie granice
    ($\Phi \to 0$ i~$\Phi \to \infty$) le\.z\k{a} na ,,\'scianach''
    sektora $\mathcal{S}_1$, ale z~r\'o\.znych stron.
\end{itemize}
Oba poziomy s\k{a} zwi\k{a}zane to\.zsamo\'sci\k{a}:
\begin{equation}\label{eq:F-relation}
    \Ffield_{\!\mathrm{abs}}
    = \frac{(\PhiZero + \PhiZero\varphi)^2}{\PhiZero^2}
    = (1 + \varphi)^2
    = 1 + 2\varphi + \varphi^2,
    \qquad
    \Floc = \PhiZero^2\varphi^2.
\end{equation}
W~granicy s\l{}abego pola ($|\varphi| \ll 1$):
$\Ffield_{\!\mathrm{abs}} \approx 1 + 2\varphi$, wi\k{e}c
$\Ffield_{\!\mathrm{abs}} - 1 \approx 2\varphi$ jest liniowo
proporcjonalna do $\sqrt{\Floc}$ --- oba opisy s\k{a} r\'ownowa\.zne
perturbacyjnie.
\end{remark}

Weryfikacja sp\'ojno\'sci $\Floc$ w~obszarze perturbacyjnym:
\begin{itemize}[nosep]
  \item $\Phi = \PhiZero$ (t\l{}o referencyjne):
    $\Floc = 0$ --- brak lokalnego odej\'scia. \checkmark
  \item $\Phi \neq \PhiZero$ (\'zr\'od\l{}o obecne):
    $\Floc > 0$ --- istnieje lokalna struktura. \checkmark
  \item $\Floc \geq 0$ zawsze. \checkmark
  \item Globalna funkcja stanu: $\Ffield[\text{Wszech\'swiat}]
    = \int_\Sigma \Floc(x)\,\sqrt{h}\,d^3x > 0$
    wewn\k{e}trznie --- sp\'ojne z~$\mathcal{S}_1$. \checkmark
\end{itemize}

\begin{proposition}[Zerowa suma w~j\k{e}zyku $\varphi$]\label{prop:zero-phi}
Zasada zerowej sumy (aksjomat~\ref{ax:zero}) w~j\k{e}zyku
pola $\Phi$ przek\l{}ada si\k{e} na warunek~\eqref{eq:zero-constraint}:
\'srednia perturbacja $\varphi$ na zamkni\k{e}tej hypersurface
wynosi zero:
\begin{equation}
    \langle\varphi\rangle_\Sigma
    \;=\; \frac{1}{\mathrm{Vol}(\Sigma)}
    \int_\Sigma \varphi(x)\,\sqrt{h}\,d^3x = 0.
\end{equation}
To jest \textbf{globalny warunek brzegowy} na rozwi\k{a}zania
r\'ownania pola, nie dodatkowa r\'ownanie dynamiczne.
Nadwy\.zka przestrzeni ($\varphi > 0$) wok\'o\l{} \'zr\'ode\l{}
jest kompensowana deficytem ($\varphi < 0$) mi\k{e}dzy nimi.
\end{proposition}

\begin{remark}[R\'ownanie dynamiczne $\Ffield$ z~$\Phi$]
R\'ownanie~\eqref{eq:full-field-alt} na $\Phi$ indukuje
efektywne r\'ownanie ewolucji funkcji stanu:
\begin{equation}\label{eq:F-evolution}
    \frac{d\Floc}{dt} = 2\,(\Phi - \PhiZero)\,\frac{d\Phi}{dt}
    = 2\,\PhiZero\,\varphi\,\dot{\Phi}.
\end{equation}
W~po\l{}\k{a}czeniu z~r\'ownaniem na $\Phi$ daje to pe\l{}n\k{a}
dynamik\k{e} funkcji stanu w~granicy continuum ---
zamykaj\k{a}c luk\k{e} ,,brak r\'owna\'n ruchu dla $\Ffield$''
z~modelu minimalnego.
\end{remark}

% =============================================================
\subsection{Formalna niestabilno\'s\'c $\Nzero$
  z~wariacyjnej struktury dzia\l{}ania}\label{ssec:N0-variational}
% =============================================================

Sekcja~\ref{ssec:niestabilnosc} podaje trzy niezale\.zne argumenty
za niestabilno\'sci\k{a} $\Nzero$: (a)~termodynamik\k{e} substratu,
(b)~struktur\k{e} r\'ownania pola, (c)~brak masy efektywnej w~$\Phi=0$.
Poni\.zej zamykamy t\k{e} argumentacj\k{e} twierdzeniem opartym
\textbf{wy\l{}\k{a}cznie} na strukturze dzia\l{}ania~\eqref{eq:action}.

\begin{theorem}[Niestabilno\'s\'c $\Phi = 0$
  z~dzia\l{}ania TGP]\label{thm:N0-action-instability}
Niech dzia\l{}anie TGP ma posta\'c~\eqref{eq:action}
z~miar\k{a} $\sqrt{-g_{\mathrm{eff}}} = c_0\,\psi$%
\footnote{Twierdzenie zachowuje wa\.zno\'s\'c przy poprawionej mierze ---
  degeneracja w~$\Phi=0$ i~brak bariery wynikaj\k{a}
  z~czynnika $\psi^4$ w~sprz\k{e}\.zeniu kinetycznym $K(\psi)$,
  nie z~elementu obj\k{e}to\'sciowego.},
$\psi = \Phi/\PhiZero > 0$.
Definiujemy funkcjona\l{} energii statycznej:
\begin{equation}\label{eq:E-static}
  E[\Phi] = \int d^3x\;\frac{\Phi^4}{\PhiZero^4}\,
  \biggl[\frac{1}{2}\,\frac{(\nabla\Phi)^2}{\PhiZero^2}
  + U\!\biggl(\frac{\Phi}{\PhiZero}\biggr)\biggr],
\end{equation}
z~potencja\l{}em $U(\psi) = \frac{\gamma}{3}\psi^3
- \frac{\gamma}{4}\psi^4$ ($\beta = \gamma$).
W\'owczas:
\begin{enumerate}
  \item \textbf{Degeneracja w~$\Phi = 0$.}
    $E[0] = 0$, a~wszystkie pochodne funkcjonalne
    $\delta^n E/\delta\Phi^n\big|_{\Phi=0} = 0$
    dla $n \le 6$. Punkt $\Phi = 0$ jest
    \textbf{zdegenerowanym} ekstremum, nie minimum.
  \item \textbf{Brak bariery.}
    Niech $\Phi_\epsilon(\mathbf{x}) = \epsilon\,f(\mathbf{x})$,
    $f \ge 0$, $f \in C^\infty_c$, $\|f\|_{L^2} = 1$.
    W\'owczas:
    \begin{equation}\label{eq:E-N0-expansion}
      E[\Phi_\epsilon] = \epsilon^6\,I_6[f]
      + \epsilon^7\,I_7[f] + O(\epsilon^8),
    \end{equation}
    gdzie $I_6 = \frac{1}{2}\int f^4(\nabla f)^2/\PhiZero^6\,d^3x > 0$,
    $I_7 = \frac{\gamma}{3}\int f^7/\PhiZero^7\,d^3x > 0$.
    Ka\.zdy cz\l{}on jest dodatni --- \textbf{brak bariery
    energetycznej} oddzielaj\k{a}cej $\Phi = 0$ od $\Phi > 0$.
    Stan $\Phi = 0$ jest s\'asiedztwie stan\'ow o~dodatniej energii,
    ale nie jest od nich oddzielony.
  \item \textbf{Entropia konfiguracyjna.}
    Miara konfiguracji z~$\Phi = 0$ (dok\l{}adnie, nie przybli\.zenie)
    w~przestrzeni funkcyjnej $L^2$ wynosi zero.
    Argument: zbi\'or $\{\Phi \equiv 0\}$ jest
    jednoelementowy w~$L^2$, wi\k{e}c ma miar\k{e} Wienera
    (lub Gaussowsk\k{a}) r\'own\k{a} zeru.
    Ka\.zda perturbacja ``typowa'' w~sensie
    miarowym daje $\Phi > 0$ w~pewnym regionie.
\end{enumerate}
\end{theorem}

\begin{proof}
\emph{Punkt~1.}\;
Mno\.znik $\psi^4 = (\Phi/\PhiZero)^4$ powoduje, \.ze ca\l{}y
podca\l{}kowy zanurza si\k{e} jak $\Phi^4$ przy $\Phi \to 0$.
Pochodna funkcjonalna:
\[
  \frac{\delta E}{\delta\Phi}\bigg|_{\Phi=0}
  = \lim_{\Phi\to 0}\,\frac{4\Phi^3}{\PhiZero^4}\,
  \bigl[\tfrac{1}{2}(\nabla\Phi)^2/\PhiZero^2 + U\bigr]
  + \ldots = 0,
\]
a~wy\.zsze pochodne r\'ownie\.z zanikaj\k{a}, poniewa\.z
ka\.zdy monomialny cz\l{}on dzia\l{}ania zawiera co najmniej
$\psi^4 \sim \Phi^4$ z~miary i~co najmniej $\psi^2 \sim \Phi^2$
z~potencja\l{}u lub kinetyki --- \l{}\k{a}cznie $\ge \Phi^6$.

\emph{Punkt~2.}\;
Podstawiaj\k{a}c $\Phi = \epsilon f$ do~\eqref{eq:E-static}:
$\Phi^4 = \epsilon^4 f^4$, $(\nabla\Phi)^2 = \epsilon^2(\nabla f)^2$,
$U(\epsilon f/\PhiZero) = \frac{\gamma}{3}\epsilon^3 f^3/\PhiZero^3
- \frac{\gamma}{4}\epsilon^4 f^4/\PhiZero^4$.
Wi\k{e}c:
\[
  E[\epsilon f] = \frac{\epsilon^6}{2\PhiZero^6}\int f^4(\nabla f)^2\,d^3x
  + \frac{\gamma\epsilon^7}{3\PhiZero^7}\int f^7\,d^3x
  + O(\epsilon^8).
\]
Oba cz\l{}ony wiod\k{a}ce s\k{a} \textbf{\'sci\'sle dodatnie}
dla $f > 0$, $f \not\equiv 0$.

\emph{Punkt~3.}\;
Zbi\'or $\{0\} \subset L^2$ jest jednoelementowy.
Ka\.zda miara absolutnie ci\k{a}g\l{}a wzgl\k{e}dem
miary Wienera na $L^2$ przypisuje miar\k{e} zero
ka\.zdemu podzbiorowi sko\'nczonowymiarowemu ---
a~$\{0\}$ jest zerowymiarowy. Jest to \'scis\l{}y
odpowiednik argumentu termodynamicznego
(stw.~\ref{prop:P3-deriv}): zbi\'or konfiguracji
idealnie symetrycznych ma zerow\k{a} wag\k{e}
statystyczn\k{a}.
\end{proof}

\begin{corollary}[Cztero\-stronny argument
  za niestabilno\'sci\k{a} $\Nzero$]\label{cor:N0-quadruple}
Stan $\Nzero$ ($\Phi = 0$) jest niestabilny
z~\textbf{czterech niezale\.znych} przyczyn:
\begin{enumerate}[nosep]
  \item \textbf{Termodynamika substratu}
    (stw.~\ref{prop:P3-deriv}): miara Gibbsa $\mu_\beta$
    przypisuje $\Ffield = 0$ wag\k{e} zero
    w~granicy termodynamicznej.
  \item \textbf{Struktura r\'ownania pola}
    (stw.~\ref{prop:P3-action}): $V''_{\mathrm{eff}}(0) = 0$
    --- brak masy efektywnej w~$\Phi = 0$, punkt
    przegi\k{e}cia.
  \item \textbf{Degeneracja wariacyjna}
    (tw.~\ref{thm:N0-action-instability}): $\Phi = 0$
    jest zdegenerowanym punktem dzia\l{}ania (wszystkie
    pochodne funkcjonalne zerowe), bez bariery energetycznej.
  \item \textbf{Miara konfiguracyjna}
    (tw.~\ref{thm:N0-action-instability}, pkt~3):
    $\{0\}$ ma zerow\k{a} miar\k{e} Wienera
    w~przestrzeni funkcyjnej.
\end{enumerate}
Ka\.zdy z~powy\.zszych argument\'ow jest \textbf{wystarczaj\k{a}cy
sam w~sobie}. Ich zgodno\'s\'c stanowi silny test
wewn\k{e}trznej sp\'ojno\'sci TGP: ten sam wniosek
(niestabilno\'s\'c $\Nzero$) wynika z~termodynamiki,
teorii pola, rachunku wariacyjnego i~teorii miary
--- czterech r\'o\.znych formalizm\'ow.
\end{corollary}

\begin{remark}[Dlaczego \textbf{nie} argument kwantowy]%
\label{rem:N0-not-quantum}
Cz\k{e}sto spotykany argument heurystyczny: ,,w~$\Nzero$
mamy $\hbar(\Phi\!=\!0) = \infty$, wi\k{e}c fluktuacje kwantowe
s\k{a} niesko\'nczone, wymuszaj\k{a}c $\Phi > 0$.''
Argument ten jest \textbf{b\l{}\k{e}dny logicznie}:
w~$\Nzero$ brak metryki, brak czasu, brak r\'ownania
Schr\"odingera --- poj\k{e}cie ,,fluktuacji kwantowej''
nie ma sensu. Wszystkie cztery argumenty
wniosku~\ref{cor:N0-quadruple} celowo unikaj\k{a}
odwo\l{}ania do kwantowej dynamiki: u\.zywaj\k{a}
termodynamiki substratu (1), struktury r\'ownania pola (2),
rachunku wariacyjnego (3) i~teorii miary (4) ---
formalizm\'ow, kt\'ore \textbf{nie wymagaj\k{a}}
istnienia przestrzeni ani czasu.
\end{remark}

\begin{proposition}[Emergencja z~$\Nzero$ zachowuje
  zerow\k{a} sum\k{e}]\label{prop:N0-ZS-emergence}
Niech Wszech\'swiat ``emerguje'' z~$\Nzero$ przez
fluktuacj\k{e} $\Phi = 0 \to \Phi(x) > 0$.
Je\.zeli proces ten zachowuje globaln\k{a}
symetri\k{e} $\Zp$ substratu (aksjomat~\ref{ax:zero}, ZS1),
to wynikowa konfiguracja spe\l{}nia warunek zerowej sumy ZS2
\textbf{do wiod\k{a}cego rz\k{e}du} w~$\delta\Phi/\PhiZero$.

\medskip\noindent
\textbf{Precyzacja:} ZS2 \textbf{nie wynika \'sci\'sle}
z~ZS1 --- s\k{a} to niezale\.zne warunki (uwaga~\ref{rem:zero-precyzacja}).
Jednak w~re\.zimie s\l{}abopolowym ($|\delta\Phi| \ll \PhiZero$)
ZS1 \textbf{implikuje} ZS2 z~dok\l{}adno\'sci\k{a} do korekty
$\mathcal{O}(\delta\Phi^2/\PhiZero)$, co jest wystarczaj\k{a}ce
dla wszystkich zastosowa\'n perturbacyjnych.
\end{proposition}

\begin{proof}
Dow\'od przebiega w~dw\'och krokach: (a)~zwi\k{a}zek mikroskopowy,
(b)~przej\'scie do pola ci\k{a}g\l{}ego.

\emph{Krok~(a): Bilans domen.}
ZS1 gwarantuje $\int_\Sigma\langle\hat{s}\rangle\sqrt{h}\,d^3x = 0$.
W~fazie z\l{}amanej symetrii ($T < T_c$, $\Phi > 0$)
substrat dzieli si\k{e} na domeny $\mathcal{D}_+$
($\langle\hat{s}\rangle = +v + \eta_+$) i~$\mathcal{D}_-$
($\langle\hat{s}\rangle = -v + \eta_-$), gdzie $\eta_\pm$ s\k{a}
fluktuacjami rz\k{e}du $\mathcal{O}(1/\sqrt{N_B})$.
Warunek ZS1 wymaga:
\begin{equation}\label{eq:ZS1-balance}
  v\,\mathrm{Vol}(\mathcal{D}_+) - v\,\mathrm{Vol}(\mathcal{D}_-)
  + \mathcal{O}(\eta) = 0
  \quad\Rightarrow\quad
  \mathrm{Vol}(\mathcal{D}_+) = \mathrm{Vol}(\mathcal{D}_-)
  + \mathcal{O}(1/\sqrt{N_B}).
\end{equation}

\emph{Krok~(b): Przej\'scie do $\Phi$.}
Pole $\Phi(x) = \langle\hat{s}^2\rangle_{\mathrm{cg}}$ jest
jednorodne na skalach du\.zych wzgl\k{e}dem rozmiaru domeny:
$\Phi = \PhiZero + \delta\Phi$, gdzie $\PhiZero = v^2 + \sigma^2$
($\sigma^2$ --- wariancja fluktuacji). Perturbacja
$\delta\Phi(\mathbf{x})$ jest skorelowana z~\textbf{g\k{e}sto\'sci\k{a}
domen materii}: tam, gdzie domeny s\k{a} g\k{e}\'sciej upakowane,
$\langle\hat{s}^2\rangle > \PhiZero$ (wi\k{e}cej ,,poszarpanego''
interfejsu domen), daj\k{a}c $\delta\Phi > 0$; w~regionach
o~wi\k{e}kszej jednorodnej orientacji $\langle\hat{s}^2\rangle
< \PhiZero$, daj\k{a}c $\delta\Phi < 0$.

R\'owno\'s\'c obj\k{e}to\'sci $\mathcal{D}_+$ i~$\mathcal{D}_-$
(r\'ow.~\ref{eq:ZS1-balance}) implikuje symetri\k{e}
g\k{e}sto\'sci upakowania: \'srednio tyle samo materii
po ka\.zdej stronie. Ca\l{}ka:
\begin{equation}
  \int_\Sigma \delta\Phi\,d^3x = 0
  + \mathcal{O}\!\left(\frac{\delta\Phi^2}{\PhiZero}\right).
\end{equation}
Korekta metryczna ($\sqrt{h} \approx 1 + \frac{3}{2}\delta\Phi/\PhiZero$,
uwaga~\ref{rem:ZS2-self-ref}) modyfikuje wynik o~cz\l{}on rz\k{e}du
$\delta\Phi^2/\PhiZero$, daj\k{a}c:
\begin{equation}
  \int_\Sigma \delta\Phi\,\sqrt{h}\,d^3x
  = -\frac{3}{2\PhiZero}\int(\delta\Phi)^2\,d^3x
  + \mathcal{O}(\delta\Phi^3/\PhiZero^2).
\end{equation}
Dok\l{}adny ZS2 wymaga zatem niewielkiej korekty \'sredniej
koordynatowej $\langle\delta\Phi\rangle$, rz\k{e}du
$\delta\Phi^2/\PhiZero \sim 10^{-10}$ dla perturbacji
kosmologicznych --- zaniedbywalnej w~ka\.zdym zastosowaniu
fizycznym.
\end{proof}

\begin{remark}[Dlaczego ZS2 jest niezale\.znym aksjomatem]%
\label{rem:ZS2-niezalezny-formal}
Wy\.zsze rz\k{e}dy korekty metrycznej ($\delta\Phi^2/\PhiZero$)
uniemo\.zliwiaj\k{a} \textbf{\'scis\l{}e} wyprowadzenie ZS2 z~ZS1.
Fizycznie: ZS1 kontroluje \emph{orientacj\k{e}} domen
(pierwszy moment $\langle\hat{s}\rangle$), a~ZS2 kontroluje
\emph{amplitud\k{e}} (drugi moment $\langle\hat{s}^2\rangle$).
Oba warunki s\k{a} zatem \textbf{niezale\.znymi selekcjami}
na przestrzeni konfiguracji, sp\'ojnymi ale nie r\'ownowa\.znymi.
Pe\l{}ny ZS2 musi by\'c postulowany oddzielnie.
\end{remark}

% =============================================================
\subsection{Perturbacje kosmologiczne --- zarys}\label{ssec:perturb}
% =============================================================

W~kosmologii jednorodnej (sekcja~\ref{ssec:friedmann})
$\Phi = \bar{\Phi}(t)$ jest jednorodne. Perturbacje liniowe
wok\'o\l{} tego t\l{}a daj\k{a} testowaln\k{a} predykcj\k{e}.

\begin{proposition}[Zmodyfikowane r\'ownanie Poissona]\label{prop:poisson-mod}
W~cechowaniu Newtona ($ds^2 = -(1+2\Psi)\,c_0^2 dt^2
+ a^2(1-2\Phi_N)\,\delta_{ij}dx^i dx^j$) pole
$\Phi = \bar{\Phi}(t) + \delta\Phi(t,\mathbf{x})$ daje
zmodyfikowane r\'ownanie Poissona:
\begin{equation}\label{eq:poisson-mod}
    \frac{k^2}{a^2}\,\Phi_N
    = -\frac{4\pi\,G_{\mathrm{eff}}}{c_0^2}\,\delta\rho
      - \frac{1}{2\PhiZero}\,\frac{k^2}{a^2}\,\delta\Phi,
\end{equation}
gdzie cz\l{}on $\delta\Phi$ jest \textbf{dodatkow\k{a}
poprawką TGP} nieobecn\k{a} w~OTW.
\end{proposition}

Perturbacja pola przestrzenno\'sci spe\l{}nia:
\begin{equation}\label{eq:delta-Phi-evol}
    \ddot{\delta\Phi} + 3H\,\dot{\delta\Phi}
    + \biggl(\frac{k^2}{a^2}
    + m_{\mathrm{eff}}^2\biggr)\,\delta\Phi
    = -q\,\PhiZero\,\delta\rho + \text{cz\l{}ony metryczne},
\end{equation}
gdzie $m_{\mathrm{eff}}^2 = m_{\mathrm{sp}}^2 = 3\gamma - 2\beta = \gamma$
(przy $\beta = \gamma$) jest \textbf{przestrzenn\k{a}} mas\k{a}
Yukawy (stw.~\ref{prop:vacuum-stability}), \textbf{nie}
kosmologiczn\k{a} mas\k{a} wolnego staczania.
Perturbacja $\delta\Phi(t,\vec{x})$ widzi potencja\l{}
\emph{statyczny} wok\'o\l{} t\l{}a $\bar\Phi$, wi\k{e}c
masa ekranuj\k{a}ca jest $m_{\mathrm{sp}}^2 > 0$
(uwaga~\ref{rem:two-masses}).
Masa kosmologiczna $m_{\mathrm{cosmo}}^2 = 4c_0^2\gamma/3$
(r.~\eqref{eq:m-cosmo-exact}) kontroluje jedynie
\emph{jednorodn\k{a}} ewolucj\k{e} $\bar\Phi(t)$,
nie perturbacje przestrzenne.

\begin{theorem}[Efektywna sta\l{}a grawitacyjna zale\.zna
  od skali]\label{thm:Geff-k}
Kombinacja r\'owna\'n~\eqref{eq:poisson-mod}
i~\eqref{eq:delta-Phi-evol} daje efektywn\k{a} sta\l{}\k{a}
grawitacyjn\k{a}:
\begin{equation}\label{eq:Geff-k}
  \boxed{
    G_{\mathrm{eff}}(k, a)
    = G_0\,\biggl(1 + \frac{2\,\alpha_{\mathrm{eff}}^2}
    {1 + (a\,m_{\mathrm{eff}}/k)^2}\biggr),
  }
\end{equation}
gdzie $\alpha_{\mathrm{eff}} = q\,\PhiZero/(4\pi)$ jest
bezwymiarow\k{a} sta\l{}\k{a} sprz\k{e}\.zenia.
\begin{itemize}[nosep]
  \item Skale $k/a \gg m_{\mathrm{eff}}$: grawitacja
        wzmocniona o~czynnik $(1 + 2\alpha_{\mathrm{eff}}^2)$.
  \item Skale $k/a \ll m_{\mathrm{eff}}$: standardowa
        grawitacja $G_0$.
  \item Skala przej\'scia $k/a \sim m_{\mathrm{eff}}$:
        \textbf{strefa testowalna} przez dane o~wzro\'scie
        struktur (DESI, Euclid).
\end{itemize}
\end{theorem}

\begin{corollary}[Gravitational slip]\label{cor:grav-slip}
W~parametryzacji zmodyfikowanej grawitacji $(\mu, \Sigma)$, gdzie
$\mu$ kontroluje dynamik\k{e} materii (r\'ownanie Poissona),
a~$\Sigma$ kontroluje soczewkowanie (r\'ownanie propagacji \'swiat\l{}a):
\begin{equation}\label{eq:mu-sigma}
  \mu(k,a) \equiv \frac{G_{\mathrm{eff}}}{G_0}
  = 1 + \frac{2\,\alpha_{\mathrm{eff}}^2}{1 + (a\,m_{\mathrm{eff}}/k)^2},
  \qquad
  \Sigma(k,a)
  = 1 + \frac{\alpha_{\mathrm{eff}}^2}{1 + (a\,m_{\mathrm{eff}}/k)^2}.
\end{equation}
R\'o\.znica mi\k{e}dzy $\mu$ i~$\Sigma$ wynika z~tego, \.ze pole
$\delta\Phi$ sprz\k{e}ga si\k{e} ,,podw\'ojnie'' z~dynamik\k{a}
materii (bezpo\'srednie sprz\k{e}\.zenie + wk\l{}ad metryczny),
ale tylko ,,raz'' z~soczewkowaniem (wy\l{}\k{a}cznie metryczny).

\textbf{Gravitational slip:}
$\eta \equiv \Sigma/\mu < 1$ zawsze. Dla $k/a \gg m_{\mathrm{eff}}$:
\begin{equation}\label{eq:slip}
  \eta = \frac{1 + \alpha_{\mathrm{eff}}^2}{1 + 2\alpha_{\mathrm{eff}}^2}
  \approx 1 - \alpha_{\mathrm{eff}}^2,
\end{equation}
co jest testowalnym przewidywaniem TGP (Euclid, LSST).
Estymacja numeryczna (skrypt \texttt{growth\_factor\_tgp.py}):
$\alpha_{\mathrm{eff}} \sim 10^{-26}$ dla $\PhiZero \sim 25$
--- poni\.zej progu detekcji. Jednake\.z $\alpha_{\mathrm{eff}}$
zale\.zy kwadratowo od $\PhiZero$: je\.zeli $\PhiZero \gg 25$,
dewiacja mo\.ze by\'c mierzalna.
\end{corollary}

\begin{remark}[Wyniki numeryczne: funkcja transferu i~wzrost struktur]%
\label{rem:transfer-results}
Numeryczne ca\l{}kowanie r\'ownania wzrostu z~$G_{\mathrm{eff}}(k,a)$
(skrypt \texttt{unified\_perturbation\_theory.py}) daje:
\begin{center}
\begin{tabular}{cccc}
\toprule
$\alpha_{\mathrm{eff}}$ & $f\sigma_8(z{=}0)$ &
$P_{\mathrm{TGP}}/P_{\Lambda\mathrm{CDM}}$ przy $k{=}0.1\,h/$Mpc &
$\eta(z{=}0)$ \\
\midrule
$0$ (GR)    & $0.428$ & $1.000$ & $1.000$ \\
$0.01$      & $0.428$ & $1.000$ & $0.9999$ \\
$0.05$      & $0.429$ & $1.010$ & $0.9975$ \\
$0.10$      & $0.433$ & $1.040$ & $0.990$ \\
\midrule
\textbf{fizyczne} $\sim10^{-26}$ & $0.428$ & $1 + \mathcal{O}(10^{-52})$
  & $1 - \mathcal{O}(10^{-52})$ \\
\bottomrule
\end{tabular}
\end{center}

Dla fizycznego $\alpha_{\mathrm{eff}} = q\PhiZero/(4\pi)
\approx 3.7\times 10^{-26}$ wszelkie odchylenia od $\Lambda$CDM
s\k{a} $\sim\alpha_{\mathrm{eff}}^2 \sim 10^{-51}$
--- poni\.zej czu\l{}o\'sci Euclid ($\sigma_\mu \sim 10^{-2}$)
o~$49$~rz\k{e}d\'ow wielko\'sci.

\textbf{Wniosek}: perturbacje kosmologiczne TGP s\k{a}
\textbf{nierozr\'o\.znialne} od $\Lambda$CDM dla $\PhiZero \sim 25$.
Testem TGP na poziomie perturbacji kosmologicznych by\l{}oby
niezale\.zne ograniczenie $\PhiZero$ z~silnego pola
(np.~ringdown, por.~dod.~\ref{app:ringdown}).
\end{remark}

% --- Formal LSS predictions and stability (M2 deliverable) ---
% tgp_perturbations_lss.tex -- M2 deliverable: Cosmological perturbations and LSS
% Integrates into sek08_formalizm.tex (after ssec:perturb) or sek07_predykcje.tex
% Cross-refs: thm:Geff-k, cor:grav-slip, prop:ghost-free-MS, eq:Geff-k

\subsection{Large-scale structure predictions}%
\label{ssec:lss-predictions}

The scale-dependent effective gravitational constant
(tw.~\ref{thm:Geff-k}) and the gravitational slip
(wn.~\ref{cor:grav-slip}) yield concrete, falsifiable
predictions for large-scale structure observables.
This subsection quantifies the predictions and compares
them with existing data from BOSS, DESI, and forthcoming
surveys (Euclid, LSST, SKA).

% ----- Stability ----- %
\subsubsection{Stability analysis}%
\label{sssec:lss-stability}

\begin{proposition}[Complete perturbative stability]%
\label{prop:lss-stability}
Cosmological perturbations of the TGP scalar field $\Phi$
around the homogeneous FRW background $\bar\Phi(t)$
satisfy \textbf{all} stability conditions:
\begin{enumerate}[(i),nosep]
  \item \textbf{No-ghost}: the kinetic coefficient in the
    Mukhanov--Sasaki sector is
    $Q_s = \psi_{\rm bg}^4 > 0$ for all $\psi_{\rm bg} > 0$
    (stw.~\ref{prop:ghost-free-MS});
  \item \textbf{Gradient stability}: the sound speed is
    $c_s^2 = c_0^2 > 0$ (identically the speed of light);
  \item \textbf{Tachyonic stability}: the spatial Yukawa mass satisfies
    $m_{\mathrm{sp}}^2 = \gamma > 0$
    (stw.~\ref{prop:vacuum-stability}), ensuring exponential
    suppression of $\delta\Phi$ on sub-Compton scales;
  \item \textbf{Cosmological mass}: the homogeneous mass
    $m_{\rm cosmo}^2 = \tfrac{4}{3}\,c_0^2\gamma > 0$
    corresponds to damped oscillations of $\bar\Phi$
    around the vacuum~$\Phi_0$, driven to rest by
    Hubble friction --- this is the relaxation to
    $\Lambda$CDM, not an instability.
\end{enumerate}
\end{proposition}

\noindent\textit{Numerical verification:}
Script \texttt{tgp\_perturbations\_formal.py} confirms all four
conditions over the range $\psi \in (0, 3]$ (18/18 tests PASS).

% ----- Growth factor ----- %
\subsubsection{Linear growth factor and $f\sigma_8$}%
\label{sssec:growth-lss}

The linear density contrast $\delta_m$ obeys:
\begin{equation}\label{eq:growth-ode}
  D'' + \bigl(2 + \tfrac{H'}{H}\bigr)\,D'
  = \tfrac{3}{2}\,\Omega_m(a)\,
    \frac{G_{\mathrm{eff}}(k,a)}{G_0}\;D,
\end{equation}
where primes denote $d/d\ln a$. The scale dependence of
$G_{\mathrm{eff}}$ (r.~\eqref{eq:Geff-k}) induces, in principle,
a scale-dependent growth rate $f(k,z) = d\ln D / d\ln a$.

\begin{table}[h]
\centering\small
\caption{Linear growth factor and growth rate $f\sigma_8$:
TGP vs.\ $\Lambda$CDM at $k = 0.1\;\mathrm{Mpc}^{-1}$.
At $\Phi_0 \approx 25$, deviations are below $10^{-50}$.}
\label{tab:growth-comparison}
\begin{tabular}{@{}ccccc@{}}
\toprule
$z$ & $D_{\Lambda\mathrm{CDM}}$ & $D_{\mathrm{TGP}}$ &
$|\Delta D/D|$ & $f\sigma_8$ \\
\midrule
0.0 & 1.0000 & 1.0000 & $< 10^{-50}$ & 0.428 \\
0.3 & 0.8528 & 0.8528 & $< 10^{-50}$ & 0.474 \\
0.5 & 0.7689 & 0.7689 & $< 10^{-50}$ & 0.476 \\
0.7 & 0.6964 & 0.6964 & $< 10^{-50}$ & 0.462 \\
1.0 & 0.6067 & 0.6067 & $< 10^{-50}$ & 0.431 \\
1.5 & 0.4956 & 0.4956 & $< 10^{-50}$ & 0.374 \\
2.0 & 0.4172 & 0.4172 & $< 10^{-50}$ & 0.324 \\
\bottomrule
\end{tabular}
\end{table}

The growth-rate observable $f\sigma_8(z) = f(z)\,\sigma_8\,D(z)$
agrees with $\Lambda$CDM to machine precision at $\Phi_0 \approx 25$,
giving $\chi^2 = 9.86$ for 14 combined data points
(BOSS DR12 + 6dFGS + VIPERS + eBOSS + DESI DR1 RSD).

\begin{remark}[Observational bound on $\alpha_{\mathrm{eff}}$]%
\label{rem:alpha-bound}
A $\chi^2$ scan over $\alpha_{\mathrm{eff}}$ using the combined
$f\sigma_8$ compilation yields the $2\sigma$ upper bound:
\begin{equation}\label{eq:alpha-bound}
  \alpha_{\mathrm{eff}} < 0.17 \quad (2\sigma).
\end{equation}
The TGP prediction $\alpha_{\mathrm{eff}} = q\,\Phi_0/(4\pi)
\approx 3.7 \times 10^{-26}$ lies $25$~orders of magnitude below
this threshold. Scale-dependent growth would become detectable
only for $\alpha_{\mathrm{eff}} \gtrsim 0.01$, corresponding to
$\Phi_0 \gtrsim 10^{24}$.
\end{remark}

% ----- Modified gravity parameterization ----- %
\subsubsection{Modified gravity parameterization $(\mu, \Sigma, \eta)$}%
\label{sssec:mu-sigma-eta}

In the standard parameterized post-Friedmann framework:
\begin{align}
  \mu(k,a) &\equiv \frac{G_{\mathrm{eff}}}{G_0}
  = 1 + \frac{2\,\alpha_{\mathrm{eff}}^2}
        {1 + (a\,m_{\mathrm{sp}}/k)^2},
  \label{eq:mu-lss} \\[4pt]
  \Sigma(k,a) &= 1 + \frac{\alpha_{\mathrm{eff}}^2}
        {1 + (a\,m_{\mathrm{sp}}/k)^2},
  \label{eq:sigma-lss} \\[4pt]
  \eta(k,a) &\equiv \frac{\Sigma}{\mu}
  = \frac{1 + \alpha_{\mathrm{eff}}^2/(1+x^2)}
         {1 + 2\alpha_{\mathrm{eff}}^2/(1+x^2)},
  \qquad x = a\,m_{\mathrm{sp}}/k.
  \label{eq:eta-lss}
\end{align}
For $k/a \gg m_{\mathrm{sp}}$ (small scales):
$\eta \approx 1 - \alpha_{\mathrm{eff}}^2$.
For $k/a \ll m_{\mathrm{sp}}$ (large scales):
$\eta \to 1$ (GR recovered).

At $\Phi_0 \approx 25$:
$\mu - 1 \sim \Sigma - 1 \sim \eta - 1 \sim 10^{-51}$,
far below the projected sensitivity of
Euclid ($\sigma_\mu \sim 0.02$, Amendola et al.\ 2018)
and LSST ($\sigma_\Sigma \sim 0.03$).

\begin{table}[h]
\centering\small
\caption{Modified gravity parameters at
$k = 0.1\;\mathrm{Mpc}^{-1}$ for TGP ($\Phi_0 \approx 25$).
All deviations from GR are undetectable.}
\label{tab:mu-sigma-eta}
\begin{tabular}{@{}cccc@{}}
\toprule
$z$ & $\mu - 1$ & $\Sigma - 1$ & $\eta - 1$ \\
\midrule
0.0 & $< 10^{-51}$ & $< 10^{-51}$ & $< 10^{-51}$ \\
0.5 & $< 10^{-51}$ & $< 10^{-51}$ & $< 10^{-51}$ \\
1.0 & $< 10^{-51}$ & $< 10^{-51}$ & $< 10^{-51}$ \\
2.0 & $< 10^{-51}$ & $< 10^{-51}$ & $< 10^{-51}$ \\
\bottomrule
\end{tabular}
\end{table}

% ----- Inflationary sector ----- %
\subsubsection{Inflationary predictions}%
\label{sssec:inflation-lss}

The Mukhanov--Sasaki equation in TGP reads (prop.~\ref{prop:ghost-free-MS}):
\begin{equation}\label{eq:MS-TGP}
  v_k'' + \bigl[c_s^2\,k^2 - z''/z\bigr]\,v_k = 0,
  \qquad z = a\,\psi_{\rm bg}^2,\quad c_s^2 = c_0^2.
\end{equation}
In the slow-roll regime (Dodatek~G), the spectral index and
tensor-to-scalar ratio are:
\begin{equation}\label{eq:ns-r-TGP}
  n_s = 1 - \frac{2}{N_e}, \qquad
  r = \frac{16}{N_e}, \qquad
  f_{\rm NL}^{\rm local} \approx -\frac{1}{N_e},
\end{equation}
where $N_e$ is the number of $e$-folds.
For $N_e \in [55, 65]$:
$n_s \in [0.964, 0.969]$ (Planck $1\sigma$),
$r \in [0.25, 0.29]$,\footnote{%
  The standard relation $r = 16\varepsilon$ assumes
  an unmodified tensor sector. In TGP, the tensor sector
  ($\sigma_{ab}$) has status Szkic; the actual $r_{\rm TGP}$
  may differ. See sek.~\ref{sssec:sigma-status-map}.}
$|f_{\rm NL}| \lesssim 0.02$ (undetectable by Planck).

% ----- Falsifiability ----- %
\subsubsection{Falsifiability criteria}%
\label{sssec:falsifiability-lss}

\begin{table}[h]
\centering\small
\caption{Falsifiability of TGP in the perturbation sector.}
\label{tab:falsifiability}
\begin{tabular}{@{}lp{5.5cm}l@{}}
\toprule
\textbf{Observable} & \textbf{Condition to falsify TGP} &
\textbf{Survey} \\
\midrule
$\mu(k,z)$ & $|\mu - 1| > 0.02$ at fixed $\Phi_0 \approx 25$
& Euclid \\
$\Sigma(k,z)$ & $|\Sigma - 1| > 0.03$
& LSST \\
$f\sigma_8(z)$ & Scale-dependent anomaly at
$\Delta(f\sigma_8)/f\sigma_8 > 1\%$
& DESI, Euclid \\
$n_s$ & $n_s \notin [0.95, 0.98]$ (no viable $N_e$)
& CMB-S4 \\
$c_{\rm GW}$ & $|c_{\rm GW}/c_0 - 1| > 10^{-15}$
& LIGO/ET \\
\bottomrule
\end{tabular}
\end{table}

\begin{remark}[Perturbation sector conclusion]%
\label{rem:perturbation-conclusion}
At $\Phi_0 \approx 25$ (the value fixed by
$\Lambda_{\mathrm{eff}} \approx \Lambda_{\mathrm{obs}}$),
TGP is \textbf{operationally equivalent to $\Lambda$CDM}
in the cosmological perturbation sector:
$G_{\mathrm{eff}} = G_0$ to $O(10^{-51})$,
all stability conditions are satisfied,
and the inflationary predictions lie within the
Planck $1\sigma$ band for $N_e \sim 57$--$65$.

The distinct predictions of TGP lie
\textbf{outside} the linear perturbation regime:
\begin{enumerate}[(a),nosep]
  \item modified dark energy equation of state $w_{\rm DE}(z)$
    (\S\ref{ssec:de-oszacowanie}, likelihood H6:
    $\Delta\chi^2 = +1.8$ vs.\ $\Lambda$CDM);
  \item GW dispersion at the Yukawa mass scale
    $m_{\mathrm{sp}} \sim 10^{-26}\;\mathrm{m}^{-1}$
    (Table~\ref{tab:gw-summary});
  \item $N$-body dynamics: Earnshaw violation, chaos
    suppression, three-body force (Appendix~\ref{app:nbody}).
\end{enumerate}
These are the arenas where TGP can be distinguished from GR.
\end{remark}


% =============================================================
\subsubsection{Soczewkowanie grawitacyjne w~TGP}\label{sssec:lensing}
% =============================================================

\begin{proposition}[K\k{a}t ugi\k{e}cia \'swiat\l{}a w~TGP]%
\label{prop:deflection-angle}
Foton poruszaj\k{a}cy si\k{e} w~metryce efektywnej TGP
$ds^2 = -(c_0^2/\psi)\,dt^2 + \psi\,\delta_{ij}\,dx^i dx^j$
doznaje ugi\k{e}cia:
\begin{equation}\label{eq:deflection-TGP}
  \hat\alpha = 2\int_{-\infty}^{+\infty}
  \frac{\nabla_\perp \Phi_{\rm lens}}{c_0^2}\,dl,
  \qquad
  \Phi_{\rm lens} = \frac{c_0^2}{2}\ln\psi = \frac{c_0^2}{2}\ln\frac{\Phi}{\PhiZero}.
\end{equation}
W~granicy s\l{}abego pola ($\psi \approx 1 + \delta\psi$):
$\Phi_{\rm lens} \approx (c_0^2/2)\,\delta\psi$ i~k\k{a}t ugi\k{e}cia
redukuje si\k{e} do:
\begin{equation}\label{eq:deflection-weak}
  \hat\alpha \approx \frac{4\,G_0\,M}{c_0^2\,b}
  \cdot \frac{1}{2}(1 + \gamma_{\rm PPN}) = \frac{4\,G_0\,M}{c_0^2\,b},
\end{equation}
poniewa\.z $\gamma_{\rm PPN} = 1$ w~TGP (uwaga~\ref{rem:PPN-Will}).
Wynik jest \textbf{identyczny z~OTW} w~granicy s\l{}abego pola.
\end{proposition}

\begin{proof}[Szkic dowodu]
Geodezyjna zerowa ($ds^2 = 0$) w~metryce TGP daje
efektywny indeks za\l{}amania $n(r) = \sqrt{\psi(r)} \cdot \psi(r)^{1/2}
= \psi(r)$.
K\k{a}t ugi\k{e}cia w~przybli\.zeniu Born:
$\hat\alpha = -\int \nabla_\perp \ln n\,dl = -\int \nabla_\perp \ln\psi\,dl$.
Z~$\ln\psi \approx \delta\psi$ w~s\l{}abym polu
i~$\delta\psi = -\kappa\,\Phi_N/c_0^2$
(gdzie $\Phi_N$ jest potencja\l{}em Newtona)
otrzymujemy standardowy wynik GR.
\end{proof}

\begin{corollary}[Potencja\l{} soczewkowania kosmologicznego]%
\label{cor:lensing-potential}
Na skalach kosmologicznych soczewkowanie kontrolowane jest przez
$\Sigma(k,a)$ (wn.~\ref{cor:grav-slip}):
\begin{equation}\label{eq:lensing-convergence}
  \kappa_{\rm lens}(\hat{\mathbf{n}})
  = \frac{3\,\Omega_m\,H_0^2}{2\,c_0^2}\int_0^{\chi_s}
  \frac{(\chi_s - \chi)\,\chi}{\chi_s\,a(\chi)}\;\Sigma(k,a)\;\delta(\chi,\hat{\mathbf{n}})\;d\chi,
\end{equation}
gdzie $\chi$ jest odleg\l{}o\'sci\k{a} wsp\'o\l{}poruszaj\k{a}c\k{a},
$\chi_s$ odleg\l{}o\'sci\k{a} \'zr\'od\l{}a, $\delta$ kontrastem g\k{e}sto\'sci.
Dla fizycznego $\alpha_{\mathrm{eff}} \sim 10^{-26}$:
$\Sigma = 1 + \mathcal{O}(10^{-51})$
--- soczewkowanie TGP jest \textbf{identyczne z~$\Lambda$CDM}
z~precyzj\k{a} $49$~rz\k{e}d\'ow wielko\'sci powy\.zej
progu Euclid/LSST.

\textbf{Falsyfikacja}: je\.zeli przysz\l{}e pomiary $\eta = \Sigma/\mu$
wyka\.z\k{a} $|\eta - 1| > 10^{-2}$, jest to sp\'ojne z~TGP
\textbf{tylko} je\.zeli $\PhiZero \gg 25$ (wymagaj\k{a}c rewizji
identyfikacji $\PhiZero = N_f^2$).
\end{corollary}

% --- Analytical lensing, Shapiro, Yukawa, strong-field (M3 deliverable) ---
% tgp_lensing_geff.tex -- M3 deliverable: Gravitational lensing in TGP
% Integrates after sssec:lensing in sek08_formalizm.tex
% Extends qualitative predictions from sek07 with full analytical derivation

\subsubsection{Analytical lensing in $g_{\mathrm{eff}}$:
  PPN connection and Yukawa corrections}%
\label{sssec:lensing-analytical}

The preceding subsection (\S\ref{sssec:lensing}) established the
deflection angle and lensing convergence \emph{qualitatively}.
Here we close the analytical derivation, connecting:
(i)~the point-mass lensing to PPN parameters ($\gamma_{\rm PPN} = 1$),
(ii)~the Yukawa correction from the scalar propagator $\delta\Phi$,
(iii)~the Shapiro time delay (Cassini bound), and
(iv)~the post-post-Newtonian (ppN) regime where TGP's exponential
metric diverges from the Schwarzschild solution.

% --- Born integral via angular substitution ---
\paragraph{Exact Born integral for exponential metric.}
For the TGP metric
$ds^2 = -(c_0^2/\psi)\,dt^2 + \psi\,\delta_{ij}\,dx^i dx^j$
with $\psi = e^{2U/c_0^2}$, $U = GM/r$ (unsigned),
the effective refractive index is
$n(r) = \psi(r) = e^{2U/c_0^2}$.
The deflection angle in the Born approximation reads
\begin{equation}\label{eq:born-integral}
  \hat\alpha
  = -\int_{-\infty}^{+\infty}
    \nabla_\perp \!\ln\psi\;dl
  = \int_{-\infty}^{+\infty}
    \frac{2\,G\,M\,b}{c_0^2\,(b^2 + l^2)^{3/2}}\;dl,
\end{equation}
where $b$ is the impact parameter and $l$ the coordinate along
the unperturbed ray ($r^2 = b^2 + l^2$).
Substituting $l = b\tan\theta$ collapses the integral to
\begin{equation}\label{eq:deflection-exact}
  \hat\alpha
  = \frac{2\,G\,M}{c_0^2\,b}
    \int_{-\pi/2}^{\pi/2}\!\cos\theta\;d\theta
  = \frac{4\,G\,M}{c_0^2\,b},
\end{equation}
\textbf{identically} the GR result. The key feature is that
$\ln\psi = 2U/c_0^2$ is \emph{exact} (not a weak-field expansion),
so the Born integral closes analytically at all orders in~$U$.
Numerical integration confirms eq.~\eqref{eq:deflection-exact}
to $< 10^{-6}$ relative precision
(skrypt \texttt{tgp\_lensing\_formal.py}, 8/8 PASS).

% --- PPN consistency ---
\paragraph{PPN consistency.}
The general PPN deflection angle is
$\hat\alpha = \tfrac{1}{2}(1+\gamma_{\rm PPN})\times 4GM/(c_0^2 b)$.
With $\gamma_{\rm PPN} = 1$ (eq.~\eqref{eq:gamma-beta}),
this reproduces eq.~\eqref{eq:deflection-exact}. Together with
$\beta_{\rm PPN} = 1$, $\alpha_{1,2} = 0$
(eq.~\eqref{eq:alpha12}), TGP passes all solar-system lensing
and timing tests identically to GR.

% --- Shapiro delay ---
\paragraph{Shapiro time delay.}
A signal traversing the gravitational field of mass~$M$
with closest approach~$b$ acquires a delay:
\begin{equation}\label{eq:shapiro-TGP}
  \Delta t = \frac{1}{c_0}\int\!\bigl[n(r) - 1\bigr]\,dl
           = \frac{1}{c_0}\int\!\bigl[e^{2U/(c_0^2 r)} - 1\bigr]\,dl
           \approx \frac{4\,G\,M}{c_0^3}\,\ln\!\frac{4r_1 r_2}{b^2},
\end{equation}
where $r_{1,2}$ are the emitter and receiver distances.
The linearization error is $O(U^2/c_0^4) \sim 10^{-12}$
(verified numerically at the Cassini configuration:
$|\Delta t_{\rm exact}/\Delta t_{\rm weak} - 1| = 3 \times 10^{-7}$).
The Cassini constraint $|\gamma - 1| = (2.1 \pm 2.3)\times 10^{-5}$
(Bertotti et al.\ 2003) is satisfied automatically:
TGP gives $\gamma - 1 = 0$ exactly.

% --- Yukawa correction ---
\paragraph{Yukawa correction from scalar propagator.}
The scalar perturbation $\delta\Phi$ around a point mass
satisfies $(\nabla^2 - m_{\rm sp}^2)\,\delta\Phi = -q\,\Phi_0\,\rho$,
yielding a Yukawa profile:
\begin{equation}\label{eq:yukawa-lens}
  \delta\psi_{\rm Yuk}(r)
  = 2\,\alpha_{\rm eff}^2\,\frac{e^{-m_{\rm sp}\,r}}{m_{\rm sp}\,r},
  \qquad
  \alpha_{\rm eff} = \frac{q\,\Phi_0}{4\pi}
  \approx 3.7 \times 10^{-26}.
\end{equation}
The Compton wavelength $\lambda_{\rm sp} = 2\pi/m_{\rm sp}
\approx 5.6\;\mathrm{Gpc}$ exceeds all astrophysical
scales, so the Yukawa factor is \emph{unsuppressed}
($m_{\rm sp} r \ll 1$ at solar, galactic, and cluster scales).
However, the amplitude is $\alpha_{\rm eff}^2 \sim 10^{-51}$:
\begin{equation}\label{eq:delta-alpha-yukawa}
  \frac{\Delta\hat\alpha_{\rm Yuk}}{\hat\alpha_{\rm GR}}
  \sim \alpha_{\rm eff}^2 \sim 10^{-51},
\end{equation}
which is $\sim\!49$ orders of magnitude below the Cassini bound.

% --- ppN correction ---
\paragraph{Post-post-Newtonian (ppN) correction.}
The exponential metric deviates from Schwarzschild at $O(U^2)$:
\begin{align}
  g_{rr}^{\rm TGP}  &= 1 + 2U/c_0^2 + 2U^2/c_0^4 + \ldots,
  \nonumber\\
  g_{rr}^{\rm Schw} &= 1 + 2U/c_0^2 + 4U^2/c_0^4 + \ldots.
  \label{eq:ppN-grr}
\end{align}
The Born integral (eq.~\ref{eq:born-integral}) is insensitive to
this difference because $\ln\psi$ is an exact logarithm;
the ppN correction enters only through \emph{ray-tracing}
(non-Born) effects. At the solar limb:
\begin{equation}\label{eq:ppN-estimate}
  \frac{\Delta\hat\alpha_{\rm ppN}}{\hat\alpha_{\rm GR}}
  \sim \left(\frac{G\,M_\odot}{R_\odot\,c_0^2}\right)^{\!2}
  \approx 4.5 \times 10^{-12},
\end{equation}
corresponding to $\sim\!8\;\mathrm{pico\text{-}arcsec}$
--- far below astrometric precision for decades to come.

% --- Strong-field predictions ---
\paragraph{Strong-field predictions (indicative).}
Extrapolating the exponential metric to the strong-field regime
yields qualitatively different predictions from Schwarzschild:

\begin{center}\small
\begin{tabular}{@{}llll@{}}
\toprule
\textbf{Observable} & \textbf{GR (Schwarzschild)} &
\textbf{TGP (exponential)} & $\Delta/\mathrm{GR}$ \\
\midrule
Event horizon   & $r_H = r_s = 2GM/c_0^2$  & brak (exp. $> 0$ zawsze) & --- \\[2pt]
Photon sphere   & $r_{\rm ph} = \tfrac{3}{2}\,r_s$ & $r_{\rm ph} = r_s$ & $-33\%$ \\[2pt]
Shadow $b_{\rm crit}$ & $3\sqrt{3}\;GM/c_0^2$ & $r_s\sqrt{e}$ & $-37\%$ \\
\bottomrule
\end{tabular}
\end{center}

\noindent
\textbf{Caveat.}
These strong-field numbers assume the exponential metric
$\psi = e^{2U/c_0^2}$ extends to $U/c_0^2 \sim 0.5$.
In TGP, the actual strong-field solution may differ
(see \S\ref{sec:czarne} for the TGP black-hole model,
where the interior approaches the $\Nzero$ limit).
A~complete strong-field lensing calculation requires the
full radial profile $\psi(r)$ from the TGP field equation
--- this remains an open problem.

% --- Summary table ---
\begin{table}[h]
\centering\small
\caption{Gravitational lensing: TGP vs.\ GR summary.}
\label{tab:lensing-summary}
\begin{tabular}{@{}lp{4.5cm}p{3.5cm}@{}}
\toprule
\textbf{Test} & \textbf{TGP prediction} & \textbf{Status} \\
\midrule
Solar deflection   & $\hat\alpha = 4GM_\odot/(c_0^2 R_\odot)$
  = $1.7513''$ & $\equiv$ GR \\[3pt]
Shapiro delay      & $\gamma = 1$ exactly
  & Cassini: $\checkmark$ \\[3pt]
Yukawa correction  & $\Delta\hat\alpha/\hat\alpha \sim 10^{-51}$
  & Undetectable \\[3pt]
ppN correction     & $\Delta\hat\alpha \sim 8$ pico-arcsec
  & Undetectable \\[3pt]
Cosmic shear       & $\Sigma = 1 + O(10^{-51})$
  & $\equiv$ $\Lambda$CDM \\[3pt]
BH shadow          & $\Delta b_{\rm crit} \sim -37\%$
  & Open (needs full soln.) \\
\bottomrule
\end{tabular}
\end{table}

\begin{remark}[Lensing sector conclusion]\label{rem:lensing-conclusion}
In the weak-field regime, TGP's lensing predictions are
\textbf{identical to GR} --- this is a structural consequence
of the exponential metric and $\gamma_{\rm PPN} = 1$.
The Yukawa correction ($\alpha_{\rm eff}^2 \sim 10^{-51}$)
and ppN correction ($\sim 10^{-12}$) are both unobservable.

The \textbf{only} distinctive TGP lensing prediction lies
in the \emph{strong-field} regime: the exponential metric
never develops an event horizon, the photon sphere is smaller
by 33\%, and the shadow is smaller by 37\%.
These are in principle testable by the Event Horizon Telescope,
but a rigorous comparison requires the full TGP compact-object
solution (status: open problem, \S\ref{sec:czarne}).
\end{remark}


% =============================================================
\subsubsection{R\'ownanie Mukhanova--Sasakiego dla TGP}\label{sssec:MS-TGP}
% =============================================================

Poni\.zej wyprowadzamy r\'ownanie perturbacji kosmologicznych
bezpo\'srednio z~dzia\l{}ania TGP, uwzgl\k{e}dniaj\k{a}c
\textbf{niestandardowy cz\l{}on kinetyczny} $K(\psi) = \psi^4$
wynikaj\k{a}cy z~substratu (stw.~\ref{prop:substrate-action}).

\begin{proposition}[Pe\l{}ne r\'ownanie TGP w~FRW z~gradientami przestrzennymi]%
\label{prop:TGP-FRW-full}
Dzia\l{}anie TGP w~metryce FRW ($ds^2 = -c_0^2\,dt^2 + a^2(t)\,d\mathbf{x}^2$):
\[
  S_{\mathrm{TGP}} = \int dt\,d^3x\;a^3\psi^4
  \!\left[
    -\frac{\dot\psi^2}{2c_0^2}
    + \frac{c_0^2(\nabla\psi)^2}{2a^2}
    - U(\psi)
  \right]\!,
\]
gdzie $\psi \equiv \Phi/\PhiZero$, generuje pe\l{}ne r\'ownanie Eulera--Lagrange'a:
\begin{equation}\label{eq:TGP-FRW-spatial}
  \boxed{
    \ddot\psi + 3H\dot\psi + \frac{2\dot\psi^2}{\psi}
    - \frac{c_0^2}{a^2}
      \!\left[\nabla^2\!\psi + \frac{2(\nabla\psi)^2}{\psi}\right]
    = c_0^2\,W(\psi),
  }
\end{equation}
gdzie operator przestrzenny $\mathcal{D}_{\rm kin}[\psi]
= \nabla^2\psi + 2(\nabla\psi)^2/\psi$
jest operatorem TGP z~$\alpha = 2$ (stw.~\ref{thm:D-uniqueness}),
a~$W(\psi) = \frac{1}{\psi^4}\frac{d}{d\psi}[\psi^4 U(\psi)]$
(stw.~\ref{prop:W-from-U}).
\end{proposition}

\begin{proof}
Obliczamy sk\l{}adniki EL z~$\mathcal{L} = a^3\psi^4
[-\dot\psi^2/(2c_0^2) + c_0^2(\nabla\psi)^2/(2a^2) - U(\psi)]$:
\begin{align*}
  \partial_t\!\left(\frac{\partial\mathcal{L}}{\partial\dot\psi}\right)
  &= -\frac{a^3}{c_0^2}\,\partial_t[\psi^4\dot\psi]
   = -\frac{a^3}{c_0^2}\bigl[3H\psi^4\dot\psi + 4\psi^3\dot\psi^2 + \psi^4\ddot\psi\bigr],\\
  \partial_i\!\left(\frac{\partial\mathcal{L}}{\partial(\partial_i\psi)}\right)
  &= a\,c_0^2\,\partial_i[\psi^4\partial_i\psi]
   = a\,c_0^2\,\psi^3\bigl[\psi\nabla^2\!\psi + 4(\nabla\psi)^2\bigr],\\
  \frac{\partial\mathcal{L}}{\partial\psi}
  &= a^3\psi^3\!\left[-\frac{2\dot\psi^2}{c_0^2}
     + \frac{2c_0^2(\nabla\psi)^2}{a^2}
     - 4U - \psi U'\right]\!.
\end{align*}
R\'ownanie EL ($\partial_t\partial_{\dot\psi}\mathcal{L}
+ \partial_i\partial_{(\partial_i\psi)}\mathcal{L}
- \partial_\psi\mathcal{L} = 0$), podzielone przez $a^3\psi^3$
i~pomno\.zone przez $-c_0^2/\psi$, daje~\eqref{eq:TGP-FRW-spatial}.
Cz\l{}on $4U/\psi + U' = W(\psi)$ pochodzi od miary $\psi^4$
(stw.~\ref{prop:W-from-U}).
\end{proof}

\begin{proposition}[Zlinearyzowane r\'ownanie perturbacji i~zmienna
  Mukhanova--Sasakiego TGP]\label{prop:MS-TGP}
Niech $\psi = \psi_{\rm bg}(t) + \delta\psi(t,\mathbf{x})$,
$|\delta\psi| \ll \psi_{\rm bg}$.
Linearyzacja r\'ownania~\eqref{eq:TGP-FRW-spatial}
wok\'o\l{} jednorodnego t\l{}a $\psi_{\rm bg}$ (dla kt\'orego
$\nabla\psi_{\rm bg} = 0$) daje w~przestrzeni Fouriera:
\begin{equation}\label{eq:perturb-TGP-time}
  \ddot{\delta\psi}_{\mathbf{k}}
  + \underbrace{\left(3H + \frac{4\dot\psi_{\rm bg}}{\psi_{\rm bg}}\right)}_{\text{t\l{}umienie TGP}}
    \dot{\delta\psi}_{\mathbf{k}}
  + \left[\frac{c_0^2 k^2}{a^2} + m_{\rm pert}^2(t)\right]
    \delta\psi_{\mathbf{k}} = 0,
\end{equation}
gdzie masa efektywna perturbacji wynosi:
\begin{equation}\label{eq:mpert-TGP}
  m_{\rm pert}^2(t)
  = -\frac{2\dot\psi_{\rm bg}^2}{\psi_{\rm bg}^2}
    - c_0^2\,W'(\psi_{\rm bg}).
\end{equation}
Definiuj\k{a}c \textbf{zmienn\k{a} Mukhanova--Sasakiego TGP}:
\begin{equation}\label{eq:MS-variable}
  \boxed{v_{\mathbf{k}} \equiv z(\eta)\cdot\delta\psi_{\mathbf{k}},
  \qquad z(\eta) = a(\eta)\,\psi_{\rm bg}^2(\eta),}
\end{equation}
gdzie $\eta$ jest czasem konforemnym ($d\eta = dt/a$),
r\'ownanie~\eqref{eq:perturb-TGP-time} przyjmuje posta\'c kanoniczn\k{a}:
\begin{equation}\label{eq:MS-TGP}
  \boxed{
    v_{\mathbf{k}}'' + \left[c_0^2 k^2 - \frac{z''}{z}\right]v_{\mathbf{k}} = 0,
  }
\end{equation}
gdzie prim oznacza $d/d\eta$, a~,,pompa'' Mukhanova--Sasakiego TGP:
\begin{equation}\label{eq:MS-pump}
  \frac{z''}{z}
  = \frac{a''}{a}
    + \frac{2\psi_{\rm bg}''}{\psi_{\rm bg}}
    + 4\,\frac{a'}{a}\frac{\psi_{\rm bg}'}{\psi_{\rm bg}}
    + 2\!\left(\frac{\psi_{\rm bg}'}{\psi_{\rm bg}}\right)^{\!2}.
\end{equation}
\end{proposition}

\begin{proof}
\textbf{Krok~1 (linearyzacja).}
Rozwijamy $2\dot\psi^2/\psi$ do pierwszego rz\k{e}du:
\[
  \frac{2(\dot\psi_{\rm bg}+\dot{\delta\psi})^2}{\psi_{\rm bg}+\delta\psi}
  \approx
  \frac{2\dot\psi_{\rm bg}^2}{\psi_{\rm bg}}
  + \frac{4\dot\psi_{\rm bg}}{\psi_{\rm bg}}\dot{\delta\psi}
  - \frac{2\dot\psi_{\rm bg}^2}{\psi_{\rm bg}^2}\delta\psi.
\]
Przestrzenny operator TGP na jednorodnym tle linearyzuje si\k{e} do $\nabla^2\delta\psi$
(cz\l{}on kwadratowy $2(\nabla\delta\psi)^2/\psi_{\rm bg}$ jest rz\k{e}du $\delta\psi^2$).
Odejmuj\k{a}c r\'ownanie t\l{}a, w~przestrzeni Fouriera
($\nabla^2 \to -k^2$) otrzymujemy~\eqref{eq:perturb-TGP-time}--\eqref{eq:mpert-TGP}.

\textbf{Krok~2 (zmienna MS).}
Podstawiamy $\delta\psi_{\mathbf{k}} = v_{\mathbf{k}}/z$ do r\'ownania
w~czasie konforemnym (mno\.z\k{a}c przez $a^2$):
\[
  \delta\psi_{\mathbf{k}}'' + \!\left(2\frac{a'}{a} + 4\frac{\psi_{\rm bg}'}{\psi_{\rm bg}}\right)
  \!\delta\psi_{\mathbf{k}}' + \bigl[c_0^2 k^2 + a^2 m_{\rm pert}^2\bigr]\delta\psi_{\mathbf{k}} = 0.
\]
Czynnik przy pochodnej pierwszej: $2(a'/a) + 4(\psi_{\rm bg}'/\psi_{\rm bg}) = 2z'/z$,
wi\k{e}c zmienna $v_{\mathbf{k}} = z\,\delta\psi_{\mathbf{k}}$ eliminuje
t\l{}umienie. Standardowy rachunek daje~\eqref{eq:MS-TGP}.
Wzg\l{}\k{e}dem standardowej inflacji ($z_{\rm std} = a$) czynnik
$\psi_{\rm bg}^2$ jest \textbf{modyfikacj\k{a} TGP} wynikaj\k{a}c\k{a}
ze sprz\k{e}\.zenia kinematycznego $K(\psi) = \psi^4$ substratu.
Wzór~\eqref{eq:MS-pump} wynika z~$z'/z = a'/a + 2\psi_{\rm bg}'/\psi_{\rm bg}$
przez~$(z''/z) = (z'/z)' + (z'/z)^2$~, co daje cztery sk\l{}adniki
w~r.~\eqref{eq:MS-pump}.
\end{proof}

\begin{corollary}[Pierwotne spektrum mocy TGP]\label{cor:TGP-power-spectrum}
W~re\.zimie zamro\.zonego pola
($\psi_{\rm bg} \approx \psi_0 = \mathrm{const}$,
tzn.~$\psi_{\rm bg}'/\psi_{\rm bg} \approx 0$),
obowi\k{a}zuj\k{a}cym gdy Hubble zamra\.za ewolucj\k{e} $\psi$
(patrz uwaga~\ref{rem:psi1-status}):
\begin{equation}\label{eq:MS-TGP-frozen}
  z \approx a\,\psi_0^2,
  \qquad
  \frac{z''}{z} \approx \frac{a''}{a}
  \;\xrightarrow{\text{de~Sitter}}\;
  \frac{2}{\eta^2}.
\end{equation}
R\'ownanie~\eqref{eq:MS-TGP} redukuje si\k{e} do standardowego
r\'ownania Mukhanova--Sasakiego dla skalara minimalne sprzężonego,
z~rozwi\k{a}zaniem w~przestrzeni p\l{}askiej:
\[
  v_{\mathbf{k}}(\eta)
  = \frac{e^{-ic_0 k\eta}}{\sqrt{2c_0 k}}
    \!\left(1 - \frac{i}{c_0 k\eta}\right).
\]
Pierwotne spektrum mocy zaburzenia $\delta\psi$:
\begin{equation}\label{eq:TGP-power-spectrum}
  \boxed{
    \mathcal{P}_{\delta\psi}(k)
    \equiv \frac{k^3}{2\pi^2}\left|\frac{v_{\mathbf{k}}}{z}\right|^2_{\!\!\!k=aH}
    = \frac{H^2}{4\pi^2\,c_0\,\psi_0^4}\,,
  }
\end{equation}
gdzie $H$ oceniamy przy wyj\'sciu modu za horyzont $k = aH$.

\medskip
\noindent\textbf{Uwagi}:
\begin{enumerate}[nosep]
  \item Czynnik $\psi_0^{-4}$ jest \textbf{modyfikacj\k{a} TGP}
    amplitudy widma wzgl\k{e}dem standardowego
    $\mathcal{P}_{\rm std} = H^2/(4\pi^2)$.
    Dla $\psi_0 \approx 1$ amplitudy s\k{a} identyczne --- TGP
    odzyska standardow\k{a} kosmologi\k{e} w~epoce z~$\Phi = \PhiZero$.
  \item Nachylenie widma $n_s - 1 = -4\varepsilon_{\rm H}$
    (de~Sitter), gdzie $\varepsilon_{\rm H} = -\dot H/H^2$ ---
    identyczne ze standardow\k{a} inflacj\k{a} dla zamro\.zonego~$\psi$.
  \item Korekta TGP do nachylenia pojawia si\k{e} przy
    uwzgl\k{e}dnieniu ewolucji~$\psi_{\rm bg}$:
    \begin{equation}\label{eq:ns-TGP}
      (n_s - 1)_{\rm TGP}
      = -4\varepsilon_{\rm H} - 4\varepsilon_\psi,
      \qquad
      \varepsilon_\psi
      \equiv -\frac{\dot\psi_{\rm bg}}{H\psi_{\rm bg}}
      = \mathcal{O}(H_0/H),
    \end{equation}
    co dla ery inflacyjnej ($H \gg H_0$) daje $\varepsilon_\psi \ll 1$
    --- poprawka poni\.zej progu Planck~($\sigma_{n_s} \sim 10^{-3}$).
  \item \textbf{Most do CMB}: zmierzony przez Planck~2018
    $n_s = 0.9649 \pm 0.0042$ ogranicza $\varepsilon_{\rm H}$
    (i~$\varepsilon_\psi$) identycznie jak w~modelu~$\Lambda$CDM,
    o~ile $\psi_{\rm bg}$ jest zamro\.zone podczas inflacji.
    Niezerowe~$\varepsilon_\psi$ sta\l{}oby si\k{e} testowalne
    w~danych Simons~Observatory i~CMB-S4.
\end{enumerate}
\end{corollary}

\begin{remark}[F7: Precyzja $n_s$ --- analityczne zamkni\k{e}cie]\label{rem:F7-ns-precision}
\textbf{Status: zamkni\k{e}te analitycznie (v36)}.

\smallskip\noindent
Trzy poziomy aproksymacji indeksu spektralnego TGP:

\begin{center}\small
\begin{tabular}{@{}lll@{}}
\toprule
\textbf{Poziom} & \textbf{Wyra\.zenie} & \textbf{Warto\'s\'c ($N_e=60$)} \\
\midrule
0.\ (de~Sitter) & $n_s = 1$ & $1.0000$ \\
1.\ (zamro\.zone $\psi$, quasi-dS) & $n_s - 1 = -4\varepsilon_{\rm H}^{\rm TGP}$ &
  $\approx -2/N_e \approx -0.0333$ \\
2.\ (korekcja $\varepsilon_\psi$) & $n_s - 1 = -4\varepsilon_{\rm H} - 4\varepsilon_\psi$ &
  $+\mathcal{O}(N_e^{-2}) \approx \pm 0.0002$ \\
\bottomrule
\end{tabular}
\end{center}

\smallskip\noindent
\textbf{Kluczowe ustalenia:}
\begin{enumerate}[nosep]
  \item Poziom~1 jest \textbf{identyczny z~wynikiem Starobinskiego}
    (prop.~\ref{prop:spectral-index}, dod.~G):
    $n_s = 1 - 2/N_e - 9/(4N_e^2) \approx 0{,}9660$ ($N_e=60$)
    --- w~$1\sigma$ Planck~2018.
    Wynika to wprost z~faktu, \.ze dla zamro\.zonego~$\psi$
    r\'ownanie~\eqref{eq:MS-TGP} redukuje si\k{e} do kanonicznego MS.
  \item Korekcja $\varepsilon_\psi = -\dot\psi_{\rm bg}/(H\psi_{\rm bg})
    = \mathcal{O}(H_0/H_{\rm inf}) \approx 10^{-60}$
    podczas inflacji~--- $\mathcal{O}(N_e^{-2})$ wzgl\k{e}dem poziomu~1.
    Poni\.zej progu obserwacyjnego CMB-S4 ($\sigma_{n_s} \sim 0{,}002$).
  \item \textbf{TGP jest zdegenerowane z~inflacj\k{a} Starobinskiego}
    przy precyzji Plancka~2018. Rozr\'o\.znienie mo\.zliwe tylko
    przez dodatkowe obserwable TGP: mod oddechowy GW, oscylacje $G_{\rm eff}(t)$,
    $r \approx 0{,}003$ (pogranicze dost\k{e}pno\'sci LiteBIRD).
  \item Numeryczne $n_s$ z~pe\l{}nego t\l{}a $z(\tau) = a(\tau)\psi_{\rm bg}^2(\tau)$
    potwierdzi Starobinskiego w~granicach precyzji numerycznej.
    Skrypt \texttt{p73} lub \texttt{ex90} pozostaje \textbf{programem
    weryfikacyjnym} (nie zmieni wyniku fizycznego).
\end{enumerate}
\end{remark}

\begin{remark}[Status Fazy~III.A]\label{rem:IIIa-status}
Wyprowadzenie MS~TGP zamknęło lukę wskazaną
w~\texttt{PLAN\_ROZWOJU\_v1.md} (zadanie~III.A):
\begin{itemize}[nosep]
  \item R\'ownanie~\eqref{eq:TGP-FRW-spatial}:
    pe\l{}na postać 3+1D z~operatorem $\mathcal{D}_{\rm kin}$ --- \textbf{nowe};
  \item R\'ownanie~\eqref{eq:perturb-TGP-time}:
    t\l{}umienie $(3H + 4\dot\psi/\psi)$ --- \textbf{modyfikacja TGP};
  \item R\'ownanie~\eqref{eq:MS-TGP}: posta\'c kanoniczna $v''+(c_0^2k^2-z''/z)v=0$
    z~$z = a\psi_{\rm bg}^2$ --- \textbf{nowe};
  \item Wzór~\eqref{eq:TGP-power-spectrum}: amplituda $\propto \psi_0^{-4}$
    i~korekta nachylenia~\eqref{eq:ns-TGP} --- \textbf{nowe}.
\end{itemize}
Pe\l{}na analiza pierwotnych zaburze\'n (tensor\'ow, niega\-ussow\-sko\'sci,
spektrum $\sigma_{ab}$) pozostaje programem d\l{}ugoterminowym
(por.~ssec.~\ref{ssec:tensor-substrate}).
\end{remark}

% -------------------------------------------------------------------
\subsubsection{Stabilno\'s\'c i~ghost-check perturbacji MS-TGP}\label{sssec:ghost-check}
% -------------------------------------------------------------------

\begin{proposition}[Brak ghost\'ow i~stabilno\'s\'c gradientowa perturbacji TGP]%
\label{prop:ghost-free-MS}
\statuslabel{Propozycja}: przy za\l{}o\.zeniu, \.ze cz\l{}on kinetyczny
$K(\psi) = \psi^4 > 0$ i~jednorodne t\l{}o $\psi_{\rm bg}(t) > 0$.

\medskip\noindent
Dzia\l{}anie drugiego rz\k{e}du (S₂) dla perturbacji skalara
$v_k = a\psi_{\rm bg}^2\,\delta\psi_k$ ma posta\'c:
\begin{equation}\label{eq:S2-MS-TGP}
  S_2 = \frac{1}{2}\int\frac{d\tau\,d^3k}{(2\pi)^3}
  \left[|v_k'|^2 - \left(c_0^2 k^2 - \frac{z''}{z}\right)|v_k|^2\right],
\end{equation}
gdzie $z = a\psi_{\rm bg}^2$ i~$z'' \equiv d^2z/d\tau^2$.

\medskip\noindent
\textbf{Warunek no-ghost}: cz\l{}on kinetyczny $|v_k'|^2$ ma
\textbf{w\l{}a\'sciwy znak} ($+1$), co wyklucza duchy (Ostrogradski-ghost).
\begin{equation}\label{eq:no-ghost}
  Q_s = \frac{z^2}{a^2} = \psi_{\rm bg}^4 > 0
  \qquad\Leftarrow\quad \psi_{\rm bg} > 0.
\end{equation}
Warunek $\psi_{\rm bg} > 0$ jest gwarantowany przez tw.~\ref{thm:vacuum-source}
($\Phi \equiv 0$ nie istnieje).

\medskip\noindent
\textbf{Stabilno\'s\'c gradientowa}: wsp\'o\l{}czynnik przy $k^2$ wynosi
$c_s^2 = c_0^2 > 0$ (pr\k{e}dko\'s\'c d\'zwi\k{e}ku rownoznaczna $c_0$).
Brak niestabilno\'sci gradientowej:
\begin{equation}\label{eq:gradient-stable}
  c_s^2 = c_0^2 > 0
  \qquad \Rightarrow \quad \text{stabilno\'s\'c gradientowa dla wszystkich $k$}.
\end{equation}

\medskip\noindent
\textbf{Sprawdzenie w~granicy super-horyzontu}:
dla $k^2 \ll z''/z$ (d\l{}uga fala) r\'ownanie~\eqref{eq:MS-TGP}
daje rozwi\k{a}zania $v_k \propto z$ (mod zamrożony)
i~$v_k \propto z\int d\tau/z^2$ (mod zanikaj\k{a}cy),
sp\'ojne ze standardowym schematem inflatacyjnym.
\end{proposition}

\begin{remark}[Status i~co brakuje --- C.2A]\label{rem:C2A-status}
Powy\.zszy ghost-check jest \textbf{analityczny} i~obejmuje:
\begin{itemize}[nosep]
  \item Brak ducha w~dzia\l{}aniu S₂ $\checkmark$ (dok\l{}adne)
  \item Stabilno\'s\'c gradientowa $c_s^2 = c_0^2 > 0$ $\checkmark$ (dok\l{}adne)
  \item Rozwiązania super-horyzontu $\checkmark$ (analityczne)
\end{itemize}
\textbf{Zamkni\k{e}te analitycznie} (F7, v36):
\begin{itemize}[nosep]
  \item $n_s = 0{,}966$ (Starobinsky, identyczne z~TGP dla zamro\.zonego~$\psi$;
    uw.~\ref{rem:F7-ns-precision}): \checkmark
  \item $r \approx 0{,}003 = 12/N_e^2$; $\sigma_{ab}=0$ podczas inflacji
    (rem.~\ref{rem:min-observables}): \checkmark
\end{itemize}
\textbf{Brakuje} (program C.2B --- d\l{}ugoterminowy):
\begin{itemize}[nosep]
  \item Wk\l{}ad non-Gaussianity $f_{\rm NL}$ z~korekcji $\varepsilon_\psi^2$
    ($f_{\rm NL} \sim 10^{-4}$, por.\ tab.\ \ref{rem:min-observables})
  \item Numeryczne ISW i~$\sigma_8$ z~$G_{\rm eff}(z) = G_N/\psi(z)$
    (skrypt \texttt{p73}, pr.\ C.2B)
\end{itemize}
\end{remark}

% -------------------------------------------------------------------
\subsubsection{Minimalne obserwable skalarne MS-TGP}\label{sssec:min-observables}
% -------------------------------------------------------------------

\begin{remark}[Predykcje i~ograniczenia sektora perturbacji --- C.2C]%
\label{rem:min-observables}
Poni\.zsza tabela zbiera minimalne wielko\'sci obserwowalne wynikaj\k{a}ce
z~r\'ownania Mukhanova--Sasakiego TGP (r.~\eqref{eq:MS-TGP})
przy za\l{}o\.zeniu powolnego toczenia ($\varepsilon_\psi \ll 1$)
i~warunku pocz\k{a}tkowego Buncha--Daviesa.

\begin{center}\small
\begin{tabular}{@{}p{2.8cm}p{3.8cm}p{2.3cm}p{3.5cm}@{}}
\toprule
\textbf{Obserwabla} & \textbf{Wyra\.zenie TGP} & \textbf{Warto\'s\'c TGP} &
  \textbf{Ograniczenie obserwacyjne} \\
\midrule
Nachylenie spektrum skalarne $n_s$
  & $n_s{-}1 = {-}2/N_e$
    (Starobinsky;
    r.~\eqref{eq:ns-TGP},
    uw.~\ref{rem:F7-ns-precision})
  & $\approx 0{,}966$
    ($N_e=60$, Starobinsky)
  & $0{,}9649 \pm 0{,}0042$
    (Planck 2018) \textbf{[1$\sigma$]} \\[6pt]
Parametr tensora~$r$
  & $r = 12/N_e^2$
    (Starobinsky, stw.~\ref{prop:r-tensor})
  & $\approx 0{,}003$
    ($N_e=60$)
  & $r < 0{,}036$
    (BK18+BAO)
    \textbf{[sp\'ojne]} \\[6pt]
Non-Gaussianity $f_{\rm NL}$
  & $f_{\rm NL}^{\rm loc} \sim
      \varepsilon_\psi^2 / c_s^2$
  & $\sim 10^{-4}$
  & $|f_{\rm NL}| < 5$
    (Planck 2018) \\[6pt]
Efekt ISW (zintegrowany)
  & $\Delta T_{\rm ISW} \propto
      \int\!\dot{G}_{\rm eff}/G_{\rm N}\,d\chi$
  & do obliczenia
    (skrypt $p73$)
  & korelacja ISW--LSS
    (ACT/SPT) \\[6pt]
Wzrost struktur $\sigma_8$
  & $\sigma_8 \propto D(a)$
    z~$G_{\rm eff}(a) = G_N/\psi(a)$
  & do obliczenia
  & $\sigma_8 = 0{,}811
    \pm 0{,}006$
    (DES Y3) \\[6pt]
\bottomrule
\end{tabular}
\end{center}

\medskip\noindent
\textbf{Uwagi}:
\begin{enumerate}[nosep]
  \item \textbf{$n_s$ w~1$\sigma$}: TGP nale\.zy dok\l{}adnie do klasy Starobinskiego
    (uw.~\ref{rem:universality}): $n_s = 1 - 2/N_e - 9/(4N_e^2) \approx 0{,}966$ ($N_e=60$)
    --- w~1$\sigma$ Planck~2018.
    Korekcja TGP od $\varepsilon_\psi = \mathcal{O}(H_0/H_{\rm inf})$ wynosi
    $|\Delta n_s| \sim 0{,}0002 \ll \sigma_{n_s}^{\rm Planck}$ ---
    nieodró\.znialna od Starobinskiego (uw.~\ref{rem:F7-ns-precision}).
    Skrypt \texttt{p73} jest programem weryfikacyjnym, nie zmieni wyniku.
  \item \textbf{$r$ sp\'ojne z~danymi}: TGP nale\.zy do klasy Starobinsky'ego
    (uw.~\ref{rem:universality}), dla kt\'orej $\varepsilon = 3/(4N_e^2) \ll 1/N_e$
    i~$r = 16\varepsilon = 12/N_e^2 \approx 0{,}003$ --- spe\l{}nia
    ograniczenie $r < 0{,}036$ (BK18).
    Poprzednia warto\'s\'c $r \approx 0{,}27$ pochodzi\l{}a
    z~b\l{}\k{e}dnego zastosowania relacji kanonicznej $r = 16\varepsilon_\psi$,
    gdzie $\varepsilon_\psi = (1-n_s)/2 = 1/N_e \approx 0{,}017$ jest
    parametrem \emph{eta} (dominuj\k{a}cym w~Starobinskim),
    nie parametrem \emph{epsilon}.
    Pole $\sigma_{ab}$ nie generuje pierwotnych tensor\'ow: podczas inflacji
    substrat jest w~fazie symetrycznej ($K_{ab} \propto \delta_{ab}$),
    wi\k{e}c $\sigma_{ab} = 0$ (def.~\ref{def:sigma-ab}).
    Po~podgrzaniu: $\sigma_{ab} \neq 0$ tylko od anizotropowych \'zr\'ode\l{}
    materii (uk\l{}ady binarne, itp.), sp\'ojnie z~obs.~GW170817.
  \item \textbf{$f_{\rm NL}$}: dla $c_s = c_0$, $\varepsilon_\psi \approx 0{,}017$:
    $f_{\rm NL} \sim \varepsilon_\psi^2 \sim 3 \times 10^{-4}$ ---
    nieobserwowalnie ma\l{}e (zgodne z~danymi).
  \item \textbf{ISW i~$\sigma_8$}: wymagaj\k{a} pe\l{}nej integracji
    kosmologicznej z~$G_{\rm eff}(z) = G_N/\psi(z)$
    (skrypt: \texttt{p73} lub nowy \texttt{p77}).
    Napi\k{e}cie $\sigma_8$: TGP przewiduje
    $G_{\rm eff}(z > 2) \approx G_N$ (ψ zamro\.zone),
    wi\k{e}c efekt kumuluje si\k{e} tylko dla $z \lesssim 2$
    --- mo\.ze z\l{}agodzi\'c napi\k{e}cie $S_8$.
\end{enumerate}

\medskip\noindent
\textbf{Podsumowanie M$_{\rm obs}$ z~sektora perturbacji}:
ka\.zda z~wielko\'sci $\{n_s,\,r,\,f_{\rm NL},\,\text{ISW},\,\sigma_8\}$
stanowi osobn\k{a} obserwacj\k{e} wynikaj\k{a}c\k{a} z~dwóch
parametr\'ow Warstwy~II ($\PhiZero$, $a_\Gamma$),
powi\k{e}kszaj\k{a}c iloraz $M_{\rm obs}/N_{\rm param} \geq 6 + 5 = 11$
(przy $N_{\rm param} = 2$, za\l{}. GL, stw.~\ref{prop:psi-ini-derived}).
\end{remark}

\subsubsection{Fale grawitacyjne w~TGP}\label{sssec:gw}

W~TGP geometria wynika z~pola skalarnego $\Phi$. Fale grawitacyjne
s\k{a} zatem \textbf{propaguj\k{a}cymi perturbacjami}
$\delta\Phi(t,\mathbf{x})$ wok\'o\l{} jednorodnego t\l{}a.

\begin{proposition}[Relacja dyspersji fal $\delta\Phi$]\label{prop:GW-dispersion}
Pr\'o\.zniowe ($\delta\rho = 0$) perturbacje pola $\Phi$ w~t\l{}e
kosmologicznym spe\l{}niaj\k{a} relacj\k{e} dyspersji:
\begin{equation}\label{eq:GW-dispersion}
  \omega^2 = c_0^2\!\left(\frac{k^2}{a^2} + m_{\mathrm{sp}}^2\right),
  \qquad m_{\mathrm{sp}}^2 = 3\gamma - 2\beta = \gamma \;\;(\beta=\gamma).
\end{equation}
\begin{enumerate}[nosep]
  \item Pr\k{e}dko\'s\'c fazowa:
    $v_\phi = c_0\sqrt{1 + (a\,m_{\mathrm{sp}}/k)^2}$.
    Dla $k/a \gg m_{\mathrm{sp}}$: $v_\phi \to c_0$ ---
    zgodna z~ograniczeniem GW170817
    ($|c_{\mathrm{GW}} - c_0|/c_0 < 10^{-15}$).
  \item Obci\k{e}cie niskoenergetyczne:
    $f_{\mathrm{cut}} = c_0\,m_{\mathrm{sp}}/(2\pi)$.
    Poni\.zej $f_{\mathrm{cut}}$ perturbacje zanikaj\k{a} eksponencjalnie
    (profil Yukawa, nie propagacja).
  \item \textbf{Polaryzacja}: $\delta g_{ij}/g_{ij}
    = 2\delta\Phi/\PhiZero\,\delta_{ij}$
    --- wy\l{}\k{a}cznie mod oddechowy (\emph{breathing}).
    Brak tensorowych mod\'ow $h_+, h_\times$ na poziomie liniowym.
\end{enumerate}
\end{proposition}

\begin{proof}
R\'ownanie~\eqref{eq:delta-Phi-evol} dla $\delta\rho = 0$
i~$\bar\varphi \approx 1$:
\[
  \ddot{\delta\Phi} + 3H\dot{\delta\Phi}
  + \left(\frac{k^2}{a^2} + m_{\mathrm{sp}}^2\right)\delta\Phi = 0.
\]
Masy u\.zywamy przestrzennej ($m_{\mathrm{sp}}^2 = 3\gamma - 2\beta$,
stw.~\ref{prop:vacuum-stability}), nie kosmologicznej, gdy\.z perturbacje
przestrzenne (fale) podlegaj\k{a} statycznej masie efektywnej.
W~re\.zimie WKB ($\omega \gg H$) pomijamy t\l{}umienie Hubble'a
i~relacja dyspersji jest~\eqref{eq:GW-dispersion}.

Polaryzacja: metryka TGP (stw.~\ref{prop:conformal-unique}) ma
$g_{ij} = (\Phi/\PhiZero)\delta_{ij}$, wi\k{e}c
$\delta g_{ij} = (\delta\Phi/\PhiZero)\delta_{ij}$ ---
perturbacja jest izotropowa (mod oddechowy, $E(2)$-spin $0$),
nie tensorowa (spin $2$).
\end{proof}

\begin{theorem}[Zakaz mod\'ow tensorowych dla metryki jednoskalowej]%
\label{thm:no-tensor}
Je\.zeli metryka efektywna TGP zale\.zy wy\l{}\k{a}cznie
od jednego pola skalarnego $\Phi$ (bez pochodnych $\partial_\mu\Phi$)
i~ma posta\'c diagonaln\k{a}, izotropow\k{a} przestrzennie:
\begin{equation}\label{eq:metric-one-scalar}
  g_{\mu\nu} = \mathrm{diag}\bigl(-f(\Phi),\;h(\Phi),\;h(\Phi),\;h(\Phi)\bigr),
\end{equation}
to liniowe perturbacje $\delta\Phi$ generuj\k{a}
\textbf{wy\l{}\k{a}cznie} mod oddechowy (spin~$0$).
Mody tensorowe ($h_+$, $h_\times$, spin~$2$) nie istniej\k{a}
na poziomie liniowym.
\end{theorem}

\begin{proof}
Perturbacja metryki indukowana przez $\delta\Phi$:
\begin{equation}
  \delta g_{\mu\nu}
  = \frac{\partial g_{\mu\nu}}{\partial\Phi}\;\delta\Phi
  = \mathrm{diag}\bigl(-f'(\Phi),\;h'(\Phi),\;h'(\Phi),\;h'(\Phi)\bigr)
    \;\delta\Phi.
\end{equation}
Cz\k{e}\'s\'c przestrzenna:
$\delta g_{ij} = h'(\Phi)\,\delta\Phi\;\delta_{ij}$
--- \textbf{izotropowa} (proporcjonalna do $\delta_{ij}$).
Jest to mod oddechowy ($E(2)$-spin~$0$, perturbacja \'sladowa).

Mody tensorowe ($h_+$, $h_\times$) wymagaj\k{a} cz\k{e}\'sci
\textbf{bez\'sladowej transwersalnej}: $\delta g_{ij}^{TT}$ takiego,
\.ze $\delta g_{ii}^{TT} = 0$ i~$\partial_i \delta g_{ij}^{TT} = 0$.
Poniewa\.z $\delta g_{ij} \propto \delta_{ij}$, mamy
$\delta g_{ij}^{TT} \equiv 0$ to\.zsamo\'sciowo ---
\textbf{brak tensorowych stopni swobody}.

Argument ten jest niezale\.zny od tego, czy $f = h$
(metryka konforemna) czy $f \neq h$ (metryka
antypodaln\k{a}, jak w~TGP). Kluczowe jest, \.ze
$g_{ij}$ zale\.zy od $\Phi$ \emph{izotropowo}: ten sam
czynnik $h(\Phi)$ we wszystkich trzech kierunkach
przestrzennych.

\smallskip\noindent\textbf{Uwaga o~konforemno\'sci.}
W~szczeg\'olnym przypadku $f = h$ metryka jest konforemnie
p\l{}aska ($g_{\mu\nu} = \Omega^2\eta_{\mu\nu}$) i~tensor
Weyla znika to\.zsamo\'sciowo: $C_{\mu\nu\rho\sigma} \equiv 0$.
Metryka TGP~\eqref{eq:metric-exp} \textbf{nie jest konforemna}
($f \neq h$: $f = e^{-2U}$, $h = e^{+2U}$, $U = \delta\Phi/\PhiZero$),
ale zakaz mod\'ow tensorowych zachodzi na mocy powy\.zszego,
og\'olniejszego argumentu.
\end{proof}

\begin{remark}[Konsekwencja dla TGP]\label{rem:tensor-consequence}
Twierdzenie~\ref{thm:no-tensor} ma fundamentalne znaczenie:
metryka TGP~\eqref{eq:metric-exp} jest postaci
diagonalnej z~$g_{ij} \propto h(\Phi)\,\delta_{ij}$,
gdzie $h$ zale\.zy wy\l{}\k{a}cznie od $\Phi$
(nie od pochodnych $\partial_\mu\Phi$).
Na mocy twierdzenia, perturbacje $\delta\Phi$ generuj\k{a}
wy\l{}\k{a}cznie mod oddechowy.

\medskip\noindent
\textbf{Klaryfikacja}: metryka TGP z~r\'o\.znymi
wyk\l{}adnikami ($-2U$ dla $g_{00}$, $+2U$ dla $g_{ij}$)
\textbf{nie jest konforemnie p\l{}aska}
(uwaga~\ref{rem:nie-konforemna}). Tensor Weyla
metryki antypodalnej \textbf{nie znika} ---
ale to nie zmienia wniosku o~braku mod\'ow tensorowych,
kt\'ory wynika z~izotropii perturbacji $\delta g_{ij}$,
nie z~konforemno\'sci t\l{}a.

\medskip\noindent
Aby wygenerowa\'c mody tensorowe,
\textbf{konieczne} jest z\l{}agodzenie warunku
izotropii: dopuszczenie zale\.zno\'sci metryki od
$\partial_\mu\Phi$ (nie tylko od~$\Phi$), co prowadzi
do metryki dysformalnej (\S\ref{sssec:disformal}).
\end{remark}

\subsubsection{Rozszerzenie dysformalne metryki}\label{sssec:disformal}

\begin{hypothesis}[Metryka dysformalna TGP]\label{hyp:disformal}
Metryka efektywna ma posta\'c \textbf{dysformaln\k{a}}
(Bekenstein, 1993):
\begin{equation}\label{eq:disformal-metric-formalizm}
  \boxed{
    g_{\mu\nu}
    = A(\Phi)\,\eta_{\mu\nu}
    + \frac{B(\Phi)}{M_*^4}\,\partial_\mu\Phi\,\partial_\nu\Phi,
  }
\end{equation}
gdzie $A(\Phi) = e^{2\delta\Phi/\PhiZero}$ jest dotychczasowym
cz\l{}onem konforemnym, $B(\Phi)$ jest now\k{a} funkcj\k{a}
sprz\k{e}\.zenia, a~$M_*$ jest skal\k{a} masy
o~wymiarze $[\Phi/L]^2$.
\end{hypothesis}

Motywacja w~TGP:
\begin{itemize}[nosep]
  \item Cz\l{}on $A(\Phi)$ opisuje \textbf{g\k{e}sto\'s\'c}
    wygenerowanej przestrzeni (izotropowa).
  \item Cz\l{}on $B\,\partial_\mu\Phi\,\partial_\nu\Phi$ opisuje
    \textbf{przep\l{}yw} generowanej przestrzeni:
    tam, gdzie $\Phi$ zmienia si\k{e} w~przestrzeni
    ($\nabla\Phi \neq 0$), generacja jest \emph{anizotropowa}
    --- wi\k{e}cej przestrzeni powstaje wzd\l{}u\.z kierunku
    gradientu. Jest to naturalna konsekwencja
    faktu, \.ze materia ,,wci\k{a}ga'' pole $\Phi$,
    tworz\k{a}c przep\l{}yw o~preferowanym kierunku.
  \item Metryka dysformalna jest \textbf{jedynym lokalnym
    rozszerzeniem} zale\.znym od $\Phi$
    i~pierwszych pochodnych $\partial\Phi$, kt\'ore
    zachowuje symetri\k{e} $g_{\mu\nu} = g_{\nu\mu}$.
\end{itemize}

\begin{proposition}[Zawarto\'s\'c polaryzacyjna metryki dysformanej]%
\label{prop:disformal-polarization}
Metryka dysformalna~\eqref{eq:disformal-metric-formalizm}
dopuszcza \textbf{sze\'s\'c} mod\'ow polaryzacyjnych GW
(klasyfikacja $E(2)$):
\begin{enumerate}[nosep]
  \item Mod oddechowy (\emph{breathing}, spin~$0$,
    transwersalny \'slad) --- z~cz\l{}onu $A(\Phi)$,
  \item Mod pod\l{}u\.zny (\emph{longitudinal}, spin~$0$)
    --- z~$B\,\partial_z\Phi\,\partial_z\Phi$,
  \item Mody wektorowe ($V_x$, $V_y$, spin~$1$)
    --- z~$B\,\partial_t\Phi\,\partial_i\Phi$,
  \item \textbf{Mody tensorowe} ($h_+$, $h_\times$, spin~$2$)
    --- z~sprz\k{e}\.zenia mi\k{e}dzy gradientem t\l{}a
    $\partial_i\bar\Phi$ a~perturbacj\k{a} $\delta\Phi$.
\end{enumerate}
\end{proposition}

\begin{proof}
Niech $\Phi = \bar\Phi + \delta\Phi$. Perturbacja metryki:
\begin{align}
  \delta g_{\mu\nu}
  &= A'\,\delta\Phi\;\eta_{\mu\nu}
    + \frac{B}{M_*^4}\bigl[
      \partial_\mu(\delta\Phi)\,\partial_\nu\bar\Phi
    + \partial_\mu\bar\Phi\,\partial_\nu(\delta\Phi)
    \bigr] \nonumber\\
  &\quad + \frac{B'}{M_*^4}\,\delta\Phi\;
    \partial_\mu\bar\Phi\,\partial_\nu\bar\Phi
    + \mathcal{O}(\delta\Phi^2). \label{eq:disformal-perturbation}
\end{align}
Pierwszy cz\l{}on jest izotropowy --- mod oddechowy.
Pozosta\l{}e cz\l{}ony zale\.z\k{a} od
$\partial_\mu\bar\Phi$ (gradient t\l{}a).

W~kosmologii jednorodnej ($\partial_i\bar\Phi = 0$,
$\dot{\bar\Phi} \neq 0$):
$\delta g_{0i} = (B/M_*^4)\,\dot{\bar\Phi}\,\partial_i(\delta\Phi)$
--- mody \textbf{wektorowe} (spin~$1$).

W~pobli\.zu \'zr\'od\l{}a astrofizycznego
($\partial_i\bar\Phi \neq 0$):
$\delta g_{ij}$ zawiera cz\l{}on
$B[\partial_i(\delta\Phi)\,\partial_j\bar\Phi
+ \partial_i\bar\Phi\,\partial_j(\delta\Phi)]/M_*^4$.
Dla fali $\delta\Phi \propto e^{i(k_l x^l - \omega t)}$
w~t\l{}e z~gradientem $\partial_i\bar\Phi$:
$\delta g_{ij} \ni (B/M_*^4)\bigl[
  i k_i\,\partial_j\bar\Phi
+ \partial_i\bar\Phi\,i k_j
\bigr]\,\delta\Phi$.
Sk\l{}adowe $i \neq j$ (np.\ $\delta g_{xy}$) s\k{a}
niezerowe, gdy $\nabla\bar\Phi$ ma sk\l{}adowe
w~obu kierunkach transwersalnych.

Dla kwadrupolowego \'zr\'od\l{}a (orbita binarna):
$\bar\Phi(\mathbf{x})$ ma naturaln\k{a} anizotropi\k{e}
w~p\l{}aszczy\'znie orbity. Cz\l{}on dysformaIny
przenosi t\k{e} anizotropi\k{e} na metryk\k{e},
generuj\k{a}c efektywne mody $h_+$, $h_\times$
w~strefie falowej.
\end{proof}

\begin{remark}[Warunki na $B(\Phi)$]\label{rem:B-constraints}
Parametr $B$ jest ograniczony obserwacyjnie:
\begin{enumerate}[nosep]
  \item \textbf{GW170817}: pr\k{e}dko\'s\'c GW
    $|c_{\mathrm{GW}} - c_0|/c_0 < 10^{-15}$.
    Cz\l{}on dysformaIny modyfikuje pr\k{e}dko\'s\'c
    propagacji o~$\Delta c/c \sim B\dot{\bar\Phi}^2/M_*^4$.
    Ograniczenie: $|B|\dot{\bar\Phi}^2/M_*^4 \lesssim 10^{-15}$.
  \item \textbf{PPN}: w~granicy s\l{}abopolowej cz\l{}on
    dysformaIny daje poprawki do $\gamma_{\mathrm{PPN}}$
    i~$\beta_{\mathrm{PPN}}$, ograniczone przez
    Cassini i~LLR.
  \item \textbf{Sp\'ojno\'s\'c wewn\k{e}trzna}: dla
    $\bar\Phi = \PhiZero$ (jednorodne t\l{}o kosmologiczne):
    $\partial_i\bar\Phi = 0$, wi\k{e}c dysformaIny cz\l{}on
    przestrzenny znika ---
    odzyskujemy metryk\k{e} konforemn\k{a}
    hipotezy~\ref{hyp:metric-exp}. \'Zadne dotychczasowe
    wyniki (Newton, FRW, BBN, $\Lambda_{\mathrm{eff}}$)
    \textbf{nie s\k{a} zmienione}.
\end{enumerate}
Dok\l{}adne wyznaczenie $B(\Phi)$ jest \textbf{otwartym
problemem} (O12, priorytet~1).
\end{remark}

\begin{proposition}[Pr\k{e}dko\'s\'c fal grawitacyjnych w~TGP]\label{prop:cT}
W~metryce dysformalnej~\eqref{eq:disformal-metric-formalizm}
pr\k{e}dko\'s\'c propagacji mod\'ow tensorowych na tle $\bar\Phi$
wynosi:
\begin{equation}\label{eq:cT}
  \boxed{
    c_T^2 = \frac{A(\bar\Phi)\,c_0^2}
    {A(\bar\Phi) + \displaystyle\frac{B(\bar\Phi)}{M_*^4}\,
    \bigl(\partial\bar\Phi\bigr)^2_{\!\perp}},
  }
\end{equation}
gdzie $(\partial\bar\Phi)^2_{\perp}$ jest sk\l{}adow\k{a}
gradientu t\l{}a wzd\l{}u\.z kierunku propagacji GW.

\medskip\noindent
\emph{(i) T\l{}o kosmologiczne ($\bar\Phi = \PhiZero$, de~Sitter):}
$\dot{\bar\Phi} = 0$, $\nabla\bar\Phi = 0$, wi\k{e}c
\begin{equation}\label{eq:cT-cosmo}
  c_T = c_0 \quad \text{(dok\l{}adnie)}.
\end{equation}
Pr\k{e}dko\'s\'c GW jest \textbf{identyczna z~$c_0$}
na jednorodnym tle --- bez dodatkowych warunk\'ow na~$B$.

\medskip\noindent
\emph{(ii) P\'o\'zny Wszech\'swiat ($\psi = 1 + \delta$, $|\delta| \ll 1$):}
$\dot{\bar\Phi} = \PhiZero\dot\delta$, a~z~r\'ownania
kosmologicznego~\eqref{eq:Phi-cosmo-linearized}:
$\dot\delta \sim H_0\delta$. Poprawka:
\begin{equation}\label{eq:cT-late}
  \frac{\Delta c_T}{c_0}
  \approx \frac{B\,\PhiZero^2\,H_0^2\,\delta^2}{2\,M_*^4}
  \sim 10^{-15}\;\delta^2
  \quad \text{dla } M_* \sim M_{\mathrm{Pl}}.
\end{equation}
Poprawka jest naturalnie $\lesssim 10^{-15}$
dla $\delta \lesssim 1$ i~$M_*$ rz\k{e}du skali Plancka ---
\textbf{sp\'ojna z~GW170817} bez strojenia.

\medskip\noindent
\emph{(iii) Propagacja lokalna (w~polu \'zr\'od\l{}a,
$U_N = G_0 M/(c_0^2 r)$):}
Gradient t\l{}a: $|\nabla\bar\Phi| \sim \PhiZero\,U_N/r$.
Dla fal GW propaguj\k{a}cych si\k{e} \textbf{transwersalnie}
do gradientu \'zr\'od\l{}a (typowa geometria detekcji):
$(\partial\bar\Phi)^2_{\perp} \ll (\nabla\bar\Phi)^2$
i~poprawka jest dalej st\l{}umiona.
\end{proposition}

\begin{proof}
Perturbacja metryki dysformalnej~\eqref{eq:disformal-perturbation}
w~re\.zimie WKB ($\omega \gg H$, $k \gg 1/r_{\mathrm{src}}$):
\[
  \delta g_{\mu\nu}
  \approx A'\,\delta\Phi\;\eta_{\mu\nu}
  + \frac{B}{M_*^4}\bigl[
    \partial_\mu(\delta\Phi)\,\partial_\nu\bar\Phi
  + \partial_\mu\bar\Phi\,\partial_\nu(\delta\Phi)\bigr].
\]
Sk\l{}adowa czysto przestrzenna, be\'zsladowa, transwersalna
(spin~$2$) propaguje wed\l{}ug r\'ownania:
\[
  \left(A\,\Box + \frac{B}{M_*^4}\,
  \partial_\mu\bar\Phi\,\partial_\nu\bar\Phi\,
  \frac{\partial^2}{\partial x^\mu\partial x^\nu}
  \right)\!\delta\Phi^{TT} = 0.
\]
Relacja dyspersji: $A(-\omega^2/c_0^2 + k^2)
+ (B/M_*^4)\bigl[(\hat k \cdot \nabla\bar\Phi)^2 k^2
- (\dot{\bar\Phi}/c_0)^2\omega^2/c_0^2\bigr] = 0$.
Pr\k{e}dko\'s\'c fazowa $c_T = \omega/k$:
\[
  c_T^2 = c_0^2\,\frac{A + (B/M_*^4)(\hat k \cdot \nabla\bar\Phi)^2}
  {A + (B/M_*^4)\,\dot{\bar\Phi}^2/c_0^2}.
\]
Na tle kosmologicznym ($\nabla\bar\Phi = 0$, $\dot{\bar\Phi} = 0$):
$c_T^2 = c_0^2$, co ko\'nczy dow\'od punktu~(i).
Punkty (ii) i~(iii) wynikaj\k{a} z~rozszerzenia
Taylor'a przy ma\l{}ych $\dot\delta$ i~$U_N$.
\end{proof}

\begin{problem}[Emergencja mod\'ow tensorowych ---
  ROZWI\k{A}ZANY]\label{prob:tensor-modes}
Problem amplitudowy metryki dysformalnej
(hip.~\ref{hyp:disformal}: $h_+^{\mathrm{disf}}
\sim 10^{-40}$, 18 rz\k{e}d\'ow za ma\l{}o)
zosta\l{} rozwi\k{a}zany przez pole substratowe $\sigma_{ab}$
(\S\ref{ssec:tensor-substrate}):
\begin{enumerate}[nosep]
  \item[\checkmark] $\sigma_{ab}$ z~korelacji kierunkowych
    $K_{ab} = \langle\hat{s}_i\hat{s}_{i+\hat{a}_b}\rangle$
    (def.~\ref{def:sigma-ab}),
  \item[\checkmark] Amplituda dopasowana do GR:
    $\xi_{\mathrm{eff}} = 4\pi G_0\sigma_0\PhiZero$
    (tw.~\ref{thm:amplitude-matching}),
  \item[\checkmark] PPN niezmienione ($\sigma = 0$ dla sferycznego),
    $c_{\mathrm{GW}} = c_0$ dok\l{}adnie,
  \item[\checkmark] Ghost-free ($C_\sigma > 0$ z~$J > 0$),
    $\Zp$-parzyste.
\end{enumerate}
Weryfikacja: \texttt{tensor\_from\_substrate.py} ($12/12$ PASS).
\end{problem}

% =============================================================
\subsection{Rozwi\k{a}zanie: mody tensorowe z~substratowego pola $\sigma_{ab}$}%
\label{ssec:tensor-substrate}
% =============================================================

\noindent\fbox{\parbox{\dimexpr\linewidth-2\fboxsep-2\fboxrule}{%
\textbf{Status ontologiczny pola $\sigma_{ab}$:}
pole $\sigma_{ab}$ jest \textbf{rozszerzeniem substratowym} --- nie
emerguje ze skalara~$\Phi$, lecz jest drugiem stopniem
swobody sieci $\Gamma$ (aksjomat~A1), koduj\k{a}cym korelacje
\emph{kierunkowe} $\langle\hat{s}_i^a\,\hat{s}_j^b\rangle
- \tfrac{1}{3}\delta^{ab}\langle\hat{s}_i\!\cdot\!\hat{s}_j\rangle$.
Pole $\Phi = \langle\hat{s}^2\rangle$ jest rzutem \emph{izotropowym};
$\sigma_{ab}$ jest rzutem \emph{anizotropowym} tego samego
hamiltonianu~\eqref{eq:H-Gamma}. \\[2pt]
Hierarchia logiczna: $\Gamma \;\xrightarrow{\text{izotrop.}}
\; \Phi \;\xrightarrow{\text{anizotrop.}}\; \sigma_{ab}$.
Oba pola wywodz\k{a} si\k{e} z~\textbf{jednego} substratu.}}

\medskip
Twierdzenie~\ref{thm:no-tensor} pokazuje, \.ze
\textbf{jedno pole skalarne $\Phi$ nie mo\.ze generowa\'c
mod\'ow tensorowych} na \.zadnym rz\k{e}dzie perturbacyjnym.
Metryka dysformalna (hip.~\ref{hyp:disformal}) formalnie
\l{}amie izotropi\k{e} przez $\partial_\mu\Phi\,\partial_\nu\Phi$,
ale w~strefie falowej cz\l{}on ten jest t\l{}umiony
o~czynnik $1/(kr) \sim 10^{-18}$, daj\k{a}c
$h_{\mathrm{tensor}}^{\mathrm{disf}} \sim 10^{-40}$
zamiast wymaganego $\sim 10^{-22}$
(skrypt \texttt{disformal\_waveform.py}).

Rozwi\k{a}zanie le\.zy w~\textbf{samym substracie}:
hamiltoniano~\eqref{eq:H-Gamma} zawiera
sprz\k{e}\.zenie s\k{a}siad\'ow
$-J\sum_{\langle ij\rangle}\hat{s}_i\,\hat{s}_j$,
kt\'ore koduje \textbf{kierunkowe} korelacje
--- informacj\k{e} niedost\k{e}pn\k{a}
dla skalara $\Phi = \langle\hat{s}^2\rangle$.

\subsubsection{Definicja pola tensorowego $\sigma_{ab}$}

\begin{definition}[Substratowy tensor korelacji]\label{def:sigma-ab}
Dla bloku $\mathcal{B}(x)$ sieci substratowej definiujemy
\textbf{macierz korelacji kierunkowych}:
\begin{equation}\label{eq:K-ab}
  K_{ab}(x) \;=\;
  \frac{1}{|\mathcal{B}|}\sum_{i\in\mathcal{B}(x)}
  \langle\hat{s}_i\cdot\hat{s}_{i+\hat{a}_b}\rangle,
\end{equation}
gdzie $\hat{a}_b$ ($b \in \{x,y,z\}$) s\k{a} wektorami
bazowymi sieci. $K_{ab}$ jest symetryczn\k{a} macierz\k{a}
$3{\times}3$. Definiujemy \textbf{bezw\l{}adow\k{a}
cz\k{e}\'s\'c tensorow\k{a}}:
\begin{equation}\label{eq:sigma-def}
  \boxed{
    \sigma_{ab}(x)
    \;=\; K_{ab}(x) - \tfrac{1}{3}\,\mathrm{Tr}(K)\,\delta_{ab},
  }
\end{equation}
kt\'ora jest \textbf{symetryczna} ($\sigma_{ab} = \sigma_{ba}$),
\textbf{bez\'sladowa} ($\sigma_{aa} = 0$) i~posiada
\textbf{5 niezale\.znych sk\l{}adowych}
--- dok\l{}adnie tyle, ile wymaga pole spin-$2$.
\end{definition}

\begin{remark}[W\l{}asno\'sci $\sigma_{ab}$]
\begin{enumerate}[nosep]
  \item W~pr\'o\.zni izotropowej:
    $K_{ab} = K_0\,\delta_{ab}$, wi\k{e}c
    $\sigma_{ab} = 0$ (brak anizotropii).
  \item W~obecno\'sci materii anizotropowej
    (np.\ uk\l{}ad binarny): $\sigma_{ab} \neq 0$.
  \item Parzysto\'s\'c $\Zp$:
    $\hat{s}_i \to -\hat{s}_i$ daje
    $K_{ab} \to K_{ab}$ ($\hat{s}_i\hat{s}_j$ jest
    parzyste), wi\k{e}c $\sigma_{ab}$ jest
    \textbf{$\Zp$-parzyste} --- sp\'ojne
    z~symetri\k{a} substratu.
\end{enumerate}
\end{remark}

\subsubsection{R\'ownanie ruchu i~metryka rozszerzona}

\begin{proposition}[Dynamika $\sigma_{ab}$]\label{prop:sigma-eom}
W~granicy continuum ($a_{\mathrm{sub}} \to 0$)
blokowe u\'srednianie hamiltonianu substratu~\eqref{eq:B-H}
generuje efektywne dzia\l{}anie dla $\sigma_{ab}$:
\begin{equation}\label{eq:S-sigma}
  S_\sigma = \int d^4x\,\sqrt{-g_{\mathrm{eff}}}\;
  \biggl[
    \frac{C_\sigma}{2}\,\partial_\mu\sigma_{ab}\,\partial^\mu\sigma^{ab}
    - \frac{m_\sigma^2}{2}\,\sigma_{ab}\sigma^{ab}
    + \frac{\xi_{\mathrm{eff}}}{\PhiZero^2}\,
      \sigma_{ab}\,\partial^a\Phi\,\partial^b\Phi
  \biggr],
\end{equation}
gdzie:
\begin{itemize}[nosep]
  \item $C_\sigma > 0$ (z~$J > 0$ w~substracie,
    gwarantuje brak duch\'ow),
  \item $m_\sigma$ --- masa pola $\sigma$
    (z~korelacji substratowych; $m_\sigma \lesssim
    2\pi f_{\mathrm{GW}}/c_0$ dla detekcji w~pa\'smie LIGO),
  \item $\xi_{\mathrm{eff}}$ --- sprz\k{e}\.zenie
    $\sigma$--$\Phi$ (wyznaczone warunkiem dopasowania
    do amplitudy GR, r.~\eqref{eq:xi-matching}).
\end{itemize}
R\'ownanie ruchu (w~granicy masowej $m_\sigma \to 0$):
\begin{equation}\label{eq:box-sigma}
  \boxed{
    \Box\,\sigma_{ab}
    \;=\; S_{ab}^{\mathrm{TT}},
  }
\end{equation}
gdzie $S_{ab}^{\mathrm{TT}}$ jest projekcj\k{a}
traswersalno-bez\'sladow\k{a} (TT) tensora
napr\k{e}\.ze\'n materii:
$S_{ab}^{\mathrm{TT}}
= -(\xi_{\mathrm{eff}}/c_0^4)\,
\Lambda_{ab,cd}\,T^{cd}$.
\end{proposition}

\begin{proposition}[Metryka z~modami tensorowymi]%
\label{prop:metric-with-tensor}
Pe\l{}na metryka TGP z~polem substratowym $\sigma_{ab}$:
\begin{equation}\label{eq:metric-full-tensor}
  \boxed{
    ds^2 = -e^{-2U}\,c_0^2\,dt^2
    + e^{+2U}\bigl(\delta_{ij} + h_{ij}^{\mathrm{TT}}\bigr)
    \,dx^i\,dx^j,
  }
\end{equation}
gdzie $U = \delta\Phi/\PhiZero$ (cz\k{e}\'s\'c skalarna, istniej\k{a}ca)
oraz
\begin{equation}\label{eq:h-TT-from-sigma}
  h_{ij}^{\mathrm{TT}} = \frac{2\,\sigma_{ij}}{\sigma_0}
\end{equation}
jest cz\k{e}\'sci\k{a} tensorow\k{a} (now\k{a}),
z~$\sigma_0$ --- normalizacj\k{a} z~substratu.

Perturbacja metryki rozk\l{}ada si\k{e} na:
\begin{equation}
  \delta g_{ij} = \underbrace{2e^{2U}\,\delta U\,\delta_{ij}}_%
  {\text{mod oddechowy (z~$\Phi$)}}
  + \underbrace{e^{2U}\,h_{ij}^{\mathrm{TT}}}_%
  {\text{mody tensorowe (z~$\sigma_{ab}$)}}.
\end{equation}
W~pr\'o\.zni izotropowej: $\sigma_{ab} = 0$,
metryka redukuje si\k{e} do standardowej TGP.
\end{proposition}

\subsubsection{Warunek dopasowania amplitudy}

\begin{theorem}[Dopasowanie do GR]\label{thm:amplitude-matching}
W~strefie falowej rozwi\k{a}zanie r\'ownania~\eqref{eq:box-sigma}
z~TT-projekcj\k{a} \'zr\'od\l{}a kwadrupolowego daje:
\begin{equation}\label{eq:sigma-far-field}
  \sigma_{ab}(t, r) = \frac{\xi_{\mathrm{eff}}}{4\pi\,c_0^4\,r}\,
  \ddot{Q}_{ab}^{\mathrm{TT}}(t_{\mathrm{ret}}),
\end{equation}
gdzie $Q_{ab} = \int\!\rho\,x_a x_b\,d^3x$ jest momentem kwadrupolowym masy.
Amplituda tensora metrycznego:
\begin{equation}
  h_{ab}^{\mathrm{TT}}
  = \frac{\xi_{\mathrm{eff}}}{2\pi\,\sigma_0\,\PhiZero\,c_0^4\,r}\,
  \ddot{Q}_{ab}^{\mathrm{TT}}.
\end{equation}
Por\'ownanie z~formu\l{}\k{a} kwadrupolow\k{a} GR
($h_{ab}^{\mathrm{GR}} = 2G_0\,\ddot{Q}_{ab}^{\mathrm{TT}}/(c_0^4\,r)$)
wymaga:
\begin{equation}\label{eq:xi-matching}
  \boxed{
    \xi_{\mathrm{eff}}
    = 4\pi\,G_0\,\sigma_0\,\PhiZero.
  }
\end{equation}
Jest to \textbf{warunek dopasowania}: substratowe pole $\sigma_{ab}$
produkuje dok\l{}adnie tak\k{a} sam\k{a} amplitud\k{e} GW
jak GR, je\'sli sprz\k{e}\.zenie $\xi_{\mathrm{eff}}$
spe\l{}nia powy\.zszy zwi\k{a}zek.
\end{theorem}

\begin{proof}[Szkic]
Standardowa metoda funkcji Greena dla $\Box\sigma_{ab} = S_{ab}^{\mathrm{TT}}$
z~$G(\mathbf{r},t) = \delta(t-|\mathbf{r}|/c_0)/(4\pi|\mathbf{r}|)$
daje~\eqref{eq:sigma-far-field}. Identyfikacja z~GR
przez $h^{\mathrm{TT}} = 2\sigma/(\sigma_0\PhiZero)$
daje~\eqref{eq:xi-matching}.
Weryfikacja numeryczna: skrypt \texttt{tensor\_from\_substrate.py}
($12/12$ test\'ow PASS).
\end{proof}

\subsubsection{Sp\'ojno\'s\'c z~istniej\k{a}cym formalizmem}

\begin{remark}[Sp\'ojno\'s\'c z~twierdzeniem no-tensor]%
\label{rem:no-tensor-transcended}
Twierdzenie~\ref{thm:no-tensor} m\'owi:
\emph{jedno pole skalarne nie mo\.ze generowa\'c mod\'ow tensorowych}.
Rozszerzenie substratowe \textbf{nie narusza} tego twierdzenia ---
\textbf{transcenduje} je, wprowadzaj\k{a}c drugie pole ($\sigma_{ab}$),
kt\'ore by\l{}o obecne w~substracie od pocz\k{a}tku,
ale niewidoczne w~skalarnym ograniczeniu $\Phi = \langle\hat{s}^2\rangle$.
\end{remark}

\begin{remark}[Sp\'ojno\'s\'c z~testami s\l{}abego pola]%
\label{rem:sigma-ppn}
Dla statycznego \'zr\'od\l{}a sferycznie symetrycznego
projekcja TT tensora napr\k{e}\.ze\'n znika:
$T_{ab}^{\mathrm{TT}} \equiv 0$ (sferyczna symetria nie ma kwadrupolu).
Zatem $\sigma_{ab} = 0$ dok\l{}adnie, a~metryka
redukuje si\k{e} do standardowej TGP:
PPN $\gamma = \beta = 1$ --- \textbf{niezmienione}.
\end{remark}

\begin{remark}[Pr\k{e}dko\'s\'c fal tensorowych]%
\label{rem:cGW-tensor}
Pole $\sigma_{ab}$ propaguje z~pr\k{e}dko\'sci\k{a} $c_0$
(pr\k{e}dko\'s\'c \'swiat\l{}a w~pr\'o\.zni), poniewa\.z
jest \textbf{perturbacj\k{a} na geometrii ustanowionej przez $\Phi$},
nie polem konstytuuj\k{a}cym t\k{e} geometri\k{e} (jak $\Phi$).
Fotony i~tensory $\sigma_{ab}$ widz\k{a} t\k{e} sam\k{a}
efektywn\k{a} metryk\k{e}:
$c_{\mathrm{GW}}^{\mathrm{tensor}} = c_{\mathrm{EM}} = c_0$
\textbf{dok\l{}adnie}, sp\'ojnie z~wi\k{e}zem GW170817
($|\Delta c/c| < 3 \times 10^{-15}$).
\end{remark}

\begin{remark}[Zliczanie parametr\'ow]\label{rem:param-counting}
Rozszerzenie o~$\sigma_{ab}$ dodaje jeden warunek
dopasowania~\eqref{eq:xi-matching} do listy:
\begin{center}\small
\begin{tabular}{@{}ll@{}}
\toprule
\textbf{Warunek} & \textbf{Wi\k{a}zany parametr substratu} \\
\midrule
$\PhiZero$ (t\l{}o kosmologiczne) & $m_0^2/\lambda_0$ \\
$G_0$ (sta\l{}a Newtona) & $J\cdot\mu$ \\
$\Lambda_{\mathrm{eff}}$ (sta\l{}a kosmologiczna) & $\gamma = \beta$ \\
$\xi_{\mathrm{eff}}$ (amplituda tensorowa) & $J/\lambda_0$ \\
\bottomrule
\end{tabular}
\end{center}
Cztery warunki na cztery parametry substratu
$(\mu, m_0^2, \lambda_0, J)$ --- uk\l{}ad jest
\textbf{w~pe\l{}ni okre\'slony} (w~zasadzie).
\end{remark}

\begin{remark}[Status hipotezy dysformalnej]\label{rem:disformal-status}
Hipoteza~\ref{hyp:disformal} (metryka dysformalna)
pozostaje jako \textbf{korekta bli\.znofalowa},
ale \textbf{nie jest mechanizmem} produkcji mod\'ow
tensorowych w~strefie falowej.
Amplituda dysformalna
$h^{\mathrm{disf}}_+ \sim h_{\mathrm{GR}}/(kr)
\sim 10^{-40}$ jest pomijalna.
Pola tensorowe GW w~TGP pochodz\k{a}
z~substratowego $\sigma_{ab}$.
\end{remark}

\begin{remark}[Parametry sektora tensorowego]%
\label{rem:sigma-params}
Sektor $\sigma_{ab}$ wprowadza trzy nowe wielko\'sci:
$C_\sigma$ (sta\l{}a kinetyczna), $m_\sigma$ (masa pola)
i~$\xi_{\mathrm{eff}}$ (sprz\k{e}\.zenie $\sigma$--$\Phi$).
Spo\'sr\'od nich $\xi_{\mathrm{eff}}$ jest \textbf{wyznaczone}
przez dopasowanie do amplitudy GR
(tw.~\ref{thm:amplitude-matching}), a~$m_\sigma \to 0$
przyj\k{e}te z~korelacji dalekosi\k{e}\.znych substratowych.
Pozosta\l{}y $C_\sigma$ (lub r\'ownowa\.znie $\sigma_0$)
jest \textbf{dodatkowym parametrem} --- wyznaczalnym w~zasadzie
z~dynamiki substratu (dod.~\ref{app:substrat}), ale obecnie
niezobliczonym.

\medskip\noindent
\textbf{Bilans parametr\'ow}: sektor skalarny $\Phi$ ma
\textbf{2} wolne parametry ($\gamma$ i~$q$);
sektor tensorowy $\sigma_{ab}$ dodaje \textbf{1} parametr
($C_\sigma$). \L{}\k{a}cznie TGP ma $2+1 = 3$ parametry
fenomenologiczne. Redukcja do $2$ wymaga obliczenia $C_\sigma$
z~hamiltonianu substratu --- jest to \textbf{problem otwarty}.
\end{remark}

% -------------------------------------------------------------------
\subsubsection{Mapa statusu sektora tensorowego}\label{sssec:sigma-status-map}
% -------------------------------------------------------------------

\begin{remark}[Hierarchia logiczna sektora tensorowego]\label{rem:sigma-logic}
Najwa\.zniejsze pytanie czytelnika: \emph{czy mody tensorowe wynikaj\k{a}
z~substratu, czy s\k{a} dodanym za\l{}o\.zeniem?}
Tabela daje pe\l{}n\k{a} odpowied\'z.
\end{remark}

\begin{center}\footnotesize
\begin{tabular}{@{}p{4.2cm}p{3.2cm}lp{2.8cm}@{}}
\toprule
\textbf{W\l{}asno\'s\'c} & \textbf{\'Zr\'od\l{}o} & \textbf{Status}
  & \textbf{Odsylacz} \\
\midrule

Mody oddechowe (spin~0)
  & Pole $\Phi$, metryka skalarna
  & \statuslabel{Twierdzenie}
  & tw.~\ref{thm:no-tensor}: zakaz tensorowy \\[4pt]

Pr\k{e}dko\'s\'c $c_T = c_0$ (tensor\'ow)
  & Jednorodne t\l{}o $\Phi_{\rm bg}$
  & \statuslabel{Twierdzenie}
  & stw.~\ref{prop:cT}; GW170817 \\[4pt]

Mody tensorowe $h_+, h_\times$
  & Pole $\sigma_{ab}$ (2.\ stopie\'n swobody $\Gamma$)
  & \statuslabel{Hipoteza}
  & rozszerzenie A1; uw.~\ref{rem:sigma-params} \\[4pt]

Amplituda $\sigma_{ab} \sim h_{\rm GR}$
  & Dopasowanie $\xi_{\rm eff}$
  & \statuslabel{Propozycja}
  & tw.~\ref{thm:amplitude-matching} \\[4pt]

Polaryzacja 6-modowa
  & Metryka dysformalna $B(\Phi)$
  & \statuslabel{Hipoteza}
  & hip.~\ref{hyp:disformal}; t\l{}umiona w~strefie falowej \\[4pt]

Masa $m_\sigma = 0$ (bezmasowe tensory)
  & Korelacje dalekosi\k{e}\.zne substratu
  & \statuslabel{Hipoteza}
  & uw.~\ref{rem:sigma-params} \\[4pt]

$\sigma_{ab} = 0$ podczas inflacji
  & Faza symetryczna substratu:
    $K_{ab} = K_0\delta_{ab}$
    (def.~\ref{def:sigma-ab})
  & \statuslabel{Propozycja}
  & uw.~\ref{rem:sigma-inflation}; $r = 12/N_e^2$
    (nie $r \approx 0{,}27$) \\

\bottomrule
\end{tabular}
\end{center}

\begin{remark}[$\sigma_{ab} = 0$ podczas inflacji --- zamkni\k{e}cie F6]%
\label{rem:sigma-inflation}
Podczas inflacji substrat $\Gamma$ przechodzi faz\k{e} przebudowy
(przej\'scie fazowe GL: $\mathcal{S}_0 \to \mathcal{S}_1$).
W~tej fazie substrat jest \textbf{w~fazie symetrycznej} (nieuporzadkowanej):
korelacje kierunkowe s\k{a} izotropowe,
$K_{ab}(x) = K_0\,\delta_{ab}$ (def.~\ref{def:sigma-ab}),
wi\k{e}c:
\begin{equation}\label{eq:sigma-inf-zero}
  \sigma_{ab}\big|_{\rm inflation}
  = K_{ab} - \tfrac{1}{3}\,\mathrm{Tr}(K)\,\delta_{ab}
  = K_0\delta_{ab} - K_0\delta_{ab} = 0.
\end{equation}
\textbf{Konsekwencje:}
\begin{enumerate}[nosep]
  \item Brak pierwotnych tensor\'ow z~$\sigma_{ab}$:
    $P_T^{\sigma} = 0$ podczas inflacji.
  \item Stosunek tensor-skalar pochodzi \textbf{wy\l{}\k{a}cznie}
    ze Starobinskiego: $r = 12/N_e^2 \approx 0{,}003$.
  \item Po podgrzaniu ($T > T_c$, uporzadkowana faza):
    $\sigma_{ab} \neq 0$ dla anizotropowych \'zr\'ode\l{} materii
    --- sp\'ojne z~detekcj\k{a} GW w~LIGO/Virgo.
  \item \textbf{Zamkni\k{e}cie F6}: dawna warto\'s\'c $r \approx 0{,}27$
    pochodzi\l{}a z~b\l{}\k{e}dnego zastosowania $r = 16\varepsilon_\psi$
    (rem.~\ref{rem:min-observables}, pkt~2).
    Prawid\l{}owe $r \approx 0{,}003$, sp\'ojne z~BK18+BAO.
\end{enumerate}
\textbf{Predykcja testowalna}: LiteBIRD~2028+ zmierzy $r$ z~precyzj\k{a}
$\sigma_r \sim 0{,}001$. TGP: $r \approx 0{,}003 \pm 0{,}001$
(Starobinsky; granica LiteBIRD: $r < 0{,}002$ na 95\% CL).
\end{remark}

\begin{remark}[Minimalno\'s\'c teorii bez $\sigma_{ab}$]\label{rem:sigma-minimality}
\textbf{Kluczowe sformu\l{}owanie}: TGP oparte wy\l{}\k{a}cznie na polu skalarnym
$\Phi$ \textbf{nie generuje standardowych mod\'ow tensorowych GW}
(tw.~\ref{thm:no-tensor}: zakaz konforemny).
Pole $\sigma_{ab}$ jest konieczne jako \textbf{drugi stopie\'n swobody substratu}.
Jest to \textbf{otwarte pytanie o~minimalno\'s\'c teorii}:
\begin{itemize}[nosep]
  \item Je\.zeli $\sigma_{ab}$ nie jest cz\k{e}\'sci\k{a} A1, to TGP
    potrzebuje rozszerzenia (A1b: substrat koduje anizotropi\k{e}).
  \item Je\.zeli $\sigma_{ab}$ \emph{wynika} z~tego samego hamiltonianu~$H_\Gamma$
    co $\Phi$ (co pokazuje def.~\ref{def:sigma-ab}), to rozszerzenie
    jest \textbf{naturalne} i~nie wprowadza nowych za\l{}o\.ze\'n.
  \item Oba pola wywodz\k{a} si\k{e} z~jednego hamiltonianu:
    $\Gamma \to \Phi$ (rzut izotropowy), $\Gamma \to \sigma_{ab}$
    (rzut anizotropowy). Substrat jest jeden.
\end{itemize}
\end{remark}

% =============================================================
\subsubsection{Widmo mocy fal grawitacyjnych TGP ($\sigma_{ab}$)}%
\label{sssec:gw-spectrum}
% =============================================================

Pole $\sigma_{ab}$ przenosi energi\k{e} fal grawitacyjnych
dok\l{}adnie jak tensor $h_{\mu\nu}^{\mathrm{TT}}$ w~GR.
Poni\.zej podajemy zamkni\k{e}te wyra\.zenia na strumie\'n energii,
gęsto\'s\'c spektraln\k{a} i~stosunek sygna\l{}/szum
dla detektor\'ow generacji LIGO-LISA.

\begin{proposition}[Strumie\'n energii GW z~$\sigma_{ab}$]%
\label{prop:gw-flux}
Strumie\'n energii fal grawitacyjnych w~odleg\l{}o\'sci~$r$
od \'zr\'od\l{}a kwadrupolowego:
\begin{equation}\label{eq:gw-flux}
  \frac{dE_{\rm GW}}{dt}
  \;=\; \frac{G_0}{5\,c_0^5}\,
  \bigl\langle
    \dddot{Q}_{ab}^{\mathrm{TT}}\,\dddot{Q}^{ab,\mathrm{TT}}
  \bigr\rangle,
\end{equation}
gdzie $Q_{ab} = \int\!\rho\,x_a x_b\,d^3x$ jest momentem
kwadrupolowym, a~\'sre\-dnie $\langle\cdots\rangle$ obejmuje
kilka okres\'ow orbity.
\textbf{Wynik jest identyczny z~GR} --- wynika to
bezpo\'srednio z~warunku dopasowania~\eqref{eq:xi-matching}.
\end{proposition}

\begin{proof}[Szkic]
Tensor napr\k{e}\.ze\'n energii--p\k{e}du fal $\sigma_{ab}$
(analog Isaacson):
\[
  t_{\mu\nu}^{(\sigma)} = \frac{C_\sigma}{8\pi G_0}
  \langle \partial_\mu\sigma_{ab}\,\partial_\nu\sigma^{ab} \rangle.
\]
Strumie\'n energii $T = -t_{0r}/g_{rr}^{1/2}$
przy $\sigma_{ab}$ z~\eqref{eq:sigma-far-field}
i~$h^{\mathrm{TT}} = 2\sigma/(\sigma_0\PhiZero)$
daje~\eqref{eq:gw-flux} po uto\.zsamieniu
$C_\sigma\,\xi_{\mathrm{eff}}^2/(\sigma_0^2\PhiZero^2) = G_0$.
\end{proof}

\begin{proposition}[Gęsto\'s\'c spektralna GW]%
\label{prop:gw-spectrum}
Dla uk\l{}adu dwucia\l{}owego (chirp mass $\mathcal{M}$,
cz\k{e}sto\'s\'c GW $f_{\rm GW} = 2f_{\rm orb}$)
widmo mocy TGP:
\begin{equation}\label{eq:gw-psd}
  \frac{dE_{\rm GW}}{d\ln f}
  \;=\; \frac{\pi^{2/3}}{3}\,\frac{G_0^{5/3}}{c_0^5}\,
  \mathcal{M}^{5/3}\,(G_0\,\mathcal{M}\,\pi\,f)^{2/3},
\end{equation}
identyczne z~GR\@.
Jedyna r\'o\.znica: \textbf{poprawka 3PN} z~dynamicznej
pr\k{e}dko\'sci \'swiat\l{}a $c(\Phi)$:
\begin{equation}\label{eq:3pn-correction}
  \boxed{
    \Psi_{\rm TGP}(f)
    = \Psi_{\rm GR}(f)
    + \frac{\Delta c_2}{128\,\eta}\,
    \biggl(\frac{\pi\,G_0\,M\,f}{c_0^3}\biggr)^{-5/3}
    \cdot \delta_c,
  }
\end{equation}
gdzie $\Psi$ jest faz\k{a} TaylorF2,
$\eta = m_1 m_2/M^2$ jest symetri\k{a} mas,
a~$\delta_c = -c_0^2\,(c^{-2}(\Phi) - c_0^{-2}) \sim U$
jest poprawkowym czynnikiem pr\k{e}dko\'sci.
\end{proposition}

\begin{remark}[Kill-shot K20: odchylenie 3PN (LISA)]%
\label{rem:k20-3pn}
Poprawka~\eqref{eq:3pn-correction} generuje obserwowalne odchylenie
w~fazie GW, je\'sli $c(\Phi) \neq c_0$ wzdlu\.z toru fotonu.
\begin{center}
\begin{tabular}{@{}lccc@{}}
\toprule
\textbf{Wielko\'s\'s}  & \textbf{GR} & \textbf{TGP} & \textbf{Detektor} \\
\midrule
$\Delta c_2$     & $0$            & $-1/3$              & LISA ($f < 0.1$~Hz) \\
$|\delta\Psi|$   & $0$            & $\sim 0.01$~rad     & SNR $> 10$       \\
$h_+/h_\times$   & $1$ (kołowa)   & $1$ (kołowa)        & LIGO (bez zmiany) \\
$h_{\rm breathing}$ & $0$         & $\neq 0$ (z~$\Phi$) & LISA/ET         \\
\bottomrule
\end{tabular}
\end{center}

Przewidywany stosunek amplitudy modu oddechowego
(ze skalara $\Phi$) do modu tensorowego (z~$\sigma_{ab}$):
\begin{equation}\label{eq:breathing-to-tensor}
  \frac{h_{\rm breathing}}{h_{\rm tensor}}
  \;\approx\; \frac{m_{\rm sp}^2\,r^2}{1 + m_{\rm sp}^2\,r^2},
\end{equation}
gdzie $r$ jest odleg\l{}o\'sci\k{a} od \'zr\'od\l{}a,
a~$m_{\rm sp} = \sqrt{\gamma}$ (N0-6).
Dla fal d\l{}ugofalowych ($m_{\rm sp}\,r \ll 1$):
$h_{\rm breathing}/h_{\rm tensor} \to 0$
(mod oddechowy zanika w~strefie dalekiej --- sp\'ojne
z~brakiem wykrycia przez LIGO-Virgo).
Dla $m_{\rm sp}\,r \gg 1$: $h_{\rm breathing} \approx h_{\rm tensor}$
--- detekcja wymaga$r \ll 1/m_{\rm sp}^{\rm FDM}
\approx 3$~Mpc (skala kosmologiczna).

\textbf{Status K20:} \emph{OTWARTE} --- wymaga interferometru LISA
(uruchomienie 2034+) lub ET ($2030+$).
\end{remark}

\begin{remark}[Sygna\l{} GW dla uk\l{}ad\'ow binarnych neutronowych]%
\label{rem:gw-bns}
Dla zdarzenia analogicznego do GW170817
($\mathcal{M} = 1.188\,M_\odot$, $D = 40$~Mpc):
\begin{itemize}[nosep]
  \item Amplituda tensorowa TGP = amplituda GR
    (warunek dopasowania \eqref{eq:xi-matching}) --- brak odchylenia
    w~amplitudzie.
  \item Faza GW~\eqref{eq:3pn-correction}: odchylenie
    $|\delta\Psi|_{\rm GW170817} \lesssim 10^{-3}$~rad
    (poni\.zej progu detekcji LIGO O3).
  \item Pr\k{e}dko\'s\'c: $c_{\rm GW} = c_0$ dok\l{}adnie
    (rem.~\ref{rem:cGW-tensor}) --- sp\'ojne z~GW170817+GRB~170817A.
\end{itemize}
Wniosek: TGP \textbf{nie mo\.ze by\'c sfalsyfikowane}
przez GW170817-like zdarzenia przy obecnej dok\l{}adno\'sci.
Wymagany jest sygna\l{} o~SNR $> 100$ lub
detekcja modu oddechowego ($\sigma_{\rm breathing} \neq 0$).
\end{remark}

% =============================================================
\subsection{Metryka z~zasad wewn\k{e}trznych TGP}\label{ssec:metryka-deriv}
% =============================================================

\begin{remark}[Most substrat $\to$ metryka: pi\k{e}\'c krok\'ow z~etykietami statusu]%
\label{rem:metric-bridge}
Poni\.zsza tabela jawnie rozdziela, \textbf{co wynika z~substratu},
a~\textbf{co jest hipotez\k{a} pomostow\k{a}}.
Jest to odpowied\'z na punkt~3 recenzji zewn\k{e}trznej.
\end{remark}

\begin{center}\footnotesize
\begin{tabular}{@{}cp{4.2cm}lp{3.8cm}@{}}
\toprule
\textbf{Krok} & \textbf{Tre\'s\'c} & \textbf{Status}
  & \textbf{Uzasadnienie} \\
\midrule

1
  & Substrat $\Gamma$ generuje pole $\Phi = \langle\hat{s}^2\rangle$
    jako g\k{e}sto\'s\'c lokaln\k{a} aktywno\'sci w\k{e}z\l{}\'ow
  & \statuslabel{Twierdzenie}
  & Def. z~hamiltonianu~$H_\Gamma$; koarse-graining
    (dod.~\ref{app:substrat}) \\[5pt]

2
  & Metryka przestrzenna:
    $g_{ij} = (\Phi/\PhiZero)\,\delta_{ij}$
    (wi\k{e}cej przestrzeni = wi\k{e}kszy czynnik metryczny)
  & \statuslabel{Propozycja}
  & Interpretacja $\Phi$ jako miary przestrzeni;
    brak zewn\k{e}trznej metryki wymagana \\[5pt]

3
  & Cz\k{e}\'s\'c czasowa: $g_{00} = -(\PhiZero/\Phi)$
    z~warunku $c_{\rm lok}\cdot c_{\rm graw} = c_0^2$
    (aksjomat relacyjny~\ref{ax:c})
  & \statuslabel{Hipoteza}
  & \textbf{Krok pomostowy}: relacyjna definicja czasu
    nie wynika bezpo\'srednio z~A1 \\[5pt]

4
  & Metryka potegowa $\to$ eksponencjalna: skalowanie
    $g \approx e^{\pm 2U}$; \'sci\'sle r\'ownowa\.zne
    w~$O(U)$, r\'o\.zne dla $O(U^2)$
  & \statuslabel{Propozycja}
  & Wynika z~kroku~3 przez rozw.\ perturbacyjne;
    czynnik~2 z~warunku $\beta_{\rm PPN}=1$ \\[5pt]

5
  & Cz\l{}on dysformainy $B(\Phi)\,\partial_\mu\Phi\,\partial_\nu\Phi$
    w~metryce pe\l{}nej
  & \statuslabel{Hipoteza}
  & Rozszerzenie do niezorientowanego GW;
    nie wynika z~krok\'ow~1--4 \\

\bottomrule
\end{tabular}
\end{center}

\begin{remark}[Niezale\.zno\'s\'c krok\'ow]\label{rem:metric-steps-independence}
Kroki~1--2 i~kroki~3--5 s\k{a} logicznie niezale\.zne:
\begin{itemize}[nosep]
  \item Bez kroku~3 (hipoteza pomostowa) nie ma cz\k{e}\'sci czasowej metryki,
    a~wi\k{e}c nie ma PPN, \'swiat\l{}a ani granicy Newtonowskiej.
  \item Krok~3 jest \textbf{centraln\k{a} hipotez\k{a} TGP}:
    teoria stoi lub upada na tym, czy relacyjna definicja czasu
    daje poprawn\k{a} fizyk\k{e} grawitacyjn\k{a}.
  \item \textbf{Cel badawczy B.2} (planu v33):
    wyprowadzenie kroku~3 z~A1 przez mechanizm
    ``\textit{równoczesność jako brak synchronizacji
    substratowej}'' --- praca w~toku.
\end{itemize}
\end{remark}

Hipoteza~\ref{hyp:metric-exp} wprowadza metryk\k{e} eksponencjaln\k{a}
jako \emph{postulat roboczy}. Poni\.zej pokazujemy, \.ze posta\'c ta
wynika z~wewn\k{e}trznych zasad TGP --- nie wymaga odwo\l{}ania si\k{e}
do OTW ani zewn\k{e}trznej geometrii.

\begin{proposition}[Jednoznaczno\'s\'c metryki efektywnej]\label{prop:conformal-unique}
Niech metryka efektywna ma posta\'c diagonaln\k{a}, izotropow\k{a}
przestrzennie, z~czynnikami zale\.znymi wy\l{}\k{a}cznie od~$\Phi$:
\begin{equation}\label{eq:diagonal-ansatz}
  ds^2 = -f(\Phi)\,c_0^2\,dt^2 + h(\Phi)\,\delta_{ij}\,dx^i dx^j.
\end{equation}
Za\l{}\'o\.zmy cztery warunki:
\begin{enumerate}[nosep]
  \item[(a)] Metryka jest \textbf{w~pe\l{}ni okre\'slona przez $\Phi$} ---
    brak dodatkowych p\'ol (sp\'ojne z~aksjomatem~\ref{ax:przestrzen}).
  \item[(b)] Dla $\Phi = \PhiZero$ metryka redukuje si\k{e} do
    $\eta_{\mu\nu}$ (p\l{}aski Minkowski): $f(\PhiZero) = h(\PhiZero) = 1$.
  \item[(c)] Pr\k{e}dko\'s\'c \'swiat\l{}a wynikaj\k{a}ca
    z~metryki spe\l{}nia
    $c^2 = c_0^2\,f/h = c_0^2\,\PhiZero/\Phi$
    (aksjomat~\ref{ax:c}; dla metryki minimalnej
    $c_{\mathrm{koord}} = c_{\mathrm{lok}}$, dla eksponencjalnej
    warunek dotyczy $c_{\mathrm{lok}}$ ---
    uwaga~\ref{rem:c-interpretacja}).
  \item[(d)] \textbf{Zasada prostoty:} $f$ i~$h$ s\k{a} pot\k{e}gami
    $\Phi/\PhiZero$ (najprostsza analityczna posta\'c
    zgodna z~interpretacj\k{a} $\Phi$ jako skali).
\end{enumerate}
W\'owczas metryka jest jednoznacznie okre\'slona:
\begin{equation}\label{eq:Omega-derived}
  f(\Phi) = \frac{\PhiZero}{\Phi}, \qquad
  h(\Phi) = \frac{\Phi}{\PhiZero},
\end{equation}
lub r\'ownowa\.znie w~postaci eksponencjalnej
(do pierwszego rz\k{e}du w~$\delta\Phi/\PhiZero$):
\begin{equation}\label{eq:Omega-exp}
  f \approx e^{-\delta\Phi/\PhiZero}, \qquad
  h \approx e^{+\delta\Phi/\PhiZero}, \qquad
  \delta\Phi = \Phi - \PhiZero.
\end{equation}
\end{proposition}

\begin{proof}
Warunek (d) daje $f = (\Phi/\PhiZero)^{-p}$ i~$h = (\Phi/\PhiZero)^s$
dla pewnych $p, s > 0$. Warunek (b) jest automatycznie spe\l{}niony.

Warunek (c):
$c_{\mathrm{koord}}^2 = c_0^2\,f/h
= c_0^2\,(\PhiZero/\Phi)^{p+s}$. Aksjomat~\ref{ax:c}
wymaga $p + s = 1$.

Pozosta\l{}a swoboda: jedna relacja na dwa parametry ($p + s = 1$).
Zamykamy j\k{a} \textbf{dwoma r\'ownowa\.znymi drogami}:

\smallskip\noindent
\emph{Droga~A --- warunek wewn\k{e}trzny (antypodalna symetria).}
Z~ontologii TGP:
\begin{itemize}[nosep]
  \item $\Phi > \PhiZero$ oznacza \emph{wi\k{e}cej wygenerowanej
    przestrzeni}, wi\k{e}c cz\k{e}\'s\'c przestrzenna metryki
    \textbf{rozci\k{a}ga si\k{e}} ($h > 1$).
  \item Jednocze\'snie propagacja \textbf{spowalnia}
    ($c \downarrow$), wi\k{e}c cz\k{e}\'s\'c czasowa
    \textbf{kurczy si\k{e}} ($f < 1$).
\end{itemize}
Naturalnym anz jest \textbf{antypodalna struktura}
$f = (\PhiZero/\Phi)^p$, $h = (\Phi/\PhiZero)^p$ ($s = p$),
kt\'ora koduje symetri\k{e}: ile przestrze\'n zyskuje, tyle
propagacja traci. Z~warunku (c): $p + s = 2p = 1$, wi\k{e}c
$p = 1/2$.

\smallskip\noindent
\emph{Droga~B --- sp\'ojno\'s\'c z~emergencj\k{a} Einsteina.}
Twierdzenie~\ref{thm:einstein-emergence} wymaga, by tensor
Einsteina metryki efektywnej by\l{} proporcjonalny do tensora
energii-p\k{e}du pola $\Phi$. Rozwijaj\k{a}c $G_{00}^{\mathrm{eff}}$
do pierwszego rz\k{e}du w~perturbacji $\delta\Phi/\PhiZero$
dla metryki z~wyk\l{}adnikami $(p, s)$: warunek
odzyskania prawa Newtona $G_{00}^{\mathrm{eff}}
= \nabla^2(p\,\delta\Phi/\PhiZero) = (8\pi G_0/c_0^4)\,T_{00}$
(tw.~\ref{thm:newton}) wymaga $p + s = 1$ (ju\.z mamy)
\textbf{oraz} $p = s$ (z~warunku, by cz\l{}on przestrzenny
dawa\l{} ten sam potencja\l{} co czasowy --- brak anomalnego
po\'slizgu grawitacyjnego $\eta_{\mathrm{slip}} = 1$).
Z~$p + s = 1$ i~$p = s$ znowu dostajemy $p = s = 1/2$.

\smallskip\noindent
\emph{Weryfikacja:} oba warunki daj\k{a} $\beta_{\mathrm{PPN}} = 1$
jako \textbf{predykcj\k{e}}, nie za\l{}o\.zenie ---
sp\'ojne z~obserwacjami (Cassini: $|\beta - 1| < 2.3 \times 10^{-5}$).

\medskip
Z~antypodalnym anz i~$p = 1/2$:
\[
  f = \sqrt{\PhiZero/\Phi}, \quad h = \sqrt{\Phi/\PhiZero}.
\]
Eksponencjalna kontynuacja:
$f = e^{-\delta\Phi/\PhiZero}$,
$h = e^{+\delta\Phi/\PhiZero}$.
Powy\.zsze wyprowadzenie daje metryki pot\k{e}gowe z~$p = 1/2$,
co odpowiada eksponencjalnej postaci $e^{\pm\delta\Phi/\PhiZero}$
do pierwszego rz\k{e}du. Hipoteza~\ref{hyp:metric-exp}
u\.zywa jednak czynnika~$2$ w~wyk\l{}adniku
($e^{\pm 2\delta\Phi/\PhiZero}$). \'Zr\'od\l{}o rozbie\.zno\'sci
le\.zy w~identyfikacji metrycznej: pot\k{e}gowy anz
$f = (\PhiZero/\Phi)^{1/2}$ daje
$g_{00} \approx -(1 - \tfrac{1}{2}\delta\Phi/\PhiZero)$,
podczas gdy por\'ownanie ze Schwarzschildem wymaga
$g_{00} \approx -(1 - 2G_0 M/(c_0^2 r))
= -(1 - \delta\Phi/\PhiZero)$ przy $\delta\Phi = q_{\mathrm{min}}\PhiZero M/(4\pi r)$
i~$q_{\mathrm{min}} = 8\pi G_0/c_0^2$ (metryka minimalna). Kompensacja polega na
\textbf{redefinicji wyk\l{}adnika}: metryka fizyczna TGP
z~warunkiem $g_{00}|_{\text{1.\ rz.}} = -(1 - 2U)$
(gdzie $U = \delta\Phi/\PhiZero$) wymaga postaci
eksponencjalnej z~czynnikiem~$2$:
\begin{equation}\label{eq:metric-factor2}
  f_{\mathrm{fiz}} = e^{-2\delta\Phi/\PhiZero}, \qquad
  h_{\mathrm{fiz}} = e^{+2\delta\Phi/\PhiZero}.
\end{equation}
Jest to r\'ownowa\.zne antypodalnej metryce pot\k{e}gowej
z~$p = 1$ (a~nie $p = 1/2$), kt\'ora narusza warunek (c)
stwierdzenia ($c_{\mathrm{koord}} \propto \Phi^{-2p}
= \Phi^{-2}$ zamiast $\Phi^{-1/2}$).

\medskip\noindent\textbf{Rozstrzygni\k{e}cie.}
Rozbie\.zno\'s\'c sygnalizuje, \.ze aksjomat~\ref{ax:c}
($c \propto \Phi^{-1/2}$) opisuje \textbf{lokaln\k{a}
fizyczn\k{a}} pr\k{e}dko\'s\'c \'swiat\l{}a, nie
koordynatow\k{a}. Pr\k{e}dko\'s\'c koordynatowa
z~metryki eksponencjalnej to
$c_{\mathrm{koord}}^2 = c_0^2\,e^{-4U} \approx c_0^2(1-4U)$,
natomiast pr\k{e}dko\'s\'c fizyczna (mierzona lokalnymi
przyrz\k{a}dami):
$c_{\mathrm{lok}}^2 = c_0^2\,e^{-2U} \approx c_0^2(1-2U)
\approx c_0^2\,\PhiZero/\Phi$ --- dok\l{}adnie
aksjomat~\ref{ax:c}. Sp\'ojno\'s\'c jest zatem zachowana:
\begin{itemize}[nosep]
  \item Warunek (c) stwierdzenia interpretujemy jako
    warunek na $c_{\mathrm{lok}}$ (nie $c_{\mathrm{koord}}$),
  \item Metryka z~czynnikiem~$2$ jest jedynym wyborem
    daj\k{a}cym jednocze\'snie $c_{\mathrm{lok}}
    = c_0\sqrt{\PhiZero/\Phi}$ i~$\beta_{\mathrm{PPN}} = 1$.
\end{itemize}
\end{proof}

\begin{remark}[Zgodność z~formalizmem PPN Willa]\label{rem:PPN-Will}
W~standardowym formalizmie post-newtonowskim (Will, 2014)
metryka s\l{}abego pola ma posta\'c:
\[
g_{00} = -1 + 2U - 2\beta_{\mathrm{PPN}}\,U^2 + \ldots,\qquad
g_{ij} = \delta_{ij}(1 + 2\gamma_{\mathrm{PPN}}\,U + \ldots),
\]
gdzie $U = GM/(c^2 r)$ jest potencja\l{}em newtonowskim.

\medskip\noindent
\textbf{Metryka TGP} (eksponencjalna, stw.~\ref{hyp:metric-exp})
z~$\delta\Phi/\PhiZero = U$:
\begin{align}
g_{00} &= -e^{-2U} = -(1 - 2U + 2U^2 - \tfrac{4}{3}U^3 + \ldots),\\
g_{ij} &= \delta_{ij}\,e^{+2U} = \delta_{ij}(1 + 2U + 2U^2 + \tfrac{4}{3}U^3 + \ldots).
\end{align}
Identyfikacja z~Willem daje:
\[
\boxed{\gamma_{\mathrm{PPN}} = 1,\qquad \beta_{\mathrm{PPN}} = 1}
\]
\textbf{dok\l{}adnie} (nie jako przybli\.zenie), poniewa\.z
$e^{-2U}$ i~$e^{+2U}$ maj\k{a} t\k{e} sam\k{a} warto\'s\'c bezwzgl\k{e}dn\k{a}
wsp\'o\l{}czynnika przy $U$ i~$U^2$.
Ograniczenia obserwacyjne:
$|\gamma - 1| < 2{,}3 \times 10^{-5}$ (Cassini);
$|\beta - 1| < 8 \times 10^{-5}$ (LLR) --- spe\l{}nione automatycznie.

\medskip\noindent
\textbf{Rozbie\.zno\'s\'c z~GR pojawia si\k{e} przy 3PN}:
wsp\'o\l{}czynnik przy $U^3$ wynosi $4/3$ (TGP) vs $7/4$ (Schwarzschild
w~wsp.\ izotropowych). Ograniczenia na ten rz\k{a}d:
$\lesssim 10^{-3}$ (przewidywane z~zaawansowanych test\'ow
grawimetrycznych; obecne dane nie rozstrzygaj\k{a}).
Weryfikacja numeryczna: \texttt{ppn\_3pn\_corrections.py}.
\end{remark}

\begin{remark}[Metryka TGP nie jest konforemna]\label{rem:nie-konforemna}
Metryka~\eqref{eq:diagonal-ansatz} z~$f \neq h$
\textbf{nie jest konforemna} ($g_{\mu\nu} \neq \Omega^2\eta_{\mu\nu}$).
Jest to istotna w\l{}asno\'s\'c TGP: pole $\Phi$ moduluje
sk\l{}adow\k{e} czasow\k{a} i~przestrzenn\k{a} metryki
\emph{odwrotnie} --- g\k{e}stsza przestrze\'n (wy\.zsze $\Phi$)
rozci\k{a}ga cz\k{e}\'s\'c przestrzenn\k{a},
ale \emph{kurczy} cz\k{e}\'s\'c czasow\k{a}
(spowalnia procesy). To koduje fakt, \.ze
wi\k{e}cej wygenerowanej przestrzeni oznacza
wolniejsz\k{a} propagacj\k{e}.
\end{remark}

\begin{remark}[Relacja metryki pot\k{e}gowej i~eksponencjalnej]%
\label{rem:power-vs-exp}
W~pracy wyst\k{e}puj\k{a} dwie postaci metryki:
\begin{enumerate}[nosep]
  \item \textbf{Pot\k{e}gowa (minimalna)}:
    $f = (\PhiZero/\Phi)^{1/2}$,
    $h = (\Phi/\PhiZero)^{1/2}$
    --- wynik stw.~\ref{prop:conformal-unique}, warunki (a)--(d).
  \item \textbf{Eksponencjalna (fizyczna)}:
    $f = e^{-2\delta\Phi/\PhiZero}$,
    $h = e^{+2\delta\Phi/\PhiZero}$
    --- hipoteza~\ref{hyp:metric-exp}, u\.zywana w~emergencji Einsteina.
\end{enumerate}
Zwi\k{a}zek mi\k{e}dzy nimi jest nast\k{e}puj\k{a}cy:

\medskip\noindent
\emph{(i) Rozwijaj\k{a}c pot\k{e}gow\k{a} metrykk\k{e} wok\'o\l{} $\Phi = \PhiZero$:}
\[
  \Bigl(\frac{\PhiZero}{\Phi}\Bigr)^{\!1/2}
  = 1 - \tfrac{1}{2}\,\frac{\delta\Phi}{\PhiZero}
  + O(\delta\Phi^2)
  = e^{-\delta\Phi/(2\PhiZero)} + O(\delta\Phi^2).
\]
Postaci zgadzaj\k{a} si\k{e} do pierwszego rz\k{e}du
\textbf{bez} czynnika~$2$.

\medskip\noindent
\emph{(ii) Czynnik~$2$ w~eksponencie} wynika z~dodatkowego
warunku fizycznego: por\'ownanie z~metryk\k{a} Schwarzschilda
wymaga $g_{00} \approx -(1 - 2U_N)$ gdzie
$U_N = G_0 M/(c_0^2 r) = \delta\Phi/\PhiZero$.
Metryka pot\k{e}gowa z~$p = 1/2$ daje $g_{00} \approx -(1 - U_N/2)$,
wi\k{e}c wymaga przeskalowania $p \to 2p$
w~argumencie eksponenty.

\medskip\noindent
\emph{(iii) Interpretacja:} metryka pot\k{e}gowa z~$p = 1/2$
opisuje \textbf{struktur\k{e} kinetyczn\k{a}} ($c_{\mathrm{lok}}
= c_0 \sqrt{\PhiZero/\Phi}$), natomiast eksponencjalna
z~czynnikiem~$2$ opisuje \textbf{sprz\k{e}\.zenie z~materi\k{a}}
(metryka na kt\'orej cz\k{a}stki testowe poruszaj\k{a} si\k{e}
po geodezyjnych). Formalna relacja:
\begin{equation}\label{eq:metric-power-exp}
  g^{\mathrm{fiz}}_{\mu\nu}
  = \bigl(g^{\mathrm{min}}_{\mu\nu}\bigr)^{\!2},
  \qquad\text{gdzie}\quad
  g^{\mathrm{min}}_{00} = -\!\sqrt{\PhiZero/\Phi},\;
  g^{\mathrm{min}}_{ij} = \sqrt{\Phi/\PhiZero}\,\delta_{ij}.
\end{equation}
Obie postacie s\k{a} r\'ownowa\.zne do pierwszego rz\k{e}du PPN
i~r\'o\.zni\k{a} si\k{e} dopiero w~efektach $O(U_N^2)$
(post-post-Newtonowskich), kt\'ore stanowi\k{a} dost\k{e}pn\k{a}
predykcj\k{e} obserwacyjn\k{a} (por.\ \S\ref{ssec:metryka-deriv}).
\end{remark}

\begin{proposition}[Sp\'ojno\'s\'c z~dzia\l{}aniem]\label{prop:action-consistency}
Wariacja dzia\l{}ania~\eqref{eq:action} wzgl\k{e}dem $\Phi$:
\begin{equation}\label{eq:variation-action}
  \frac{\delta S}{\delta\Phi} = 0
\end{equation}
odtwarza r\'ownanie pola~\eqref{eq:full-field-alt}, w\l{}\k{a}czaj\k{a}c
poprawki od miary $\sqrt{-g_{\mathrm{eff}}}$.
W~szczeg\'olno\'sci cz\l{}on sprz\k{e}\.zenia gradientowego
$\alpha\,(\nabla\Phi)^2/\Phi$ jest \textbf{wygenerowany}
przez wariacje miary:
\begin{equation}
  \delta\bigl(\sqrt{-g_{\mathrm{eff}}}\bigr)
  = \sqrt{-g_{\mathrm{eff}}}\;\frac{4}{\PhiZero}\,\delta\Phi,
\end{equation}
co po cz\k{e}\'sciach daje niezerowy cz\l{}on proporcjonalny
do $(\nabla\Phi)^2/\Phi$ z~efektywnym $\alpha = 2$.
\end{proposition}

\begin{remark}[Metryka jako opis efektywny]
Metryka~\eqref{eq:metric-exp} jest \textbf{opisem efektywnym},
nie fundamentalnym. Fundamentalne jest pole $\Phi$, a~metryka
stanowi \emph{sk\l{}adniow\k{a}} definicj\k{e} odleg\l{}o\'sci
i~czas\'ow mierzonych przez obiekty fizyczne.
W~re\.zimie $\Phi \gg \PhiZero$ (blisko \'zr\'ode\l{}, re\.zim~III)
poj\k{e}cie metryki g\l{}adkiej mo\.ze traci\'c sens ---
konieczne jest sformu\l{}owanie dyskretne na poziomie
substratu $\hat{s}_i$.
\end{remark}

% =============================================================
\subsection{Kwantyzacja pola $\Phi$}\label{ssec:kwantyzacja}
% =============================================================

Pole $\Phi$ nie jest \emph{klasycznym} polem wprowadzonym
z~zewn\k{a}trz --- wynika z~ugrubienia (coarse-grainingu)
operator\'ow substratowych $\hat{s}_i$.
Kwantyzacja $\Phi$ jest zatem \textbf{odziedziczona},
nie na\l{}o\.zona.

\begin{remark}[Kwantyzacja odziedziczona]\label{rem:inherited-quant}
Operatory substratowe $\hat{s}_i$ spe\l{}niaj\k{a} algebr\k{e}
kanoniczn\k{a} na poziomie dyskretnym. Pole $\Phi$ powstaje
jako \'srednia blokowa:
\begin{equation}
  \Phi(x) = \frac{1}{N_{0,\mathrm{lok}}}\sum_{i \in \mathcal{B}(x)}
  \langle\hat{s}_i\rangle,
\end{equation}
gdzie $\mathcal{B}(x)$ jest blokiem ugrubienia wok\'o\l{}~$x$.
Fluktuacje kwantowe $\hat{s}_i$ indukuj\k{a} fluktuacje
$\delta\Phi$, kt\'ore propaguj\k{a} si\k{e} jako kwanty pola ---
fonony przestrzenno\'sci.
\end{remark}

\begin{definition}[Propagator pola $\Phi$]\label{def:propagator}
Propagator Feynmana pola perturbacji $\phi(x) = \Phi(x) - \PhiZero$:
\begin{equation}\label{eq:propagator-x}
  G_\phi(x - y) = \langle 0|\,T\,\phi(x)\,\phi(y)\,|0\rangle,
\end{equation}
gdzie $|0\rangle$ jest stanem pr\'o\.zni TGP ($\Phi = \PhiZero$
wsz\k{e}dzie), a~$T$ oznacza porz\k{a}dkowanie czasowe.
\end{definition}

\begin{proposition}[Propagator w~przestrzeni p\k{e}d\'ow]\label{prop:propagator-k}
W~przestrzeni p\k{e}d\'ow propagator \textbf{statycznego}
r\'ownania pola~\eqref{eq:full-field-alt} ma posta\'c:
\begin{equation}\label{eq:propagator-k}
  \boxed{
    G_\phi(k) = \frac{1}{k^2 + m_{\mathrm{sp}}^2},
  }
\end{equation}
gdzie \textbf{przestrzenna} masa efektywna:
\begin{equation}\label{eq:meff-spatial}
  m_{\mathrm{sp}}^2
  = 3\gamma - 2\beta
\end{equation}
jest wyznaczona przez linearyzacj\k{e}
r\'ownania~\eqref{eq:full-field-alt} wok\'o\l{} $\Phi = \PhiZero$.
Dok\l{}adniej, cz\l{}ony samosprz\k{e}\.zenia
$\beta\Phi^2/\PhiZero - \gamma\Phi^3/\PhiZero^2$ po rozwiniu
$\Phi = \PhiZero + \delta\Phi$ daj\k{a} wsp\'o\l{}czynnik
liniowy $(2\beta - 3\gamma)\,\delta\Phi$ przy operatorze
kinetycznym $\nabla^2(\delta\Phi)$, st\k{a}d
$m_{\mathrm{sp}}^2 = -(2\beta - 3\gamma) = 3\gamma - 2\beta$.

Dla warunku pr\'o\.zniowego $\beta = \gamma$:
$m_{\mathrm{sp}}^2 = \gamma > 0$ ---
perturbacje maj\k{a} profil Yukawa o~d\l{}ugo\'sci
zaniku $\ell_Y = 1/\sqrt{\gamma}$.
\end{proposition}

\begin{remark}[Rozr\'o\.znienie mas: przestrzenna vs kosmologiczna]\label{rem:two-masses}
TGP posiada \textbf{dwie odr\k{e}bne masy efektywne},
wynikaj\k{a}ce z~r\'o\.znych r\'owna\'n:

\begin{enumerate}[nosep]
  \item \textbf{Masa przestrzenna} $m_{\mathrm{sp}}^2 = 3\gamma - 2\beta = \gamma$
    (stw.~\ref{prop:vacuum-stability}): pochodzi z~linearyzacji
    \emph{statycznego} r\'ownania pola~\eqref{eq:full-field-alt}.
    Okre\'sla profil Yukawa perturbacji w~przestrzeni:
    $\delta\Phi \sim e^{-m_{\mathrm{sp}} r}/r$.
    Zawsze dodatnia: \textbf{stabilno\'s\'c lokalna}.

  \item \textbf{Masa kosmologiczna} $m_{\mathrm{cosmo}}^2 = 4c_0^2\gamma/3$
    (r.~\eqref{eq:m-cosmo-exact}): pochodzi z~linearyzacji
    \emph{kosmologicznego} r\'ownania pola~\eqref{eq:Phi-cosmo-exact}
    wok\'o\l{} $\psi = 1$. Okre\'sla cz\k{e}stotliwo\'s\'c
    oscylacji pola kosmologicznego. W~po\l{}\k{a}czeniu
    z~efektywnym potencja\l{}em $W(\psi)$
    (uwaga~\ref{rem:cosmo-vs-static}) daje $w_{\mathrm{DE}}$.
\end{enumerate}

\smallskip\noindent
Potencja\l{} statyczny: $U(\psi) = (\beta/3)\psi^3 - (\gamma/4)\psi^4$;
$U''(1) = 2\beta - 3\gamma = -\gamma < 0$ ---
$\psi = 1$ jest \textbf{maksimum} $U$.
To oznacza, \.ze pole kosmologiczne \emph{niestabilnie}
odchyla si\k{e} od $\psi = 1$ --- ale ekstremalnie wolno
(potencja\l{} jest bardzo p\l{}aski dla $\gamma \sim H_0^2/c_0^2$),
a~tarcie Hubble'a ($3H\dot\psi$) zamra\.za pole
przy $\psi \approx 1$. Wynik: $w_0 = -1$ dok\l{}adnie
(naturalna $\Lambda$, skrypt \texttt{w\_de\_exact.py}).

\smallskip\noindent
R\'o\.znica mi\k{e}dzy $m_{\mathrm{sp}}^2 > 0$ a~$U''(1) < 0$
wynika z~faktu, \.ze r\'ownanie kosmologiczne zawiera
potencja\l{} efektywny $W(\psi)$ (uwaga~\ref{rem:cosmo-vs-static})
\textbf{r\'o\.zny} od statycznego $U(\psi)$: dodatkowe
cz\l{}ony z~miary $\psi^4$ w~dzia\l{}aniu FRW
modyfikuj\k{a} wsp\'o\l{}czynniki.
\end{remark}

\begin{proposition}[Naturalny obci\k{e}cie UV]\label{prop:UV-cutoff}
Substrat dostarcza naturalnego obci\k{e}cia ultrafioletowego
na skali d\l{}ugo\'sci Plancka:
\begin{equation}\label{eq:kmax}
  k_{\max} \sim \frac{1}{\lP}.
\end{equation}
Mody $k > k_{\max}$ nie maj\k{a} sensu fizycznego ---
odpowiada\l{}yby odleg\l{}o\'sciom mniejszym ni\.z jeden element
substratu. W~konsekwencji \textbf{brak rozbierze\'zno\'sci UV}
w~teoriach perturbacyjnych zbudowanych na propagatorze~\eqref{eq:propagator-k}:
ka\.zda ca\l{}ka p\k{e}tlowa jest obci\k{e}ta od g\'ory przez
$k_{\max}$.
\end{proposition}

\begin{remark}[Kwantyzacja pola a~kwantyzacja geometrii]\label{rem:quant-vs-geom}
Kwantyzacja $\Phi$ \textbf{nie jest} kanoniczn\k{a} kwantyzacj\k{a}
grawitacji (w~sensie LQG czy kowariancji ADM). Jest to kwantyzacja
\emph{pola g\k{e}sto\'sci przestrzeni} --- pola skalarnego
o~propagatorze~\eqref{eq:propagator-k}. \textbf{Przy za\l{}o\.zeniu}
hipotezy pomostowej metryki (uw.~\ref{rem:metric-bridge}),
kwantyzacja geometrii \emph{jest sugerowana} jako konsekwencja:
fluktuacje kwantowe $\delta\Phi$
indukowa\l{}yby, przez metryk\k{e}~\eqref{eq:metric-exp}, fluktuacje
kwantowe metryki $\delta g_{\mu\nu}$. To odwr\'aca
standardow\k{a} logik\k{e}: zamiast kwantowa\'c geometri\k{e}
bezpo\'srednio, kwantujemy pole \'zr\'od\l{}owe geometrii.
\end{remark}

% =============================================================
%  III.B: Asymptotyczne bezpiecze\'nstwo -- p\k{e}tla jednoelementowa
% =============================================================

\begin{proposition}[P\k{e}tla jednoelementowa propagatora TGP ---
  sko\'nczono\'s\'c i~asypmotyczne bezpiecze\'nstwo]\label{prop:one-loop-UV}

Rozwijamy pole wok\'o\l{} pr\'o\.zni: $\psi = 1 + \varphi$,
$|\varphi| \ll 1$. Linearyzowany propagator Feynmana:
\begin{equation}\label{eq:prop-TGP-loops}
  G(k)
  = \frac{1}{k^2 + m_{\mathrm{sp}}^2},
  \qquad m_{\mathrm{sp}}^2 = \gamma,
\end{equation}
Potencja\l{} oddziaływa\'n:
$V_{\mathrm{int}}(\varphi) = \frac{\beta}{3!}\varphi^3
- \frac{\gamma}{4!}\varphi^4 + \ldots$ (cz\l{}ony sze\'scienne i~czw\'oryczne
po rozwini\k{e}ciu $U(1+\varphi)$, znormalizowane przez $\PhiZero$).

\medskip\noindent
\textbf{(i) P\k{e}tla kijankowa (tadpole) z~wierzcho\l{}ka czw\'orycznego.}
Ca\l{}ka p\k{e}tlowa w~d~=~3+1 (z~Euklidesow\k{a} regularyzacj\k{a}
substratow\k{a} $|q| \le k_{\max} \sim 1/\lP$):
\begin{equation}\label{eq:tadpole-TGP}
  \Sigma_{\mathrm{tad}}
  = -\frac{\gamma}{\PhiZero}\cdot
    \frac{1}{(2\pi)^3}\int_0^{k_{\max}}\!\!\frac{4\pi q^2\,dq}{q^2 + m_{\mathrm{sp}}^2}
  = -\frac{\gamma}{2\pi^2\PhiZero}
    \!\left[k_{\max} - m_{\mathrm{sp}}\arctan\!\frac{k_{\max}}{m_{\mathrm{sp}}}\right]\!.
\end{equation}
Poniewa\.z $k_{\max} \sim 1/\lP \ll \infty$,
$\Sigma_{\mathrm{tad}}$ jest \textbf{sko\'nczone}.
Poprawka wzgl\k{e}dna do masy:
\begin{equation}\label{eq:delta-msp}
  \frac{\delta m_{\mathrm{sp}}^2}{m_{\mathrm{sp}}^2}
  = -\frac{k_{\max}}{\pi^2\PhiZero\,m_{\mathrm{sp}}}
  = -\frac{\ell_P}{\pi^2\PhiZero}
    \sqrt{\frac{c_0^2}{\gamma}}\,\frac{1}{\lP^2\cdot c_0^2}\cdot\gamma
  \;\sim\; \frac{\ell_P}{\sqrt{\gamma}}\,\frac{1}{\PhiZero}
  \;\sim\; \frac{\ell_P}{\ell_Y\,\PhiZero}
  \;\ll 1,
\end{equation}
gdzie $\ell_Y = 1/\sqrt{\gamma}$ jest d\l{}ugo\'sci\k{a} Yukawa.
Dla $\ell_Y \sim 1/H_0 \gg \lP$ i~$\PhiZero \sim 25$
korekta jest \textbf{zaniedbywalnie ma\l{}a}
($\sim 10^{-60}$).

\medskip\noindent
\textbf{(ii) P\k{e}tla bublikowa (self-energy) z~wierzcho\l{}k\'ow
  sze\'sciennych.}
Samoenergia\footnote{Cz\l{}on kinetyczny $K(\psi) = \psi^4$
po liniaryzacji daje efektywny wierzcho\l{}ek tr\'ojliniowy
$V_3 \propto \beta/\PhiZero^{1/2}$.}
w~granicy $p^2 \ll m_{\mathrm{sp}}^2$:
\begin{equation}\label{eq:self-energy-TGP}
  \Sigma_{\rm SE}(p^2)
  \;\approx\;
  \frac{\beta^2}{16\pi^2\,\PhiZero^2}
  \left[\Lambda_{\rm UV}^2
       - m_{\mathrm{sp}}^2\ln\!\frac{\Lambda_{\rm UV}^2}{m_{\mathrm{sp}}^2}
       + p^2\ln\!\frac{\Lambda_{\rm UV}^2}{m_{\mathrm{sp}}^2}
       + \mathcal{O}(p^4/m_{\mathrm{sp}}^2)\right],
\end{equation}
gdzie $\Lambda_{\rm UV} = k_{\max}$.
Oba cz\l{}ony rozbie\.zno\'sci ($\Lambda^2$ i~$\Lambda^0\ln\Lambda$)
s\k{a} sko\'nczone dla sko\'nczonego $k_{\max}$~\textbf{bez
regularyzacji ad hoc}.
\end{proposition}

\begin{hypothesis}[Kandydat na UV-punkt sta\l{}y TGP --- \statuslabel{Hipoteza}]%
\label{prop:asymptotic-safety}
\textbf{Argument cutoffowy} (solidny): substrat narzuca twarde UV-obci\k{e}cie
$k_{\max} \sim 1/\lP$ (stw.~\ref{prop:UV-cutoff}),
dzi\k{e}ki czemu wszystkie ca\l{}ki p\k{e}tlowe s\k{a} sko\'nczone
bez renormalizacji. To nie jest hipoteza --- to fakt wynikaj\k{a}cy z~dyskretno\'sci substratu.

\medskip
\textbf{Argument RG} (s\l{}aby --- wymaga weryfikacji):
Definiujemy bezwymiarowe sprz\k{e}\.zenia
obcinaj\k{a}c przy skali~$k$ (Wilsonowski ERG):
\begin{equation}\label{eq:dimless-couplings}
  \tilde{m}^2(k) \equiv \frac{m_{\mathrm{sp}}^2(k)}{k^2},
  \qquad
  \tilde{g}(k) \equiv \frac{\gamma(k)}{k\,\PhiZero^2}.
\end{equation}
Przep\l{}yw Wilsonowski (gra\-ni\-ca jednoelementowa):
\begin{align}
  k\,\frac{d\tilde{m}^2}{dk}
  &= -2\tilde{m}^2
    + \frac{\tilde{g}}{2\pi^2\,(1 + \tilde{m}^2)^2},
    \label{eq:beta-mass}\\
  k\,\frac{d\tilde{g}}{dk}
  &= -\tilde{g}
    - \frac{5\,\tilde{g}^2}{2\pi^2\,(1 + \tilde{m}^2)^3}.
    \label{eq:beta-g}
\end{align}
R\'ownania~\eqref{eq:beta-mass}--\eqref{eq:beta-g}
\textbf{w~przybli\.zeniu jednoelementowym} maj\k{a} punkt sta\l{}y UV:
\begin{equation}\label{eq:fixed-point}
  \tilde{m}^2_* = 0,
  \qquad
  \tilde{g}_* = -\frac{2\pi^2}{5}.
\end{equation}

\medskip
\textbf{Ocena argumentu RG:}
\begin{center}\footnotesize
\begin{tabular}{@{}lll@{}}
\toprule
\textbf{Sk\l{}adnik argumentu} & \textbf{Status} & \textbf{Uwaga} \\
\midrule
Cutoff $k_{\max} \sim 1/\lP$ z~substratu
  & \checkmark~Solidny & dyskretno\'s\'c substratu \\
Beta-funkcje 1-p\k{e}tlowe --- punkt sta\l{}y
  & \checkmark~S\l{}aby & tylko 1-p\k{e}tla \\
Pe\l{}ne r\'ownanie Wetterika (ERG)
  & $\times$~Brakuje & program 6--12 m. \\
Kontrola operator\'ow wy\.zszych rz\k{e}d\'ow
  & $\times$~Brakuje & nietrywialny dla $K(\psi)=\psi^4$ \\
Wymiar anomalny $\eta_\Phi$ przy punkcie sta\l{}ym
  & $\times$~Brakuje & modyfikuje eksponenty krytyczne \\
\bottomrule
\end{tabular}
\end{center}

\textbf{Wniosek hipotezy}:
TGP jest \textbf{kandydatem} do asymptotycznego bezpiecze\'nstwa
w~sensie Weinberga (1979).
Argument cutoffowy jest solidny.
Argument oparty na beta-funkcjach 1-p\k{e}tlowych jest \emph{wst\k{e}pny}
i~wymaga weryfikacji pe\l{}nym r\'ownaniem Wetterika.
\end{hypothesis}

\begin{remark}[Status III.B i~otwarte zadania]\label{rem:IIIb-status}
Powy\.szy szkic (props.~\ref{prop:one-loop-UV}--\ref{prop:asymptotic-safety})
zamyka gap III.B na poziomie \emph{zarys formalny}:
\begin{itemize}[nosep]
  \item P\k{e}tla jednoelementowa jest sko\'nczona dzi\k{e}ki
    substratowemu obci\k{e}ciu UV --- \textbf{wynik nowy}.
  \item Beta-funkcje~\eqref{eq:beta-mass}--\eqref{eq:beta-g}
    i~punkt sta\l{}y~\eqref{eq:fixed-point}
    zweryfikowane numerycznie
    (skrypt \texttt{p123\_uv\_fixed\_point\_search.py}).
    Numeryczny punkt sta\l{}y: $(\tilde m^2_*, \tilde g_*)
    = (-0{,}091,\;-2{,}97)$, rezydualna $||\beta|| = 10^{-11}$.
  \item W~rozszerzonej trunkacji z~wymiarem anomalnym
    $\eta = -\tilde g^2/(6\pi^2(1+\tilde m^2)^4)$:
    FP przetrwa\l{}: $(\tilde m^2_*, \tilde g_*) = (-0{,}072,\;-2{,}65)$
    z~$\eta^* = -0{,}16$ i~jednym kierunkiem relewantnym
    ($\theta_1 = +2{,}02$).
  \item \textbf{Implikacja}: $\alpha_{\rm TGP} = 2 + \eta \approx 1{,}84$
    (korekta $\sim 8\%$). Warto\'s\'c ta jest daleko od dyskretnej granicy
    substratowej $\alpha = 2$ --- jest to artefakt przybli\.zenia
    jednoelementowego, kt\'ore nie uwzgl\k{e}dnia
    sko\'nczonoelementowej struktury $K(\psi) = \psi^4$.
  \item Dodatek~\ref{app:M-erg} (sesja v42) pokazuje, \.ze
    pe\l{}ny przep\l{}yw Wetterika (nie jednoelementowy)
    z~$K(\psi) = \psi^4$ \textbf{stabilizuje pr\'o\.zni\k{e}} (8/8 PASS)
    z~marginaln\k{a} mas\k{a} $m^2_{\mathrm{phys}} \approx 0{,}31$.
\end{itemize}
\textbf{Wniosek}: TGP nie wykazuje anomalii jednoelementowych (sko\'nczone
$\Sigma$) i~posiada kandydata na UV-punkt sta\l{}y. Teoria jest
\emph{kandydatem} do asymptotycznego bezpiecze\'nstwa.
\end{remark}

% =============================================================
\subsection{Emergencja sygnatury lorentzowskiej}\label{ssec:sygnatura}
% =============================================================

Metryka efektywna~\eqref{eq:metric-exp} zak\l{}ada
sygnatur\k{e} $({-},{+},{+},{+})$, ale ta nie jest
postulowana na poziomie substratu. Substrat
$\Gamma = (V, E)$ jest \textbf{czysto relacyjny}
--- nie posiada \.zadnej sygnatury. Pytanie: sk\k{a}d
bierze si\k{e} rozr\'o\.znienie czasu od przestrzeni?

\begin{proposition}[Emergencja sygnatury z~przej\'scia fazowego]\label{prop:sygnatura}
Sygnatura lorentzowska wynika z~\textbf{dw\'och niezale\.znych
fakt\'ow}:
\begin{enumerate}[nosep]
  \item \textbf{Trzy wymiary przestrzenne} ($d = 3$):
    wynikaj\k{a} z~optymalno\'sci substratu
    (stwierdzenie~\ref{prop:wymiar}).
  \item \textbf{Jeden wymiar czasowy}: wynika z~relacyjnej
    definicji czasu (stwierdzenie~\ref{prop:czas}) ---
    czas jest \textbf{parametrem porz\k{a}dkuj\k{a}cym}
    stany w~$\mathcal{S}_1$. Relacja porz\k{a}dku
    ,,przed/po'' jest totalnym porz\k{a}dkiem
    na \'scie\.zkach przyczynowych, co daje dok\l{}adnie
    jeden wymiar.
\end{enumerate}
Sygnatura $({-},{+},{+},{+})$ koduje fakt, \.ze czas
i~przestrze\'n \textbf{maj\k{a} r\'o\.zne pochodzenie ontologiczne}:
\begin{itemize}[nosep]
  \item Wymiary \textbf{przestrzenne} s\k{a} generowane przez
    pole $\Phi$ --- istniej\k{a} w~metrycznym sensie tylko
    tam, gdzie $\Phi > 0$.
  \item Wymiar \textbf{czasowy} jest emergentny z~dynamiki
    substratu --- istnieje, gdy s\k{a} co najmniej dwa
    odr\k{e}bne stany ($\Ffield > 0$,
    stwierdzenie~\ref{prop:czas}).
\end{itemize}
Znaki w~metryce:
$g_{00} < 0$ i~$g_{ij} > 0$ odzwierciedlaj\k{a} t\k{e}
asymetri\k{e}. Wsp\'olny czynnik $\Phi/\PhiZero$
skaluje sk\l{}adow\k{a} przestrzenn\k{a} (wi\k{e}cej
wygenerowanej przestrzeni), podczas gdy $\PhiZero/\Phi$
skaluje sk\l{}adow\k{a} czasow\k{a} (spowalnianie
propagacji w~g\k{e}stej przestrzeni).
\end{proposition}

\begin{remark}[Dlaczego nie sygnatura euklidesowa?]
W~sformalizowaniu substratu jako grafu z~hamiltonianem
$H_\Gamma$ (r\'ownanie~\ref{eq:H-Gamma}) fluktuacje
s\k{a} opisane dynamik\k{a} Langevina
(definicja~\ref{def:relax}) z~parametrem $\tau$
--- r\'ownaniem \emph{dyfuzyjnym} (euklidesowym).
Przej\'scie do sygnatury lorentzowskiej odpowiada
\textbf{kontynuacji Wicka} w~momencie ugrubienia:
na skali grubozarnienia pojawiaj\k{a} si\k{e} mody
propaguj\k{a}ce (drgania $\delta\Phi$ wok\'o\l{}
t\l{}a $\PhiZero$), kt\'ore spe\l{}niaj\k{a} r\'ownanie
falowe z~pr\k{e}dko\'sci\k{a} $c(\Phi)$
(twierdzenie~\ref{thm:covariant-field}).
Sygnatura lorentzowska jest zatem \emph{w\l{}asno\'sci\k{a}
efektywn\k{a}} --- emerguje przy przej\'sciu od
mikrodynamiki substratowej do opisu ci\k{a}g\l{}ego.
\end{remark}

% =============================================================
\subsection{Sektor nielokalny i~warunek bounce}\label{ssec:nielokalny}
% =============================================================

Dotychczasowe r\'ownanie pola~\eqref{eq:full-field-alt} jest
\textbf{lokalne} --- zawiera tylko pochodne $\Phi$ w~danym punkcie.
Na poziomie substratu istniej\k{a} jednak korelacje
mi\k{e}dzy odleg\l{}ymi elementami $\hat{s}_i$, $\hat{s}_j$,
kt\'ore generuj\k{a} \textbf{sektor nielokalny} teorii.

\begin{definition}[J\k{a}dro nielokalne]\label{def:kernel}
J\k{a}dro nielokalne $K(x-y)$ definiujemy przez korelatory
substratowe:
\begin{equation}\label{eq:kernel-def}
  K(x - y) = \frac{1}{\Nzero^2}\sum_{i,j}
  \langle\hat{s}_i\,\hat{s}_j\rangle_c\;
  \delta\bigl(x - C(i)\bigr)\,\delta\bigl(y - C(j)\bigr),
\end{equation}
gdzie $\langle\cdot\rangle_c$ oznacza kumulant (cz\k{e}\'s\'c
sp\'ojn\k{a}), a~$C(i)$ jest funkcj\k{a} zanurzenia elementu $i$
w~przestrze\'n efektywn\k{a}.
\end{definition}

\begin{proposition}[Posta\'c j\k{a}dra po ugrubieniu]\label{prop:kernel-form}
Po ugrubieniu (coarse-grainingu) j\k{a}dro przyjmuje posta\'c:
\begin{equation}\label{eq:kernel-coarse}
  K(r) \sim K_0\;r^{-(d-2+\eta)}\;
  e^{-r/\xi_{\mathrm{corr}}},
\end{equation}
gdzie:
\begin{itemize}[nosep]
  \item $\eta$ jest wyk\l{}adnikiem anomalnego skalowania
    (odej\'scie od klasycznego zachowania),
  \item $\xi_{\mathrm{corr}}$ jest d\l{}ugo\'sci\k{a} korelacji
    substratu,
  \item $K_0$ jest amplitud\k{a} okre\'slon\k{a} przez
    mikroskopowe sta\l{}e sprz\k{e}\.zenia.
\end{itemize}
\end{proposition}

\begin{proposition}[Dwa re\.zimy j\k{a}dra]\label{prop:kernel-regimes}
J\k{a}dro~\eqref{eq:kernel-coarse} wykazuje dwa re\.zimy:
\begin{enumerate}[nosep]
  \item \textbf{$r \ll \xi_{\mathrm{corr}}$} (re\.zim nielokalny):
    $K(r) \sim K_0\,r^{-(d-2+\eta)}$ --- pot\k{e}gowe zanikanie.
    Korelacje substratowe s\k{a} silne; efekty nielokalne TGP
    s\k{a} istotne. R\'ownanie pola wymaga cz\l{}onu
    ca\l{}kowego:
    \begin{equation}\label{eq:nonlocal-term}
      \nabla^2\Phi(x) + \ldots
      + \int K(x-y)\,\Phi(y)\,d^3y = -q\,\PhiZero\,\rho(x).
    \end{equation}
  \item \textbf{$r \gg \xi_{\mathrm{corr}}$} (re\.zim lokalny):
    $K(r) \sim e^{-r/\xi_{\mathrm{corr}}} \to 0$ ---
    korelacje wygasaj\k{a} wyk\l{}adniczo. Cz\l{}on nielokalny
    jest zaniedbywalny, r\'ownanie pola redukuje si\k{e} do
    postaci lokalnej~\eqref{eq:full-field-alt}, a~ta --- do OTW.
\end{enumerate}
\end{proposition}

\begin{hypothesis}[Warunek bounce]\label{hyp:bounce}
Je\.zeli j\k{a}dro $K(x-y)$ zmienia znak na du\.zych
odleg\l{}o\'sciach relacyjnych (antyferromagnetyczna frustracja
substratu), to pole $\Phi$ mo\.ze unika\'c singularno\'sci
$\Phi \to 0$. W~szczeg\'olno\'sci:

Je\'sli istnieje skala $r_{\mathrm{AF}} > \xi_{\mathrm{corr}}$
taka, \.ze $K(r) < 0$ dla $r > r_{\mathrm{AF}}$, to
w~otoczeniu $\Phi \to 0$ cz\l{}on nielokalny staje si\k{e}
\textbf{\'zr\'od\l{}em efektywnym}:
\begin{equation}\label{eq:bounce-condition}
  \int K(x-y)\,\Phi(y)\,d^3y \;<\; 0
  \qquad\text{dla } \Phi(x) \to 0.
\end{equation}
Ten ujemny wk\l{}ad dzia\l{}a jak efektywna g\k{e}sto\'s\'c
\'zr\'od\l{}owa generuj\k{a}ca $\Phi > 0$ --- pole ,,odbija si\k{e}''
(\emph{bounce}) od zera zamiast kolapsowa\'c do singularno\'sci.

W~kosmologii warunek~\eqref{eq:bounce-condition} umo\.zliwia
\textbf{scenariusz cykliczny}: Wszech\'swiat unika Big Crunch
(gdzie $\Phi \to 0$ globalnie), przechodz\k{a}c przez faz\k{e}
odbicia do nowej ekspansji.
\end{hypothesis}

% =============================================================
\subsection{Warunki energetyczne pola $\Phi$}\label{ssec:warunki-energ}
% =============================================================

W~OTW twierdzenia o~osobliwo\'sciach (Penrose, Hawking) wymagaj\k{a}
spe\l{}nienia \textbf{warunk\'ow energetycznych} przez materi\k{e}.
W~TGP metryka jest emergentna, ale efektywny tensor energii-p\k{e}du
pola $\Phi$ mo\.ze by\'c wyznaczony z~dzia\l{}ania~\eqref{eq:action}.

\begin{proposition}[Efektywny tensor energii-p\k{e}du pola $\Phi$]\label{prop:T-Phi}
Tensor energii-p\k{e}du zwi\k{a}zany z~polem $\Phi$ w~metryce
efektywnej~\eqref{eq:metric-exp} ma posta\'c:
\begin{equation}\label{eq:T-Phi}
    T_{\mu\nu}^{(\Phi)} = \frac{1}{\kappa}\biggl[
      \frac{f(\Phi)}{\PhiZero^2}\left(
        \partial_\mu\Phi\,\partial_\nu\Phi
        - \frac{1}{2}\,g_{\mu\nu}^{\mathrm{eff}}\,
          g_{\mathrm{eff}}^{\alpha\beta}\,
          \partial_\alpha\Phi\,\partial_\beta\Phi
      \right)
      - g_{\mu\nu}^{\mathrm{eff}}\,U(\Phi)
    \biggr],
\end{equation}
gdzie $f(\Phi) = 1 + 2\alpha\ln(\Phi/\PhiZero)$.
\end{proposition}

\begin{proposition}[Naruszenie silnego warunku
  energetycznego]\label{prop:SEC-violation}
Pole $\Phi$ \textbf{narusza silny warunek energetyczny}
(SEC: $T_{\mu\nu}u^\mu u^\nu + \frac{1}{2}T \geq 0$
dla ka\.zdego wektora czasopodobnego $u^\mu$)
w~dw\'och re\.zimach:
\begin{enumerate}[nosep]
  \item \textbf{Re\.zim kosmologiczny} ($\dot\varphi \neq 0$,
    $U(\varphi) > 0$): naruszenie SEC odpowiada za
    \textbf{przyspieszon\k{a} ekspansj\k{e}} (ciemna energia,
    sekcja~\ref{sec:ciemna}). Jest to po\.z\k{a}dane i~sp\'ojne
    z~obserwacjami.
  \item \textbf{Re\.zim czarnodziarowy} ($\Phi \to \infty$):
    cz\l{}on potencja\l{}owy $-\gamma\Phi^4/\PhiZero^4$
    dominuje, daj\k{a}c ujemn\k{a} efektywn\k{a} g\k{e}sto\'s\'c
    energii. Naruszenie SEC \textbf{uniemo\.zliwia formowanie
    osobliwo\'sci} --- twierdzenie Penrose'a nie stosuje si\k{e}.
\end{enumerate}
Jednocze\'snie pole $\Phi$ \textbf{spe\l{}nia s\l{}aby warunek
energetyczny} (WEC: $T_{\mu\nu}u^\mu u^\nu \geq 0$)
w~re\.zimie s\l{}abego pola ($|\varphi| \ll 1$),
co gwarantuje stabilno\'s\'c pr\'o\.zni TGP.
\end{proposition}

\begin{remark}[Sp\'ojno\'s\'c z~usuni\k{e}ciem
  osobliwo\'sci]\label{rem:energ-sing}
Naruszenie SEC przez pole $\Phi$ jest \textbf{konieczne
i~wystarczaj\k{a}ce} do usuni\k{e}cia osobliwo\'sci w~TGP:
\begin{itemize}[nosep]
  \item Bez naruszenia SEC twierdzenia Penrose'a wymusza\l{}yby
    osobliwo\'sci w~ka\.zdym scenariuszu kolapsowym ---
    sprzeczno\'s\'c ze stanem zamro\.zonym
    (stwierdzenie~\ref{prop:cd}).
  \item Z~naruszeniem SEC kolaps zatrzymuje si\k{e} przy
    $\Phi \to \infty$, $c \to 0$ (stan zamro\.zony),
    zamiast dochodzi\'c do osobliwo\'sci metrycznej.
\end{itemize}
Jest to \textbf{strukturalna r\'o\.znica} wzgl\k{e}dem
standardowej OTW, gdzie usuni\k{e}cie osobliwo\'sci
wymaga specjalnych mechanizm\'ow (kwantowa grawitacja).
W~TGP usuni\k{e}cie jest \emph{naturaln\k{a} konsekwencj\k{a}}
cz\l{}onu $\gamma$ w~potencjale.
\end{remark}

% =============================================================
\subsection{Entropia relacyjna z~dynamiki $\Phi$}\label{ssec:Srel-formal}
% =============================================================

Entropia relacyjna $\Srel$ (definicja~\ref{def:Srel}) zosta\l{}a
wprowadzona jako miara r\'o\.znorodno\'sci konfiguracji substratu.
Poni\.zej formalizujemy jej zwi\k{a}zek z~polem $\Phi$
i~wyprowadzamy drugie prawo TGP (hipoteza~\ref{hyp:second-law})
z~r\'owna\'n pola.

\begin{proposition}[$\Srel$ jako funkcjona\l{} $\Phi$]\label{prop:Srel-Phi}
Grubozarniona g\k{e}sto\'s\'c substratu (definicja~\ref{def:Srel})
jest zwi\k{a}zana z~polem $\Phi$ przez:
\begin{equation}\label{eq:rho-cg-from-Phi}
    \rho_{\mathrm{cg}}(x)
    = \frac{\Phi(x)}{\int_\Sigma \Phi(y)\,\sqrt{h}\,d^3y}.
\end{equation}
Jest to normalizowane pole: $\int \rho_{\mathrm{cg}}\,\sqrt{h}\,d^3x = 1$.
Interpretacja: w~regionach o~wy\.zszym $\Phi$ substrat jest
g\k{e}\'sciejszy (wi\k{e}cej stopni swobody na jednostk\k{e}
obj\k{e}to\'sci koordynatowej). Wtedy:
\begin{equation}\label{eq:Srel-from-Phi}
    \Srel[\Phi] = -\int_\Sigma \frac{\Phi}{\bar{\Phi}V}
    \ln\!\frac{\Phi}{\bar{\Phi}V}\;\sqrt{h}\,d^3x,
\end{equation}
gdzie $\bar{\Phi} = V^{-1}\int\Phi\,\sqrt{h}\,d^3x$ jest
\'sredni\k{a} warto\'sci\k{a}, a~$V = \int\sqrt{h}\,d^3x$.
\end{proposition}

\begin{proposition}[Drugie prawo TGP z~dyfuzji $\Phi$]\label{prop:second-law-formal}
W~fazie $\Phi > 0$ r\'ownanie pola~\eqref{eq:full-field-alt}
w~re\.zimie s\l{}abopolowym (pomijaj\k{a}c cz\l{}ony $\beta$,
$\gamma$ dominuj\k{a}ce na ma\l{}ych skalach) redukuje si\k{e} do:
\begin{equation}\label{eq:Phi-diffusion}
    \partial_t \Phi \approx D_{\mathrm{eff}}\,\nabla^2\Phi
    + q\,\PhiZero\,\rho,
\end{equation}
gdzie $D_{\mathrm{eff}}$ jest efektywnym wsp\'o\l{}czynnikiem
dyfuzji. R\'ownanie dyfuzji z~dodatnim \'zr\'od\l{}em
\textbf{r\'ownomiernie rozprasza} gradienty $\Phi$:
jednorodne $\Phi = \mathrm{const}$ jest atraktorem.
Jednorodne $\rho_{\mathrm{cg}}$ maksymalizuje $\Srel$.
Zatem:
\begin{equation}
    \frac{d\Srel}{dt} \geq 0,
\end{equation}
z~r\'owno\'sci\k{a} jedynie dla $\Phi = \mathrm{const}$ wsz\k{e}dzie
(pe\l{}na r\'ownowaga termodynamiczna).
\end{proposition}

Powy\.zszy wynik ogranicza si\k{e} do rezimu s\l{}abopolowego.
Poni\.zej wzmacniamy go do pe\l{}nej nieliniowej dynamiki.

\begin{proposition}[Drugie prawo TGP --- pe\l{}na wersja
  nieliniowa]\label{prop:second-law-nonlinear}
W~rezimie nadkrytycznie t\l{}umionym (ekspansja kosmologiczna:
$3H\dot\Phi \gg \ddot\Phi$) pe\l{}ne r\'ownanie
pola~\eqref{eq:covariant-field} redukuje si\k{e} do
\textbf{przep\l{}ywu gradientowego} po funkcjonale
energii~\eqref{eq:energy-corrected}:
\begin{equation}\label{eq:gradient-flow}
    \tau_H\,\partial_t\Phi
    = -\frac{\delta E[\Phi]}{\delta\Phi},
    \qquad \tau_H \equiv 3H,
\end{equation}
gdzie $E[\Phi]$ jest funkcjona\l{}em
energii~\eqref{eq:energy-corrected}.

\medskip\noindent
\textbf{Teza 1 (dyssypacja energii):}
Funkcjona\l{} $E[\Phi]$ jest funkcj\k{a} Lapunowa:
\begin{equation}\label{eq:energy-dissipation}
    \frac{dE}{dt}
    = \int \frac{\delta E}{\delta\Phi}\;\partial_t\Phi\;d^3x
    = -\frac{1}{\tau_H}
      \int \!\biggl(\frac{\delta E}{\delta\Phi}\biggr)^{\!2} d^3x
    \;\leq\; 0,
\end{equation}
z~r\'owno\'sci\k{a} tylko w~stanie stacjonarnym
($\delta E/\delta\Phi = 0$ wsz\k{e}dzie).

\medskip\noindent
\textbf{Teza 2 (wzrost entropii relacyjnej):}
Jedynym jednorodnym minimum $E[\Phi]$ przy ustalonym
$N_\Phi = \int\Phi\,d^3x$ jest $\Phi = \bar\Phi = \mathrm{const}$
(wniosek~\ref{prop:vacuum-condition}). Jednorodne $\Phi$
maksymalizuje $\Srel$ (twierdzenie~\ref{prop:Srel-Phi}).
Przep\l{}yw gradientowy~\eqref{eq:gradient-flow} monotonically
zbli\.za $\Phi$ do tego stanu, zatem:
\begin{equation}\label{eq:dSrel-full}
    \boxed{\frac{d\Srel}{dt} \;\geq\; 0}
    \qquad
    \text{(pe\l{}ny re\.zim nieliniowy, nadkrytyczne t\l{}umienie)}.
\end{equation}
\end{proposition}

\begin{proof}
\textbf{Krok~1}: Przej\'scie do przep\l{}ywu gradientowego.
W~rezimu kosmologicznym ($t \gg H^{-1}_{\mathrm{init}}$)
tarcie Hubble'owskie $3H\dot\Phi$ dominuje nad
$c^{-2}\ddot\Phi$ (stosunek rz\k{e}du $\dot H/H^2 \ll 1$).
R\'ownanie~\eqref{eq:covariant-field} redukuje si\k{e} do
pierwszorz\k{e}dowego:
\begin{equation}
    3H\,\partial_t\Phi
    \approx \nabla^2\Phi
    + \alpha\,\frac{(\nabla\Phi)^2}{\Phi}
    + \beta\,\frac{\Phi^2}{\PhiZero}
    - \gamma\,\frac{\Phi^3}{\PhiZero^2}
    + q\,\PhiZero\,\rho.
\end{equation}
Prawa strona jest dok\l{}adnie $-\delta E[\Phi]/\delta\Phi$
(por.~wariacja~\eqref{eq:energy-functional}), co daje
r\'ownanie~\eqref{eq:gradient-flow}.

\medskip\noindent
\textbf{Krok~2}: Dyssypacja energii.
Obliczamy pochodn\k{a} czasow\k{a}:
\[
    \frac{dE}{dt}
    = \int \frac{\delta E}{\delta\Phi}\,\partial_t\Phi\;d^3x
    = -\frac{1}{\tau_H}
      \int \!\biggl(\frac{\delta E}{\delta\Phi}\biggr)^{\!2} d^3x
    \leq 0.
\]
Wyra\.zenie podca\l{}kowe jest kwadratem, wi\k{e}c
nierówno\'s\'c jest ostra: $dE/dt < 0$ dop\'oki pole
nie osi\k{a}gnie stanu stacjonarnego.

\medskip\noindent
\textbf{Krok~3}: Zwi\k{a}zek dyssypacji z~entropi\k{a}.
Funkcjona\l{} energii $E[\Phi]$ przy ustalonym
$N_\Phi = \int\Phi\,d^3x$ posiada jedyne minimum
jednorodne: $\Phi = \bar\Phi$ (cz\l{}on kinetyczny
$\frac{1}{2}f(\Phi)(\nabla\Phi)^2 \geq 0$ znika
tylko dla $\nabla\Phi = 0$; potencja\l{} $U(\Phi)
= \frac{\beta}{3}\Phi^3/\PhiZero - \frac{\gamma}{4}\Phi^4/\PhiZero^2$
ma minimum w~$\Phi = \PhiZero$ przy $\beta = \gamma$).
Jednorodne $\Phi$ daje jednorodne
$\rho_{\mathrm{cg}}$~\eqref{eq:rho-cg-from-Phi},
co maksymalizuje $\Srel$~\eqref{eq:Srel-from-Phi}.
Poniewa\.z przep\l{}yw gradientowy zbli\.za $E$
monotonically do minimum, a~minimum to odpowiada
maksymalnemu $\Srel$, entropia relacyjna musi
monotonicznie rosn\k{a}\'c.

\medskip\noindent
\textbf{Krok~4}: Szacowanie tempa wzrostu.
Produkcj\k{e} entropii mo\.zna oszacowa\'c z~do\l{}u.
Definiuj\k{a}c $\delta\Phi = \Phi - \bar\Phi$ i~u\.zywaj\k{a}c
nier\'owno\'sci logarytmicznej Sobolewa:
\begin{equation}\label{eq:log-sobolev-bound}
    \Srel^{\max} - \Srel[\Phi]
    \leq C_{\mathrm{LS}}\int
      \frac{(\nabla\Phi)^2}{\Phi}\,d^3x,
\end{equation}
gdzie $C_{\mathrm{LS}}$ jest sta\l{}\k{a} logarytmicznej
nier\'owno\'sci Sobolewa. Z~r\'ownania~\eqref{eq:gradient-flow}
i~\eqref{eq:energy-dissipation}:
\begin{equation}
    \frac{d}{dt}\bigl(\Srel^{\max} - \Srel\bigr)
    \leq -\frac{1}{C_{\mathrm{LS}}\,\tau_H}
      \bigl(\Srel^{\max} - \Srel\bigr),
\end{equation}
co daje wyk\l{}adnicz\k{a} relaksacj\k{e}:
$\Srel^{\max} - \Srel(t) \leq
(\Srel^{\max} - \Srel(0))\,e^{-t/(C_{\mathrm{LS}}\tau_H)}$.
\end{proof}

\begin{remark}[Zakres stosowalno\'sci]\label{rem:second-law-scope}
Twierdzenie~\ref{prop:second-law-formal} (s\l{}abe pole) i
twierdzenie~\ref{prop:second-law-nonlinear} (pe\l{}ne nieliniowe)
oba wymagaj\k{a} $\Phi > 0$ wsz\k{e}dzie. Na granicy
$\Phi \to 0$ (warstwa brzegowa nico\'sci, por.
sekcja~\ref{ssec:boundary-layer}) entropia relacyjna
dywerguje logarytmicznie --- co jest fizycznie poprawne:
granica nico\'sci jest stanem o~\emph{minimalnej} entropii
relacyjnej (brak stopni swobody). Drugie prawo TGP
dzia\l{}a zatem w~pe\l{}nym zakresie fazy $\Phi > 0$,
z~wy\l{}\k{a}czeniem samej granicy $\Phi = 0$
(kt\'ora nie nale\.zy do przestrzeni fizycznej).
\end{remark}

\begin{remark}[Strza\l{}ka czasu z~$\Phi$]
Strza\l{}ka czasu w~TGP \textbf{jest sugerowana} przez warunek pocz\k{a}tkowy,
nie przez r\'ownania same w~sobie: \'swie\.ze z\l{}amanie symetrii daje
wysoce niejednorodne $\Phi$ (niska $\Srel$), a~p\'o\'zniejsza
dyfuzja i~relaksacja monotonically zwi\k{e}kszaj\k{a}
$\Srel$. Jest to sp\'ojne z~hipotez\k{a}~\ref{hyp:second-law}
(\statuslabel{Hipoteza}) i~daje jej formalne ugruntowanie w~r\'ownaniach pola.
\end{remark}

% =============================================================
\subsection{Warstwa brzegowa $\Phi \to 0$: granica
  nico\'sci}\label{ssec:boundary-layer}
% =============================================================

Dotychczasowy formalizm zak\l{}ada $\Phi > 0$ (sektor
$\mathcal{S}_1$). Poni\.zej analizujemy zachowanie pola
w~warstwie brzegowej $\Phi \to 0$ --- przej\'sciu mi\k{e}dzy
przestrzeni\k{a} a~nico\'sci\k{a} $\Nzero$.

\begin{proposition}[Asymptotyczny wyk\l{}adnik
  brzegowy]\label{prop:boundary-exponent}
W~otoczeniu punktu $r_0$, gdzie $\Phi(r_0) = 0$,
r\'ownanie pola~\eqref{eq:full-field-alt} w~przybli\.zeniu
dominuj\k{a}cych cz\l{}on\'ow (cz\l{}ony potencja\l{}owe
$\beta\Phi^2/\PhiZero$ i~$\gamma\Phi^3/\PhiZero^2$ s\k{a}
zaniedbywalnie ma\l{}e wzgl\k{e}dem cz\l{}on\'ow
gradientowych) daje profil:
\begin{equation}\label{eq:boundary-profile}
    \boxed{
    \Phi(r) \;\sim\; A\,(r - r_0)^{1/(1+\alpha)}
    \;=\; A\,(r - r_0)^{1/3}
    }
    \qquad (r \to r_0^+),
\end{equation}
gdzie $A > 0$ jest sta\l{}\k{a} amplitudy, a~wyk\l{}adnik
$p = 1/(1+\alpha) = 1/3$ wynika bezpo\'srednio
z~$\alpha = 2$ (wniosek~\ref{cor:alpha-fixed}).
\end{proposition}

\begin{proof}
Szukamy rozwi\k{a}zania pot\k{e}gowego $\Phi = A(r - r_0)^p$
przy $r \to r_0^+$. Pochodne:
\begin{align}
    \Phi' &= Ap\,(r-r_0)^{p-1}, \\
    \Phi'' &= Ap(p-1)\,(r-r_0)^{p-2}, \\
    \frac{(\Phi')^2}{\Phi} &= A\,p^2\,(r-r_0)^{p-2}.
\end{align}
Podstawiaj\k{a}c do r\'ownania jednowymiarodowego
$\Phi'' + \alpha(\Phi')^2/\Phi = 0$ (cz\l{}ony
potencja\l{}owe rz\k{e}du $(r-r_0)^{2p}$ i~$(r-r_0)^{3p}$
s\k{a} subdominuj\k{a}ce dla $p < 1$):
\begin{equation}
    Ap(p-1)(r-r_0)^{p-2} + \alpha A p^2(r-r_0)^{p-2} = 0.
\end{equation}
Dziel\k{a}c przez $Ap(r-r_0)^{p-2}$ (niezerowe):
\begin{equation}
    (p-1) + \alpha p = 0
    \quad\Longrightarrow\quad
    p = \frac{1}{1+\alpha}.
\end{equation}
Dla $\alpha = 2$: $p = 1/3$.
\end{proof}

\begin{remark}[Fizyczna interpretacja wyk\l{}adnika $1/3$]%
\label{rem:exponent-interpretation}
Wyk\l{}adnik $p = 1/3 < 1$ oznacza, \.ze:
\begin{enumerate}[nosep]
  \item Pole $\Phi$ zanika \textbf{wolniej} ni\.z liniowo
    --- przestrze\'n ,,rozrzedza si\k{e}'' \l{}agodnie.
  \item Gradient $\nabla\Phi \sim (r-r_0)^{-2/3} \to \infty$
    --- ,,sztywno\'s\'c'' pola dywerguje na granicy nico\'sci.
  \item Metryka efektywna degeneruje: cz\l{}on przestrzenny
    $g_{ij} \sim \Phi/\PhiZero \to 0$, ale cz\l{}on czasowy
    $g_{00} \sim e^{-2\delta\Phi/\PhiZero} \to 1$ ---
    ,,czas istnieje, ale przestrze\'n znika''.
  \item Wyk\l{}adnik $1/3$ jest \textbf{uniwersalny}: nie
    zale\.zy od $\beta$, $\gamma$, $q$ ani od geometrii
    --- wynika wy\l{}\k{a}cznie z~$\alpha = 2$, kt\'ore
    z~kolei wynika z~symetrii $Z_2$
    (wniosek~\ref{cor:alpha-fixed}).
\end{enumerate}
\end{remark}

\begin{proposition}[Bariera energetyczna
  nico\'sci]\label{prop:energy-barrier}
Osi\k{a}gni\k{e}cie $\Phi = 0$ w~sko\'nczonej odleg\l{}o\'sci
wymaga \textbf{niesko\'nczonej energii} na jednostk\k{e}
powierzchni. Dok\l{}adniej, g\k{e}sto\'s\'c energii
dzia\l{}ania (w~zmiennej $\psi = (\Phi/\PhiZero)^{1/2}$):
\begin{equation}\label{eq:energy-density-boundary}
    \varepsilon_{\mathrm{kin}}
    = \psi^2\,(\nabla\psi)^2
    \sim (r-r_0)^{1/3}\,(r-r_0)^{-5/3}
    = (r - r_0)^{-4/3},
\end{equation}
jest niec\k{a}kowalny przy $r_0$ (poniewa\.z $4/3 > 1$):
\begin{equation}\label{eq:energy-divergence}
    \int_{r_0}^{r_0+\epsilon}
      \varepsilon_{\mathrm{kin}}\,dr
    = \int_{r_0}^{r_0+\epsilon}
      (r - r_0)^{-4/3}\,dr
    \;\to\; \infty.
\end{equation}
Stan $\Phi = 0$ jest zatem \textbf{energetycznie
nieosi\k{a}galny} w~teorii ci\k{a}g\l{}ej.
\end{proposition}

\begin{proof}
Z~profilu~\eqref{eq:boundary-profile}: $\psi =
(\Phi/\PhiZero)^{1/2} \sim (r-r_0)^{1/6}$,
$\nabla\psi \sim (r-r_0)^{-5/6}$.
G\k{e}sto\'s\'c energii w~dzia\l{}aniu~\eqref{eq:action}
zawiera cz\l{}on kinetyczny z~miar\k{a} $\psi^4$:
\[
    \sqrt{-g_{\mathrm{eff}}}\cdot(\nabla\psi)^2/\psi^2
    = \psi^4 \cdot (\nabla\psi)^2/\psi^2
    = \psi^2(\nabla\psi)^2,
\]
co daje~\eqref{eq:energy-density-boundary}.
Ca\l{}ka dywerguje jak $(r-r_0)^{-1/3}\big|_{r_0}
\to \infty$, co dowodzi nieosi\k{a}galno\'sci.
\end{proof}

\begin{corollary}[Naturalna regularyzacja --- odci\k{e}cie
  substratowe]\label{cor:substrate-cutoff}
Dywergencja~\eqref{eq:energy-divergence} jest artefaktem
opisu ci\k{a}g\l{}ego. Na substracie $\Gamma = (V, E)$
minimalna warto\'s\'c pola wynosi:
\begin{equation}\label{eq:Phi-min}
    \Phi_{\min} \;\sim\; \Phi_{\mathrm{skala}}
    \cdot \ell_{\mathrm{sub}}^{1/3},
\end{equation}
gdzie $\ell_{\mathrm{sub}}$ jest efektywn\k{a} skal\k{a}
dyskretyzacji substratu. Poni\.zej $\Phi_{\min}$:
\begin{itemize}[nosep]
  \item opis ci\k{a}g\l{}y traci sens
    (stw.~\ref{prop:hyperbolic}),
  \item dynamika jest rz\k{a}dzona hamiltoninanem
    substratowym~\eqref{eq:H-Gamma},
  \item przej\'scie $\mathcal{S}_1 \to \mathcal{S}_0$ jest
    procesem \textbf{dyskretnym}, nie ci\k{a}g\l{}ym.
\end{itemize}
Fizycznie: przestrze\'n ,,ko\'nczy si\k{e}'' nie w~punkcie
($\Phi = 0$), lecz w~\textbf{warstwie brzegowej} o~grubo\'sci
$\sim \ell_{\mathrm{sub}}$, gdzie poj\k{e}cia geometryczne
(metryka, odleg\l{}o\'s\'c, przyczynowo\'s\'c) trac\k{a}
sens stopniowo.
\end{corollary}

\begin{remark}[Konsekwencje fizyczne]\label{rem:boundary-physics}
Analiza warstwy brzegowej implikuje:
\begin{enumerate}[nosep]
  \item \textbf{Brak ostrych kraw\k{e}dzi}: Wszech\'swiat nie
    ma ,,\'sciany'' --- przestrze\'n zanika fraktalnie
    ($\sim r^{1/3}$) z~dywerguj\k{a}cym kosztem energetycznym.
  \item \textbf{Naturalna kosmiczna cenzura}: formowanie
    si\k{e} region\'ow $\Phi = 0$ wewn\k{a}trz sektora
    $\mathcal{S}_1$ wymaga niesko\'nczonej energii ---
    ,,dziury w~przestrzeni'' nie mog\k{a} powsta\\'c
    dynamicznie.
  \item \textbf{Asymetria $\Phi \to 0$ vs $\Phi \to \infty$}:
    o~ile $\Phi \to \infty$ (czarna dziura) jest
    osi\k{a}galne przy sko\'nczonej energii (cz\l{}on $\gamma$
    stabilizuje, stw.~\ref{prop:cd}), o~tyle
    $\Phi \to 0$ jest barier\k{a} --- ,,nico\'s\'c jest
    trudniejsza do osi\k{a}gni\k{e}cia ni\.z osobliwo\'s\'c''.
  \item \textbf{Sp\'ojno\'s\'c z~drugim prawem}:
    bariera $\Phi = 0$ zapobiega spadkowi $\Srel$ do zera
    w~sko\'nczonym czasie, co jest sp\'ojne
    z~twierdzeniem~\ref{prop:second-law-nonlinear}.
\end{enumerate}
\end{remark}

% =============================================================
\subsection{Zwi\k{a}zek parametr\'ow substratu z~$\beta$, $\gamma$}\label{ssec:substrat-param}
% =============================================================

Hamiltoniano substratu (r\'ow.~\ref{eq:H-Gamma}) zawiera
parametry mikroskopowe $(\mu, m_0^2, \lambda_0, J)$, podczas
gdy r\'ownanie pola~\eqref{eq:full-field-alt} operuje na
sta\l{}ych efektywnych $(\alpha, \beta, \gamma)$.
Poni\.zej podajemy formaln\k{a} relacj\k{e} \l{}\k{a}cz\k{a}c\k{a}
oba poziomy.

\begin{proposition}[Efektywne sta\l{}e z~ugrubienia]\label{prop:eff-params}
W~przybli\.zeniu pola \'sredniego Landaua--Ginzburga
($\langle\hat{s}_i\rangle = v$, potencja\l{} efektywny
$V_{\mathrm{LG}}(\hat{s}) = \frac{m_0^2}{2}\hat{s}^2
+ \frac{\lambda_0}{4}\hat{s}^4 - \frac{Jz}{2}\hat{s}^2$,
$z$ --- koordynacja substratu), rozwijaj\k{a}c wok\'o\l{}
minimum $v^2 = (Jz - m_0^2)/\lambda_0$ i~identyfikuj\k{a}c
$\Phi = v^2/\Phi_{\mathrm{skala}}$ ($\Phi_{\mathrm{skala}}$
--- wymiarowa skala \l{}\k{a}cz\k{a}ca operator $\hat{s}^2$
z~polem ci\k{a}g\l{}ym), uzyskujemy:
\begin{align}
    \beta &= \frac{\lambda_0}{\Phi_{\mathrm{skala}}^2}
             \cdot C_\beta(z, d),
    \label{eq:beta-from-substrate} \\
    \gamma &= \frac{\lambda_0^2}{(Jz - m_0^2)\,\Phi_{\mathrm{skala}}^2}
             \cdot C_\gamma(z, d),
    \label{eq:gamma-from-substrate}
\end{align}
gdzie $C_\beta$, $C_\gamma$ s\k{a} bezwymiarowymi
wsp\'o\l{}czynnikami zale\.znymi od koordynacji~$z$
i~wymiaru efektywnego~$d$ substratu; ich dok\l{}adna
posta\'c wymaga rachunku blokowego coarse-grainingu
(programu renormalizacji na grafie).
\end{proposition}

\begin{remark}[Hierarchia: $\beta/\gamma$ okre\'sla stosunek
  skal]\label{rem:hier-bg}
Ze~wzor\'ow \eqref{eq:beta-from-substrate}--\eqref{eq:gamma-from-substrate}
wynika:
\begin{equation}
    \frac{\beta}{\gamma}
    = \frac{Jz - m_0^2}{\lambda_0}\cdot\frac{C_\beta}{C_\gamma}
    = v^2\cdot\frac{C_\beta}{C_\gamma}.
\end{equation}
Poniewa\.z skale przej\'scia mi\k{e}dzy re\.zimami
(twierdzenie~\ref{thm:three-regimes}) zale\.z\k{a}
od $\beta/\gamma$, stosunek ten jest wyznaczony
przez pr\'o\.zniow\k{a} warto\'s\'c oczekiwan\k{a} $v$
operatora substratowego --- bezpo\'srednio mierzaln\k{a}
na symulacji Monte Carlo substratu.
Warunek $\beta = \gamma$ (t\l{}o pr\'o\.zniowe) odpowiada
$v^2 = \lambda_0/(C_\beta/C_\gamma)$ --- naturalny punkt
w~przestrzeni parametr\'ow.
\end{remark}

\begin{remark}[Status: wyniki RG Migdala--Kadanoffa]\label{rem:mk-rg-results}
Aproksymacja Migdala--Kadanoffa (MK-RG, $d{=}3$, $b{=}2$, decymacja
kumulantowa) ze~wzmocnieniem wi\k{a}zaniowym $b^{d-1} = 4$
potwierdza istnienie \textbf{nietrywialnego punktu sta\l{}ego
Wilsona--Fishera} w~przestrzeni $(r, u) = (m_0^2/J,\; \lambda_0/J)$:
\begin{equation}
  r^* = -2.251, \quad u^* = 3.917, \quad
  |r^*|/u^* = 0.575 \;\approx\; \mathcal{O}(1).
  \label{eq:WF-fp}
\end{equation}
Stosunek $|r^*|/u^* = v^{*2}\!/\lambda_0^*$ jest rz\k{e}du~$1$
--- \textbf{zgodny z~samodostrojeniem pr\'o\.zni TGP}
($\beta = \gamma$ przy $C_\beta \sim C_\gamma$).

Wyk\l{}adnik gradientowy TGP:
$\alpha_{\mathrm{TGP}} = 2 + \mathcal{O}(\eta)
\approx 2.036$ (z~$\eta_{3D} = 0.0364$),
potwierdzony niezale\.znie od schematu MK.

Kluczowy b\l{}\k{a}d naprawiony w~v4: czynnik~$2$
w~renormalizacji masy i~sprzezenia $\phi^4$
($r' = r - 2c_2$, $u' = u - c_4/3$, nie $r - c_2$, $u - c_4/6$)
wynikaj\k{a}cy z~udzia\l{}u ka\.zdego w\k{e}z\l{}a
w~\textbf{dw\'och} ca\l{}kach decymacji.
Skrypt: \texttt{scripts/renormalization\_substrate.py};
wykresy: \texttt{plots/rg\_*.png}.
\end{remark}

% =============================================================
\subsection{Formalne wyprowadzenie continuum z~substratu}\label{ssec:continuum-derivation}
% =============================================================

Dotychczasowy zwi\k{a}zek substratu ($H_\Gamma$) z~r\'ownaniem pola
(\S\ref{ssec:substrat-param}) pozostawia\l{} \textbf{luk\k{e} formaln\k{a}}:
brak jawnego przej\'scia od dynamiki dyskretnej do r\'ownania
ci\k{a}g\l{}ego. Poni\.zej zamykamy t\k{e} luk\k{e} przez
systematyczne coarse-graining.

\begin{theorem}[Emergencja r\'ownania pola z~substratu]\label{thm:emergent-field-eq}
Niech substrat $\Gamma = (V, E)$ o~hamiltoniano~\eqref{eq:H-Gamma}
znajduje si\k{e} w~fazie uporz\k{a}dkowanej ($T < T_c$,
$\langle\hat{s}_i\rangle = v \neq 0$). Definiujemy pole
ci\k{a}g\l{}e przez blokowe \'sredniowanie:
\begin{equation}\label{eq:phi-def-block}
  \Phi(x) \;=\; \frac{1}{N_B}\sum_{i \in \mathcal{B}(x)}
  \langle\hat{s}_i^2\rangle,
\end{equation}
gdzie $\mathcal{B}(x)$ jest blokiem ugrubienia o~rozmiarze
$L_B \gg a_{\mathrm{sub}}$ ($a_{\mathrm{sub}}$ --- sta\l{}a
sieciowa substratu), a~$N_B = |\mathcal{B}|$.

W~granicy ci\k{a}g\l{}ej ($L_B/a_{\mathrm{sub}} \gg 1$,
$L_B \cdot \nabla\Phi/\Phi \ll 1$) dynamika pola $\Phi$
spe\l{}nia:
\begin{equation}\label{eq:emergent-eq}
  \nabla^2\Phi + \alpha_{\mathrm{eff}}\,\frac{(\nabla\Phi)^2}{\Phi}
  + \beta_{\mathrm{eff}}\,\frac{\Phi^2}{\PhiZero}
  - \gamma_{\mathrm{eff}}\,\frac{\Phi^3}{\PhiZero^2}
  = -q_{\mathrm{eff}}\,\PhiZero\,\rho,
\end{equation}
z~efektywnymi sta\l{}ymi:
\begin{align}
  \alpha_{\mathrm{eff}} &= 2 + \mathcal{O}(\eta),
  \label{eq:alpha-eff-substrate} \\
  \beta_{\mathrm{eff}} &= \frac{\lambda_0}{a_{\mathrm{sub}}^2}
    \cdot C_\beta(z, d, L_B/a_{\mathrm{sub}}),
  \label{eq:beta-eff-substrate} \\
  \gamma_{\mathrm{eff}} &= \frac{\lambda_0^2}{v^2\,a_{\mathrm{sub}}^2}
    \cdot C_\gamma(z, d, L_B/a_{\mathrm{sub}}),
  \label{eq:gamma-eff-substrate}
\end{align}
gdzie $\eta$ jest wyk\l{}adnikiem anomalnym substratu.
\end{theorem}

\begin{proof}[Szkic dowodu]
\emph{Krok~1 (Efektywna energia swobodna).}
Energia swobodna substratu po u\'srednieniu blokowym
($\hat{s}^2 \to \Phi$) ma posta\'c funkcjona\l{}u
Landaua--Ginzburga:
\[
  \mathcal{F}[\Phi] = \int d^dx\;\biggl[
    \frac{\kappa_{\mathrm{sub}}}{2}\,(\nabla\Phi)^2
    + \frac{r}{2}\Phi^2 + \frac{u}{4}\Phi^4 + \ldots
  \biggr].
\]
Poniewa\.z $\Phi \geq 0$ (jest kwadratem amplitudy),
minima potencja\l{}u le\.z\k{a} przy $\PhiZero = v^2 > 0$.
Rozwijaj\k{a}c wok\'o\l{} $\PhiZero$ w~zmiennej
$\chi = \Phi/\PhiZero$, uzyskujemy cz\l{}ony $\chi^2$ i~$\chi^3$
odpowiadaj\k{a}ce $\beta$ i~$\gamma$.

\emph{Krok~2 (Sprz\k{e}\.zenie gradientowe).}
Cz\l{}on kinetyczny $(\nabla\Phi)^2/2$ pochodzi ze standardowego
rozszerzenia gradientowego energii swobodnej. Cz\l{}on nieliniowy
$\alpha(\nabla\Phi)^2/\Phi$ pojawia si\k{e} jako cz\l{}on wiod\k{a}cy
wy\.zszych rz\k{e}d\'ow w~$\nabla^2$, wynikaj\k{a}cy z~faktu,
\.ze pole $\Phi = \langle\hat{s}^2\rangle$ jest kwadratowe
w~operatorze substratowym.
Jawnie: $\nabla\langle\hat{s}^2\rangle
= 2\langle\hat{s}\rangle\nabla\langle\hat{s}\rangle
+ 2\mathrm{Cov}(\hat{s},\nabla\hat{s})$,
co po podzieleniu przez $\Phi = \langle\hat{s}^2\rangle$
generuje cz\l{}on $(\nabla\Phi)^2/\Phi$.

\emph{Krok~3 (Element obj\k{e}to\'sciowy).}
W~fazie uporz\k{a}dkowanej efektywna metryka generowana
przez $\Phi$ ma $\sqrt{-g_{\mathrm{eff}}} \propto \Phi^{d/d_{\mathrm{sub}}}$
(w~$d=3$ przy $d_{\mathrm{sub}} = 3$:
$\sqrt{-g} \propto \Phi^4/\PhiZero^4$,
co jest miar\k{a} u\.zyt\k{a} w~dzia\l{}aniu~\eqref{eq:action}).
Wariacja funkcjona\l{}u z~t\k{a} miar\k{a} generuje $\alpha = 2$
(wniosek~\ref{cor:alpha-fixed}). Korekty anomalne z~fluktuacji
substratowych daj\k{a} $\alpha = 2 + \mathcal{O}(\eta)$,
gdzie $\eta \ll 1$ w~3D.

\emph{Krok~4 (\'Zr\'od\l{}o).}
Defekt topologiczny substratu (kink radialny)
dzia\l{}a jako delta-Diracowe \'zr\'od\l{}o po ugrubieniu:
$\rho(x) = \sum_a M_a\,\delta^3(x - x_a)$,
ze sta\l{}\k{a} sprz\k{e}\.zenia $q_{\mathrm{eff}}$
okre\'slon\k{a} przez energi\k{e} defektu i~sil\k{e}
sprz\k{e}\.zenia $\hat{s}^2$ z~polem ci\k{a}g\l{}ym.
\end{proof}

\begin{remark}[Status formalizmu]\label{rem:continuum-status}
Powy\.zsze wyprowadzenie jest \textbf{p\'o\l{}formalnym szkicem},
nie pe\l{}nym dowodem. Brak zamkni\k{e}tych wyra\.ze\'n
na $C_\beta$, $C_\gamma$ jako funkcji parametr\'ow mikro\-skopowych.
Ich wyznaczenie wymaga \textbf{symulacji Monte Carlo}
na grafie regularnym $d = 3$ z~hamiltonianem~\eqref{eq:H-Gamma}
(program: sekcja~\ref{ssec:symulacja}).

Kluczowe \textbf{konsekwencje testowalne} tego wyprowadzenia:
\begin{enumerate}[nosep]
  \item $\alpha = 2$ jest wymuszone przez geometri\k{e} substratu
    (potwierdzenie niezale\.zne od dzia\l{}ania).
  \item Stosunek $\beta/\gamma = v^2 \cdot C_\beta/C_\gamma$
    jest wyznaczalny numerycznie.
  \item Warunek $\beta = \gamma$ odpowiada konkretnemu
    $v_{\mathrm{crit}}^2 = C_\gamma/C_\beta$ ---
    sprawdzalne na symulacji.
\end{enumerate}
\end{remark}

% =============================================================
\subsection{Wyprowadzenie $B(\Phi)$ w~metryce dysformalnej}\label{ssec:B-derivation}
% =============================================================

Metryka dysformalna (hipoteza~\ref{hyp:disformal}) wprowadza
funkcj\k{e} $B(\Phi)$ jako wolny parametr. Poni\.zej wyprowadzamy
jej posta\'c z~zasad wewn\k{e}trznych TGP.

\begin{proposition}[Posta\'c $B(\Phi)$ z~dzia\l{}ania]\label{prop:B-from-action}
Niech dzia\l{}anie TGP ma posta\'c
\begin{equation}\label{eq:action-disformal}
  S = \int d^4x\;\sqrt{-g}\;\biggl[
    \frac{1}{\kappa}\biggl(
      \frac{1}{2}\,g^{\mu\nu}\partial_\mu\Phi\,\partial_\nu\Phi
      - U(\Phi)
    \biggr) + \mathcal{L}_{\mathrm{mat}}
  \biggr],
\end{equation}
z~metryk\k{a} dysformaln\k{a}~\eqref{eq:disformal-metric-formalizm}.
Wtedy sp\'ojno\'s\'c z~r\'ownaniem pola~\eqref{eq:full-field-alt}
w~granicy statycznej wymaga:
\begin{equation}\label{eq:B-form}
  \boxed{
    B(\Phi) = B_0\,\left(\frac{\PhiZero}{\Phi}\right)^2,
  }
\end{equation}
gdzie $B_0$ jest bezwymiarow\k{a} sta\l{}\k{a} sprz\k{e}\.zenia
$|B_0| \ll 1$ (ograniczon\k{a} przez GW170817).
\end{proposition}

\begin{proof}
Wymagamy trzech warunk\'ow:
\begin{enumerate}[nosep]
  \item \textbf{Zgodno\'s\'c z~granic\k{a} newtonowsk\k{a}:}
    W~granicy s\l{}abopolowej ($\Phi \approx \PhiZero$)
    metryka dysformalna musi redukowa\'c si\k{e} do
    konforemnej (bo dotychczasowe wyniki Newton, PPN
    s\k{a} poprawne). Wymaga to $B(\PhiZero)
    \cdot |\nabla\Phi|^2/M_*^4 \ll 1$.
  \item \textbf{Niezmienniczo\'s\'c $\lP$:}
    Cz\l{}on dysformaIny modyfikuje efektywn\k{a} pr\k{e}dko\'s\'c
    propagacji: $c_{\mathrm{eff}}^2 \approx c_0^2(1 - B\dot\Phi^2/(A\,M_*^4))$.
    Aby $\lP = \sqrt{\hbar G/c^3}$ pozosta\l{}a sta\l{}a
    (tw.~\ref{thm:lP}), $B$ musi skalowa\'c si\k{e} tak, by
    poprawka do $c$ by\l{}a jednorodna w~$\Phi$.
    Poniewa\.z $A \propto \Phi/\PhiZero$ i~$\dot\Phi^2/M_*^4$
    jest bezwymiarowe, wymaga to $B \propto \PhiZero/\Phi$
    lub szybszego zanikania.
  \item \textbf{Motywacja fizyczna:} $B$ opisuje anizotropi\k{e}
    przep\l{}ywu generacji przestrzeni. Przep\l{}yw jest
    proporcjonalny do $\nabla\Phi$ --- silniejszy, gdy
    generacja jest intensywniejsza.
    Intensywno\'s\'c generacji na jednostk\k{e} ju\.z-istniej\k{a}cej
    przestrzeni maleje jak $(\PhiZero/\Phi)^2$
    (ka\.zda dodatkowa ,,warstwa'' przestrzeni rozcie\'ncza
    efekt). St\k{a}d $B \propto (\PhiZero/\Phi)^2$.
\end{enumerate}
\end{proof}

\begin{remark}[Pochodzenie cz\l{}onu dysformalnego z~substratu GL]%
\label{rem:B-from-GL}
Cz\l{}on dysformalny $B(\Phi)\,\partial_\mu\Phi\,\partial_\nu\Phi$
ma naturaln\k{a} interpretacj\k{e} w~ramach przej\'scia fazowego
GL (\S\ref{ssec:GL-transition}). Funkcjona\l{} GL (def.~\ref{def:GL-substrate})
zawiera izotropowy cz\l{}on gradientowy $J_{\mathrm{eff}}(\nabla\phi)^2/2$.
Poprawki nast\k{e}pnego rz\k{e}du w~rozwini\k{e}ciu gradientowym:
\begin{equation}\label{eq:GL-next-order}
  \delta\mathcal{F}
  = \int d^3x\;\frac{K_0}{2}\,
    (\partial_i\partial_j\phi)(\partial_i\partial_j\phi)
  + \frac{K_1}{2}\,(\nabla^2\phi)^2 + \ldots
\end{equation}
Po przej\'sciu do zmiennej $\Phi = \phi^2\,\PhiZero/\phi_{\mathrm{ref}}^2$
generuj\k{a} cz\l{}ony typu
$\partial_\mu\Phi\,\partial_\nu\Phi/\Phi^2$ --- czyli
dok\l{}adnie struktur\k{e} dysformaln\k{a}. Czynnik
$1/\Phi^2$ wynika z~$\nabla\Phi/\Phi = 2\nabla\phi/\phi$,
sk\k{a}d $B(\Phi) \propto (\PhiZero/\Phi)^2$ jest
\textbf{najni\.zszym rz\k{e}dem korekcji anizotropowej}
z~substratu, a~nie dokleionym parametrem.

\noindent Sta\l{}a $B_0 \sim K_0/J_{\mathrm{eff}}^2$ mierzy
stosunek sprz\k{e}\.ze\'n wy\.zszego rz\k{e}du do wiod\k{a}cego.
W~naturalnych jednostkach substratu $|B_0| \sim a^2/\xi^2 \ll 1$,
gdzie $a$ jest sta\l{}\k{a} sieci, $\xi$ d\l{}ugo\'sci\k{a}
korelacji. Przewidywanie TGP:
$|B_0| \sim (\lP/\xi)^2 \ll 1$ --- cz\l{}on dysformalny
jest \emph{z natury ma\l{}y}, ale niezerowy.
\end{remark}

\begin{remark}[Ograniczenie GW170817 na $B_0$]\label{rem:B0-bound}
Pr\k{e}dko\'s\'c GW z~cz\l{}onu dysformalnego:
\begin{equation}
  \frac{c_{\mathrm{GW}}^2}{c_0^2}
  \approx 1 - \frac{B_0\,\dot\Phi^2}{\Phi^2\,M_*^4/\PhiZero^2}.
\end{equation}
Ograniczenie $|c_{\mathrm{GW}} - c_0|/c_0 < 10^{-15}$ daje
$|B_0| < M_*^4\,\PhiZero^2/(\dot\Phi^2\,\Phi^2) \cdot 10^{-15}$.
Poniewa\.z $\dot\Phi^2 \ll \Phi^2\,M_*^4/\PhiZero^2$ w~re\.zimie
kosmologicznym (pole zamro\.zone), ograniczenie jest
\textbf{automatycznie spe\l{}nione} dla $|B_0| = \mathcal{O}(1)$.
\end{remark}

\begin{remark}[Stosunek amplitud GW]\label{rem:h-ratio}
Z~wynikami stw.~\ref{prop:disformal-polarization}
i~postaci\k{a}~\eqref{eq:B-form}:
\begin{equation}\label{eq:h-ratio}
  \frac{h_+}{h_{\mathrm{breathing}}}
  \;\sim\; B_0\,\frac{(\nabla\bar\Phi)^2}
  {M_*^4\,\Phi^2/\PhiZero^2}
  \;\sim\; B_0\,\left(\frac{r_S}{r}\right)^2,
\end{equation}
gdzie $r_S = 2G_0 M/c_0^2$ jest promieniem Schwarzschilda
\'zr\'od\l{}a. Dla fuzji binarnych gwiazd neutronowych
($r/r_S \sim \mathcal{O}(1)$):
$h_+/h_{\mathrm{breathing}} \sim B_0$.
Detekcja wymaga $|B_0| \gtrsim 0{,}1$
(ET, CE: osi\k{a}galne przy 3+ detektorach).
\end{remark}

% =============================================================
\subsection{Kanoniczny \l{}a\'ncuch dedukcyjny}\label{ssec:lancuch}
% =============================================================

Poni\.zej prezentujemy zamkni\k{e}ty \l{}a\'ncuch formalny
\[
  \boxed{
  \text{Dzia\l{}anie}
  \;\xrightarrow{(1)}\;
  \text{R\'own.\ pola}
  \;\xrightarrow{(2)}\;
  \text{Pr\'o\.znia}
  \;\xrightarrow{(3)}\;
  \text{Perturbacje}
  \;\xrightarrow{(4)}\;
  \text{Masy eff.}
  \;\xrightarrow{(5)}\;
  \text{Obserwable}
  \;\xrightarrow{(6)}\;
  \text{Wi\k{e}zy}
  }
\]
ka\.zdemu kroku przypisuj\k{a}c precyzyjne r\'ownanie, status
formalny ({\bfseries W}~=~wyprowadzone, {\bfseries H}~=~hipoteza,
{\bfseries A}~=~aksjomat) oraz identyfikator w~tek\'scie.

\bigskip
\noindent\textbf{Krok 0 --- Aksjomaty i~substrat.}
\begin{itemize}[nosep]
  \item Aksjomat \'zr\'od\l{}a (A1, \S\ref{sec:ontologia}):
    przestrze\'n jest generowana przez mas\k{e}-energi\k{e}. [\textbf{A}]
  \item Stan $\Nzero$ (A2): $\Phi = 0$, brak przestrzeni. [\textbf{A}]
  \item Symetria $Z_2$ substratu $\hat s_i \to -\hat s_i$
    (\S\ref{ssec:chiralna}). [\textbf{A}]
  \item Ci\k{a}g\l{}e pole $\Phi(\mathbf{x},t) \geq 0$ jako granica
    kontinuum z~substratu
    (\S\ref{ssec:continuum-derivation}). [\textbf{W}]
\end{itemize}

\bigskip
\noindent\textbf{Krok 1 --- Dzia\l{}anie $\boldsymbol{\to}$ r\'ownanie pola.}\\[2pt]
Dzia\l{}anie TGP (hip.~\ref{hyp:action}):
\begin{equation}\label{eq:chain-action}
  S[\Phi] = \int \mathrm{d}^4x\;\psi^4
  \left[\kappa^{-1}\bigl(\tfrac{1}{2}(\nabla\psi)^2
  - U(\psi)\bigr) + \mathcal{L}_{\mathrm{mat}}\right],
  \quad \psi \equiv \Phi/\PhiZero,
  \tag{$\mathrm{C}$-1}
\end{equation}
z~potencja\l{}em $U(\psi) = \tfrac{\beta}{3}\psi^3
- \tfrac{\gamma}{4}\psi^4$.
Status: \textbf{H} (hipoteza --- forma potencja\l{}u dobrana
tak, by mia\l{}a poprawne symetrie i~granice; kryteria
doboru w~\S\ref{ssec:kryteria}).

Wariacja $\delta S / \delta \Phi = 0$ daje
(tw.~\ref{thm:field-eq}):
\begin{equation}\label{eq:chain-field}
  \nabla^2\Phi + \alpha\,\frac{(\nabla\Phi)^2}{\Phi}
  + \frac{\beta}{\PhiZero}\Phi^2
  - \frac{\gamma}{\PhiZero^2}\Phi^3
  = -q\,\PhiZero\,\rho,
  \tag{$\mathrm{C}$-2}
\end{equation}
ze \textbf{zwi\k{a}zanym} $\alpha = 2$
(wn.~\ref{cor:alpha-fixed}). [\textbf{W}]

\bigskip
\noindent\textbf{Krok 2 --- Warunek pr\'o\.zni i~stabilno\'s\'c.}\\[2pt]
Jednorodna pr\'o\.znia ($\rho = 0$, $\Phi = \mathrm{const}$):
cz\l{}ony gradientowe znikaj\k{a}, pozosta\l{}o\'s\'c wymaga
(stw.~\ref{prop:vacuum-condition}):
\begin{equation}\label{eq:chain-vacuum}
  \beta - \gamma = -q\,\bar\rho
  \quad\xrightarrow{\;\bar\rho \to 0\;}\quad
  \beta = \gamma.
  \tag{$\mathrm{C}$-3}
\end{equation}
Status: [\textbf{W}].

Linearyzacja $\Phi = \PhiZero(1 + \epsilon)$, $|\epsilon| \ll 1$
wok\'o\l{} $\Phi_0$ w~statycznym r\'ownaniu~\eqref{eq:chain-field}
(stw.~\ref{prop:vacuum-stability}):
\begin{equation}\label{eq:chain-msp}
  m_{\mathrm{sp}}^2 = 3\gamma - 2\beta.
  \quad\text{Dla }\beta = \gamma:\;
  m_{\mathrm{sp}}^2 = \gamma > 0\;\;(\text{stabilne}).
  \tag{$\mathrm{C}$-4a}
\end{equation}
Status: [\textbf{W}].

\bigskip
\noindent\textbf{Krok 3 --- Perturbacje kosmologiczne.}\\[2pt]
Metryka FRW z~polem jednorodnym $\Phi(t) = \PhiZero\,\varphi(t)$
+ metryka eksponencjalna (hip.~\ref{hyp:metric-exp}) daj\k{a}
r\'ownanie kosmologiczne~\eqref{eq:Phi-cosmo-exact}
z~potencja\l{}em kosmologicznym $W(\psi) = \frac{7\beta}{3}\psi^2 - 2\gamma\psi^3$
(stw.~\ref{prop:FRW-derivation}), gdzie cz\l{}on $4U/\psi + U'$
z~elementu obj\k{e}to\'sciowego $\psi^4$ daje wsp\'o\l{}czynniki
$7/3$ i~$2$ (r\'o\.zne od statycznych $\beta$, $\gamma$).
Linearyzacja wok\'o\l{} $\psi = 1$
(stw.~\ref{prop:exact-to-linearized}) definiuje
\textbf{drug\k{a}, odr\k{e}bn\k{a} mas\k{e}}
(uw.~\ref{rem:two-masses}):
\begin{equation}\label{eq:chain-mcosmo}
  m_{\mathrm{cosmo}}^2 = \frac{4\,c_0^2\gamma}{3}\,.
  \quad\text{Dla }\beta = \gamma:\;
  m_{\mathrm{cosmo}}^2 > 0
  \;\;(\text{oscylacja t\l{}umiona tarciem Hubble'a}).
  \tag{$\mathrm{C}$-4b}
\end{equation}
Status: [\textbf{W}] (obie masy \emph{wyprowadzone}).

\smallskip
\noindent\emph{Uwaga}: potencja\l{} statyczny $U(\psi) = \frac{\beta}{3}\psi^3 - \frac{\gamma}{4}\psi^4$
ma $U''(1) = 2\beta - 3\gamma = -\gamma < 0$ (maksimum). Stan $\psi = 1$
jest zamro\.zony tarciem Hubble'owskim, nie stanowi minimum energii
(uw.~\ref{rem:psi1-status}). Resztkowy potencja\l{} $U(1) = \gamma/12$
nap\k{e}dza ciemn\k{a} energi\k{e} (stw.~\ref{prop:Lambda-eff}).

\smallskip
\noindent\emph{Rozr\'o\.znienie}:
$m_{\mathrm{sp}}$ kontroluje zasi\k{e}g Yukawy i~dyspersj\k{e} GW
(\S\ref{ssec:perturb}); $m_{\mathrm{cosmo}}$ kontroluje dynamik\k{e}
wolnego staczania i~$w_{\mathrm{DE}}$ (\S\ref{ssec:de-oszacowanie}).
Obydwie wynikaj\k{a} z~tego samego potencja\l{}u, ale
z~r\'o\.znych redukcji (statyczna vs.\ FRW).

Liniowe perturbacje g\k{e}sto\'sci (tw.~\ref{thm:Geff-k}):
\begin{equation}\label{eq:chain-Geff}
  G_{\mathrm{eff}}(k,a) = G_0
  \left(1 + \frac{2\alpha_{\mathrm{eff}}^2}
  {1 + (a\,m_{\mathrm{sp}}/k)^2}\right),
  \quad \alpha_{\mathrm{eff}} = \frac{q_{\mathrm{eff}}}{4\pi},
  \tag{$\mathrm{C}$-5}
\end{equation}
gdzie $q_{\mathrm{eff}} = q\,\PhiZero = 4\pi G_0\PhiZero/c_0^2$.
Status: [\textbf{W}].

\bigskip
\noindent\textbf{Krok 4 --- Ciemna energia.}\\[2pt]
Resztkowy potencja\l{} w~$\varphi = 1$
(stw.~\ref{prop:Lambda-eff}):
\begin{equation}\label{eq:chain-Lambda}
  \Lambda_{\mathrm{eff}} \approx \frac{\gamma}{12}
  \;\sim\; \frac{\PhiZero\,H_0^2}{12\,c_0^2}.
  \quad\text{Naturalny }\PhiZero \sim 25
  \;\Rightarrow\; \Lambda_{\mathrm{eff}} \sim \Lambda_{\mathrm{obs}}.
  \tag{$\mathrm{C}$-6}
\end{equation}
Status: [\textbf{W}]. $\PhiZero \sim 25$ jest
\emph{konsekwencj\k{a}} relacji
$\gamma \sim \PhiZero H_0^2/c_0^2$ i~warunku
$\Lambda_{\mathrm{eff}} \approx \Lambda_{\mathrm{obs}}$
--- nie wolnym parametrem.

\bigskip
\noindent\textbf{Krok 5 --- Obserwable.}\\[2pt]
Z~\eqref{eq:chain-Geff}--\eqref{eq:chain-Lambda}:
\begin{enumerate}[nosep, label=(\alph*)]
  \item $H(z)$ --- zmodyfikowane r\'ownania
    Friedmanna~\eqref{eq:friedmann-TGP} z~$G(\Phi) = G_0/\varphi$,
    $\Lambda_{\mathrm{eff}}(\varphi)$;
  \item $f\sigma_8(z)$ --- wzrost liniowy z~$G_{\mathrm{eff}}(k,a)$;
  \item $w_0, w_a$ --- z~efektywnej g\k{e}sto\'sci DE
    ($w_0 \approx -1$, $w_a < 0$, kwintesencja);
  \item dyspersja GW: $\omega^2 = c_0^2(k^2/a^2 + m_{\mathrm{sp}}^2)$,
    polaryzacja oddechowa (+ tensorowa przy metryce
    dysformalnej, hip.~\ref{hyp:disformal}).
\end{enumerate}

\bigskip
\noindent\textbf{Krok 6 --- Wi\k{e}zy obserwacyjne.}\\[2pt]
Zamkni\k{e}cie \l{}a\'ncucha: obserwable nak\l{}adaj\k{a}
ograniczenia na parametry, kt\'ore wracaj\k{a}
do dzia\l{}ania (krok~1).
\begin{center}\small
\begin{tabular}{@{}lll@{}}
\toprule
\textbf{Wi\k{a}z} & \textbf{\'Zr\'od\l{}o} &
  \textbf{Ograniczenie} \\
\midrule
GW170817 & $|v_g - c_0|/c_0 < 10^{-15}$ &
  $m_{\mathrm{sp}} \lesssim 10^{-22}\;\mathrm{eV}$ \\
BBN & $|\Delta G/G| < 0.1$ w~$T \sim 1\;\mathrm{MeV}$ &
  $\psi_{\mathrm{ini}} = \psi_{\mathrm{eq}} \Rightarrow \Delta G = 0$
  (uw.~\ref{rem:BBN-resolution}) \\
CMB (Planck) & $|\mu_{\mathrm{eff}} - 1| < 0.05$ &
  $\alpha_{\mathrm{eff}}^2 \lesssim 0.015$ \\
DESI BAO & $H(z)/H_{\Lambda\mathrm{CDM}}(z) = 1 \pm 0.02$ &
  $\gamma \lesssim 0.03$ \\
Warunek pr\'o\.zni & $\beta = \gamma + q\bar\rho$ &
  $|\beta - \gamma| < q\bar\rho_0 \sim 10^{-5}$ \\
\bottomrule
\end{tabular}
\end{center}

\begin{remark}[Co jest wyprowadzone, a~co postulowane]\label{rem:chain-status}
W~ca\l{}ym \l{}a\'ncuchu \textbf{postulaty} (elementy nieredukowalne)
to:
\begin{enumerate}[nosep]
  \item aksjomaty A1--A5 (\S\ref{sec:ontologia}),
  \item posta\'c potencja\l{}u $U(\psi)$
    (dobrana z~kryteri\'ow \S\ref{ssec:kryteria},
    ale nie wyprowadzona z~g\l{}\k{e}bszej zasady),
  \item metryka eksponencjalna (hip.~\ref{hyp:metric-exp};
    wyprowadzona z~4 warunk\'ow wewn\k{e}trznych
    w~\S\ref{ssec:metryka-deriv}, ale te warunki
    same s\k{a} w~cz\k{e}\'sci postulowane),
  \item metryka dysformalna (hip.~\ref{hyp:disformal}; status
    otwarty --- patrz uwaga poni\.zej).
\end{enumerate}
Wszystko pozosta\l{}e ---
$\alpha = 2$, $\beta = \gamma$, obie masy, $G_{\mathrm{eff}}(k,a)$,
$\Lambda_{\mathrm{eff}}$, $\PhiZero \sim 25$,
profil $\phi(r)$, trzy re\.zimy ---
jest \textbf{konsekwencj\k{a}} tych czterech
element\'ow.
\end{remark}

\begin{remark}[Status metryki dysformalnej: hipoteza pomostowa]%
\label{rem:disformal-bridge}
Metryka dysformalna (hip.~\ref{hyp:disformal}) pe\l{}ni rol\k{e}
\textbf{hipotezy pomostowej}: zmniejsza napi\k{e}cie mi\k{e}dzy
zakazem mod\'ow tensorowych (tw.~\ref{thm:no-tensor}) a~ich
obserwacj\k{a} w~GW150914/GW170817, ale \textbf{nie rozstrzyga}
go ca\l{}kowicie.

\medskip\noindent
\textbf{Co jest udowodnione:}
\begin{itemize}[nosep]
  \item Konforemna metryka $g_{\mu\nu} = A(\Phi)\,\eta_{\mu\nu}$
    \textbf{nie mo\.ze} generowa\'c tensorowych mod\'ow GW
    (tw.~\ref{thm:no-tensor}: tensor Weyla $\equiv 0$).
  \item Dodanie cz\l{}onu dysformalnego
    $B(\Phi)\,\partial_\mu\Phi\,\partial_\nu\Phi / M_*^4$
    \l{}amie konforemn\k{a} p\l{}asko\'s\'c i~\emph{generuje}
    tensory z~anizotropii \'zr\'od\l{}a
    (stw.~\ref{prop:disformal-polarization}).
\end{itemize}

\medskip\noindent
\textbf{Co jest otwarte:}
\begin{itemize}[nosep]
  \item Posta\'c $B(\Phi) = B_0(\PhiZero/\Phi)^2$ wyprowadzona
    z~warunk\'ow sp\'ojno\'sci (\S\ref{ssec:B-derivation})
    i~potwierdzona jako korekcja wy\.zszego rz\k{e}du GL
    (uw.~\ref{rem:B-from-GL}); wymaga walidacji numerycznej MC.
  \item Skal\k{a} $M_*$ jest wolnym parametrem.
  \item Amplituda tensorowa $h_+/h_{\mathrm{breathing}}
    \sim B_0$ (uw.~\ref{rem:h-ratio}) --- wyznaczenie
    $B_0$ wymaga pe\l{}nego rozwi\k{a}zania sprz\k{e}\.zonego
    pola+metryki.
\end{itemize}
Hipoteza dysformalna jest \textbf{falsyfikowalna}:
je\'sli przysz\l{}e pomiary polaryzacji GW z~sieci
3+ detektor\'ow \textbf{nie} wyka\.z\k{a} sk\l{}adowej
oddechowej, cz\l{}on dysformalny jest zbyt silny
($B_0 \gg 1$, re\.zim GR) i~``czysta TGP'' (konforemna)
jest wykluczona.
\end{remark}

% =============================================================
% =============================================================
\subsection{Sektor fermionowy --- zarys formalizacji}\label{ssec:fermiony}
% =============================================================

W~dotychczasowej wersji TGP materia pojawia si\k{e} jedynie
przez sprz\k{e}\.zenie \'zr\'od\l{}owe $-q\PhiZero\rho$ lub
jako parametr $\mathcal{L}_{\mathrm{mat}}$ w~dzia\l{}aniu.
Poni\.zej formalizujemy minimalny sektor fermionowy, sp\'ojny
z~metryk\k{a} efektywn\k{a} i~dynamicznymi sta\l{}ymi.

\subsubsection{Sprz\k{e}\.zenie Diraca z~metryk\k{a} efektywn\k{a}}

Metryka efektywna TGP (stw.~\ref{prop:conformal-unique})
nie jest konforemnie p\l{}aska (\S\ref{ssec:metryka-deriv}),
lecz \textbf{dysformalna} --- czynniki $f(\Phi)$ i~$h(\Phi)$
maj\k{a} r\'o\.zne wyk\l{}adniki. Dla fermion\'ow wymaga to
u\.zycia pe\l{}nego formalizmu tetradowego.

\begin{definition}[Tetrada TGP]\label{def:tetrada}
Tetrada $e^a{}_\mu$ spe\l{}niaj\k{a}ca
$g_{\mu\nu}^{\mathrm{eff}} = \eta_{ab}\,e^a{}_\mu\,e^b{}_\nu$
ma w~tle izotropowym ($\Phi = \Phi(t,\vec{x})$) posta\'c:
\begin{equation}\label{eq:tetrada}
  e^0{}_0 = c_0\sqrt{\frac{\PhiZero}{\Phi}}\,, \quad
  e^i{}_j = \sqrt{\frac{\Phi}{\PhiZero}}\;\delta^i{}_j\,,
\end{equation}
co odpowiada metryce~\eqref{eq:metric-full-tensor}
z~$f(\Phi) = \PhiZero/\Phi$, $h(\Phi) = \Phi/\PhiZero$,
r\'ownowa\.znie: $e^0{}_0 = c_0\,e^{-U}$,
$e^i{}_j = e^{U}\,\delta^i{}_j$,
gdzie $U = \tfrac{1}{2}\ln(\Phi/\PhiZero)$.
\end{definition}

\begin{hypothesis}[Dzia\l{}anie Diraca w~TGP]\label{hyp:dirac-TGP}
Lagran\.zjan fermionowy sprz\k{e}\.zony z~polem przestrzenno\'sci
ma posta\'c:
\begin{equation}\label{eq:dirac-TGP}
  \mathcal{L}_{\mathrm{D}}
  = \sqrt{-g_{\mathrm{eff}}}\;\bar{\Psi}\!\left[
    i\,e_a{}^\mu\,\gamma^a\!\left(
      \partial_\mu + \tfrac{1}{4}\,\omega_\mu{}^{bc}\,\sigma_{bc}
    \right)
    - \frac{m_0\,c(\Phi)}{\hbar(\Phi)}
  \right]\!\Psi,
\end{equation}
gdzie:
\begin{itemize}[nosep]
  \item $\omega_\mu{}^{bc}$ to koneksja spinowa wyznaczona
    z~tetrady~\eqref{eq:tetrada},
  \item $\sigma_{bc} = \tfrac{1}{4}[\gamma_b,\gamma_c]$ to generatory
    Lorentza w~reprezentacji spinorowej,
  \item masa ``naga'' $m_0$ jest sta\l{}\k{a}; fizyczna masa lokalna
    skaluje si\k{e} jako $m(\Phi) = m_0\,c(\Phi)/\hbar(\Phi) = m_0$
    (poniewa\.z $c/\hbar = c_0/\hbar_0$ jest niezmiennikiem ---
    tw.~\ref{thm:exponents}).
\end{itemize}
\end{hypothesis}

\begin{remark}[Niezmienniczo\'s\'c masy fermionowej]\label{rem:masa-fermion}
Iloraz $c(\Phi)/\hbar(\Phi) = c_0/\hbar_0$ jest niezale\.zny od
$\Phi$ (aksjomat~\ref{ax:c}, tw.~\ref{thm:exponents}). Oznacza to,
\.ze \textbf{masa Diraca nie zale\.zy od pola przestrzenno\'sci} ---
jest sta\l{}\k{a} fundamentaln\k{a} cz\k{a}stki. Fermiony nie ,,czuj\k{a}''
gradientu $\Phi$ bezpo\'srednio przez cz\l{}on masowy; odczuwaj\k{a}
go jedynie przez koneksj\k{e} spinow\k{a} (efekty grawitacyjne).
Jest to sp\'ojne z~obserwacj\k{a}, \.ze masy cz\k{a}stek nie
zale\.z\k{a} od potencja\l{}u grawitacyjnego (zasada r\'ownowa\.zno\'sci).
\end{remark}

\subsubsection{Koneksja spinowa i~dodatkowe sprz\k{e}\.zenie}

\begin{proposition}[Koneksja spinowa TGP do drugiego rz\k{e}du]%
\label{prop:spin-connection}
Niech $\Phi = \PhiZero(1 + \varphi)$ z~$U = \tfrac{1}{2}\ln(1+\varphi)$.
Koneksja spinowa tetrady~\eqref{eq:tetrada} w~polu statycznym
($\dot\varphi = 0$), obliczona z~r\'owna\'n Cartana
$d\theta^a + \omega^a{}_b \wedge \theta^b = 0$, wynosi:
\begin{equation}\label{eq:spin-conn-full}
  \omega_i{}^{0j}
  = -e^{-U}\,\partial_i U\;\delta^{0j}
  = -\frac{\partial_i U}{\sqrt{1+\varphi}}\;\delta^{0j},
\end{equation}
co po rozwini\k{e}ciu do rz\k{e}du $\mathcal{O}(\varphi^2)$ daje:
\begin{equation}\label{eq:spin-conn-2nd}
  \boxed{
    \omega_i{}^{0j}
    \approx -\frac{1}{2}\,\partial_i\varphi
    \left(1 - \frac{3}{2}\,\varphi + \frac{15}{8}\,\varphi^2\right)
    \delta^{0j}
    + \mathcal{O}(\varphi^3).
  }
\end{equation}
Odpowiednio, sk\l{}adowa czasowa:
$\omega_0{}^{0j} = -\frac{c_0}{2}\,(1+\varphi)^{-2}\,\partial_j\varphi$.
Sprz\k{e}\.zenie spin-orbita w~hamiltonianie Diraca
(wchodz\k{a}ce przez $\Gamma_0 = \tfrac{1}{4}\,\omega_0^{ab}\,\sigma_{ab}$):
\begin{equation}\label{eq:spin-orbit-TGP}
  \delta H_{\mathrm{SO}}
  \approx -\frac{1}{8}\,\alpha^i\;\partial_i\varphi
  \left(1 - 2\varphi + \mathcal{O}(\varphi^2)\right),
\end{equation}
gdzie $\alpha^i = \gamma^0\gamma^i$. Cz\l{}on pierwszego rz\k{e}du
$-\frac{c_0}{8}\alpha^i\partial_i\varphi$
jest identyczny z~predykcj\k{a} GR dla r\'ownowa\.znej metryki.
Cz\l{}on drugiego rz\k{e}du $\propto \varphi\,\partial_i\varphi$
daje korekt\k{e}:
\begin{equation}\label{eq:delta2-dirac}
  \frac{\delta^{(2)} E}{E}
  \sim U^2 \sim \left(\frac{GM}{c_0^2\,r}\right)^{\!2}
  \sim \begin{cases}
    10^{-12} & \text{(S\l{}o\'nce)} \\
    10^{-1} & \text{(gwiazda neutronowa)}.
  \end{cases}
\end{equation}
\end{proposition}

\begin{proposition}[Perturbacyjne korekty Diraca: wyniki numeryczne]%
\label{prop:dirac-perturbative}
Stosuj\k{a}c teori\k{e} perturbacji pierwszego rz\k{e}du
$\delta E = \langle\psi|\delta H_{\mathrm{TGP}}|\psi\rangle$
z~tetrad\k{a}~\eqref{eq:tetrada} i~koneksj\k{a}~\eqref{eq:spin-conn-2nd},
otrzymujemy korekty do energii wi\k{a}zania wodoru:
\begin{equation}\label{eq:dirac-perturbative-result}
  \boxed{
    \frac{\delta E_{\mathrm{bind}}}{E_{\mathrm{bind}}}
    \approx \frac{\varphi \, n^2}{\alpha^2}
    \left(1 + \mathcal{O}(\varphi)\right)
    + \underbrace{\mathcal{O}(10^{-21})}_{\text{koneksja spinowa}}
  }
\end{equation}
Wyniki numeryczne (skrypt \texttt{dirac\_tgp\_corrected.py}):
\begin{center}
\begin{tabular}{lcccc}
\toprule
  Scenariusz & $\varphi$ & $\delta E/E$ (tetrada) &
  spin-conn/tetrada & 2.~rz\k{a}d/tetrada \\
\midrule
  Ziemia (pow.) & $1.4\times 10^{-9}$
    & $2.6\times 10^{-5}$ & $10^{-24}$ & $10^{-14}$ \\
  GN (pow.) & $8.3\times 10^{-1}$
    & $1.6\times 10^{4}$ & $10^{-21}$ & $-0.21$ \\
\bottomrule
\end{tabular}
\end{center}
\textbf{Kluczowy wniosek:}
Wiod\k{a}ca korekta tetradowa jest \emph{identyczna} z~predykcj\k{a}
OTW (przesuni\k{e}cie grawitacyjne energii spoczynkowej wzgl\k{e}dem
ma\l{}ej energii wi\k{a}zania $E_{\mathrm{bind}} \sim \alpha^2 mc^2$).
Efekty specyficzne dla TGP (koneksja spinowa) s\k{a} nieobserwowalne
($\lesssim 10^{-21}$ razy mniejsze).
Korekty drugiego rz\k{e}du ($\sim\varphi^2$) s\k{a} istotne
jedynie na powierzchni gwiazdy neutronowej ($\sim 20\%$ tetrady).
\end{proposition}

\subsubsection{Materia jako defekty: od kinka do fermion\'ow}

\begin{remark}[Czterostopniowa hierarchia sektor\'u fermionowego TGP]\label{rem:fermion-4-steps}
Sekcja ta wprowadza sektor fermionowy w~czterech \textbf{oddzielnych} krokach,
ka\.zdy z~w\l{}asnym statusem logicznym. Czytelnik mo\.ze zaakceptowa\'c
kroki~D.1a i~D.1b bez~D.1c--D.1d. Kroki~D.1c--D.1d to otwarte problemy badawcze.
\end{remark}

\begin{center}\footnotesize
\begin{tabular}{@{}cp{4.0cm}lp{4.2cm}@{}}
\toprule
\textbf{Krok} & \textbf{Tre\'s\'c} & \textbf{Status}
  & \textbf{Uzasadnienie / odsylacz} \\
\midrule

D.1a
  & Kink topologiczny $\Phi(r)$ jako obiekt nierozk\l{}adalny;
    winding number $n \in \mathbb{Z}$
  & \statuslabel{Twierdzenie}
  & Wynika z~$\pi_1(\mathbb{Z}_2\text{-instanton}) = \mathbb{Z}$;
    stabilno\'s\'c energetyczna (stw.~\ref{prop:vacuum-stability}) \\[5pt]

D.1b
  & Efektywny spin~$\tfrac{1}{2}$: ZS1 $\Rightarrow$ $Q_{\rm eff} = \tfrac{1}{2}$,
    obrót $2\pi$ nadaje faz\k{e} $(-1)$
  & \statuslabel{Propozycja}
  & Topologiczny argument $\mathbb{RP}^2$;
    tw.~\ref{thm:zs1-spin}; wymaga klasyfikacji repr. \\[5pt]

D.1c
  & Pe\l{}ne statystyki Fermiego: wymiana dw\'och kink\'ow
    $\Psi(2,1) = -\Psi(1,2)$ (zasada Pauliego)
  & \statuslabel{Szkic}
  & Wymaga wielocz\k{a}stkowej przestrzeni Fock dla kink\'ow;
    stw.~\ref{prop:fermi-statistics}: holonomia Berry \\[5pt]

D.1d
  & Chiralno\'s\'c: $U(\varphi) \neq U(2-\varphi)$ $\Rightarrow$
    asymetria kink/antykink $\Rightarrow$ $m_L \neq m_R$;
    po\l{}\k{a}czenie z~sektorem Dirac/Weyl
  & \statuslabel{Hipoteza}
  & Mechanizm wskazany; brak sformalizowanego
    odpowiednika Higgsa; hip.~\ref{hyp:spin-chirality} \\

\bottomrule
\end{tabular}
\end{center}

\begin{remark}[Hierarchia formalizacji materii]\label{rem:materia-hierarchia}
TGP oferuje dwa komplementarne opisy materii:
\begin{enumerate}[nosep]
  \item \textbf{Opis efektywny} (g\'orny): pole $\Psi$
    na tle metryki efektywnej $g_{\mu\nu}^{\mathrm{eff}}(\Phi)$
    --- hipoteza~\ref{hyp:dirac-TGP}. Odpowiada standardowemu QFT
    na zakrzywionej przestrzeni z~dodatkowym polem skalarnym.
  \item \textbf{Opis fundamentalny} (dolny): cz\k{a}stka jest
    radialnym kinkiem pola~$\Phi$ (def.~\ref{def:kink},
    \S\ref{ssec:defekty}). Spin-$\frac{1}{2}$ wynika z~topologii
    defektu w~sektorze chiralnym $\Delta(x)$.
\end{enumerate}
Pe\l{}ne po\l{}\k{a}czenie tych dw\'och opis\'ow wymaga
pokazania, \.ze kwantowy kink pola $\Phi$ reprodukuje
propagator Diraca w~granicy d\l{}ugofalowej. Jest to
\textbf{otwarty problem} analogiczny do derywacji
fermion\'ow z~modelu Skyrme'a.
\end{remark}

\begin{hypothesis}[Spin z~asymetrii chiralnej substratu]%
\label{hyp:spin-chirality}
Pole asymetrii chiralnej $\Delta(x) = \langle\hat{s}\rangle$ (ze znakiem)
posiada dodatkow\k{a} struktur\k{e} --- w~otoczeniu defektu pola $\Phi$
pole $\Delta$ tworzy konfiguracje z~nietrywialn\k{a} topologi\k{a}:
\begin{equation}\label{eq:Delta-hedgehog}
  \Delta_a(\vec{x}) = \Delta_0(r)\,\hat{x}_a,
  \qquad a \in \{1,2,3\},
\end{equation}
gdzie $\hat{x}_a = x_a/r$ to jednostkowy wektor radialny
(\emph{konfiguracja je\.za}). Takie odwzorowanie
$S^2 \to S^2$ nosi \l{}adunek topologiczny (stopie\'n
odwzorowania), kt\'ory identyfikujemy ze spinem.

Poniewa\.z $\pi_2(S^2) = \mathbb{Z}$, mo\.zliwe s\k{a}
konfiguracje o~dowolnym \l{}adunku ca\l{}kowitym. Jednak
dla profilu $\Delta_0(r)$ spe\l{}niaj\k{a}cego warunek
zerowej sumy (aksjomat~\ref{ax:zero}, ZS1:
$\int\Delta\sqrt{h}\,d^3x = 0$), minimalna stabilna
konfiguracja odpowiada \l{}adunkowi $\pm 1/2$
--- spinowi $1/2$ fermion\'ow.
\end{hypothesis}

\begin{theorem}[ZS1 $\Rightarrow$ spin-$1/2$]\label{thm:zs1-spin}
Niech $(\Phi, \Delta)$ b\k{e}dzie konfiguracj\k{a} defektow\k{a} TGP,
gdzie:
\begin{enumerate}[nosep]
  \item[(i)] $\Phi(r)$ jest radialnym kinkiem (def.~\ref{def:kink}),
  \item[(ii)] $\Delta_a(\vec{x}) = \Delta_0(r)\,\hat{x}_a$
    jest je\.zem z~odwzorowaniem $\hat{n}\colon S^2 \to S^2/\Zp = \mathbb{RP}^2$
    (symetria $\Zp$ substratu: $\hat{s} \to -\hat{s}$),
  \item[(iii)] ZS1 zachodzi:
    $\int\Delta_0(r)\,r^2\,(\Phi/\PhiZero)^{3/2}\,dr = 0$.
\end{enumerate}
W\'owczas:
\begin{enumerate}[nosep]
  \item[(A)] Efektywny \l{}adunek topologiczny $Q_{\mathrm{eff}}$
    przyjmuje warto\'sci w~$\tfrac{1}{2}\mathbb{Z}$
    (p\'o\l{}ca\l{}kowite), z~energi\k{a} k\k{a}tow\k{a}
    $E_{\mathrm{ang}} = Q_{\mathrm{eff}}(Q_{\mathrm{eff}}+1)$.
  \item[(B)] ZS1 wymusza \emph{kompensowany} profil
    $\Delta_0(r)$ (musi zmienia\'c znak),
    dodaj\k{a}c koszt energetyczny $E_{\mathrm{rad}} > 0$.
  \item[(C)] Energia ca\l{}kowita $E = E_{\mathrm{rad}} + E_{\mathrm{ang}}$
    jest minimalizowana przy $Q_{\mathrm{eff}} = 1/2$:
    \begin{equation}\label{eq:E-ang-half}
      \boxed{E_{\mathrm{ang}}(1/2) = \tfrac{3}{4}
        < E_{\mathrm{ang}}(1) = 2.}
    \end{equation}
    $Q_{\mathrm{eff}} = 1/2$ jest \textbf{niepodzielny}
    w~$\mathbb{RP}^2$ (stabilno\'s\'c topologiczna).
  \item[(D)] Przy obrocie uk\l{}adu o~$2\pi$, konfiguracja
    z~$Q_{\mathrm{eff}} = 1/2$ nabywa faz\k{e} $(-1)$
    --- spe\l{}nia prawo transformacji \textbf{spinora}.
\end{enumerate}
\textbf{Wniosek:} minimalny stabilny defekt topologiczny
w~TGP ma spin~$1/2$ (fermion).
\end{theorem}

\begin{proof}
\textbf{(A)} Przestrze\'n docelowa je\.za to
$\mathbb{RP}^2 = S^2/\Zp$ (z~symetrii $\Zp$ substratu),
z~grup\k{a} homotopii $\pi_2(\mathbb{RP}^2) = \mathbb{Z}$.
Odwzorowanie $\hat{n}\colon S^2 \to \mathbb{RP}^2$
o~stopniu $n \in \mathbb{Z}$ odpowiada efektywnemu
\l{}adunkowi $Q_{\mathrm{eff}} = n/2$ na $S^2$
(poniewa\.z $\mathbb{RP}^2$ jest dwukrotnym ilorazem $S^2$).

\textbf{(B)} ZS1 $\int\Delta_0\,r^2\,h^{3/2}\,dr = 0$
z~warunkiem brzegowym $\Delta_0(\infty) = 0$
wymusza co najmniej jedno przej\'scie przez zero ---
profil musi by\'c kompensowany.

\textbf{(C)} Energie k\k{a}towe: $E_{\mathrm{ang}}(1/2) = 3/4$,
$E_{\mathrm{ang}}(1) = 2$, $E_{\mathrm{ang}}(3/2) = 15/4$.
Dla $Q_{\mathrm{eff}} \geq 1$: rozszczepienie na dwa defekty
$Q_{\mathrm{eff}}/2$ zmniejsza energi\k{e} k\k{a}tow\k{a}
o~$\Delta E = Q(Q+1) - 2(Q/2)(Q/2+1) = Q^2/2 > 0$,
wi\k{e}c $Q_{\mathrm{eff}} \geq 1$ jest niestabilny.
$Q_{\mathrm{eff}} = 1/2$ nie mo\.ze si\k{e} rozszczepi\'c
(minimalny \l{}adunek w~$\mathbb{RP}^2$).

\textbf{(D)} Obr\'ot o~$2\pi$ na $S^2$ dzia\l{}a na
odwzorowanie $\hat{n}\colon S^2 \to \mathbb{RP}^2$
jako $\hat{n} \to -\hat{n}$ (z~identyfikacji $\Zp$).
Dla $Q_{\mathrm{eff}} = 1/2$: mno\.znik fazowy wynosi
$e^{i\pi \cdot 2Q_{\mathrm{eff}}} = e^{i\pi} = -1$ ---
transformacja spinorowa (weryfikacja numeryczna:
\texttt{zs1\_spin\_half\_proof.py}).
\end{proof}

\begin{remark}[Status sektora fermionowego]\label{rem:fermion-status}
Powy\.zszy zarys formalizacji fermion\'ow ma status
\textbf{zamkni\k{e}ty}:
\begin{enumerate}[nosep]
  \item[\checkmark] jawne obliczenie koneksji spinowej do drugiego
    rz\k{e}du w~$\varphi$ (stw.~\ref{prop:spin-connection})
    --- sp\'ojne z~PPN do rz\k{e}du 2PN
    (skrypt \texttt{spin\_connection\_2nd\_order.py}),
  \item[\checkmark] poprawiona tetrada (def.~\ref{def:tetrada})
    z~wyk\l{}adnikami $1/2$ (nie $1/4$)
    --- metryka $g_{00} = -c_0^2 \PhiZero/\Phi$,
    $g_{ij} = (\Phi/\PhiZero)\delta_{ij}$,
  \item[\checkmark] perturbacyjne rozwi\k{a}zanie r\'ownania Diraca
    z~TGP-tetrad\k{a} w~polu sferycznym
    (skrypt \texttt{dirac\_tgp\_corrected.py}):
    dominuje przesuni\k{e}cie tetradowe
    $\delta E_{\mathrm{bind}}/E_{\mathrm{bind}}
    \approx \varphi \cdot n^2/\alpha^2$,
    identyczne z~GR (zasada r\'ownowa\.zno\'sci);
    koneksja spinowa daje korekt\k{e}
    $\sim 10^{-21}$ razy mniejsz\k{a} od tetrady,
  \item[\checkmark] formalny dow\'od ZS1 $\Rightarrow$ spin~$1/2$
    (tw.~\ref{thm:zs1-spin}): $\mathbb{RP}^2 = S^2/\Zp$
    daje $Q_{\mathrm{eff}} = 1/2$ jako minimalny stabilny defekt
    (skrypt \texttt{zs1\_spin\_half\_proof.py}),
  \item[\checkmark] pe\l{}ne sprz\k{e}\.zenie
    $\Psi$--$\Phi$--$\Delta$: iteracja samospojna zbiega si\k{e}
    wyk\l{}adniczo (skrypt \texttt{psi\_phi\_delta\_coupling.py}).
    Fizyczna reakcja zwrotna $|\Psi|^2 \to \Phi$ wynosi
    $q_{\mathrm{phys}} = Gm_e/(c^2\lambda_C) \approx 1{,}75\times 10^{-45}$
    --- \textbf{zero numeryczne} w~float64.
    Zamkni\k{e}cie ontologiczne TGP:
    $\Psi$ generuje $\Phi$ przez $\rho = |\Psi|^2$,
    $\Phi$ kszta\l{}tuje $\Psi$ przez $g_{\mathrm{eff}}(\Phi)$
    + koneksj\k{e} spinow\k{a},
    $\Delta_0$ jest zdeterminowane przez ZS1 na tle $\Phi$.
    System jest \textbf{formalnie zamkni\k{e}ty}, cho\'c
    efekt fizyczny le\.zy 35+ rz\k{e}d\'ow poni\.zej wszystkich
    znanych korekt (QED, hyperfine, rozmiar j\k{a}dra).
\end{enumerate}
\end{remark}

% =============================================================
\subsubsection{D.1d: Chiralno\'s\'c i~asymetria kink/antykink}
\label{sssec:chirality-D1d}
% =============================================================

\begin{hypothesis}[Asymetria potencja\l{}u a~chiralno\'s\'c fermion\'ow]%
\label{hyp:chirality-kink}
\statuslabel{Hipoteza}: mechanizm wskazany, brak sformalizowanego
odpowiednika Higgsa.

\medskip
\textbf{Obserwacja}: Potencja\l{} TGP
$U(\psi) = \tfrac{\beta}{3}\psi^3 - \tfrac{\gamma}{4}\psi^4$ (przy $\beta=\gamma$)
nie jest symetryczny wzgl\k{e}dem inwersji o~pr\'o\.zni\k{e} $\psi = 1$:
\begin{equation}\label{eq:U-asymmetry}
  U(2 - \psi) \neq U(\psi).
\end{equation}
Przyk\l{}adowo, przy $\beta = \gamma$ i~$\psi = 1/2$ (centrum kinka):
\begin{equation}
  U\!\left(\tfrac{1}{2}\right)
  = \frac{\gamma}{24} - \frac{\gamma}{64}
  = \frac{5\gamma}{192}
  \;\neq\;
  U\!\left(\tfrac{3}{2}\right)
  = \frac{\gamma\cdot 27}{24} - \frac{\gamma\cdot 81}{64}
  = -\frac{9\gamma}{64}.
\end{equation}
Og\'olnie: asymetria $\Delta U(\psi) \equiv U(2-\psi) - U(\psi)$ znika
jedynie w~punkcie pr\'o\.zni $\psi = 1$ oraz w~$\psi = 0$
(centrum kinka).
Dla wszystkich pozosta\l{}ych $\psi \in (0,1)$: $\Delta U(\psi) \neq 0$.

\medskip
\textbf{Interpretacja fizyczna}:
\begin{itemize}[nosep]
  \item \textbf{Kink}: profil radialny $\psi(r)$: $0 \to 1$ (od centrum kinka do pr\'o\.zni).
    Cz\k{a}stka ,,'schodzi'' z~centrum ($\psi \approx 0 \approx \Nzero$)
    do pr\'o\.zni ($\psi = 1$).
  \item \textbf{Antykink}: profil $\psi(r)$: $2-0 \to 2-1 = 1$ (inwersja).
    Gdyby symetria $\psi \to 2-\psi$ by\l{}a spe\l{}niona,
    kink i~antykink by\l{}yby energetycznie r\'ownowa\.zne.
  \item \textbf{Asymetria}: $U(2-\psi) \neq U(\psi)$
    $\Rightarrow$ profile kink i~antykink maj\k{a} r\'o\.zne energie i~kszta\l{}ty.
    Jest to \textbf{wbudowana} asymetria kink/antykink w~TGP
    --- nie wymaga dodatkowego pola.
\end{itemize}

\medskip
\textbf{Zwi\k{a}zek z~chiralno\'sci\k{a}}:
W~opisie efektywnym (Dirac na tle $g_{\mu\nu}^{\mathrm{eff}}(\Phi)$,
hip.~\ref{hyp:dirac-TGP}), lewostronne $\Psi_L$ i~prawostronne
$\Psi_R$ komponenty spinora sprz\k{e}gaj\k{a} si\k{e}
z~$\Phi$ przez t\k{e}trad\k{e}:
\begin{align}
  \mathcal{L}_{\mathrm{Dirac}} &=
    \bar\Psi\,\bigl(i\gamma^\mu\nabla_\mu - m_{\rm eff}(\Phi)\bigr)\,\Psi,
    \label{eq:Dirac-TGP-chiral} \\
  m_{\rm eff}(\Phi) &= m_0\,\bigl(\Phi/\PhiZero\bigr)^n,
\end{align}
gdzie wyk\l{}adnik $n$ zale\.zy od profilu sprz\k{e}\.zenia.
Asymetria profilu kink/antykink przek\l{}ada si\k{e} na
\textbf{r\'o\.zne efektywne masy}:
\begin{equation}\label{eq:mL-neq-mR}
  m_L \equiv \langle\Psi_L|\,m_{\rm eff}(\Phi)\,|\Psi_L\rangle
  \;\neq\;
  m_R \equiv \langle\Psi_R|\,m_{\rm eff}(\Phi)\,|\Psi_R\rangle,
\end{equation}
gdy profil $\Phi(r)$ kinka ($\psi: 0\to 1$) r\'o\.zni si\k{e} od profilu
antykinka ($\psi: 2\to 1$) --- co jest zagwarantowane przez
$U(2-\psi) \neq U(\psi)$.
\end{hypothesis}

\begin{remark}[Por\'ownanie z~mechanizmem Higgsa]\label{rem:TGP-vs-Higgs}
W~Modelu Standardowym chiralno\'s\'c jest z\l{}amana przez
pole Higgsa: $m_f = y_f v / \sqrt{2}$ gdzie $v$~to VEV
i~$y_f$~to sprz\k{e}\.zenie Yukawa. W~TGP:
\begin{center}\small
\begin{tabular}{@{}p{3.5cm}p{5.5cm}p{5cm}@{}}
\toprule
\textbf{Element} & \textbf{Model Standardowy} & \textbf{TGP} \\
\midrule
\'Zr\'od\l{}o z\l{}amania chiralno\'sci
  & Pole Higgsa $H$ (nowe pole)
  & Asymetria $U(\psi) \neq U(2-\psi)$ (z~d.\ TGP) \\[3pt]
Masa fermion\'ow
  & $m_f = y_f v/\sqrt{2}$ (parametryczne)
  & $m_{L/R}$ ze \'srednich po profilu kinka (prognostycznie) \\[3pt]
Chiralny prze\l{}om
  & $\langle H \rangle = v \neq 0$ (z. spontaniczne)
  & $U(2-\psi) \neq U(\psi)$ (z~substratu) \\[3pt]
Status
  & Eksperymentalnie potwierdzony
  & \statuslabel{Hipoteza}: masy niewyliczone \\[3pt]
\bottomrule
\end{tabular}
\end{center}
TGP \emph{nie} zast\k{e}puje Higgsa; wskazuje, \.ze chiralne z\l{}amanie
mo\.ze mie\'c g\l{}\'ebsze \'zr\'od\l{}o substratowe.
\end{remark}

\begin{remark}[Otwarte pytania D.1d]\label{rem:D1d-open}
\begin{enumerate}[nosep]
  \item \textbf{Obliczenie $m_L, m_R$}: wymaga numerycznego rozwi\k{a}zania
    profilu kinka $\psi(r)$ i~obliczenia
    $\langle\psi^n_{\rm kink}\rangle$ i~$\langle\psi^n_{\rm antykink}\rangle$
    (skrypt docelowy: \texttt{p75\_chiral\_mass\_split.py}).
  \item \textbf{Zachowanie liczby leptonowej}: czy asymetria kink/antykink
    t\l{}umaczy baryon\k{e} asymetri\k{e} Wszech\'swiata?
    ($N_B - N_{\bar{B}} \neq 0$ z~$U(2-\psi) \neq U(\psi)$).
  \item \textbf{Masy neutrin}: czy $m_\nu = 0$ wynika z~dodatkowej
    symetrii substratu przy \l{}adunku topologicznym $n_Q = 0$?
  \item \textbf{Renormalizacja $m_{L/R}$}: jak ewoluuj\k{a} ze skal\k{a}
    w~ramach przep\l{}ywu Wilsonowskiego TGP?
\end{enumerate}
Kroki D.1d s\k{a} \statuslabel{Hipoteza} $\to$ docelowo \statuslabel{Propozycja}
po domkni\k{e}ciu OP-D1d (obliczenie $m_{L/R}$).
\end{remark}

% =============================================================
\subsubsection{D.3: Hierarchia mas --- rekonstrukcja vs.\ predykcja}%
\label{sssec:mass-hierarchy-D3}
% =============================================================

\begin{remark}[Co jest predykcj\k{a}, co dopasowaniem --- sektor mas]\label{rem:mass-hierarchy}
\statuslabel{Uwaga}. Poni\.zsza tabela precyzuje status ka\.zdego wyniki
sektora masowego TGP:
\begin{center}\small
\begin{tabular}{@{}p{4.5cm}p{2.5cm}p{3.5cm}p{3.5cm}@{}}
\toprule
\textbf{Wynik} & \textbf{Typ} & \textbf{Co wyznaczono} & \textbf{Co dopasowano} \\
\midrule
Formu\l{}a Koidego $Q = 2/3$
  & Warunek bifurkacji
  & --- (def.\ $Q_{\mathrm{TGP}}$)
  & definicja \\[4pt]
Trzy generacje ($n=0,1,2$)
  & Wyprowadzone z~WKB
  & $\alpha_K$, $a_\Gamma$
  & --- \\[4pt]
$r_{21} = 206{,}77$ ($m_\mu/m_e$)
  & \textbf{Dopasowanie wej\'sciowe}
  & $\alpha_K \approx 8{,}56$
  & $r_{21,\mathrm{PDG}}$ \\[4pt]
$a_\Gamma = 0{,}040049$
  & \textbf{Zdefiniowane} jako $a_c(\alpha_K)$
  & --- (identyczno\'s\'c)
  & $\alpha_K$ \\[4pt]
$m_e, m_\mu, m_\tau$
  & Odtwarzane
  & ---
  & $r_{21}$ \\[4pt]
Masy kwork\'ow
  & Analogia (Szkic)
  & $\alpha_q \neq \alpha_K$
  & dopasowanie $\alpha_q$ \\[4pt]
\bottomrule
\end{tabular}
\end{center}
\medskip
\textbf{Wniosek}: W~obecnej wersji $\alpha_K$ jest \textbf{wej\'sciem} dopasowanym
do $r_{21,\mathrm{PDG}}$, nie predykcj\k{a}. Sektor mas jest rekonstrukcj\k{a},
nie predykcj\k{a} --- a\.z do domkni\k{e}cia~OP-3
(prob.~\ref{prob:OP-3} poni\.zej).
\end{remark}

\begin{remark}[Status kanoniczny $\alpha_K$ po sesjach v45--v47b]%
\label{rem:alphaK-legacy-status}
W~sektorze leptonowym wsp\'o\l{}istniej\k{a} obecnie dwa j\k{e}zyki opisu:
\begin{enumerate}[nosep]
  \item \textbf{J\k{e}zyk legacy}:
    $\alpha_K^{\rm old}\approx 8{,}56$ oraz relacje
    $r_{21} \leftrightarrow \alpha_K \leftrightarrow a_\Gamma$,
    u\.zywane w~starszych skanach bifurkacyjnych ODE.
  \item \textbf{J\k{e}zyk kanoniczny}:
    $g_0^e$, $g_0^\mu$, $g_0^\tau$ oraz punkt sta\l{}y
    $\varphi$-FP $g_0^*$ (Formulacja~A, $K(g)=g^4$),
    w~kt\'orym relacja Koidego i~stosunek $r_{21}$ s\k{a} kodowane
    bez bezpo\'sredniego odwo\l{}ania do $\alpha_K$.
\end{enumerate}
W~tej wersji manuskryptu \textbf{kanoniczny} jest drugi j\k{e}zyk.
Parametr $\alpha_K^{\rm old}$ pozostaje u\.zytecznym obiektem
translacyjnym mi\k{e}dzy formulacjami, ale jego status epistemiczny jest
obecnie s\l{}abszy ni\.z status punktu sta\l{}ego $g_0^*$ i~$\varphi$-FP.
Dlatego OP-3 rozbijamy logicznie na dwa podproblemy:
\begin{enumerate}[nosep,label=(\alph*)]
  \item \textbf{OP-3A}: czy $\alpha_K^{\rm old}$ wynika z~Warstwy~I--II
    bez u\.zycia $r_{21}$ jako wej\'scia?
  \item \textbf{OP-3B}: czy istnieje jawna i~jednoznaczna mapa
    $\alpha_K^{\rm old}\leftrightarrow(g_0^*,g_0^e,\varphi\text{-FP})$,
    tak aby stara formulacja sta\l{}a si\k{e} pochodn\k{a} nowej?
\end{enumerate}
Zamkni\k{e}cie OP-3 mo\.ze zatem nast\k{a}pi\'c dwiema drogami:
przez pe\l{}ne wyprowadzenie $\alpha_K^{\rm old}$ albo przez formaln\k{a}
supersesj\k{e} tego parametru przez kanoniczny j\k{e}zyk $g_0^*$.
\end{remark}

\begin{problem}[OP-3: Wyznaczenie $\alpha_K$ z~substratu --- ZAMKNI\k{E}TY]%
\label{prob:OP-3}
\statuslabel{CZ. ZAMK.}
(zamkni\k{e}ty przez tw.~\ref{thm:Z-alpha-degeneracy}: $\alpha_K$ jest artefaktem
koordynatowym; kanoniczny j\k{e}zyk u\.zywa $g_0^*$ z~$\varphi$-FP).
Parametr $\alpha_K \approx 8{,}56$ jest obecnie dopasowywany do
$r_{21,\mathrm{PDG}} = 206{,}77$.
Celem jest pokazanie, \.ze $\alpha_K$ \textbf{wynika} z~danych Warstwy~I--II
bez odwo\l{}ania do $r_{21}$.

\medskip
\textbf{\'Scie\.zka~1 (topologiczna)}: Warunek minimalno\'sci energii solitonu:
\begin{equation}\label{eq:OP3-path1}
  \alpha_K \;=\; \arg\min_{\alpha}\, E_{\rm sol}(\alpha)
  \quad\text{przy}\quad
  Q_{\mathrm{TGP}}\!\bigl(K_1(\alpha),K_2(\alpha),K_3(\alpha)\bigr) = \tfrac{2}{3},
\end{equation}
gdzie $K_n(\alpha)$ to energie kinku $n$-tej generacji ($n=0,1,2$) obliczone
z~profilu WKB przy danym~$\alpha$, a~$Q_{\mathrm{TGP}} = 2/3$
to warunek Koidego jako selektor bifurkacyjny.
Je\'sli r\'ownanie~\eqref{eq:OP3-path1} ma jednoznaczne rozwi\k{a}zanie
$\alpha_K^*$, to $\alpha_K$ staje si\k{e} Warstw\k{a}~III
i~$r_{21}$ jest \textbf{predykcj\k{a}}.

\medskip
\textbf{\'Scie\.zka~2 (relacyjna)}: Warunek bifurkacji z~$a_\Gamma$:
\begin{equation}\label{eq:OP3-path2}
  a_c(\alpha_K) \;=\; a_\Gamma
  \quad \Leftrightarrow \quad
  \alpha_K = a_c^{-1}(a_\Gamma),
\end{equation}
gdzie $a_c(\alpha)$ to krytyczna warto\'s\'c parametru substratu,
przy kt\'orej bifurkuje $n$-te rozwi\k{a}zanie WKB.
Je\'sli $a_c^{-1}$ jest obliczalna numerycznie, to $\alpha_K$
wynika z~$a_\Gamma$ (Warstwa~II).
W~po\l{}\k{a}czeniu z~\textbf{za\l{}o\.zeniem~GL} (\#18,
stw.~\ref{prop:psi-ini-derived}): $N_{\rm param} = 2 \to 2$ (masy leptonowe
przechodz\k{a} do Warstwy~III, $M_{\rm obs}$ ro\'snie).

\medskip
\textbf{\'Scie\.zka~3 (stabilno\'s\'c --- nova via, p80):} [\textbf{ZAMKNI\k{E}TA~--- artefakt, p95}]
Warunek bifurkacji siod\l{}o-w\k{e}ze\l{} ODE solitonu:
\begin{equation}\label{eq:OP3-path3}
  \alpha_K \;=\; \alpha_{\rm max} \cdot (1 - \varepsilon),
  \qquad
  \varepsilon \;=\; 3{,}51\%,
\end{equation}
gdzie $\alpha_{\rm max} = 8{,}8734$ to granica wyznaczona przy konwencji $r_{\rm max}=40$
(uw.~\ref{rem:alpha-max-bifurcation}).
Numerycznie: $\alpha_K/\alpha_{\rm max} = 0{,}9649 = n_s^{\rm Planck}$
(r\'ownanie~\eqref{eq:ns-coincidence}).
\emph{Skrypt p95 wykaza\l{}, \.ze $K_{\rm nan,R}$ (definiuj\k{a}cy $\alpha_{\rm max}$)
nie istnieje jako fizyczna granica przy $r_{\rm max}\geq 60$~--- jest artefaktem
fazy oscylacyjnego ogona w~punkcie $r=40$.
Koincydencja $\alpha_K/\alpha_{\rm max} = n_s$ jest zatem COFNI\k{E}TA.}
Status OP-3~\'Scie\.zka~3: \textbf{ZAMKNI\k{E}TA (artefakt p94--p95)}.
Pozostaje \'Scie\.zka~1: $\alpha_K \approx \arg\min E^*(\alpha)$ (p82--p83)
[\textbf{COFNI\k{E}TA p96, patrz ni\.zej}].

\medskip
\textbf{Status~\'Scie\.zki~1 (p82--p84, p96):}
Precyzyjny skan $\alpha\in[8{,}50;\,8{,}71]$ (p82, krok~$0{,}01$)
oraz \'sledzenie ga\l{}\k{e}zi (p83) wykaza\l{}y pozorne minimum
$E^*(\alpha,\, r_{\rm max}=40)$ przy $\alpha\approx 8{,}65$.
Skrypt~\texttt{p96\_Estar\_rmax\_convergence.py} (skan $\alpha\in[8{,}20;\,8{,}70]$
z~$r_{\rm max}\in\{40,60,80\}$, 2/5~PASS) zbada\l{} zbiezno\'s\'c tego minimum:
\begin{center}\small
\begin{tabular}{@{}lccc@{}}
\toprule
$r_{\rm max}$ & $\arg\min E^*$ & $\delta\!\equiv\!\arg\min-\alpha_K$ \\
\midrule
$40$ & $8{,}646$ & $+0{,}085$ \\
$60$ & $\geq 8{,}70$ & $\geq +0{,}138$ \\
$80$ & $\geq 8{,}70$ & $\geq +0{,}138$ \\
$\infty$ (extrap) & $\geq 8{,}70$ & $\geq +0{,}138$ \\
\bottomrule
\end{tabular}
\end{center}
Przy $r_{\rm max}\geq 60$ ekstrapolowane $K^*(\alpha,\infty)$ jest
\textbf{monotonicznie malej\k{a}ce} w~ca\l{}ym zakresie $[8{,}20;\,8{,}70]$
--- brak minimum fizycznego.
Pozorne minimum przy $r_{\rm max}=40$ ($\alpha\approx 8{,}65$) to kolejny
artefakt fazy oscylacyjnego ogona (analogiczny do $K_{\rm nan,R}$ i~Ga\l{}\k{a}zi~L).

\textbf{\'Scie\.zka~1: ZAMKNI\k{E}TA (p96)}:
$K^*(\alpha,r_{\rm max}\to\infty)$ jest monotonicznie malej\k{a}ce.
Brak minimum energii solitonu TGP wskazuj\k{a}cego na~$\alpha_K$.

\medskip
\textbf{Bilans \'Scie\.zek (stan 2026-03-26):}
\begin{itemize}[nosep]
  \item \'Scie\.zka~1 (argmin $E^*$): \textbf{ZAMKNI\k{E}TA} --- $K^*(\alpha,\infty)$
    monotonicznie malej\k{a}ce, brak minimum energii przy~$\alpha_K$ (p96).
    \emph{Podwariant Koide (p98--p99)}: $K^*_2$ jest fizyczny ($r_{\rm max}=80$),
    ale $K^*_2/K^*_1(\infty)\approx21$--$25 \neq 206{,}77$~---
    standardowy $V_{\rm mod}$ nie spe\l{}nia warunku Koidego;
    wymaga modyfikacji (p15).
  \item \'Scie\.zka~2 (relacyjna): \textbf{ZAMKNI\k{E}TA (nie potwierdzona, p101)}.
    Skrypt~\texttt{p101\_K2\_agamma\_corrected.py}
    (skan $a_\Gamma\in[0{,}010;\,0{,}060]$, $r_{\rm max}=80$,
    w\k{a}skie okna $\psi$:
    $K^*_1$ z~$\psi\in[1{,}001;\,1{,}8]$, $K^*_2$ z~$\psi\in[3{,}5;\,7{,}0]$,
    5/5~PASS) da\l{} jednoznaczny wynik:

    \medskip
    \emph{Odkrycia p101}:
    \begin{itemize}[nosep]
      \item \textbf{Universalno\'s\'c $\psi_0$}:
        $\psi_0(K^*_1) = 1{,}2419 = \mathrm{const}$%
        \footnote{Warto\'s\'c kanoniczna: $g_0^e = 0{,}86941$ (Formulacja~A, $K = g^4$). Warto\'s\'c $1{,}24$ odpowiada Formulacji~B ($f = 1+4\ln g$).}
        dla wszystkich $a_\Gamma\in[0{,}010;\,0{,}060]$ ---
        warto\'s\'c pocz\k{a}tkowa solitonu nie zale\.zy od $a_\Gamma$!
      \item $K^*_1 \propto a_\Gamma$ (skalowanie Schwarzschilda,
        odch.\ $\leq 3{,}3\%$ bias $r_{\rm max}=80$);
      \item $K^*_2$ (ga\l{}\k{a}\'z $\psi_0\approx4{,}1$) \textbf{nie istnieje}
        dla $a_\Gamma \leq 0{,}015$~--- krytyczna warto\'s\'c
        $a_c \in (0{,}015;\,0{,}020]$;
      \item $K^*_2/K^*_1 \approx 16{,}7$ (sta\l{}e dla $a_\Gamma\in[0{,}020;\,0{,}060]$,
        oba $\propto a_\Gamma$);
      \item $a_c \approx 0{,}017 \neq 0{,}040 = a_\Gamma^{\rm phys}$~---
        \textbf{warunek $a_c(\alpha_K) = a_\Gamma$ nie spe\l{}niony}.
    \end{itemize}
    \textbf{Wniosek}: OP-3 \'Scie\.zka~2 nie jest potwierdzona
    w~zakresie dostepno\'sci numerycznej.
    Uwaga: $K^*_2$ z~p101 ($\psi_0\approx4{,}1$) i~z~p99 ($\psi_0\approx5{,}0$)
    to \emph{r\'o\.zne} ga\l{}\k{e}zie solitonu;
    bogaty spektrum wzbudze\'n $K^*_n$ wymaga dalszej identyfikacji.
  \item \'Scie\.zka~3: \textbf{ZAMKNI\k{E}TA} --- $K_{\rm nan,R}$ artefakt
    $r_{\rm max}=40$; koincydencja $\alpha_K/\alpha_{\rm max}=n_s$ cofni\k{e}ta (p94--p95).
\end{itemize}
\textbf{Wniosek (stan 2026-03-26):} $\alpha_K = 8{,}5616$ pozostaje parametrem
Warstwy~II (obserwacyjnym --- z~warunku $K^*_2/K^*_1 = r_{21}$).
\begin{itemize}[nosep]
  \item \'Scie\.zka~1 (argmin $E^*$): \textbf{ZAMKNI\k{E}TA} (p96).
  \item \'Scie\.zka~2 (relacyjna): \textbf{ZAMKNI\k{E}TA --- nie potwierdzona} (p101):
    $a_c(\alpha_K)\approx0{,}017 \neq 0{,}040$.
  \item \'Scie\.zka~3 ($n_s$-koincydencja): \textbf{ZAMKNI\k{E}TA} (p94--p95, artefakt).
\end{itemize}
\textbf{Aktualizacja statusu (v47b).}
Bilans z~sesji v45--v47b pokazuje, \.ze trzy dawne \'scie\.zki selekcji
$\alpha_K$ (argmin~$E^*$, warunek $a_c(\alpha_K)=a_\Gamma$,
koincydencja z~$n_s$) nie daj\k{a} zamkni\k{e}cia OP-3A.
Najbardziej obiecuj\k{a}ca bezpo\'srednia droga nie polega ju\.z na
szukaniu ,,wyr\'o\.znionej liczby'' w~skanach $\alpha$, lecz na
przepisaniu problemu w~postaci implicitnej
$g_0(K)+\alpha\,g_1(K)=0$ i~pr\'obie identyfikacji $\alpha$ jako funkcji
obiekt\'ow Warstwy~I--II.

\textbf{R\'ownoleg\l{}a droga OP-3B}: formalna supersesja j\k{e}zyka legacy
przez j\k{e}zyk kanoniczny $g_0^*$ oraz $\varphi$-FP.
W~tym wariancie $\alpha_K^{\rm old}$ staje si\k{e} parametrem
\emph{recovered/legacy}, a~rzeczywist\k{a} luk\k{a} pozostaje most
substrat $\to g_0^e \to \varphi\text{-FP} \to r_{21},Q_K,m_\tau$.
\emph{Nowe odkrycie (p101)}: $\psi_0(K^*_1) = 1{,}2419 = \mathrm{const}$%
\footnote{Warto\'s\'c kanoniczna: $g_0^e = 0{,}86941$ (Formulacja~A, $K = g^4$). Warto\'s\'c $1{,}24$ odpowiada Formulacji~B ($f = 1+4\ln g$).}
dla wszystkich $a_\Gamma\in[0{,}010;\,0{,}060]$ --- universalno\'s\'c warunku pocz\k{a}tkowego solitonu.
\end{problem}

% =============================================================
\subsection{Pola cechowania i~\l{}adunek: zarys emergencji}%
\label{ssec:gauge-emergence}
% =============================================================

TGP w~obecnej postaci opisuje grawitacj\k{e} (pole~$\Phi$)
i~zarys fermion\'ow (defekty $\Phi$ + pole chiralne $\Delta$).
Brakuje formalnego mechanizmu emergencji p\'ol cechowania
(elektromagnetyzmu, oddzia\l{}ywa\'n s\l{}abych i~silnych).
Poni\.zej przedstawiamy zarys dw\'och komplementarnych
scenariuszy, kt\'ore s\k{a} sp\'ojne z~duchem TGP.

\subsubsection{Scenariusz~I: Pola cechowania z~symetrii substratu}

\begin{hypothesis}[Lokalne symetrie z~substratu]%
\label{hyp:gauge-substrate}
Hamiltonjan substratu $H_\Gamma$ (r.~\eqref{eq:H-Gamma})
posiada lokalne symetrie wewn\k{e}trzne, kt\'ore
\textbf{nie s\k{a} z\l{}amane} przez przej\'scie fazowe
generuj\k{a}ce $\Phi$. Je\.zeli substrat ma wewn\k{e}trzn\k{a}
symetri\k{e} grupy $G$ ($G \supseteq \mathrm{SU}(3)\times
\mathrm{SU}(2)\times \mathrm{U}(1)$), to po coarse-grainingu
powstaj\k{a} efektywne pola cechowania $A_\mu^a(x)$
jako d\l{}ugofalowe stopnie swobody odpowiadaj\k{a}ce
generatorom $G$.

Mechanizm jest analogiczny do emergencji fonon\'ow
i~faz Goldstone'a w~fizyce materii skondensowanej.
Kluczowa r\'o\.znica: symetria $Z_2$ substratu ($\hat{s}_i
\to -\hat{s}_i$) jest z\l{}amana spontanicznie (generuj\k{a}c
$\Phi$), ale symetrie \emph{wewn\k{e}trzne}
mog\k{a} pozosta\'c niez\l{}amane i~manifestowa\'c si\k{e}
jako lokalne symetrie cechowania.
\end{hypothesis}

\subsubsection{Scenariusz~II: \L{}adunek jako topologia defektu}

\begin{hypothesis}[\L{}adunek elektryczny z~defekt\'ow $\Phi$]%
\label{hyp:charge-topology}
Radialny kink pola $\Phi$ (def.~\ref{def:kink}) wraz
z~konfiguracj\k{a} chiraln\k{a} $\Delta$ (hip.~\ref{hyp:spin-chirality})
posiada dodatkowe stopnie swobody topologiczne.
Pole $\Delta_a(\vec{x})$ na sferze $S^2$ otaczaj\k{a}cej
defekt odwzorowuje $S^2 \to S^2$ (stopie\'n = spin)
i~jednocze\'snie mo\.ze posiada\'c \emph{faz\k{e}} $e^{i\theta}$
zwi\k{a}zan\k{a} z~orientacj\k{a} wewn\k{e}trzn\k{a} substratu:
\begin{equation}\label{eq:charge-phase}
  \Delta_a(\vec{x}) = |\Delta_0(r)|\,\hat{x}_a\,
  e^{i\,n_Q\,\theta(\vec{x})},
  \qquad n_Q \in \mathbb{Z}.
\end{equation}
Liczba ca\l{}kowita $n_Q$ odpowiada \l{}adunkowi
elektrycznemu w~jednostkach $e$. Kwantyzacja
\l{}adunku wynika z~topologii ($\pi_1(\mathrm{U}(1)) = \mathbb{Z}$),
analogicznie do monopoli Diraca.
\end{hypothesis}

\subsubsection{Szkic wyprowadzenia: od substratu do U(1)}

Przedstawiamy formalny zarys emergencji pola cechowania U(1)
z~substratu. Uog\'olnienie na grupy nieabelowe (SU($N$))
przebiega analogicznie.

\begin{proposition}[Emergencja U(1) z~fazy substratu]%
\label{prop:U1-emergence}
Za\l{}\'o\.zmy, \.ze operator substratowy $\hat{s}_i$ jest
\textbf{zespolony}: $\hat{s}_i = |\hat{s}_i|\,e^{i\theta_i}$,
a~hamiltonjan substratu $H_\Gamma$ (r.~\eqref{eq:H-Gamma})
jest niezmienniczy wzgl\k{e}dem \textbf{globalnych} obrot\'ow
fazy $\theta_i \to \theta_i + \alpha$ (symetria $\mathrm{U}(1)$).

Po przej\'sciu fazowym ($T < T_c$) pole modu\l{}u
$|\hat{s}|$ kondensuje do $v > 0$ (generuj\k{a}c $\Phi$),
ale faza $\theta(x)$ pozostaje mi\k{e}kkim stopniem swobody
(mod Goldstone'a). Efektywne dzia\l{}anie fazy
po coarse-grainingu ma posta\'c:
\begin{equation}\label{eq:S-phase}
  S_\theta = \int d^4x\,\sqrt{-g_{\mathrm{eff}}}\;
  \frac{v^2}{2}\,g_{\mathrm{eff}}^{\mu\nu}\,
  \partial_\mu\theta\,\partial_\nu\theta\,.
\end{equation}
Definiuj\k{a}c \textbf{pole cechowania}
$A_\mu \equiv -v\,\partial_\mu\theta$ i~tensor pola
$F_{\mu\nu} = \partial_\mu A_\nu - \partial_\nu A_\mu$,
otrzymujemy w~d\l{}ugofalowej granicy
(rozwijaj\k{a}c wok\'o\l{} wolno zmiennej fazy):
\begin{equation}\label{eq:L-Maxwell-emergent}
  \mathcal{L}_{\mathrm{em}} =
  -\frac{1}{4\mu_0}\,F_{\mu\nu}F^{\mu\nu},
  \qquad
  \mu_0^{-1} = v^2/\Phi_0\,,
\end{equation}
gdzie $\mu_0$ jest przenikalno\'sci\k{a} magnetyczn\k{a},
a~zale\.zno\'s\'c od $\Phi$ koduje sprz\k{e}\.zenie
elektromagnetyzmu z~grawitacj\k{a}.
\end{proposition}

\begin{proof}[Szkic]
Standardowe wyprowadzenie pola cechowania z~modu Goldstone'a.
Gradient fazy $\partial_\mu\theta$ jest jedynym niezmienniczym
cechowaniowo obiektem liniowym w~$\theta$.
Energia kinetyczna~\eqref{eq:S-phase} jest najni\.zszym
cz\l{}onem w~rozwini\k{e}ciu gradientowym.
Kolejny cz\l{}on ($\sim (\partial\partial\theta)^2$) daje
kinetyczny lagran\.zjan Maxwella po podstawieniu $A_\mu$.
R\'ownowa\.zno\'s\'c z~sieci\k{a}: na sieci substratowej
r\'o\.znica faz $\theta_i - \theta_j$ mi\k{e}dzy s\k{a}siadami
jest odpowiednikiem po\l{}\k{a}czenia (Wilson link),
a~p\k{e}tla wok\'o\l{} plakietki daje $F_{\mu\nu}$
--- dok\l{}adnie jak w~kratowej teorii cechowania Wilsona.
\end{proof}

\begin{remark}[Uog\'olnienie na SU($N$)]\label{rem:SUN-gauge}
Je\.zeli operator substratowy ma $N$ sk\l{}adowych
wewn\k{e}trznych (np.\ $\hat{s}_i \in \mathbb{C}^N$)
z~symetri\k{a} SU($N$), to po kondensacji modu\l{}u
$|\hat{s}|$ powstaje $N^2 - 1$ mod\'ow Goldstone'a,
kt\'ore po coarse-grainingu daj\k{a} pola cechowania
$A_\mu^a$ ($a = 1, \ldots, N^2-1$) z~lagran\.zjanem
Yanga--Millsa:
\begin{equation}\label{eq:L-YM-emergent}
  \mathcal{L}_{\mathrm{YM}} =
  -\frac{1}{4g_s^2}\,
  \mathrm{Tr}(F_{\mu\nu}F^{\mu\nu}),
  \qquad
  F_{\mu\nu} = \partial_\mu A_\nu - \partial_\nu A_\mu
  + [A_\mu, A_\nu].
\end{equation}
Sta\l{}a sprz\k{e}\.zenia $g_s$ jest okre\'slona
przez parametry substratu ($v$, koordynacj\k{e} $z$,
temperatur\k{e} $\beta_c$).

Dla reprodukcji Modelu Standardowego potrzebny jest
substrat z~wewn\k{e}trzn\k{a} struktur\k{a}
SU(3) $\times$ SU(2) $\times$ U(1), co nak\l{}ada
minimalne wymagania:
\begin{itemize}[nosep]
  \item $\hat{s}_i$ musi mie\'c co najmniej $3 + 2 + 1 = 6$
    wewn\k{e}trznych sk\l{}adowych (lub $3 \times 2 = 6$
    z~tensora iloczynowego),
  \item SU(2) musi by\'c \l{}amana spontanicznie
    (mechanizm Higgsa) --- w~TGP odpowiada to
    dodatkowej kondensacji substratu,
  \item SU(3) musi pozosta\'c niez\l{}amana
    (uwi\k{e}zienie kolor\'ow).
\end{itemize}
\end{remark}

% ---------------------------------------------------------------
% Sesja v32, 2026-03-24 --- zadanie II.E planu \texttt{PLAN\_ROZWOJU\_v1}
% ---------------------------------------------------------------

\begin{proposition}[SU(2) z~dw\'och ga\l{}\k{e}zi substratu --- szkic]%
\label{prop:SU2-from-substrate}
Niech substrat ma dwie wewn\k{e}trzne ga\l{}\k{e}zie
$(\hat{s}_i^\uparrow, \hat{s}_i^\downarrow)$
(np.\ dwa podsublattice lub dwie sk\l{}adowe spinowe):
\begin{equation}\label{eq:substrate-doublet}
  \boldsymbol{\varphi}_i
  \;=\;
  \begin{pmatrix} \varphi_i^\uparrow \\ \varphi_i^\downarrow \end{pmatrix},
  \qquad
  \varphi_i^{\uparrow,\downarrow} = \langle(\hat{s}_i^{\uparrow,\downarrow})^2\rangle^{1/2}.
\end{equation}
Swobodna energia GL z~symetri\k{a} globalneg SU(2):
\begin{equation}\label{eq:F-SU2}
  \mathcal{F}_{\mathrm{SU(2)}} = \int d^3x\;\Bigl[
    \frac{K}{2}\,(\nabla\boldsymbol{\varphi})^\dagger(\nabla\boldsymbol{\varphi})
    + V(\boldsymbol{\varphi}^\dagger\boldsymbol{\varphi})
  \Bigr].
\end{equation}
W\'owczas:
\begin{enumerate}[(i)]
  \item \textbf{Emergencja dubletu izospinowego.}
    Kondensacja $\boldsymbol{\varphi} = (\varphi_0, 0)^T + \delta\boldsymbol{\varphi}$
    daje trzy mody Goldstone'a odpowiadaj\k{a}ce
    generatorom SU(2) $\{T_+, T_-, T_3\}$.
    Po coarse-grainingu: pola cechowania
    $W_\mu^\pm$ i~$Z_\mu$ (bez foton\'ow --- U(1)$_Y$ oddzielnie).

  \item \textbf{Warunek symetrii SU(2) granicy continuum.}
    Transformacja $({\varphi}^\uparrow, {\varphi}^\downarrow)
    \to U \cdot ({\varphi}^\uparrow, {\varphi}^\downarrow)$
    (gdzie $U \in SU(2)$) jest symetri\k{a}
    $\mathcal{F}_{\mathrm{SU(2)}}$
    wtedy i~tylko wtedy, gdy potencja\l{}
    $V(\boldsymbol{\varphi}^\dagger\boldsymbol{\varphi})$
    zale\.zy wy\l{}\k{a}cznie od $|\boldsymbol{\varphi}|^2
    = (\varphi^\uparrow)^2 + (\varphi^\downarrow)^2$.

  \item \textbf{Spontaniczne \l{}amanie SU(2) $\to$ U(1).}
    Minimum $|\boldsymbol{\varphi}|^2 = v_0^2$ z~niezerowym
    $\langle\varphi^\uparrow\rangle = v_0$, $\langle\varphi^\downarrow\rangle = 0$
    \l{}amie SU(2) do U(1)$_{T_3}$:
    daje 2 mody Goldstone'a (absorbowane przez $W^\pm$)
    i~1 ciężki mod (masa Higgsa~$m_H^2 = V''(v_0^2)\cdot 2v_0^2$).

  \item \textbf{Sta\l{}a sprz\k{e}\.zenia.}
    Analogicznie do U(1) (prop.~\ref{prop:U1-emergence}):
    \begin{equation}\label{eq:g2-from-substrate}
      g_2^2 = \frac{K\,v_0^2\,a_{\mathrm{sub}}^{d-2}}{N_B}.
    \end{equation}
\end{enumerate}
\end{proposition}

\begin{remark}[SU(2): status i~co brakuje]%
\label{rem:SU2-status}
Prop.~\ref{prop:SU2-from-substrate} jest \textbf{szkicem formalizmu}
(Faza II.E planu~\texttt{PLAN\_ROZWOJU\_v1}).
Otwarte problemy:
\begin{enumerate}[nosep]
  \item \emph{Chiralność}: fermiony L i~R mają r\'o\.zne sprz\k{e}\.zenia
    z~$W^\pm$ --- w~TGP pochodzi to z~asymetrii kink--antykink
    (prop.~\ref{prop:kink-chirality}); powi\k{a}zanie z~$SU(2)_L$
    wymaga jawnego rachunku;
  \item \emph{\L{}amanie Higgsa}: masa $m_H$ i~masy $W, Z$ z~parametr\'ow
    substratu ($v_0, K, a_{\mathrm{sub}}$) --- bez wolnych parametr\'ow;
  \item \emph{Mieszanie Weinberga}: $\theta_W$ z~$g_2/g_1$ okre\'slone
    przez stosunek st\k{a}\l{}ych sprz\k{e}\.ze\'n
    dwu-ga\l{}\k{e}ziowego (SU(2)) i~jedno-ga\l{}\k{e}ziowego (U(1)) substratu.
\end{enumerate}
\end{remark}

\subsubsection{\L{}adunek, kwantyzacja i~zerowa suma}

\begin{proposition}[Kwantyzacja \l{}adunku z~topologii
  defektu]\label{prop:charge-quantization}
Niech defekt pola $\Phi$ (kink radialny) posiada
konfiguracj\k{e} fazow\k{a} $\theta(\vec{x})$
z~hip.~\ref{hyp:charge-topology}. Ca\l{}ka fazowa
po \l{}ukach $\gamma$ otaczaj\k{a}cych defekt spe\l{}nia:
\begin{equation}\label{eq:charge-winding}
  n_Q = \frac{1}{2\pi}\oint_\gamma d\theta
  \in \mathbb{Z}
\end{equation}
(z~$\pi_1(\mathrm{U}(1)) = \mathbb{Z}$).
\L{}adunek elektryczny defektu wynosi $Q = n_Q\,e_0$,
gdzie $e_0$ jest elementarnym \l{}adunkiem okre\'slonym
przez sta\l{}\k{a} sprz\k{e}\.zenia:
\begin{equation}\label{eq:e0-from-substrate}
  e_0^2 = 4\pi\alpha_{\mathrm{em}}\,\hbar_0\,c_0
  = \frac{v^2}{\PhiZero}\,\frac{4\pi}{\mu_0 c_0}\,.
\end{equation}
\end{proposition}

\begin{remark}[Zerowa suma \l{}adunek\'ow]\label{rem:charge-zero-sum}
Aksjomat zerowej sumy (A3, ZS1) gwarantuje
$\int\Delta\sqrt{h}\,d^3x = 0$ (chiralno\'s\'c).
Analogiczny argument dla \l{}adunku:
je\.zeli faza $\theta(x)$ jest g\l{}adka w~ca\l{}ej
przestrzeni z~wyj\k{a}tkiem punkt\'ow defekt\'ow,
to \textbf{suma \l{}adunk\'ow topologicznych}
wszystkich defekt\'ow musi zerowa\'c si\k{e}:
\begin{equation}\label{eq:charge-ZS}
  \sum_{\mathrm{defekty}} n_Q = 0,
\end{equation}
poniewa\.z ca\l{}ka $\oint d\theta$ po sferze
otaczaj\k{a}cej \emph{wszystkie} defekty jest
trywialna (brak osobliwo\'sci na niesko\'nczono\'sci).
Zerowa suma \l{}adunk\'ow wynika wi\k{e}c
z~topologii, nie jest dodatkowym za\l{}o\.zeniem.
\end{remark}

\begin{remark}[\L{}adunki frakcyjne kwarks\'ow]\label{rem:fractional-charges}
Kwantyzacja~\eqref{eq:charge-winding} daje \l{}adunki ca\l{}kowite.
\L{}adunki frakcyjne kwarks\'ow ($\pm 1/3$, $\pm 2/3$)
wymagaj\k{a} struktury SU(3)$_c$: kwarki nie istniej\k{a}
jako izolowane defekty, lecz jedynie w~konfigurach
tr\'ojkowych (bariony, $n_Q = 1/3 + 1/3 + 1/3 = 1$)
lub par ($n_Q + \bar{n}_Q = 0$). Efektywny \l{}adunek
frakcyjny jest artefaktem konwencji --- fizyczny (obserwowalny)
\l{}adunek izolowanego obiektu jest zawsze ca\l{}kowity.
Jest to sp\'ojne z~uwi\k{e}zieniem kolor\'ow.
\end{remark}

\begin{remark}[Status sektora cechowania --- tabela por\'ownawcza]\label{rem:gauge-status}
Poni\.zsza tabela podaje \textbf{etapowy status} ka\.zdej symetrii cechowania
oddzielnie --- co jest najwa\.zniejsz\k{a} informacj\k{a} dla czytelnika
znaj\k{a}cego Standardowy Model.
\end{remark}

\begin{center}\footnotesize
\begin{tabular}{@{}p{1.4cm}p{2.8cm}lp{2.4cm}p{2.6cm}@{}}
\toprule
\textbf{Symetria} & \textbf{\'Zr\'od\l{}o w~TGP} & \textbf{Status}
  & \textbf{Wolne param.} & \textbf{Otwarte pytanie} \\
\midrule

U(1)$_{\rm em}$
  & Faza $\varphi \to \varphi\,e^{i\theta}$ (subst.\ wielosk\l{}adnikowy)
  & \statuslabel{Szkic}
  & $g_1$ z~substratu; niobl.
  & Czy $g_1^2 = e^2/\epsilon_0$ dok\l{}adnie? \\[5pt]

SU(2)$_L$
  & Doublet $(\varphi^\uparrow, \varphi^\downarrow)$; spontaniczne
    \l{}amanie $\to$ U(1)
  & \statuslabel{Szkic}
  & $g_2$, $v_0$ z~substratu; niobl.
  & Masy $W^\pm, Z$; chiralność $L\neq R$? \\[5pt]

SU(3)$_c$
  & Tr\'ojkolor $(\varphi_r, \varphi_g, \varphi_b)$;
    topologia defektu
  & \statuslabel{Hipoteza}
  & $g_3$ --- brak
  & Mechanizm uwi\k{e}zienia kolor\'ow \\[5pt]

Uwi\k{e}zienie
  & Potencja\l{} $V(r) \sim r$ dla $\epsilon$-mod\'ow
  & \statuslabel{Szkic}
  & ---
  & Symulacja Monte Carlo na $\Gamma$ \\

\bottomrule
\end{tabular}
\end{center}

\begin{remark}[Co oddziela SU(2) od SU(3) w~TGP]\label{rem:SU2-vs-SU3}
SU(2) ma szkic formalny (prop.~\ref{prop:SU2-from-substrate}) z~wyp\l{}ywa\-j\k{a}c\k{a}
ask\k{a}$g_2$; dla SU(3) nie ma nawet formalnego szkicu --- status to
\statuslabel{Hipoteza} (mechanizm wskazany, brak wyprowadzenia).
\textbf{Kluczowe otwarte pytanie dla SU(2)}:
czy emergentna algebra SU(2) daje \textbf{pe\l{}ne s\l{}abe oddzia\l{}ywanie}
(masywne bozony $W^\pm$, $Z$) przez \l{}amanie Higgsa z~dwu-sk\l{}adnikowego substratu,
czy tylko topologiczne cechowanie bez masy?
Odpowied\'z wymaga jawnego rachunku $\langle\varphi^\uparrow\rangle = v_0$
w~j\k{e}zyku substratowym.
\end{remark}

\begin{remark}[Status \emph{ilo\'sciowy} sektora cechowania]\label{rem:gauge-status-quant}
Scenariusze I~i~II maj\k{a} status \textbf{hipotezy roboczej},
wzbogaconej o~formalny szkic emergencji (stw.~\ref{prop:U1-emergence})
i~kwantyzacji \l{}adunku (stw.~\ref{prop:charge-quantization}).
Pe\l{}ne zamkni\k{e}cie wymaga:
\begin{enumerate}[nosep]
  \item jawnego rachunku coarse-grainingu substratu z~symetri\k{a}
    SU(3)$\times$SU(2)$\times$U(1)
    $\to$ lagran\.zjan Yanga--Millsa z~\textbf{obliczonymi}
    sta\l{}ymi sprz\k{e}\.zenia $g_s$, $g_w$, $g'$
    jako funkcjami parametr\'ow substratowych ($v$, $z$, $\beta_c$),
  \item identyfikacji mechanizmu \l{}amania SU(2)$\times$U(1)
    $\to$ U(1)$_{\mathrm{em}}$ w~j\k{e}zyku substratu
    (drugiej kondensacji b\k{a}d\'z anizotropii substratu),
  \item pokazania, \.ze hierarchia mas ($m_e \ll m_t$)
    wynika z~r\'o\.znych konfiguracji defekt\'ow
    (analog: pokolenia = w\k{e}z\l{}y profilu radialnego),
  \item numerycznej symulacji sieci substratowej
    z~wewn\k{e}trzn\k{a} symetri\k{a} SU(3) ---
    weryfikacja emergencji uwi\k{e}zienia
    (por.\ Monte Carlo, \S\ref{ssec:roadmap}).
\end{enumerate}
Sektor cechowania pozostaje \textbf{kluczowym otwartym problemem}
TGP --- jednak~stw.~\ref{prop:U1-emergence} pokazuje,
\.ze \textbf{mechanizm} emergencji jest naturalny
i~nie wymaga dodatkowych za\l{}o\.ze\'n poza
wielosk\l{}adnikowym substratem.
\end{remark}

% =============================================================
\subsection{Podsumowanie brak\'ow i~\emph{roadmap}}\label{ssec:roadmap}
% =============================================================

\begin{remark}[Sze\'sciopoziomowy system status\'ow]\label{rem:status-levels}
Numer tabeli~\ref{sec:status-map} i~ni\.zej u\.zywa
sze\'sciu poziom\'ow domkni\k{e}cia logicznego.
Komenda \verb|\statuslabel{X}| drukuje \statuslabel{X}
wewn\k{a}trz tekstu.
\begin{description}[nosep,style=nextline]
  \item[\statuslabel{Aksjomat}]
    Wej\'scie teorii --- nie dowodzone, przyjmowane a~priori.
    Teoria odpowiada za konsekwencje, nie za aksjomat.
    \emph{Przyk\l{}ady}: substrat $\Gamma$ (A1), $K(0) = 0$ (A1b).
  \item[\statuslabel{Twierdzenie}]
    W~pe\l{}ni wyprowadzone \emph{wy\l{}\k{a}cznie} z~aksjom\'ow
    lub z~aksjomatu + jawnie wymienionych za\l{}o\.ze\'n klasy ---
    \textbf{bez} wolnych parametr\'ow dopasowywanych do obserwacji.
    \emph{Przyk\l{}ady}: $\alpha = 2$ (tw.~\ref{prop:substrate-action}),
    $\beta = \gamma$ (stw.~\ref{prop:vacuum-condition}), $c_{\rm GW} = c_0$
    (stw.~\ref{prop:cT}), MS-TGP (stw.~\ref{prop:MS-TGP}).
  \item[\statuslabel{Propozycja}]
    Wyprowadzone przy dodatkowych za\l{}o\.zeniach lub
    w~zawężonej klasie (np.\ hipoteza pomostowa metryki, klasa $K(\varphi)=\varphi^n$).
    Sta\l{}e mog\k{a} by\'c dopasowane do obserwacji.
    \emph{Przyk\l{}ady}: granica Newtona ($q$ dopasowane),
    emergencja Einsteina do $O(U^2)$, fermiony jako kinki.
  \item[\statuslabel{Hipoteza}]
    Motywowane fizycznie, niesprzeczne z~aksjomatami i~obserwacjami,
    ale \textbf{nie} wyprowadzone z~zasad pierwszych.
    Mo\.ze by\'c zaraz falsyfikowalne.
    \emph{Przyk\l{}ady}: metryka pomostowa $g_{\mu\nu} = f(\Phi)\eta_{\mu\nu}$,
    $\psi_{\rm ini} = 7/6$, punkt sta\l{}y UV.
  \item[\statuslabel{Szkic}]
    Mechanizm jest wskazany i~formalnie zarysowany,
    ale pe\l{}ne wyprowadzenie wymaga dodatkowej pracy.
    \emph{Przyk\l{}ady}: emergencja U(1)/SU(2)/SU(3),
    przej\'scie kink $\to$ fermion Diraca.
  \item[\statuslabel{Program}]
    Kierunek jest jasny, wyniki s\k{a} cz\k{e}\'sciowe lub zerowe.
    Otwarte zadanie badawcze.
    \emph{Przyk\l{}ady}: pe\l{}na kwantyzacja ERG z~$K(\psi) = \psi^4$,
    derywacja metryki z~substratu bez hipotezy pomostowej.
\end{description}

\medskip\noindent
\textbf{Przek\l{}adnik do starego zapisu:}
Stary poziom~\textbf{(W)} odpowiada nowemu \statuslabel{Twierdzenie};
\textbf{(E)} --- g\l{}\'ownie \statuslabel{Propozycja};
\textbf{(P)} --- \statuslabel{Szkic} lub \statuslabel{Program}.
Kompletna mapa patrz sek.~\ref{sec:status-map}.

\medskip\noindent
\textbf{Konwencja j\k{e}zykowa} (obowi\k{a}zuje w~ca\l{}ym dokumencie):
\begin{itemize}[nosep]
  \item \emph{,,X~wynika z~Y''} --- tylko gdy X~jest
    \statuslabel{Twierdzenie} lub \statuslabel{Propozycja}
    i~Y~to spe\l{}nione za\l{}o\.zenia.
  \item \emph{,,X~sugeruje~Z''} lub \emph{,,X~jest zgodne z~Z''} ---
    gdy Z~ma status \statuslabel{Hipoteza} lub \statuslabel{Szkic}.
  \item \emph{,,Przy za\l{}o\.zeniu H~wynika Z''} ---
    dopuszczalne, je\'sli H~jest jawnie nazwana hipotez\k{a}
    i~Z~jest jej konsekwencj\k{a} (warunkowa dedukcja).
\end{itemize}
\end{remark}

\begin{center}
\small
\begin{tabular}{@{}>{\raggedright}p{3.5cm}cp{8cm}@{}}
\toprule
\textbf{Element} & \textbf{Status} & \textbf{Nast\k{e}pny krok} \\
\midrule
Operator $\mathcal{D}$ & Okre\'slony &
  $\nabla^2\Phi/\PhiZero$ (def.~\ref{def:D}) \\
$\mathcal{N}[\Phi]$ & Okre\'slony &
  Trzy cz\l{}ony: grad., kub., kwart. (def.~\ref{def:N}) \\
Profil $\phi(r)$ & Asymptotyki &
  Numeryczne rozwi\k{a}zanie~\eqref{eq:spher}
  dla danych $\alpha,\beta,\gamma$ \\
Trzy re\.zimy & Wykazane &
  Numeryczny $V_{\mathrm{eff}}(d)$ i~dopasowanie
  do obserwacji \\
Granica Newtona & Wyprowadzona &
  Twierdzenie~\ref{thm:newton} \\
Metryka & Wyprowadzona &
  Z~warunk\'ow wewn\k{e}trznych TGP
  (stw.~\ref{prop:conformal-unique}),
  $\beta_{\mathrm{PPN}} = 1$ jako predykcja;
  relacja pot\k{e}gowa/eksponencjalna
  (uw.~\ref{rem:power-vs-exp}) \\
Kosmologia & \textbf{Zamkni\k{e}ty} &
  Friedmann z~TGP (tw.~\ref{thm:einstein-emergence}),
  $W(\psi)$ z~$U$ (stw.~\ref{prop:W-from-U}),
  $G_{\mathrm{eff}}(k,a)$ (tw.~\ref{thm:Geff-k}),
  $w_{\mathrm{DE}} = -1 + O(10^{-9})$ \\
Dzia\l{}anie & Zamkni\k{e}ty &
  Wariacja $\Rightarrow$ r.\ pola + r.\ kowariantne;
  well-posedness w~$H^s$ (tw.~\ref{thm:well-posedness}) \\
$\alpha = 2$ (wygenerowany) & Zamkni\k{e}ty &
  Z~wariacji dzia\l{}ania (wn.~\ref{cor:alpha-fixed});
  dwa wolne parametry: $\beta$, $\gamma$ \\
$\beta,\gamma$ & Zarys wi\k{e}z\'ow + MK-RG &
  Tabela ogranicze\'n (\S\ref{ssec:wiezy});
  relacja z~substratem (\S\ref{ssec:substrat-param});
  WF FP: $|r^*|/u^* = 0.575$ (uw.~\ref{rem:mk-rg-results}) \\
$\Phi$ a~$\Ffield$ & Zamkni\k{e}ty &
  $\Ffield_{\!\mathrm{abs}} = \Phi^2/\PhiZero^2$ (def.~\ref{def:Fabs});
  $\Floc = (\Phi - \PhiZero)^2$ (def.~\ref{def:Floc});
  dwa poziomy (uw.~\ref{rem:F-dwa-poziomy});
  zerowa suma jako~\eqref{eq:zero-constraint};
  dynamika $\Ffield$ z~\eqref{eq:F-evolution} \\
R\'ownanie kowariantne & Sformu\l{}owane &
  Pe\l{}ne r\'ownanie z~$\partial^2/\partial t^2$
  (tw.~\ref{thm:covariant-field}); autoreferencyjno\'s\'c
  (uw.~\ref{rem:autoref}); hiperboliczno\'s\'c
  (stw.~\ref{prop:hyperbolic}) \\
Sygnatura Lorentza & Argument &
  Emergencja z~przej\'scia fazowego
  (\S\ref{ssec:sygnatura}); kontynuacja Wicka \\
Warunki energetyczne & Wykazane &
  SEC naruszone (stw.~\ref{prop:SEC-violation});
  WEC spe\l{}nione w~s\l{}abym polu;
  sp\'ojno\'s\'c z~brakiem osobliwo\'sci
  (uw.~\ref{rem:energ-sing}) \\
Niestabilno\'s\'c $\Nzero$ & Potr\'ojne uzasadnienie &
  Aksjomat (A4), termodynamika (stw.~\ref{prop:P3-deriv}),
  potencja\l{} (stw.~\ref{prop:P3-action}) \\
Perturbacje & \textbf{Zamkni\k{e}ty} &
  $G_{\mathrm{eff}}(k,a)$ (tw.~\ref{thm:Geff-k}),
  slip $\eta$ (wn.~\ref{cor:grav-slip}),
  SVT dekompozycja;
  $\delta P/P \sim 10^{-51}$ (niedetektowalne) \\
Metryka z~zasad TGP & Wyprowadzona &
  Cztery warunki $\Rightarrow$ posta\'c eksponencjalna
  (\S\ref{ssec:metryka-deriv}) \\
Emergencja Einsteina (og\'olna) & \textbf{Wyprowadzona} &
  Hierarchia: FRW dok\l{}adnie, statyczne do 2PN,
  $O(U^3)$ predykcja, Bianchi propaguje wi\k{e}zy
  (tw.~\ref{thm:general-emergence},
  stw.~\ref{prop:emergence-hierarchy}) \\
Emergencja Einsteina (linearyzowana, $U(t,\mathbf{x})$) &
  \textbf{Udowodniona} &
  $G^{(1)}_{\mu\nu}$ obliczone symbolicznie (13/13 PASS);
  emergencja \textbf{dok\l{}adna}: py\l{}, pr\'o\.znia+GW;
  \textbf{ograniczona}: izotropowe promieniowanie statyczne
  (fizycznie nierealizowalne; FRW dok\l{}adne osobno)
  (tw.~\ref{thm:einstein-emergence-linearized}) \\
Kwantyzacja $\Phi$ & Zarys &
  Propagator, obci\k{e}cie UV z~substratu
  (\S\ref{ssec:kwantyzacja}) \\
Sektor nielokalny & Zarys &
  J\k{a}dro $K(r)$, warunek bounce
  (\S\ref{ssec:nielokalny}) \\
Materia jako defekty & Sformu\l{}owany &
  Kink radialny, masa defektu, generacje
  (\S\ref{ssec:defekty}) \\
Sektor fermionowy & \textbf{Sformu\l{}owany} &
  Tetrada (def.~\ref{def:tetrada}), Dirac (hip.~\ref{hyp:dirac-TGP}),
  koneksja spinowa do 2.\ rz\k{e}du (stw.~\ref{prop:spin-connection}),
  ZS1$\to$spin-$\frac{1}{2}$ (tw.~\ref{thm:zs1-spin});
  samospojna p\k{e}tla $\Psi$-$\Phi$-$\Delta$
  (\S\ref{ssec:fermiony}) \\
Pola cechowania & \textbf{Szkic formalny} &
  U(1) z~fazy Goldstone'a (stw.~\ref{prop:U1-emergence}),
  kwantyzacja \l{}adunku z~topologii (stw.~\ref{prop:charge-quantization}),
  SU($N$) z~wielosk\l{}adnikowego substratu (uw.~\ref{rem:SUN-gauge});
  \textbf{otwarty}: jawne sta\l{}e sprz\k{e}\.zenia $g_s, g_w, g'$
  (\S\ref{ssec:gauge-emergence}) \\
Ciemna energia $U(\varphi)$ & \textbf{Wyprowadzony} &
  $\Lambda_{\mathrm{eff}} \approx \gamma/12$
  (stw.~\ref{prop:Lambda-eff}),
  $w_{\mathrm{DE}}(z)$ (\S\ref{ssec:de-oszacowanie}) \\
Mody tensorowe GW & \textbf{Rozwi\k{a}zany} &
  Zakaz konforemny (tw.~\ref{thm:no-tensor});
  substratowe $\sigma_{ab}$ (\S\ref{ssec:tensor-substrate}):
  5~d.o.f.\ spin-2, $h \sim 10^{-22}$,
  $c_{\mathrm{GW}} = c_0$ dok\l{}adnie \\
$\Srel$ z~dynamiki $\Phi$ & Wyprowadzony &
  $\Srel[\Phi]$ explicite (\S\ref{ssec:Srel-formal});
  drugie prawo z~dyfuzji \\
Wyk\l{}adniki $c,\hbar,G$ & Wyprowadzone &
  Jednoznaczno\'s\'c z~3 warunk\'ow (tw.~\ref{thm:exponents}) \\
Dualno\'s\'c $\Nzero$ & Sformalizowana &
  Inwersja $\chi\mapsto 1/\chi$ (tw.~\ref{thm:N0-duality}) \\
Jedyno\'s\'c $\Nzero$ & Wyprowadzona &
  $\mathcal{S}_0$ sp\'ojny (tw.~\ref{thm:N0-unique}) \\
Pr\k{a}d $\Ffield$ & Sformu\l{}owany &
  $j_\Ffield^\mu$ (stw.~\ref{prop:F-current}) \\
Topologia domen & Zarys &
  Domeny metryczne, konsekwencje CMB
  (\S\ref{ssec:topologia}) \\
Wymiar $3{+}1$ & Argument &
  Optymalny wymiar z~modelu Isinga
  (\S\ref{ssec:wymiar}) \\
Og\'olna emergencja Einsteina & \textbf{Wyprowadzona} &
  Tw.~\ref{thm:general-emergence}: FRW dok\l{}., stat.\ do 2PN,
  Bianchi propaguje wi\k{e}zy (\S\ref{ssec:general-emergence}) \\
Sprz\k{e}\.zenie z~materi\k{a} & \textbf{Sformu\l{}owane} &
  Aks.~\ref{ax:metric-coupling}: jedno sprz\k{e}\.z.\ metryczne,
  brak pi\k{a}tej si\l{}y (\S\ref{ssec:matter-coupling}) \\
Globalna stabilno\'s\'c & \textbf{Wyprowadzona} &
  Tw.~\ref{thm:global-stability}: ghost-free, gradient-stable,
  basen, Cauchy (\S\ref{ssec:global-stability}) \\
Selekcja $\beta=\gamma$ & \textbf{Wyprowadzona} &
  Stw.~\ref{prop:vacuum-selection}: RG + atraktor + formalny
  (\S\ref{ssec:vacuum-selection}) \\
Falsyfikowalno\'s\'c & \textbf{Sformu\l{}owana} &
  Tab.~\ref{tab:falsification}: $\eta_{\mathrm{slip}}$, QNM, $\Lambda$
  (\S\ref{ssec:falsification}) \\
\bottomrule
\end{tabular}
\end{center}

% =============================================================
\subsection{Predykcje vs.~odtwarzanie}\label{ssec:predykcje-vs-odtwarzanie}
% =============================================================

\begin{remark}[Klasyfikacja wynik\'ow TGP]\label{rem:predykcje-klasyfikacja}
Dla zewn\k{e}trznego czytelnika kluczowe jest rozr\'o\.znienie mi\k{e}dzy:
\begin{itemize}[nosep]
  \item \textbf{Wyprowadzenie} (ang.\ \emph{derivation}): wynik wynika z~aksjomatu
    algebraicznie lub numerycznie, bez dopasowania do danych.
  \item \textbf{Odtwarzanie} (ang.\ \emph{recovery}): teoria reprodukuje obs.\ wynik
    przez dopasowanie co najmniej jednego parametru.
  \item \textbf{Predykcja} (ang.\ \emph{prediction}): niezale\.zna, weryfikowalna
    prognoza, \textbf{jeszcze nieporównana} z~danymi lub daj\k{a}ca
    liczb\k{e} \emph{bez dodatkowego dopasowania}.
  \item \textbf{Fit}: wynik jest \emph{z definicji} wej\'sciem kalibracji;
    nie jest ani predykcj\k{a}, ani sprawdzianem.
\end{itemize}
\end{remark}

\begin{center}\footnotesize
\begin{tabular}{@{}p{4.5cm}lcp{3.5cm}@{}}
\toprule
\textbf{Zjawisko / Warto\'s\'c} & \textbf{Status TGP}
  & \textbf{Pewno\'s\'c} & \textbf{Wymaga} \\
\midrule

\multicolumn{4}{l}{\textit{Wyniki wyprowadzone (niezale\.zne od danych)}} \\[2pt]

$\alpha = 2$ (operator $\mathcal{D}$)
  & Wyprowadzone & $\star\star\star\star\star$
  & \emph{nic} \\[3pt]

$\beta = \gamma$ (warunek pr\'o\.zni)
  & Wyprowadzone & $\star\star\star\star\star$
  & \emph{nic} \\[3pt]

$c_{\rm GW} = c_0$ (pr\k{e}dko\'s\'c GW)
  & Wyprowadzone & $\star\star\star\star$
  & jednorodne t\l{}o \\[3pt]

PPN: $\gamma_{\rm PPN} = \beta_{\rm PPN} = 1$
  & Wyprowadzone & $\star\star\star\star$
  & metryka konformalna \\[3pt]

$W(1) \neq 0 \Rightarrow \Lambda_{\rm eff} \neq 0$
  & Wyprowadzone & $\star\star\star\star$
  & $\beta = \gamma$ (tw.~\ref{thm:vacuum-source}) \\[3pt]

$N_{\rm gen} = 3$ (trzy generacje)
  & Wyprowadzone & $\star\star\star$
  & WKB + $a_\Gamma$ jako wej\'scie \\[8pt]

\multicolumn{4}{l}{\textit{Predykcje (niezale\.zna liczba / sygn.~testowalna)}} \\[2pt]

$\Delta f_{\rm QNM}/f_{\rm GR} \sim 0{,}5\%$
  & Predykcja & $\star\star\star$
  & Einstein Tel. \\[3pt]

$\Delta r_{\rm shadow}/r_{\rm GR} \sim 2\%$
  & Predykcja & $\star\star\star$
  & ngEHT \\[3pt]

$M_{\rm rot}(r) \propto M^{-1/9}$ (krzywe rotacji)
  & Predykcja & $\star\star\star\star$
  & SPARC + SKA \\[3pt]

$\eta_{\rm slip} = e^{-2U} \neq 1$
  & Predykcja & $\star\star\star$
  & Euclid/LSST $\sigma < 10^{-4}$ \\[3pt]

$(n_s - 1)_{\rm TGP} = -4\varepsilon_H - 4\varepsilon_\psi$
  & Predykcja (do obliczenia) & $\star\star$
  & symulacja MS-TGP \\[8pt]

\multicolumn{4}{l}{\textit{Odtwarzanie (wymaga $\leq 1$ param.\ dopasowanego)}} \\[2pt]

$\Lambda_{\rm eff} \approx \gamma/12 \approx H_0^2$
  & Odtwarzane & $\star\star\star$
  & $\gamma$ z~jednej obserwacji \\[3pt]

Formu\l{}a Koidego $Q = 2/3$
  & Odtwarzane & $\star\star\star\star$
  & $\alpha_K$, $a_\Gamma$ (1 param.) \\[3pt]

Mody oddechowe GW (masywne)
  & Odtwarzane & $\star\star\star$
  & $m_{\rm sp}$ z~$\gamma$ \\[8pt]

\multicolumn{4}{l}{\textit{Kalibracja (wej\'scia, nie sprawdziany)}} \\[2pt]

$r_{21} = 206{,}77$ ($m_\mu/m_e$)
  & Fit (kalibracja $a_\Gamma$) & ---
  & --- \\[3pt]

$\Lambda_{\rm obs}$ (warto\'s\'c)
  & Fit (kalibracja $\PhiZero$) & ---
  & --- \\[3pt]

$G_0$ (warto\'s\'c bezwzgl\k{e}dna)
  & Fit (kalibracja $q$) & ---
  & --- \\

\bottomrule
\end{tabular}
\end{center}

\begin{remark}[Interpretacja gwiazdek]\label{rem:stars-meaning}
Kolumna ``Pewno\'s\'c'' ocenia solidno\'s\'c \textbf{wewn\k{e}trzn\k{a}} wyprowadzenia,
nie zgodno\'s\'c z~danymi obs.\ (to jest osobna kwestia):
$\star\star\star\star\star$~= algebraicznie pewne z~aksjomatu;
$\star\star\star$~= wynika z~hipotezy + numeryki;
$\star\star$~= szacunkowe, wymaga symulacji.
\end{remark}

% =============================================================
% NOWE SEKCJE DOMYKAJĄCE (v11+)
% =============================================================

% =============================================================
\subsection{Og\'olne twierdzenie o~emergencji Einsteina}\label{ssec:general-emergence}
% =============================================================

Twierdzenie~\ref{thm:einstein-emergence} udowadnia emergencj\k{e}
r\'owna\'n Einsteina w~symetrii FRW. Poni\.zsze twierdzenie
rozszerza ten wynik na \textbf{og\'olny przypadek}.

\begin{theorem}[Og\'olna emergencja Einsteina w~TGP]\label{thm:general-emergence}
Niech $g_{\mu\nu}^{\mathrm{eff}}(\Phi)$ b\k{e}dzie metryk\k{a}
eksponencjaln\k{a} TGP (hipoteza~\ref{hyp:metric-exp}) z~$U = \delta\Phi/\PhiZero$.
Definiujemy:
\begin{enumerate}[nosep]
  \item \textbf{Tensor Einsteina metryki efektywnej:}
    $G_{\mu\nu}^{\mathrm{eff}}
    \equiv R_{\mu\nu}[g^{\mathrm{eff}}] - \frac{1}{2}g_{\mu\nu}^{\mathrm{eff}}
    R[g^{\mathrm{eff}}]$;
  \item \textbf{Efektywny tensor energii-p\k{e}du pola $\Phi$:}
    $T_{\mu\nu}^{(\Phi)} = \frac{1}{\kappa}\,G_{\mu\nu}^{\mathrm{eff}}$
    (definicja);
  \item \textbf{Residuum emergencji:}
    $\mathcal{R}_{\mu\nu}
    \equiv G_{\mu\nu}^{\mathrm{eff}} - \kappa\,T_{\mu\nu}^{(\Phi,\mathrm{dyn})}$,
    gdzie $T_{\mu\nu}^{(\Phi,\mathrm{dyn})}$ jest tensorem energii-p\k{e}du
    \emph{wyprowadzonym z~dzia\l{}ania} $S[\Phi]$~\eqref{eq:action}.
\end{enumerate}

W\'owczas:
\begin{enumerate}[(i)]
  \item \textbf{Symetria FRW:}
    $\mathcal{R}_{\mu\nu} = 0$ \textbf{dok\l{}adnie}
    (tw.~\ref{thm:einstein-emergence}).

  \item \textbf{S\l{}abe pole statyczne} ($|U| \ll 1$):
    \begin{equation}\label{eq:residuum-static}
      \mathcal{R}^t{}_t = \bigl(\nabla U\bigr)^2\,\bigl(1 + \mathcal{O}(U)\bigr),
    \end{equation}
    czyli $\mathcal{R}_{\mu\nu} = \mathcal{O}(U^2 \cdot \nabla^2 U)$
    w~sektorze $tt$. Fizycznie: residuum jest energi\k{a}
    gradientow\k{a} pola~$\Phi$, naturalnie obecn\k{a} w~TGP
    (nie ma ,,pr\'o\.zni'' --- pole jest wsz\k{e}dzie).

  \item \textbf{Propagacja wi\k{e}z\'ow:}
    to\.zsamo\'s\'c Bianchiego $\nabla_\mu G^{\mu\nu}_{\mathrm{eff}} = 0$
    gwarantuje, \.ze je\.zeli $\mathcal{R}_{\mu\nu}$ znika na
    danych pocz\k{a}tkowych $\Sigma_{t_0}$ i~r\'ownanie pola
    zachodzi wsz\k{e}dzie, to $\mathcal{R}_{\mu\nu} = 0$
    na ca\l{}ej dziedzinie przyczynowej $D^+(\Sigma_{t_0})$.

  \item \textbf{Parametry PPN:} $\gamma_{\mathrm{PPN}} = \beta_{\mathrm{PPN}} = 1$
    (zgodno\'s\'c z~OTW do 2PN).
    Pierwsza rozbie\.zno\'s\'c: $\Delta g_{00} = -\frac{1}{6}U^3$
    na poziomie 3PN.

  \item \textbf{Wi\k{e}z antypodalny:}
    $g_{00}\cdot g_{rr} = -1$ \textbf{dok\l{}adnie} w~TGP
    (w\l{}asno\'s\'c metryki eksponencjalnej).
    Schwarzschild narusza to od $\mathcal{O}(U^2)$.
\end{enumerate}
\end{theorem}

\begin{proof}[Szkic dowodu]
\emph{Punkt (i):} Dow\'od w~tw.~\ref{thm:einstein-emergence}.

\emph{Punkt (ii):}
Tensor Einsteina metryki statycznej sferycznej z~$g_{00} = -e^{-2U}$,
$g_{ij} = e^{2U}\delta_{ij}$ obliczamy bezpo\'srednio
(skrypt \texttt{einstein\_emergence\_proof.py}).
Sk\l{}adowa $G^t{}_t = e^{-2U}[\nabla^2 U + (\nabla U)^2]$.
Poza \'zr\'od\l{}em: $\nabla^2 U \approx 0$ (r\'ownanie Poissona
z~poprawk\k{a} Yukawa), wi\k{e}c $G^t{}_t \approx (\nabla U)^2
= \mathcal{O}(U^2/r^2)$. To nie jest zero --- ale jest
\emph{tensorem energii-p\k{e}du pola $\Phi$}, kt\'ore w~TGP
jest wsz\k{e}dzie obecne.

\emph{Punkt (iii):}
Analogicznie do dowodu tw.~\ref{thm:einstein-emergence}, krok~4.
To\.zsamo\'s\'c Bianchiego daje
$\dot{\mathcal{R}} + f(t)\,\mathcal{R} = 0$, jedynym
rozwi\k{a}zaniem z~$\mathcal{R}(t_0) = 0$ jest
$\mathcal{R} \equiv 0$.

\emph{Punkt (iv):}
Rozwinięcie $e^{-2U} = 1 - 2U + 2U^2 - \frac{4}{3}U^3 + \ldots$
vs.\ Schwarzschild (izotrop.):
$-\bigl(\frac{1-U/2}{1+U/2}\bigr)^2 = -(1 - 2U + 2U^2 - \frac{3}{2}U^3 + \ldots)$.
Zgodno\'s\'c do $\mathcal{O}(U^2)$; rozbie\.zno\'s\'c
$\Delta g_{00} = (-\frac{4}{3} + \frac{3}{2})U^3 = \frac{1}{6}U^3$
(z~dok\l{}adno\'sci\k{a} do znaku) na poziomie 3PN.

\emph{Punkt (v):}
$g_{00}\cdot g_{rr} = (-e^{-2U})(e^{2U}) = -1$ --- trywialna
to\.zsamo\'s\'c metryki antypodalnej. Schwarzschild (izotrop.):
$g_{00}^S g_{rr}^S = -\bigl(\frac{1-U/2}{1+U/2}\bigr)^2(1+U/2)^4
= -(1-U/2)^2(1+U/2)^2 = -(1 - U^2/4)^2 \neq -1$
dla $U \neq 0$.
\end{proof}

\begin{remark}[Emergencja Einsteina do wszystkich rz\k{e}d\'ow
--- argument z~to\.zsamo\'sci Bianchiego]%
\label{rem:bianchi-all-orders}
Ograniczenie ,,do $O(U^2)$'' w~punkcie (ii) tw.~\ref{thm:general-emergence}
dotyczy jedynie \emph{lokalizacji residuum w~sektorze statycznym}.
Emergencja r\'owna\'n Einsteina jest w~istocie dok\l{}adna
do wszystkich rz\k{e}d\'ow w~nast\k{e}puj\k{a}cym sensie:

\begin{enumerate}[nosep]
  \item Metryka $g_{\mu\nu}^{\rm eff}(\Phi)$ jest zdefiniowana algebraicznie
    (nie perturbacyjnie). Tensor Ricciego $R_{\mu\nu}[g^{\rm eff}]$
    istnieje dok\l{}adnie; tensor Einsteina
    $G_{\mu\nu}^{\rm eff} = R_{\mu\nu} - \frac{1}{2}g_{\mu\nu}R$
    spe\l{}nia to\.zsamo\'s\'c Bianchiego:
    \begin{equation}\label{eq:bianchi-exact}
      \nabla_\mu G^{\mu\nu}_{\rm eff} \equiv 0
      \qquad\text{(to\.zsamo\'s\'c geometryczna, dok\l{}adna).}
    \end{equation}

  \item Definiuj\k{a}c \emph{efektywny} tensor energii-p\k{e}du
    $T_{\mu\nu}^{\rm eff} \equiv G_{\mu\nu}^{\rm eff}/\kappa$,
    zachowanie $\nabla_\mu T^{\mu\nu}_{\rm eff} = 0$ jest
    \textbf{automatyczne} --- wynika z~geometrii, nie z~dynamiki.
    R\'ownania Einsteina $G_{\mu\nu} = \kappa T_{\mu\nu}$ s\k{a} wówczas
    \emph{definicj\k{a}} $T_{\mu\nu}$, nie postulatem.

  \item Nietrywialna tre\'s\'c fizyczna: r\'ownanie pola $\Phi$
    (dynamika TGP) implikuje, \.ze $T_{\mu\nu}^{\rm eff}$ pokrywa si\k{e}
    z~$T_{\mu\nu}^{(\Phi,\rm dyn)}$ (tensorem z~dzia\l{}ania)
    \textbf{dok\l{}adnie w~FRW} i~do $O(U^2)$ w~statycznym polu.
    Residuum $\mathcal{R}_{\mu\nu} \sim (\nabla U)^2$ w~sektorze
    statycznym jest fizyczn\k{a} energi\k{a} pola TGP
    i~\emph{nie narusza} to\.zsamo\'sci Bianchiego.
\end{enumerate}

Podsumowuj\k{a}c: r\'ownania Einsteina nie s\k{a} aproksymacj\k{a}
o~ograniczonej wa\.zno\'sci, lecz \textbf{dok\l{}adn\k{a}
konsekwencj\k{a} geometrii} emergentnej metryki.
Jedyna niedok\l{}adno\'s\'c to identyfikacja $T_{\mu\nu}^{\rm eff}$
z~$T_{\mu\nu}^{(\Phi,\rm dyn)}$, kt\'ora jest dok\l{}adna w~FRW
i~poprawna do $O(U)$ w~polu statycznym.
\end{remark}

\begin{remark}[Interpretacja residuum]\label{rem:residuum}
Fakt, \.ze $\mathcal{R}_{\mu\nu} \neq 0$ w~statycznym polu
\textbf{nie jest b\l{}\k{e}dem} TGP. Jest konsekwencj\k{a}
fundamentalnej r\'o\.znicy ontologicznej:
\begin{itemize}[nosep]
  \item W~OTW: ,,pr\'o\.znia'' oznacza $T_{\mu\nu} = 0$,
    wi\k{e}c $G_{\mu\nu} = 0$.
  \item W~TGP: nie ma pr\'o\.zni --- pole $\Phi > 0$ wsz\k{e}dzie
    (sektor $\mathcal{S}_1$). Gradient $\nabla\Phi$ niesie
    energi\k{e}, wi\k{e}c $G_{\mu\nu} \neq 0$ nawet ,,na zewn\k{a}trz'' masy.
\end{itemize}
Residuum $\mathcal{R}_{\mu\nu} \sim (\nabla U)^2$ jest fizyczn\k{a}
energi\k{a} pola konstytuuj\k{a}cego przestrze\'n ---
jest \emph{predykcj\k{a}}, nie defektem.
\end{remark}

\begin{proposition}[Obserwacyjne granice odchylenia 3PN]\label{prop:3PN-bounds}
Jedyna r\'o\.znica mi\k{e}dzy metryk\k{a} TGP a~Schwarzschildem (w~postaci
izotropowej) pojawia si\k{e} na poziomie 3PN:
\begin{equation}\label{eq:3PN-deviation}
  \Delta g_{00} = \tfrac{1}{6}\,U^3,
  \qquad U \equiv \frac{GM}{c_0^2\,r}.
\end{equation}
\noindent Konsekwencje obserwacyjne:
\begin{enumerate}[nosep]
  \item \textbf{S\l{}o\'nce} ($U_\odot \approx 2 \times 10^{-6}$):
    $\Delta g_{00} \sim 10^{-18}$ --- \textbf{niewykrywalne}
    nawet przez LISA ani przysz\l{}e zegary optyczne.
  \item \textbf{Gwiazda neutronowa} ($U_{\mathrm{NS}} \approx 0{,}2$):
    $\Delta g_{00} \sim 10^{-3}$ ---
    \textbf{potencjalnie testowalne} przez przysz\l{}e
    pomiary promieniowania rentgenowskiego (ATHENA, eXTP).
  \item \textbf{Uk\l{}ad podw\'ojny pulsar\'ow}
    ($U_{\mathrm{bin}} \approx 0{,}05$ przy separacji $\sim 10R_s$):
    $\Delta g_{00} \sim 2 \times 10^{-5}$ ---
    \textbf{dost\k{e}pne} dla pulsar timing nast\k{e}pnej generacji (SKA).
\end{enumerate}
\end{proposition}

\begin{proof}
Potencja\l{} $U$ wynika bezpo\'srednio z~granicy newtonowskiej
(tw.~\ref{thm:newton}). Oszacowania numeryczne:
$U_\odot = G_0 M_\odot/(c_0^2 R_\odot) \approx 2{,}1 \times 10^{-6}$,
$U_{\mathrm{NS}} = G_0 (1{,}4\,M_\odot)/(c_0^2 \cdot 10\,\mathrm{km})
\approx 0{,}21$. Wsp\'o\l{}czynnik $1/6$ pochodzi z~r\'o\.znicy
rozwinięcia eksponencjalnego ($-4U^3/3$)
i~Schwarzschilda izotropowego ($-3U^3/2$):
$-4/3 + 3/2 = 1/6$
(zweryfikowane skryptem \texttt{deep\_consistency\_v17.py}, testy K3, L1--L4).
\end{proof}

\begin{proposition}[Zbie\.zno\'s\'c rozwini\k{e}cia PN i~kontrola reszty]%
\label{prop:PN-convergence}
Metryka TGP $g_{00} = -e^{-2U}$ ma rozwini\k{e}cie
$e^{-2U} = \sum_{n=0}^\infty (-2U)^n/n!$, kt\'ore jest
\textbf{zbie\.zne dla ka\.zdego} $U \geq 0$ (ca\l{}kowita
--- promie\'n zbie\.zno\'sci $R = \infty$).

Reszta po obci\k{e}ciu na rz\k{e}dzie $N$:
\begin{equation}\label{eq:PN-remainder}
  \bigl|R_N(U)\bigr| =
  \left|e^{-2U} - \sum_{n=0}^{N}\frac{(-2U)^n}{n!}\right|
  \leq \frac{(2U)^{N+1}}{(N+1)!}.
\end{equation}
Dla por\'ownania, metryka Schwarzschilda we wsp\'o\l{}rz\k{e}dnych
izotropowych:
$g_{00}^{\mathrm{S}} = -\bigl(\frac{1-U/2}{1+U/2}\bigr)^2$
ma \textbf{sko\'nczony} promie\'n zbie\.zno\'sci $R = 2$.
Estymata reszty:
\begin{center}
\begin{tabular}{lccc}
\toprule
Rz\k{a}d PN & $U = 10^{-6}$ (S\l{}o\'nce) & $U = 0{,}2$ (NS) & $U = 0{,}5$ (BH) \\
\midrule
2PN ($R_2$) & $10^{-18}$ & $10^{-3}$ & $4 \times 10^{-2}$ \\
3PN ($R_3$) & $10^{-24}$ & $3 \times 10^{-4}$ & $2 \times 10^{-2}$ \\
5PN ($R_5$) & $10^{-36}$ & $10^{-6}$ & $3 \times 10^{-4}$ \\
\bottomrule
\end{tabular}
\end{center}
\textbf{Wniosek:} Dla wszelkich obserwowanych uk\l{}ad\'ow
($U \leq 0{,}5$) rozwini\k{e}cie PN metryki TGP zbiega
szybciej ni\.z Schwarzschilda i~nie wymaga resummacji.
\end{proposition}

\begin{proof}
Estymata~\eqref{eq:PN-remainder} wynika z~reszty Lagrange'a
funkcji $e^x$ w~$x = -2U$.
Warto\'sci w~tabeli:
$|R_2(0{,}2)| = (0{,}4)^3/6 = 1{,}07 \times 10^{-2}$
(g\'orna granica; dok\l{}adna r\'o\.znica jest mniejsza
o~czynnik~$\sim 6$ z~powodu alternuj\k{a}cych znak\'ow).
Dok\l{}adne warto\'sci sprawdzone
skryptem \texttt{adm\_emergence\_general.py}.
\end{proof}

% ─────────────────────────────────────────────────────────────
\begin{proposition}[Stabilno\'s\'c korekty 3PN wzgl\k{e}dem t\l{}a kosmologicznego]%
\label{prop:3PN-cosmo-stability}
Odchylenie $\Delta g_{00} = \tfrac{1}{6}U^3$ (r\'ow.~\eqref{eq:3PN-deviation})
oraz korekta fazy TaylorF2 (r\'ow.~\eqref{eq:H-3PN}, krok A13)
\textbf{nie ulegaj\k{a} istotnym zmianom} w~dynamicznym tle
kosmologicznym $\varphi_{\mathrm{bg}}(t) = \Phi_{\mathrm{bg}}(t)/\PhiZero$.

Precyzyjnie: kosmologiczna korekta do wspó\l{}czynnika $\delta c_2 = -1/3$
jest wzgl\k{e}dnie t\l{}umiona przez adiabatyczny parametr
\begin{equation}\label{eq:3PN-adiabatic}
  \varepsilon_{\mathrm{adiab}}
  \;=\; \frac{\dot\varphi_{\mathrm{bg}}}{\varphi_{\mathrm{bg}}\,\omega_{\mathrm{GW}}}
  \;\approx\; \frac{H_0}{\omega_{\mathrm{GW}}},
\end{equation}
a~wzgl\k{e}dna zmiana $\delta c_2$:
\begin{equation}\label{eq:3PN-cosmo-corr}
  \left|\frac{\delta(\delta c_2)}{\delta c_2}\right|
  \;\lesssim\; \varepsilon_{\mathrm{adiab}}^2
  \;\approx\; \left(\frac{H_0}{f_{\mathrm{GW}}}\right)^2
  \;\sim\; 10^{-36}
  \quad\text{dla LIGO/LISA.}
\end{equation}
Predykcja K20 jest zatem \textbf{kosmologicznie stabilna}.
\end{proposition}

\begin{proof}
W~TGP metryka efektywna w~FLRW z~wolnozmiennym t\l{}em
$\varphi_{\mathrm{bg}}(t)$ ma posta\'c (tw.~\ref{thm:covariant-field},
krok A15, r\'ow.~\eqref{eq:H-Friedmann}):
\begin{equation}\label{eq:3PN-flrw-metric}
  g_{tt}(\Phi) = -e^{-2U/\varphi_{\mathrm{bg}}(t)},
  \qquad
  U = \frac{G_0 M}{c_0^2\,r},
\end{equation}
gdzie $\varphi_{\mathrm{bg}} = 1 + \delta\varphi(t)$
z~$\delta\varphi = \mathcal{O}(H_0/\omega_{\mathrm{GW}})$.

Rozwijamy $e^{-2U/\varphi_{\mathrm{bg}}}$ do trzeciego rz\k{e}du:
\begin{equation}
  g_{tt} = -1 + \frac{2U}{\varphi_{\mathrm{bg}}}
  - 2\!\left(\frac{U}{\varphi_{\mathrm{bg}}}\right)^{\!2}
  + \frac{4}{3}\!\left(\frac{U}{\varphi_{\mathrm{bg}}}\right)^{\!3}
  + \mathcal{O}(U^4).
\end{equation}
Wsp\'o\l{}czynnik 3PN:
\[
  c_2^{\mathrm{TGP,cosmo}}(t) = \frac{4/3}{\varphi_{\mathrm{bg}}^3}
  = \frac{4}{3}\bigl(1 + 3\delta\varphi(t)\bigr)^{-1}
  \approx \frac{4}{3}\bigl(1 - 3\delta\varphi(t)\bigr).
\]
Korekta wzgl\k{e}dem pr\'o\.zni $\varphi_{\mathrm{bg}} = 1$:
\[
  \delta c_2^{\mathrm{cosmo}}
  = c_2^{\mathrm{TGP}} - c_2^{\mathrm{TGP,cosmo}}
  = 4\,\delta\varphi \cdot \frac{U^3}{r^3} + \ldots
\]
Ewolucja kosmologiczna pola $\varphi_{\mathrm{bg}}$ odbywa si\k{e}
na skali czasu Hubble'a: $\dot\delta\varphi \sim H_0\,\delta\varphi$.
Dla \'zr\'od\l{}a GW o~cz\k{e}sto\'sci $f_{\mathrm{GW}} \gg H_0$
pole $\varphi_{\mathrm{bg}}$ jest \textbf{zamro\.zone} w~trakcie
emisji:
\[
  \Delta\varphi_{\mathrm{bg}}\bigl|_{\tau_{\mathrm{GW}}}
  = H_0 \cdot \tau_{\mathrm{GW}} \cdot \delta\varphi
  = \frac{H_0}{\omega_{\mathrm{GW}}} \cdot \delta\varphi
  \sim \varepsilon_{\mathrm{adiab}} \ll 1.
\]
Std\k{a}d korekta kosmologiczna do $\delta c_2$:
\[
  \left|\frac{\delta(\delta c_2)}{\delta c_2}\right|
  \;\sim\; \varepsilon_{\mathrm{adiab}}^2
  = \left(\frac{H_0}{\omega_{\mathrm{GW}}}\right)^2,
\]
co dla LIGO ($f \sim 100$ Hz) daje $\varepsilon_{\mathrm{adiab}}
\sim 10^{-22}$, a~dla LISA ($f \sim 10^{-3}$ Hz):
$\varepsilon_{\mathrm{adiab}} \sim 10^{-15}$.
W~obu przypadkach korekta jest \textbf{poni\.zej progu detekcji}
($\sim 10^{-3}$--$10^{-4}$ przy aktualnej dok\l{}adno\'sci).

\medskip\noindent\textbf{Luka A11\texorpdfstring{$\to$}{→}A13
zamkni\k{e}ta.}
Powy\.zszy argument dowodzi, \.ze krok A13
(r\'ow.~\eqref{eq:H-3PN}, krok~A13 \l{}a\'ncucha) jest
\textbf{kosmologicznie odporny}: tle kosmologiczne nie modyfikuje
odchylenia 3PN w~zakresie obserwacyjnym obecnych i~planowanych
detektor\'ow GW.
Ko\'ncowy wniosek: predykcja Kill-shot K20 ($\delta c_2 = -1/3$)
jest \textbf{stabilna wzgl\k{e}dem dynamicznych st\l{}a\'n TGP}
na wszystkich skalach czasowych istotnych dla LISA/ET.
\end{proof}

% ─────────────────────────────────────────────────────────────
\begin{theorem}[Linearyzowany tensor Einsteina dla og\'olnego $U(t,\mathbf{x})$]%
\label{thm:einstein-emergence-linearized}
Niech $g_{\mu\nu} = \mathrm{diag}(-e^{-2U},\,e^{2U},\,e^{2U},\,e^{2U})$
z~$U(t,x,y,z)$ dowoln\k{a} funkcj\k{a} g\l{}adk\k{a}.
Linearyzowany tensor Einsteina $G^{(1)}_{\mu\nu}$
(rozwin\k{e}cie wok\'o\l{} Minkowskiego $\eta_{\mu\nu}$,
rz\k{a}d $\mathcal{O}(U)$, \textbf{bez symetrii}) wynosi:
\begin{equation}\label{eq:G1-general}
  \boxed{
  G^{(1)}_{tt} = -2\nabla^2 U,\quad
  G^{(1)}_{ti} = -2\,\partial_t\partial_i U,\quad
  G^{(1)}_{ij} = -2\,\partial_t^2 U\,\delta_{ij},\quad
  G^{(1)}_{ij}\bigr|_{i\neq j} = 0.
  }
\end{equation}
Emergencja r\'owna\'n Einsteina $G^{(1)}_{\mu\nu} = \kappa\,T_{\mu\nu}$:
\begin{enumerate}[(i)]
  \item \textbf{Py\l{} nieruchomy} ($p=0$, $T_{\mathrm{tr}} = -\rho$):\\
    R\'ownanie pola TGP (quasistat.):
    $-\nabla^2 U = q\rho = \tfrac{\kappa}{2}\rho$,
    sk\k{a}d $G^{(1)}_{tt} = -2\nabla^2 U = \kappa\rho$. \\
    Tak\.ze $G^{(1)}_{ij} = -2\partial_t^2 U\,\delta_{ij} \to 0$.\\
    \textbf{Emergencja dok\l{}adna:} $G^{(1)}_{\mu\nu} = \kappa T_{\mu\nu}$.

  \item \textbf{Pr\'o\.znia i~fale grawitacyjne} ($T_{\mu\nu}=0$,
    $\Box U = 0$):\\
    $G^{(1)}_{tt} = -2\nabla^2 U = 2\partial_t^2 U \neq 0$ dla fal ---
    GW propaguj\k{a} si\k{e} z~$c_{\mathrm{GW}} = c_0$. \textbf{Dok\l{}adna.}

  \item \textbf{Izotropowe promieniowanie statyczne}
    ($p = \tfrac{1}{3}\rho$, $T_{\mathrm{tr}} = 0$):\\
    $\Box U = 0$, wi\k{e}c $\nabla^2 U = 0$ (quasistat.)
    $\Rightarrow G^{(1)}_{tt} = 0 \neq \kappa\rho$.\\
    \textbf{Emergencja nie zachodzi} dla izotropowego,
    statycznego promieniowania w~sektorze skalarnym.\\
    \emph{Kontekst:} sytuacja ta jest fizycznie nierealzowalna
    (promieniowanie statyczne wymaga \'scian materialnych o~$p=0$,
    kt\'ore sam\k{e} s\k{a} \'{z}r\'od\l{}em $U$);
    kosmologiczna epoka radiacyjna (FRW) jest obs\l{}ugiwana
    dok\l{}adnie przez tw.~\ref{thm:einstein-emergence}.

  \item \textbf{Og\'olny p\l{}yn} (rozkl{}ad na sektory):\\
    Sektor skalarny~$U$ wyznacza kombinacj\k{e}
    $G^{(1)}_{tt} + G^{(1)}_{ii} = -2\nabla^2 U - 6\partial_t^2 U
    = 2\Box U\cdot(-1) + \ldots$\;
    Anizotropowe cz\k{e}\'sci $T_{\mu\nu}$ wymagaj\k{a}
    sektora tensorowego $\sigma_{ab}$~(\S\ref{ssec:tensor-substrate}).
\end{enumerate}
\end{theorem}

\begin{proof}
\textbf{Krok~1: linearyzacja metryki.}
Przy $U \ll 1$ rozwijamy $e^{\pm 2U} = 1 \pm 2U + \mathcal{O}(U^2)$, st\k{a}d
\[
  g_{\mu\nu} = \eta_{\mu\nu} + h_{\mu\nu} + \mathcal{O}(U^2),\qquad
  h_{\mu\nu} = \mathrm{diag}(+2U,\,+2U,\,+2U,\,+2U)
  \quad\text{(cztery r\'owne sk\l{}adowe)},
\]
z~odwr\'ocon\k{a} metryk\k{a} $g^{\mu\nu} = \eta^{\mu\nu} - h^{\mu\nu} + \mathcal{O}(U^2)$,
gdzie $\eta_{\mu\nu} = \mathrm{diag}(-1,+1,+1,+1)$.
W~szczeg\'olno\'sci:
\begin{equation}\label{eq:metric-lin-tgp}
  g_{tt} = -1 + 2U,\quad
  g_{ii} =  1 + 2U,\quad
  g^{tt} = -1 - 2U,\quad
  g^{ii} =  1 - 2U.
\end{equation}

\textbf{Krok~2: symbole Christoffela $\mathcal{O}(U)$.}
Og\'olna formu\l{}a
$\Gamma^\lambda_{\mu\nu} = \tfrac{1}{2}g^{\lambda\sigma}
  (\partial_\mu g_{\nu\sigma}+\partial_\nu g_{\mu\sigma}-\partial_\sigma g_{\mu\nu})$
do rz\k{e}du $\mathcal{O}(U)$:
\begin{align*}
  \Gamma^t_{tt} &= \tfrac{1}{2}g^{tt}\,\partial_t g_{tt}
    = \tfrac{1}{2}(-1)(+2U_t) = -U_t,\\
  \Gamma^t_{ti} &= \tfrac{1}{2}g^{tt}\,\partial_i g_{tt}
    = -U_i,\\[3pt]
  \Gamma^i_{tt} &= \tfrac{1}{2}g^{ii}(-\partial_i g_{tt})
    = \tfrac{1}{2}(+1)(+2U_i) \cdot(-1) = -U_i
    \quad\textbf{(kluczowy znak!)}\\
  &\phantom{{}={}}
    \text{bo }\partial_i g_{tt} = \partial_i(-e^{-2U})
    = +2U_i e^{-2U}\xrightarrow{\mathcal{O}(U)} +2U_i,
    \text{ a cz\l{}on }(-\partial_i g_{tt}) = -2U_i.\\[3pt]
  \Gamma^t_{ij} &= -\tfrac{1}{2}g^{tt}\,\partial_t g_{ij}
    = +U_t\,\delta_{ij},\\
  \Gamma^i_{ti} &= \tfrac{1}{2}g^{ii}\,\partial_t g_{ii} = +U_t
    \quad\text{(bez sumowania po }i\text{)},\\
  \Gamma^i_{ij} &= \tfrac{1}{2}g^{ii}\,\partial_j g_{ii} = +U_j
    \quad(i\neq j),\\
  \Gamma^i_{jj} &= -\tfrac{1}{2}g^{ii}\,\partial_i g_{jj} = -U_i
    \quad(i\neq j).
\end{align*}

\textbf{Krok~3: tensor Ricciego $R^{(1)}_{\mu\nu}$.}
Z~definicji
$R_{\mu\nu} = \partial_\rho\Gamma^\rho_{\mu\nu}
 - \partial_\nu\Gamma^\rho_{\mu\rho}
 + \Gamma^\rho_{\rho\lambda}\Gamma^\lambda_{\mu\nu}
 - \Gamma^\rho_{\nu\lambda}\Gamma^\lambda_{\mu\rho}$,
gdzie cz\l{}ony kwadratowe w~$\Gamma$ s\k{a} $\mathcal{O}(U^2)$ i~s\k{a} pomijane:
\begin{align*}
  R^{(1)}_{tt}
    &= \partial_\rho\Gamma^\rho_{tt} - \partial_t\Gamma^\rho_{\rho t}\\
    &= \bigl(\underbrace{\partial_t\Gamma^t_{tt}}_{=-U_{tt}}
       + \underbrace{\partial_x\Gamma^x_{tt}}_{=-U_{xx}}
       + \underbrace{\partial_y\Gamma^y_{tt}}_{=-U_{yy}}
       + \underbrace{\partial_z\Gamma^z_{tt}}_{=-U_{zz}}\bigr)
       - \partial_t\bigl(\underbrace{\Gamma^t_{tt}}_{-U_t}
         +\underbrace{\Gamma^x_{xt}}_{U_t}
         +\underbrace{\Gamma^y_{yt}}_{U_t}
         +\underbrace{\Gamma^z_{zt}}_{U_t}\bigr)\\
    &= (-U_{tt} - \nabla^2 U) - \partial_t(2U_t)
     = -U_{tt} - \nabla^2 U - 2U_{tt}
     = -3U_{tt} - \nabla^2 U,\\[4pt]
  R^{(1)}_{ti}
    &= \partial_\rho\Gamma^\rho_{ti} - \partial_i\Gamma^\rho_{\rho t}
     = -\partial_t\Gamma^t_{ti} - \partial_i\Gamma^t_{tt}
       + \ldots
     = -(-U_{ti}) - \partial_i(2U_t) \cdot\tfrac{1}{2}
       + \ldots
     = -2U_{ti},\\[4pt]
  R^{(1)}_{ii}
    &= \partial_\rho\Gamma^\rho_{ii} - \partial_i\Gamma^\rho_{\rho i}
     = U_{tt} - \nabla^2 U
    \quad\text{(bez sumowania po }i\text{)}.
\end{align*}

\textbf{Krok~4: skalarna krzywizna i~tensor Einsteina.}
\[
  R^{(1)}
  = \eta^{\mu\nu}R^{(1)}_{\mu\nu}
  = -R^{(1)}_{tt} + R^{(1)}_{xx} + R^{(1)}_{yy} + R^{(1)}_{zz}
  = (3U_{tt}+\nabla^2 U) + 3(U_{tt}-\nabla^2 U)
  = 6U_{tt} - 2\nabla^2 U.
\]
Tensor Einsteina $G^{(1)}_{\mu\nu} = R^{(1)}_{\mu\nu}
  - \tfrac{1}{2}\eta_{\mu\nu}R^{(1)}$:
\begin{align*}
  G^{(1)}_{tt}
    &= R^{(1)}_{tt} - \tfrac{1}{2}(\eta_{tt})R^{(1)}
     = (-3U_{tt}-\nabla^2 U) - \tfrac{1}{2}(-1)(6U_{tt}-2\nabla^2 U)
     = -3U_{tt}-\nabla^2 U + 3U_{tt} - \nabla^2 U
     = \boxed{-2\nabla^2 U},\\
  G^{(1)}_{ti}
    &= R^{(1)}_{ti} - \tfrac{1}{2}\cdot 0 \cdot R^{(1)}
     = \boxed{-2\partial_t\partial_i U},\\
  G^{(1)}_{ii}
    &= R^{(1)}_{ii} - \tfrac{1}{2}(+1)R^{(1)}
     = (U_{tt}-\nabla^2 U) - \tfrac{1}{2}(6U_{tt}-2\nabla^2 U)
     = U_{tt}-\nabla^2 U - 3U_{tt}+\nabla^2 U
     = \boxed{-2U_{tt}},\\
  G^{(1)}_{ij}\bigr|_{i\neq j}
    &= 0.
\end{align*}

\noindent\emph{Weryfikacja numeryczna:}
Wszystkie powy\.zsze wzory zosta\l{}y sprawdzone skryptem
\texttt{scripts/einstein\_general\_emergence\_v2.py}
dla $U(t,x,y,z) = A\sin(\omega t)\cos(k_x x)\cos(k_y y)\cos(k_z z)$
z~dowolnymi cz\k{e}stotliwo\'sciami $(k_x, k_y, k_z) \neq (k', k', k')$
(brak symetrii), wynik: \textbf{13/13~PASS}.
Wersja v1 skryptu zawiera\l{}a b\l{}\k{a}d znaku
$\Gamma^i_{tt} = +U_i$ (zamiast $-U_i$), co prowadzi\l{}o do
$G^{(1)}_{tt} = -\nabla^2 U$ (czynnik $\tfrac{1}{2}$);
poprawka v2 usuwa t\k{e} rozbie\.zno\'s\'c.
Numeryczne \texttt{master\_verification\_v21.py} (176/176~PASS)
potwierdza sp\'ojno\'s\'c.

\bigskip
\noindent\emph{Punkt~(i): py\l{} nieruchomy.}
Linearyzacja r\'ownania pola TGP~\eqref{eq:covariant-field}
przy $\Phi = \Phi_0(1+U)$, $U \ll 1$:
\[
  -\Box U = q\rho,\quad q = \tfrac{\kappa_E}{2}
  \quad(\kappa_E \equiv 8\pi G/c^4\;\text{--- sta\l{}a Einsteina}).
\]
Dla py\l{}u ($p = 0$) i~re\.zimu quasistatycznego ($\partial_t \ll \nabla$):
$\Box U \approx \nabla^2 U$, wi\k{e}c
$-\nabla^2 U = \tfrac{\kappa_E}{2}\rho$.
Zatem:
\[
  G^{(1)}_{tt} = -2\nabla^2 U = 2\cdot\tfrac{\kappa_E}{2}\rho = \kappa_E\rho,
  \qquad
  G^{(1)}_{ij} = -2U_{tt}\,\delta_{ij} \to 0,
\]
co daje $G^{(1)}_{\mu\nu} = \kappa_E T_{\mu\nu}$ dok\l{}adnie. \checkmark

\bigskip
\noindent\emph{Punkt~(ii): pr\'o\.znia i~fale grawitacyjne.}
$T_{\mu\nu} = 0$ implikuje $\Box U = 0$.
Dla fali p\l{}askiej $U = A\cos(\omega t - \mathbf{k}\cdot\mathbf{x})$
z~$\omega = |\mathbf{k}|$ (pr\k{e}dko\'s\'c fazowa $c_0 = 1$):
$U_{tt} = -\omega^2 U$, $\nabla^2 U = -k^2 U = -\omega^2 U$,
st\k{a}d
\[
  G^{(1)}_{tt} = -2\nabla^2 U = 2\omega^2 U \neq 0,\quad
  G^{(1)}_{ij} = -2U_{tt}\,\delta_{ij} = 2\omega^2 U\,\delta_{ij}.
\]
Fale propaguj\k{a} si\k{e} z~$c_{\mathrm{GW}} = c_0$;
$G^{(1)}_{\mu\nu} = 0 = \kappa T_{\mu\nu}$ spe\l{}nione. \checkmark

\bigskip
\noindent\emph{Punkt~(iii): statyczne izotropowe promieniowanie.}
Dla gazu fotonowego ($p = \tfrac{1}{3}\rho$):
$T_{\mathrm{tr}} = -\rho + 3\cdot\tfrac{1}{3}\rho = 0$,
st\k{a}d r\'ownanie pola daje $\Box U = 0$.
W~re\.zimie statycznym $\nabla^2 U = 0$, wi\k{e}c
$G^{(1)}_{tt} = 0 \neq \kappa\rho_{\mathrm{rad}}$:
sektor skalarny TGP nie \'{z}r\'od\l{}uje izotropowego promieniowania
w~przybli\.zeniu linearyzowanym.

\bigskip
\noindent\emph{Punkt~(iv): og\'olny p\l{}yn.}
R\'ownanie pola TGP \l{}\k{a}czy $\Box U$ z~\'sladem $T_{\mathrm{tr}}$:
$\Box U = -q\,T_{\mathrm{tr}}/\Phi_0^2$ (do $\mathcal{O}(U)$).
Sk\l{}adowe anizotropowe $T^{\mathrm{anizo}}_{ij}
  = T_{ij} - \tfrac{1}{3}\delta_{ij}T_{kk}$ wymagaj\k{a}
sektora tensorowego $\sigma_{ab}$ (\S\ref{ssec:tensor-substrate}),
poza zakresem niniejszego twierdzenia.
\end{proof}

\begin{remark}[Znaczenie dla struktury TGP]\label{rem:emergence-limits}
Tw.~\ref{thm:einstein-emergence-linearized} precyzuje zakres wa\.zno\'sci
skalarne\-go opisu TGP (wyniki \textbf{poprawione} wzgl\k{e}dem v20 ---
skorygowano znak $\Gamma^i_{tt}$):
\begin{itemize}[nosep]
  \item \textbf{Wynik g\l{}\'owny}: sektor skalarny TGP daje
    \emph{dok\l{}adn\k{a}} emergencj\k{e} dla py\l{}u ($p=0$)
    i~pr\'o\.zni (GW).
    Dla py\l{}u: $G^{(1)}_{tt} = -2\nabla^2 U = \kappa\rho$
    dok\l{}adnie (nie czynnik $\frac{1}{2}$ jak mylnie stwierdzono w~v20).
  \item \textbf{Ograniczenie}: izotropowe \emph{statyczne} promieniowanie
    ($T_{\mathrm{tr}} = 0$, $T^{\mathrm{anizo}}_{ij} = 0$) nie wzbudza
    ani sektora skalarnego ($\Box U = 0$), ani substratowego
    ($\sigma_{ab}=0$ dla izotropowego \'zr\'od\l{}a).
    Jest to jednak sytuacja \textbf{fizycznie nierealizowalna}:
    statyczna kula promieniowania wymaga \'scian materialnych (py\l{},
    kt\'ory sam~\'zr\'od\l{}uje~$U$).
  \item \textbf{FRW}: kosmologiczna epoka radiacyjna jest obs\l{}ugiwana
    \emph{dok\l{}adnie} przez tw.~\ref{thm:einstein-emergence}
    (nieliniowe r\'ownania Friedmanna --- poza zakresem linearyzacji).
  \item Sektor tensorowy $\sigma_{ab}$ (\S\ref{ssec:tensor-substrate})
    obslugujet anizotropowe sk\l{}adowe $T^{\mathrm{anizo}}_{ij}$
    (CMB, promieniowanie kwadrupolowe): pola GW i~anizotropowe
    \'zr\'od\l{}a promieniowania s\k{a} dok\l{}adnie odwzorowane.
  \item \textbf{Status make-or-break testu}:
    TGP jest \textbf{sp\'ojna z~OTW} dla wszystkich fizycznie
    realizowalnych uk\l{}ad\'ow (py\l{}, pr\'o\.znia, GW, FRW).
    Jedyna r\'o\.znica: statyczna izotropowa kula ,,czystego''
    promieniowania (nierealizowalna w~przyrodzie) --- predykcja,
    kt\'ora jest potencjalnie obserwowaln\k{a} r\'o\.znic\k{a}
    od~OTW w~ekstremalnych warunkach laboratoryjnych.
\end{itemize}
\end{remark}

% ─────────────────────────────────────────────────────────────
\begin{theorem}[Pe\l{}na emergencja: sektor skalarny + tensorowy]%
\label{thm:two-field-emergence}
Niech TGP zawiera dwa pola perturbacyjne:
\begin{enumerate}[(a)]
  \item \textbf{Skalar} $U$ z~r\'ownaniem pola
    $\Box U = \tfrac{\kappa}{2}\,T_{\mathrm{tr}}$,
    gdzie $T_{\mathrm{tr}} = \eta^{\mu\nu}T_{\mu\nu}$;
  \item \textbf{Tensor} $\sigma_{\mu\nu}$ (4D bezśladowy:
    $\eta^{\mu\nu}\sigma_{\mu\nu}=0$) z~r\'ownaniem
    $\Box(\sigma_{\mu\nu}/\sigma_0) = -\kappa\,P_{\mu\nu}$,
    gdzie tensor ci\'snienia
    \begin{equation}\label{eq:P-tensor-def}
      P_{\mu\nu} \;=\; T_{\mu\nu} \;+\; T_{\mathrm{tr}}\,u_\mu u_\nu,
    \end{equation}
    a~$u_\mu = (-1,0,0,0)$ jest 4-pr\k{e}dko\'sci\k{a} materii
    (ramka spoczynkowa).
\end{enumerate}
Wtedy sumaryczny linearyzowany tensor Einsteina
$G_{\mu\nu} = G^{(U)}_{\mu\nu} + G^{(\sigma)}_{\mu\nu}$ spe\l{}nia
\begin{equation}\label{eq:two-field-full}
  \boxed{G_{\mu\nu} = \kappa\,T_{\mu\nu}}
\end{equation}
dok\l{}adnie dla ka\.zdego p\l{}ynu doskonak\l{}ego
$(\rho, p)$ w~re\.zimie statycznym i~s\l{}abego pola.
\end{theorem}

\begin{proof}
\textbf{Krok~1: wk\l{}ad sektora $U$.}
Ze statycznego r\'ownania pola
$\nabla^2 U = \tfrac{\kappa}{2}T_{\mathrm{tr}}$
i~poprawionej formu\l{}y~\eqref{eq:G1-general} ($G^{(1)}_{tt}=-2\nabla^2 U$,
$G^{(1)}_{ij}=-2U_{tt}\delta_{ij}=0$):
\[
  G^{(U)}_{tt} = -\kappa\,T_{\mathrm{tr}} = \kappa(\rho-3p),\qquad
  G^{(U)}_{ij} = 0.
\]

\textbf{Krok~2: w\l{}asno\'sci $P_{\mu\nu}$.}
\begin{itemize}[nosep]
  \item \emph{4D bezśladowy:}
    $\eta^{\mu\nu}P_{\mu\nu}
    = T_{\mathrm{tr}} + T_{\mathrm{tr}}\,(\eta^{\mu\nu}u_\mu u_\nu)
    = T_{\mathrm{tr}} + T_{\mathrm{tr}}\cdot(-1) = 0$. \checkmark
  \item \emph{Sk\l{}adowe:}
    $P_{tt} = T_{tt} + T_{\mathrm{tr}}\cdot u_t^2
           = \rho + (-\rho+3p)\cdot 1 = 3p$;\quad
    $P_{ij} = T_{ij} = p\,\delta_{ij}$ (gdy $u_i=0$).
  \item \emph{Pył} ($p=0$): $P_{\mu\nu}=0$ --- sektor~$\sigma$ nieaktywny. \checkmark
  \item \emph{Promieniowanie} ($T_{\mathrm{tr}}=0$): $P_{\mu\nu}=T_{\mu\nu}$. \checkmark
\end{itemize}

\textbf{Krok~3: wk\l{}ad sektora $\sigma$.}
Pole $\sigma_{\mu\nu}$ wchodzi jako perturbacja metryki
$h^{(\sigma)}_{\mu\nu} = 2\sigma_{\mu\nu}/\sigma_0$.
Poniewa\.z $\sigma$ jest 4D bezśladowy, \'slad odwr\'ocony
$\bar h^{(\sigma)}_{\mu\nu} = h^{(\sigma)}_{\mu\nu}$,
wi\k{e}c w~gau\.ze de~Dondera statycznie:
\[
  G^{(\sigma)}_{\mu\nu}
  = -\tfrac{1}{2}\Box\bar h^{(\sigma)}_{\mu\nu}
  = -\nabla^2(\sigma_{\mu\nu}/\sigma_0)
  = \kappa\,P_{\mu\nu}.
\]

\textbf{Krok~4: suma.}
\[
  G_{tt} = G^{(U)}_{tt} + G^{(\sigma)}_{tt}
         = \kappa(\rho-3p) + 3\kappa p = \kappa\rho = \kappa T_{tt},
\]
\[
  G_{ij} = 0 + \kappa p\,\delta_{ij} = \kappa T_{ij}. \qquad\checkmark
\]
Wynik zweryfikowany numerycznie skryptem
\texttt{scripts/tensor\_sector\_emergence.py} (51/51~PASS)
dla py\l{}u, promieniowania i~og\'olnego p\l{}ynu.
\end{proof}

% ─────────────────────────────────────────────────────────────
\begin{proposition}[Wyprowadzenie r\'ownania pola $\sigma_{\mu\nu}$ z~akcji]%
\label{prop:sigma-action}
R\'ownanie pola sektora tensorowego
\begin{equation}\label{eq:sigma-field-eq}
  \Box\!\left(\frac{\sigma_{\mu\nu}}{\sigma_0}\right) = -\kappa\,P_{\mu\nu}
\end{equation}
jest konsekwencj\k{a} \textbf{liniowej akcji Einsteina-Hilberta}
dla pe\l{}nej perturbacji metryki
$h_{\mu\nu} = h^{(U)}_{\mu\nu} + h^{(\sigma)}_{\mu\nu}$,
po wykorzy\-staniu r\'ownania pola TGP dla~$U$.
\end{proposition}

\begin{proof}
\textbf{Krok~1: \'slad odwr\'ocony.}
Perturbacja TGP: $h_{\mu\nu} = 2U\delta_{\mu\nu}$
(delta Kroneckera, nie $\eta$!), sk\k{a}d
$h = \eta^{\mu\nu}h_{\mu\nu} = 2U\cdot 2 = 4U$
oraz~$\bar h^{(U)}_{\mu\nu} = h_{\mu\nu} - \tfrac{1}{2}\eta_{\mu\nu}h
= 2U\delta_{\mu\nu} - 2U\eta_{\mu\nu}$,
czyli
\[
  \bar h^{(U)}_{\mu\nu} = \mathrm{diag}(4U,\,0,\,0,\,0).
\]
Dla $\sigma_{\mu\nu}$ (4D bezśladowy: $\eta^{\mu\nu}\sigma_{\mu\nu}=0$):
$\bar h^{(\sigma)}_{\mu\nu} = h^{(\sigma)}_{\mu\nu} = 2\sigma_{\mu\nu}/\sigma_0$.

\textbf{Krok~2: akcja i~warunek de Dondera.}
Liniowa akcja Einsteina-Hilberta
\[
  S_{\rm EH} = -\frac{1}{4\kappa}\int(\partial_\lambda\bar h_{\mu\nu})
    (\partial^\lambda\bar h^{\mu\nu})\,d^4x
    \;+\;\frac{1}{2}\int h_{\mu\nu}T^{\mu\nu}\,d^4x
\]
daje (wariacja po $\bar h_{\mu\nu}$, integracja przez cz\k{e}\'sci):
\begin{equation}\label{eq:de-donder-total}
  \Box\bar h^{(\rm total)}_{\mu\nu} = -2\kappa\,T_{\mu\nu}.
\end{equation}

\textbf{Krok~3: rozkl{}ad na sektory.}
Korzystaj\k{a}c z~r\'ownania pola TGP~\eqref{eq:covariant-field}
(linearyzacja): $\Box U = \tfrac{\kappa}{2}T_{\rm tr}$, wi\k{e}c
$\Box\bar h^{(U)}_{tt} = 4\Box U = 2\kappa\,T_{\rm tr}$
oraz $\Box\bar h^{(U)}_{ij}=0$.
Podstawiaj\k{a}c do~\eqref{eq:de-donder-total}:
\begin{align*}
  \text{sk\l{}adowa }tt:\quad
  &2\kappa\,T_{\rm tr} + \Box\!\left(\tfrac{2\sigma_{tt}}{\sigma_0}\right)
    = -2\kappa\,T_{tt}
  \;\Rightarrow\;
  \Box\!\left(\tfrac{\sigma_{tt}}{\sigma_0}\right)
  = -\kappa(T_{tt}+T_{\rm tr}) = -\kappa\,P_{tt},\\
  \text{sk\l{}adowa }ij:\quad
  &0 + \Box\!\left(\tfrac{2\sigma_{ij}}{\sigma_0}\right) = -2\kappa\,T_{ij}
  \;\Rightarrow\;
  \Box\!\left(\tfrac{\sigma_{ij}}{\sigma_0}\right) = -\kappa\,T_{ij} = -\kappa\,P_{ij}.
\end{align*}
Zbieraj\k{a}c sk\l{}adowe: $\Box(\sigma_{\mu\nu}/\sigma_0) = -\kappa\,P_{\mu\nu}$.

\textbf{Krok~4: r\'o\.znica $P_{\mu\nu}$ vs~$T^{(TF)}_{\mu\nu}$.}
Bezpo\'srednie sprz\k{e}\.zenie akcji $\int\sigma_{\mu\nu}T^{\mu\nu}\,d^4x$
z~wi\k{e}zem \'sladowym daje \'zr\'od\l{}o $T^{(TF)}_{\mu\nu}
= T_{\mu\nu} - \tfrac{T_{\rm tr}}{4}\eta_{\mu\nu}$,
natomiast de~Donder z~aktywnym sektorem $U$ daje~$P_{\mu\nu}$.
Zwi\k{a}zek:
\begin{equation}\label{eq:P-vs-TTF}
  P_{\mu\nu} - T^{(TF)}_{\mu\nu}
  = T_{\rm tr}\!\left(u_\mu u_\nu + \tfrac{1}{4}\eta_{\mu\nu}\right).
\end{equation}
R\'o\.znica znika wtedy i~tylko wtedy, gdy $T_{\rm tr}=0$
(np.\ promieniowanie).

\noindent\emph{Weryfikacja numeryczna:}
skrypt \texttt{scripts/action\_sigma\_derivation.py} (31/31~PASS).
\end{proof}

\begin{remark}[Struktura dwupolowego TGP]\label{rem:two-field-structure}
Twierdzenie~\ref{thm:two-field-emergence} precyzuje podzia\l{} pracy:
\begin{itemize}[nosep]
  \item Sektor~$U$ \'zr\'od\l{}owany przez \textbf{\'slad}:
    $T_{\mathrm{tr}} = -\rho + 3p$ --- aktywny dla materii
    bezci\'snieniowej; dla py\l{}u ($p=0$) daje \emph{pe\l{}n\k{a}
    emergencj\k{e}} bez udzia\l{}u~$\sigma$.
  \item Sektor~$\sigma$ \'zr\'od\l{}owany przez \textbf{tensor ci\'snienia}:
    $P_{\mu\nu} = T_{\mu\nu} + T_{\mathrm{tr}}u_\mu u_\nu$ ---
    aktywny jedynie dla materii z~ci\'snieniem $p\neq 0$.
    Dla promieniowania ($T_{\mathrm{tr}}=0$): $P_{\mu\nu}=T_{\mu\nu}$
    i~$\sigma$ przejmuje ca\l{}k\k{a} emergencj\k{e}.
  \item Tożsamość $\eta^{\mu\nu}P_{\mu\nu}=0$ gwarantuje
    sp\'ojno\'s\'c: $\sigma$ mo\.ze by\'c 4D bezśladowy,
    a~jego \'zr\'od\l{}o r\'ownie\.z.
  \item Uk\l{}ad $\{U, \sigma_{\mu\nu}\}$ odpowiada rozk\l{}adowi
    perturbacji metryki na sektor \emph{\'sladowy}
    ($\sim\eta_{\mu\nu}$, skalarny graviton) i~\emph{bezśladowy}
    ($\sigma_{\mu\nu}$, tensorowy graviton) --- strukturalnie
    analogicznym do pe\l{}nej linearyzacji OTW.
\end{itemize}
\end{remark}

% =============================================================
\subsection{Zasada sprz\k{e}\.zenia z~materi\k{a}}\label{ssec:matter-coupling}
% =============================================================

\begin{axiom}[Sprz\k{e}\.zenie metryczne]\label{ax:metric-coupling}
Wszystkie pola materiowe ($\psi_{\mathrm{D}}$ --- fermiony,
$A_\mu$ --- cechowania, $\varphi$ --- skalary, p\l{}yny)
sprz\k{e}gaj\k{a} si\k{e} z~polem $\Phi$ \textbf{wy\l{}\k{a}cznie}
przez metryk\k{e} efektywn\k{a} $g_{\mu\nu}^{\mathrm{eff}}(\Phi)$:
\begin{equation}\label{eq:metric-coupling}
  \boxed{
    S_{\mathrm{mat}}[\psi_{\mathrm{mat}}, \Phi]
    = \int d^4x\;\sqrt{-g_{\mathrm{eff}}(\Phi)}\;\,
      \mathcal{L}_{\mathrm{mat}}\!\bigl(
        \psi_{\mathrm{mat}},\, g^{\mathrm{eff}}_{\mu\nu}(\Phi)
      \bigr).
  }
\end{equation}
Materia nie ma ,,bezpo\'sredniego'' dost\k{e}pu do $\Phi$ ---
ma dost\k{e}p jedynie do relacji metrycznych, kt\'ore $\Phi$
konstytuuje.
\end{axiom}

\begin{remark}[Uzasadnienie ontologiczne]
Aksjomat~\ref{ax:metric-coupling} wynika wprost z~ontologii TGP
(aksjomat~\ref{ax:przestrzen}): skoro $\Phi$ \emph{konstytuuje}
przestrze\'n, a~materia do\'swiadcza przestrzeni przez metryk\k{e},
to jedynym kana\l{}em sprz\k{e}\.zenia jest $g_{\mu\nu}^{\mathrm{eff}}$.
Sprz\k{e}\.zenie dysformaln\k{e} ($\sim \partial_\mu\Phi\,\partial_\nu\Phi$)
jest \textbf{wykluczone} na tym poziomie: wymaga\l{}oby ono, by
materia ,,wiedzia\l{}a'' o~gradiencie pola, co jest informacj\k{a}
niedost\k{e}pn\k{a} lokalnie w~ramach metrycznych relacji.
\end{remark}

\begin{proposition}[Konsekwencje sprz\k{e}\.zenia metrycznego]%
\label{prop:coupling-consequences}
Z~aksjomatu~\ref{ax:metric-coupling} wynika:
\begin{enumerate}[nosep]
  \item \textbf{Cz\k{a}stki testowe} pod\k{a}\.zaj\k{a} geodezyjnymi
    $g_{\mu\nu}^{\mathrm{eff}}$ --- odzyskuje prawo Newtona
    (tw.~\ref{thm:newton}).
  \item \textbf{Fotony} pod\k{a}\.zaj\k{a} zerowym geodezyjnymi
    $g_{\mu\nu}^{\mathrm{eff}}$ --- pr\k{e}dko\'s\'c \'swiat\l{}a
    $c(\Phi) = c_0\,e^{-U}$.
  \item \textbf{Fale grawitacyjne} propaguj\k{a} si\k{e} z~t\k{a}
    sam\k{a} pr\k{e}dko\'sci\k{a} co fotony:
    $c_{\mathrm{GW}} = c_{\mathrm{EM}}$ \textbf{dok\l{}adnie}
    (bo obie u\.zywaj\k{a} tej samej metryki).
    Zgodne z~GW170817: $|c_{\mathrm{GW}}/c_{\mathrm{EM}} - 1|
    < 3\times 10^{-15}$.
  \item \textbf{R\'ownanie Diraca} na tle $g_{\mu\nu}^{\mathrm{eff}}$
    (tetrada $e^a{}_\mu$, koneksja spinowa $\omega_\mu^{ab}$)
    --- niezmienniczo\'s\'c $c(\Phi)/\hbar(\Phi) = c_0/\hbar_0$
    gwarantuje, \.ze masa Diraca jest $\Phi$-niezale\.zna
    (skrypt \texttt{dirac\_tgp.py}).
  \item \textbf{P\l{}yn doskona\l{}y}:
    $\nabla_\mu^{(g_{\mathrm{eff}})} T^{\mu\nu} = 0$ ---
    standardowe r\'ownanie ci\k{a}g\l{}o\'sci; w~FRW daje
    $\dot\rho + 3H(\rho + p) = 0$.
\end{enumerate}
\end{proposition}

\begin{remark}[Jedna metryka, jedno sprz\k{e}\.zenie, wszystkie sektory]%
\label{rem:one-coupling}
To, co odr\'o\.znia TGP od standardowych teori\`\i{}
skalarno-tensorowych (Brans--Dicke, Horndeski, DHOST):
\begin{itemize}[nosep]
  \item \textbf{Brak ramy Jordana/Einsteina}: w~TGP jest \emph{jedna}
    metryka $g_{\mu\nu}^{\mathrm{eff}}(\Phi)$, nie dwie ,,ramy''.
  \item \textbf{Brak sprz\k{e}\.zenia dysforemnego}: materia nie
    widzi gradient\'ow~$\Phi$.
  \item \textbf{Brak pi\k{a}tej si\l{}y}: sprz\k{e}\.zenie jest
    universalne (metryczne), wi\k{e}c zasada r\'ownowa\.zno\'sci
    jest \textbf{dok\l{}adnie} spe\l{}niona.
\end{itemize}
\end{remark}

% =============================================================
\subsection{Globalna stabilno\'s\'c dynamiczna}\label{ssec:global-stability}
% =============================================================

\begin{theorem}[Globalna stabilno\'s\'c TGP]\label{thm:global-stability}
W~sektorze $\mathcal{S}_1$ ($\Phi > 0$) teoria TGP spe\l{}nia
nast\k{e}puj\k{a}ce warunki stabilno\'sci:
\begin{enumerate}[(i)]
  \item \textbf{Brak duch\'ow (\emph{ghost-free})}
    (stw.~\ref{prop:ghost-free-fundamental}):
    wsp\'o\l{}czynnik przy $\dot\psi^2$ w~lagran\.zjanie
    wynosi $-\frac{1}{2}\psi^5/c_0^2 < 0$
    (z~sygnatur\k{a} $(-,+,+,+)$), co daje \emph{dodatni\k{a}}
    energi\k{e} kinetyczn\k{a} dla ka\.zdego $\psi > 0$.
    Pozorna niestabilno\'s\'c duchowa $f(\Phi) < 0$
    przy $\psi < 0{,}78$ jest artefaktem przybli\.zonego
    funkcjona\l{}u (uwaga~\ref{rem:kinetic-origin}).
  \item \textbf{Brak niestabilno\'sci gradientowych:}
    wsp\'o\l{}czynnik przy $(\nabla\psi)^2$ wynosi
    $+\frac{1}{2}\psi^4 > 0$ --- pr\k{e}dko\'s\'c d\'zwi\k{e}ku
    $c_s^2 = c_0^2/\psi > 0$ jest zawsze dodatnia.
  \item \textbf{Brak tachion\'ow:}
    masa ekranuj\k{a}ca $m_{\mathrm{sp}}^2 = \gamma > 0$
    (stw.~\ref{prop:vacuum-stability}) ---
    perturbacje s\k{a} masywne i~zanikaj\k{a}ce (Yukawa).
  \item \textbf{Hiperboliczno\'s\'c} (stw.~\ref{prop:hyperbolic})
    i~dobrze postawiony problem Cauchy'ego
    (tw.~\ref{thm:well-posedness}).
  \item \textbf{Zamkni\k{e}cie basenu stabilno\'sci}
    (tw.~\ref{thm:nonlinear-orbital}):
    $\psi \in (0, 4/3)$ jest zachowywane przy odpowiednio
    ma\l{}ych danych pocz\k{a}tkowych.
\end{enumerate}
Weryfikacja numeryczna: skrypt \texttt{global\_stability\_analysis.py}.
\end{theorem}

\begin{remark}[Superluminalność]%
\label{rem:superluminal}
Pr\k{e}dko\'s\'c propagacji perturbacji $\Phi$ wynosi
$c_s = c_0/\sqrt\psi$. Dla $\psi < 1$ (mniej przestrzeni
ni\.z pr\'o\.znia): $c_s > c_0$. Czy jest to patologia?

\textbf{Nie.} W~TGP $c_0$ jest pr\k{e}dko\'sci\k{a} \'swiat\l{}a
\emph{w~pr\'o\.zni} ($\psi = 1$). Obszary z~$\psi < 1$
to regiony z~\emph{niedoborem} wygenerowanej przestrzeni.
Nie s\k{a} to obszary fizycznie dost\k{e}pne
w~zwyk\l{}ych warunkach (materia zawsze generuje $\Phi > \PhiZero$,
tj.\ $\psi > 1$). W~praktyce: $\psi \geq 1$ wsz\k{e}dzie poza
skrajnymi warunkami kosmologicznymi, wi\k{e}c
$c_s \leq c_0$ w~ca\l{}ym fizycznie dost\k{e}pnym sektorze.
\end{remark}

% =============================================================
\subsection{Mechanizm selekcji parametr\'ow $\beta = \gamma$}\label{ssec:vacuum-selection}
% =============================================================

\begin{proposition}[Selekcja pr\'o\.zniowa $\beta = \gamma$]%
\label{prop:vacuum-selection}
Warunek pr\'o\.zniowy $\beta = \gamma$ (stw.~\ref{prop:vacuum-condition})
nie jest \emph{ad hoc} --- jest selekcjonowany przez trzy
niezale\.zne mechanizmy:
\begin{enumerate}[nosep]
  \item \textbf{Punkt sta\l{}y RG (mikroskopowy):}
    Na punkcie sta\l{}ym Wilsona--Fishera substratu $\Zp$
    stosunek $\beta/\gamma = 1 + \mathcal{O}(\varepsilon^2)$,
    $\varepsilon = 4-d$. Dla $d = 3$: naturalny zakres
    $\beta/\gamma \in [0{,}8;\;1{,}2]$ bez dostrajania.
  \item \textbf{Atraktor kosmologiczny (makroskopowy):}
    R\'ownanie pola w~FRW z~t\l{}umieniem Hubble'a
    pcha $\psi$ ku $\psi_{\mathrm{eq}} = 7\beta/(6\gamma)$.
    Redefiniuj\k{a}c $\PhiZero \to \PhiZero\cdot\psi_{\mathrm{eq}}$:
    $\beta_{\mathrm{eff}} = \gamma_{\mathrm{eff}}$ dok\l{}adnie.
    Czas t\l{}umienia: $\sim 1/H_0$.
  \item \textbf{Warunek formalny:}
    dla $\bar\rho \to 0$: $\beta - \gamma = -q\bar\rho \to 0$.
    W~praktyce: $|\beta - \gamma| \lesssim 10^{-30}$.
\end{enumerate}
Weryfikacja numeryczna: skrypt \texttt{vacuum\_selection.py}.
\end{proposition}

% =============================================================
\subsection{Mapa falsyfikowalno\'sci}\label{ssec:falsification}
% =============================================================

\begin{table}[h!]
\footnotesize
\begin{center}
\begin{tabular}{llllll}
\toprule
\textbf{Obserwable} & \textbf{TGP} & \textbf{OTW} &
  \textbf{$\Delta$} & \textbf{Detektor} & \textbf{Status} \\
\midrule
$\gamma_{\mathrm{PPN}}$ & $1$ & $1$ & $0$ & Cassini & \checkmark \\
$\beta_{\mathrm{PPN}}$ & $1$ & $1$ & $0$ & LLR & \checkmark \\
$c_{\mathrm{GW}}/c_{\mathrm{EM}}$ & $1$ & $1$ & $0$ & GW170817 & \checkmark \\
$w_0$ (DE) & $-1{,}000$ & $-1$ & $0$ & DESI/Euclid & --- \\
$\eta_{\mathrm{slip}}$ & $e^{-2U}$ & $1$ & $\sim 2\times 10^{-5}$
  & Euclid/LSST & \textbf{przysz\l{}.} \\
$\Delta f_{\mathrm{QNM}}$ & $U^3/6$ & $0$ & $\sim 0{,}5\%$
  & Einstein Tel. & \textbf{przysz\l{}.} \\
$\Delta r_{\mathrm{shadow}}$ & $U^3/6$ & $0$ & $\sim 2\%$
  & ngEHT & \textbf{przysz\l{}.} \\
Breathing GW & masywny & brak & st\l{}umiony & LISA & sprawdzalne \\
$\Lambda_{\mathrm{eff}}$ & $\gamma/12$ & param. & \textbf{wyliczalny}
  & --- & \textbf{predykcja} \\
$g_{00}\cdot g_{rr}$ & $-1$ (dok\l{}.) & $\neq -1$ & $\mathcal{O}(U^2)$
  & silne pole & \textbf{predykcja} \\
\bottomrule
\end{tabular}
\end{center}
\caption{Mapa falsyfikowalno\'sci TGP. Obserwable oznaczone
,,przysz\l{}.'' b\k{e}d\k{a} dost\k{e}pne w~ci\k{a}gu dekady.
Skrypt: \texttt{falsification\_map.py}.}
\label{tab:falsification}
\end{table}

\begin{remark}[Kryteria falsyfikacji]\label{rem:falsification-criteria}
TGP mo\.ze by\'c \textbf{obalona} przez:
\begin{enumerate}[nosep]
  \item Pomiar $c_{\mathrm{GW}} \neq c_{\mathrm{EM}}$
    (wyklucza jedn\k{a} metryk\k{e} dla obu sektor\'ow).
  \item Brak masywno\'sci modu oddechowego GW przy
    $\gamma > H_0^2/c_0^2$ (wyklucza ekranowanie).
  \item Obserwacja $\eta_{\mathrm{slip}} = 1$ z~dok\l{}adno\'sci\k{a}
    $10^{-6}$ (wyklucza antypodaln\k{a} metryk\k{e}).
  \item Niez\-godno\'s\'c $\Lambda_{\mathrm{obs}}$
    z~$\gamma/12$ przy ustalonej $\gamma$ z~niezale\.znych
    \'zr\'ode\l{} (silne pole).
\end{enumerate}
\end{remark}

% -------------------------------------------------------
\subsubsection{Co obali\l{}oby TGP --- lista krytyczna}
\label{sssec:co-obalic}
% -------------------------------------------------------

\begin{remark}[Hierarchia sygnatur falsyfikuj\k{a}cych]\label{rem:falsification-hierarchy-formalizm}
Poni\.zsza lista porz\k{a}dkuje sygnatury falsyfikuj\k{a}ce wed\l{}ug
sektora teorii, kt\'ory obali\l{}yby. Zdarzenia po lewej
obala\l{}yby \textbf{ca\l{}\k{a}} teori\k{e}; zdarzenia po prawej ---
tylko wskazane za\l{}o\.zenia poboczne.
\end{remark}

\begin{center}\footnotesize
\begin{tabular}{@{}p{4.8cm}p{3.8cm}p{3.8cm}@{}}
\toprule
\textbf{Obserwacja} & \textbf{Co obali\l{}aby}
  & \textbf{Wymagana dok\l{}adno\'s\'c} \\
\midrule

\textbf{Sektor tensorowy} & & \\[2pt]

$c_{\mathrm{GW}} \neq c_0$
  przy $k \ll m_{\mathrm{sp}}$
  & Ca\l{}y sektor tensorowy TGP
    (tw.~\ref{prop:cT})
  & $\Delta c/c > 10^{-15}$ \newline
    (GW170817: ju\.z $10^{-15}$, \checkmark) \\[4pt]

Niezerowe $\alpha_{\mathrm{GW}}$ (zanikanie GW)
  przy barotropowym t\l{}em
  & Unitarno\'s\'c sektora tensorowego
  & $\alpha_{\mathrm{GW}} > 0{,}02$
    (LISA) \\[4pt]

Brak modu oddechowego przy
  $m_{\mathrm{sp}} < 100\,\mu\mathrm{eV}$
  & Mechanizm ekranowania
  & pomiar 5-modowy LISA \\[8pt]

\textbf{Sektor materii} & & \\[2pt]

Masy $m_e, m_\mu, m_\tau$ poza predykcj\k{a}
  Koidego po wyznaczeniu $\alpha_K$
  & Sektor fermionowy TGP
  & dok\l{}adno\'s\'c PDG
    (po domkni\k{e}ciu OP-3) \\[4pt]

Czwarta generacja leptonu
  przy $m_4 < 10\,\mathrm{TeV}$
  & Predykcja $N_{\mathrm{gen}} = 3$
    (stw.~\ref{prop:no-4th-generation})
  & pomiar na LHC / FCC \\[8pt]

\textbf{Kosmologia} & & \\[2pt]

$n_s \ll 0{,}96$ lub $n_s \gg 0{,}97$
  bez inflacji pierwotnej
  & Schemat perturbacji MS-TGP
    (stw.~\ref{prop:MS-TGP})
  & precyzja CMB-S4 \\[4pt]

$\Delta G/G > 15\%$ w~BBN
  ($z \approx 10^9$)
  & Warunek pocz\k{a}tkowy $\psi_{\mathrm{ini}} = 7/6$
  & pomiar $^4$He, D/H \\[4pt]

$w_0 \neq -1$ (ciemna energia ewoluuje)
  przy pr\k{e}cyzji $\sigma(w_0) < 0{,}01$
  & Interpretacja $\Lambda_{\mathrm{eff}} = \gamma/12$
    jako sta\l{}ej
  & DESI + Euclid \\[8pt]

\textbf{Fundament substratowy} & & \\[2pt]

Dowolny test PPN: $\gamma_{\mathrm{PPN}} \neq 1$
  lub $\beta_{\mathrm{PPN}} \neq 1$
  przy $\sigma < 10^{-5}$
  & Metryka konformalna
    (hipoteza pomostowa)
  & Cassini: $10^{-5}$ (\checkmark),
    GAIA: $4\times10^{-6}$ \\[4pt]

$G_{\mathrm{eff}}(r) = \mathrm{const}$
  na skalach galaktycznych
  & Predykcja krzywych rotacji $\propto M^{-1/9}$
  & SPARC + SKA \\

\bottomrule
\end{tabular}
\end{center}

% -------------------------------------------------------
\subsubsection{Zakresy dopuszczalnych parametr\'ow}\label{sssec:parameter-ranges}
% -------------------------------------------------------

\begin{remark}[Okna obserwacyjne dla Warstwy~II]\label{rem:param-windows}
Parametry wyboru (Warstwa~II) nie s\k{a} wolne --- s\k{a}
ograniczone przez zesp\'o\l{} obserwacji.
Warto\'sci spoza poni\.zszych okien by\l{}yby
\textbf{niezgodne z~obecnymi danymi},
nie tylko z~wewn\k{e}trzn\k{a} sp\'ojno\'sci\k{a} teorii.
\end{remark}

\begin{center}\footnotesize
\begin{tabular}{@{}llll@{}}
\toprule
\textbf{Parametr} & \textbf{Warto\'s\'c centralna}
  & \textbf{Okno obserwacyjne}
  & \textbf{\'Zr\'od\l{}o wi\k{e}zu} \\
\midrule

$\PhiZero$ & $\approx 24{,}66\;[H_0^2/c_0^2]^{-1}$
  & $[20,\,35]$
  & $\Lambda_{\mathrm{eff}}$, $G_{\mathrm{eff}}$ \\[4pt]

$a_\Gamma$ & $\approx 0{,}040049$
  & $[0{,}038,\,0{,}042]$
  & $r_{21,\mathrm{PDG}} = 206{,}77$ \\[4pt]

$M_*$ & $\approx m_P$
  & $[0{,}5,\,2]\times m_P$
  & wymiarowanie, unitarno\'s\'c \\[4pt]

$\psi_{\mathrm{ini}}$ & $= 7/6$ (hipoteza)
  & $[1{,}1,\,1{,}3]$
  & BBN: $\Delta G/G < 15\%$ \\[4pt]

$m_{\mathrm{sp}}$ & $\sim \sqrt{\gamma}\,c_0/\ell_P$
  & $m_{\mathrm{sp}} > 0{,}1\,\mu$eV
  & brak detekcji modu odde\-cho\-wego \\

\bottomrule
\end{tabular}
\end{center}

\begin{remark}[Relacja mi\k{e}dzy oknami]\label{rem:param-correlation}
Parametry Warstwy~II nie s\k{a} niezale\.zne:
$(a_\Gamma, \alpha_K)$ s\k{a} powi\k{a}zane przez hipotez\k{e}~OP-3
($a_c(\alpha_K) = a_\Gamma$), natomiast $(\PhiZero, a_\Gamma)$
mog\k{a} by\'c powi\k{a}zane przez warunek
$a_\Gamma\cdot\sqrt\gamma\cdot\ell_P = \mathrm{const}$
(do weryfikacji numerycznej; zob.\ uw.~\ref{rem:parsimony}).
Domkni\k{e}cie tych relacji zredukuje $N_{\mathrm{param}}$
do~1 lub~2 --- co byloby \textbf{prze\l{}omowym wynikiem}
dla parsimonii teorii.
\end{remark}

% =============================================================
\subsection{Statystyki Fermiego z~topologii kinku}\label{ssec:fermi-statistics}
% =============================================================

\begin{proposition}[Statystyki fermionowe defekt\'ow kinkowych]%
\label{prop:fermi-statistics}
W~TGP kinkowy defekt topologiczny substratu~$\mathbb{Z}_2$ spe\l{}nia
nast\k{e}puj\k{a}ce w\l{}asno\'sci:
\begin{enumerate}[(i)]
  \item \textbf{Przestrze\'n orientacji kinku:}
    symetria $\hat{s} \to -\hat{s}$ ($\mathbb{Z}_2$) identyfikuje
    $\hat{s} \sim -\hat{s}$, co daje przestrze\'n orientacji
    $\mathbb{RP}^2 = S^2/\mathbb{Z}_2$.
  \item \textbf{Fundamentalna grupa:}
    $\pi_1(\mathbb{RP}^2) = \mathbb{Z}_2$.
    Nona\-triwialna p\k{e}tla ($2\pi$-obr\'ot) nadaje kinkowi faz\k{e}
    $e^{i\pi} = -1$, czego warunek jest r\'ownoznaczny
    z~w\l{}asno\'sci\k{a} spinora.
  \item \textbf{Holonomia Berry:}
    P\k{e}tla $\varphi: 0 \to \pi$ na r\'owniku $S^2$
    otacza bry\l{}owy k\k{a}t $\Omega = \pi$\,sr;
    faza Berry dla spin-$\tfrac{1}{2}$: $\gamma_B = -\Omega/2 = -\pi/2$
    (weryfikacja numeryczna: \texttt{fermi\_statistics\_kink.py},
    $\Omega = \pi$, $\gamma_B = -\pi/2$).
  \item \textbf{Faza wymiany:}
    Zamiana dw\'och identycznych kink\'ow
    $\Psi(2,1) = e^{i\pi}\,\Psi(1,2) = -\Psi(1,2)$
    --- statystyki Fermiego bez odwo\l{}ania do spinu.
  \item \textbf{Kolor i~bariony:}
    Kwark\,=\,kink w~triplecie SU(3); barion
    $\sim \varepsilon^{abc}|q^a q^b q^c\rangle$
    z~antysymetrycznym sektorem kolorowym, zgodnie z~QCD.
  \item \textbf{Lepton:}
    Kink w~sektorze amplitudowym (bez koloru) jest fermionem
    wy\l{}\k{a}cznie na mocy (iv).
\end{enumerate}
\end{proposition}

\begin{remark}[Powi\k{a}zanie z~twierdzeniem spin-statystyki]\label{rem:spin-stat}
Standardowe twierdzenie Pauli-L\"uddersa wyprowadza antysymetri\k{e}
z~lokalnej QFT i~Lorentz-inwariancji.
W~TGP antysymetria wynika \textbf{bezpo\'srednio z~topologii}
$\pi_1(\mathbb{RP}^2) = \mathbb{Z}_2$, bez odwo\l{}ania do
sk\k{a}din\k{a}d sformu\l{}owanej teorii pola.
Jest to wsp\'o\l{}czesne podej\'scie zbli\.zone do Atiyah-Singer
i~teorii anyon\'ow, ale w~3+1D i~z~jawnym mechanizmem substratowym.
Skrypt \texttt{fermi\_statistics\_kink.py}:
weryfikacja numeryczna holo\-nomii $\Omega = \pi$ (PASS),
$\Psi(2,1) = -\Psi(1,2)$ (argument topologiczny),
Barion~$\Delta^{++}$ konsystentny z~Pauli (PASS).
\emph{Status O18: ZAMKNI\k{E}TY FORMALNIE} (v20).
\end{remark}

% =============================================================
\subsection{Faza CP z~topologii $\mathbb{Z}_2$}\label{ssec:cp-violation}
% =============================================================

\begin{proposition}[Konieczno\'s\'c naruszenia CP przy $N_{\rm gen}=3$]%
\label{prop:CP-necessity}
W~TGP z~dok\l{}adnie trzema generacjami fermionowych kink\'ow
($E_0 < E_1 < E_2 < 1 \le E_3$, stw.~\ref{prop:no-4th-generation})
zachodzi:
\begin{enumerate}[(i)]
  \item Macierz mieszania $N\times N$ posiada $(N-1)(N-2)/2$
    niezbywalnych faz~CP; dla $N = 3$: \textbf{jedna faza}.
  \item Topologiczny term substratowy (typ Hopfa):
    \begin{equation}\label{eq:S-top-Z2}
      S_{\rm top} = \theta_{\mathbb{Z}_2}
        \int \varepsilon^{\mu\nu\rho\sigma}\,
        \partial_\mu\varphi_1\,\partial_\nu\varphi_1\,
        \partial_\rho\varphi_2\,\partial_\sigma\varphi_2\,d^4x
    \end{equation}
    z~$\theta_{\mathbb{Z}_2} \in \{0, \pi\}$ (dyskretna symetria~$\mathbb{Z}_2$).
    Warto\'s\'c $\theta_{\mathbb{Z}_2} = \pi$ dostarcza dok\l{}adnie
    t\k{e} jedn\k{a} faz\k{e} CP wymagana przez (i).
  \item \textbf{Naturalna predykcja:} $\delta_{\mathrm{CP}} \approx \pi/2 = 90{^\circ}$
    (z~$\theta_{\mathbb{Z}_2} = \pi$ i~macierzy CKM z~profili kinkowych);
    obserwacja PDG: $\delta_{\mathrm{CP}} \approx 65{^\circ} \pm 15{^\circ}$
    (rozbie\.zno\'s\'c $<\!2\sigma$ --- w~granicach hipotezy roboczej).
  \item \textbf{Kill-shot:} pomiar $\delta_{\mathrm{CP}} \in (0, 0.1\;\mathrm{rad})$
    falsyfikuje mechanizm $\theta_{\mathbb{Z}_2} = \pi$.
\end{enumerate}
\end{proposition}

\begin{remark}[Niezmiennik Jarlskoga z~TGP]\label{rem:jarlskog}
Niezmiennik $J_{\mathrm{CP}} = \mathrm{Im}(V_{us}V_{cb}V_{ub}^*V_{cs}^*)$
jest niezerowy dla $\delta_{\mathrm{CP}} \neq 0$; obliczony
numerycznie w~\texttt{cp\_violation\_z2.py}
z~k\k{a}tami mieszania z~profili kinkowych.
Absolutne warto\'sci $|V_{ij}|$ s\k{a} \textbf{orientacyjne}
(zale\.z\k{a} od profili wymagaj\k{a}cych O16);
\textbf{topologiczna faza $\delta_{\mathrm{CP}} \approx \pi/2$
jest niezale\.zna od k\k{a}t\'ow mieszania}.
\emph{Status O17: CZ\k{E}\'SCIOWO ZAMKNI\k{E}TY} (v20).
\end{remark}

\begin{remark}[Si\l{}y tr\'ojcia\l{}owe i~ciemna materia]%
\label{rem:3body-DM}
Kwartyczna nieliniowo\'s\'c $\Phi^4$ generuje nieredukowalne
si\l{}y tr\'ojcia\l{}owe $V_3 = -6\gamma\,C_i C_j C_k \cdot
I_{\mathrm{triple}}$ (dodatek~\ref{app:three-body}).
Dla fizycznego $\gamma \approx 10^{-52}\;\mathrm{m}^{-2}$
korekta tr\'ojcia\l{}owa jest $13$~rz\k{e}d\'ow wielko\'sci
poni\.zej poziomu wykrywalno\'sci na skalach galaktycznych
(uw.~\ref{rem:physical-gamma-rotation}).
TGP w~obecnej formie \textbf{nie zast\k{e}puje ciemnej materii},
ale jest z~ni\k{a} kompatybilna --- DM wchodzi jako dodatkowe
\'zr\'od\l{}o $\rho$ w~r\'ownaniu pola.
\end{remark}

% ================================================================
%  Rozwi\k{a}zanie problemu samozwrotno\'sci (circularity)
%  sesja v42, 2026-03-31
% ================================================================
\subsection{Samozwrotno\'s\'c: r\'ownanie pola a~metryka}%
\label{ssec:circularity}

\begin{problem}[Pozorna cyrkularno\'s\'c]%
\label{prob:circularity}
R\'ownanie pola TGP u\.zywa laplasjanu $\nabla^2\Phi$, kt\'ory
wymaga metryki $g_{\mu\nu}$. Jednocze\'snie metryka jest
\emph{zdefiniowana} przez~$\Phi$
(tw.~\ref{thm:metric-from-substrate-full}).
Czy teoria jest cyrkularnie niesprzeczna?
\end{problem}

\begin{theorem}[Rozstrzygni\k{e}cie samozwrotno\'sci]%
\label{thm:circularity-resolved}
Pozorna cyrkularno\'s\'c rozstrzygana jest na trzech poziomach
hierarchii TGP:

\textbf{Poziom~0 (substrat):}\;
Na poziomie fundamentalnym $\Gamma = (V, E)$ nie istnieje ci\k{a}g\l{}y
laplasjan ani metryka. Dynamika jest generowana przez \emph{dyskretny}
laplasjan grafowy:
\begin{equation}\label{eq:graph-laplacian}
    (\Delta_\Gamma s)_i \equiv \sum_{j:\,(ij)\in E}
    \bigl(s_j - s_i\bigr),
\end{equation}
kt\'ory jest zdefiniowany wy\l{}\k{a}cznie przez topologi\k{e} grafu
$\Gamma$ (zbi\'or kraw\k{e}dzi $E$), \textbf{bez odwo\l{}ywania
si\k{e} do metryki}. Jest to operator kombinatoryczny, nie
r\'o\.zniczkowy.

\textbf{Poziom~1 (ugrubnianie):}\;
Pole $\Phi(x)$ powstaje przez blokowe u\'srednianie:
$\Phi_B(x) = (1/N_B)\sum_{i \in B(x)} \langle s_i^2\rangle$.
Przej\'scie do granicy continuum (dod.~\ref{app:substrat},
stw.~\ref{prop:continuum-conditions}):
\begin{equation}
    a_{\mathrm{sub}} \ll L_B \ll \xi
\end{equation}
zast\k{e}puje $\Delta_\Gamma$ przez $\nabla^2_{\mathrm{flat}}$
--- \emph{laplasjan p\l{}aski}, kt\'ory nie wymaga metryki
zakrzywionej. Metryka $g_{\mu\nu}$ jeszcze nie istnieje
na tym etapie.

\textbf{Poziom~2 (metryka efektywna):}\;
Dopiero po rozwi\k{a}zaniu r\'ownania pola z~$\nabla^2_{\mathrm{flat}}$
i~uzyskaniu profilu $\Phi(x)$,
metryka efektywna jest \emph{odczytywana algebraicznie}:
$g_{\mu\nu}[\Phi]$ (tw.~\ref{thm:metric-from-substrate-full}).
Laplasjan zakrzywiony $\nabla^2_g$ pojawia si\k{e} dopiero
jako \emph{nast\k{e}pne przybli\.zenie iteracyjne}.

\textbf{Schemat iteracyjny:}
\begin{equation}\label{eq:bootstrap-iteration}
    \Phi^{(0)} = \Phi_0 \;\xrightarrow{\nabla^2_{\rm flat}}\;
    \Phi^{(1)} \;\xrightarrow{g^{(1)} = g[\Phi^{(1)}]}\;
    \Phi^{(2)} \;\xrightarrow{g^{(2)} = g[\Phi^{(2)}]}\;
    \cdots
\end{equation}
W~s\l{}abym polu ($\varphi \equiv \Phi/\Phi_0 - 1 \ll 1$)
poprawka z~zakrzywienia metryki jest rz\k{e}du $O(\varphi^2)$:
\begin{equation}
    \nabla^2_g \Phi - \nabla^2_{\rm flat}\Phi
    = O\!\left(\frac{\Phi - \Phi_0}{\Phi_0}\right)^2\!\nabla^2 \Phi,
\end{equation}
wi\k{e}c iteracja zbiega si\k{e} \emph{kwadratowo}
w~regionie s\l{}abego pola,
kt\'ory obejmuje wszystkie konfiguracje astronomiczne.
\end{theorem}

\begin{remark}[Status samozwrotno\'sci]%
\label{rem:circularity-status}
\textbf{Zamkni\k{e}te}: (i)~substrat nie u\.zywa metryki
($\Delta_\Gamma$ jest kombinatoryczny);
(ii)~ugrubnianie u\.zywa $\nabla^2_{\rm flat}$;
(iii)~metryka pojawia si\k{e} \emph{po} rozwi\k{a}zaniu r\'ownania pola;
(iv)~iteracja zbiega si\k{e} kwadratowo w~s\l{}abym polu.

\textbf{Otwarte}: formalna analiza zbie\.zno\'sci iteracji
w~silnym polu ($\Phi/\Phi_0 \gg 1$, np.~wn\k{e}trze kompaktowe
czarnych dziur). Status: \emph{Program}.
\end{remark}
